\documentclass[11pt,a4paper]{article}
\usepackage[T1]{fontenc}
\usepackage[utf8]{inputenc}
\usepackage{lmodern}
\usepackage[margin=25mm,headheight=15pt]{geometry}
\usepackage{amsmath,amssymb,amsthm,mathtools,mathrsfs}
\usepackage{booktabs,tabularx,array,longtable}
\usepackage{enumitem,microtype,xcolor,fancyhdr}
\usepackage[hyphens]{url}
\usepackage{hyperref}
\definecolor{inkblue}{RGB}{27,48,72}
\hypersetup{colorlinks=true,linkcolor=inkblue,citecolor=inkblue,urlcolor=inkblue,
 pdftitle={Euler's Constant: Complete Integrated Research Manuscript},
 pdfauthor={KEIO},pdfsubject={Cumulative synthesis of foundations, positive moments, infinite families, and arithmetic transfer}}
\newcommand{\Q}{\mathbb Q}
\newcommand{\R}{\mathbb R}
\newcommand{\Z}{\mathbb Z}
\newcommand{\C}{\mathbb C}
\newcommand{\Qbar}{\overline{\mathbb Q}}
\newcommand{\Ga}{\mathbb G_a}
\newcommand{\Gm}{\mathbb G_m}
\newcommand{\Ein}{\operatorname{Ein}}
\newcommand{\E}{\mathbb E}
\newcommand{\Var}{\operatorname{Var}}
\newcommand{\den}{\operatorname{den}}
\newcommand{\ord}{\operatorname{ord}}
\newcommand{\Gal}{\operatorname{Gal}}
\newcommand{\tr}{\operatorname{tr}}
\DeclareMathOperator{\rank}{rank}
\DeclareMathOperator{\lcm}{lcm}
\newtheorem{theorem}{Theorem}[section]
\newtheorem{proposition}[theorem]{Proposition}
\newtheorem{lemma}[theorem]{Lemma}
\newtheorem{corollary}[theorem]{Corollary}
\theoremstyle{definition}
\newtheorem{definition}[theorem]{Definition}
\newtheorem{problem}[theorem]{Problem}
\newtheorem{example}[theorem]{Example}
\newtheorem{conjecture}[theorem]{Conjecture}
\theoremstyle{remark}
\newtheorem{remark}[theorem]{Remark}
\numberwithin{equation}{section}
\setlist[enumerate]{topsep=4pt}
\setlength{\emergencystretch}{3em}
\setlength{\parskip}{2pt}
\allowdisplaybreaks[2]
\pagestyle{fancy}
\fancyhf{}
\fancyhead[L]{\small KEIO}
\fancyhead[R]{\small Euler's constant: integrated research manuscript}
\fancyfoot[C]{\thepage}
\renewcommand{\headrulewidth}{0.3pt}
\setcounter{tocdepth}{1}
\title{\bfseries Euler's Constant:\\[4pt]
Complete Integrated Research Manuscript\\[9pt]
\large Exponential-Integral Foundations, Positive Moments,\\
Infinite Families, Local Arithmetic, and Transfer Obstructions}
\author{KEIO}
\date{Cumulative edition --- 10 September 2026}
\begin{document}
\maketitle
\begin{abstract}
This cumulative research manuscript brings together the gamma-related
foundations of companion manuscripts [001] and [002], the earlier
35-page positive-moment paper, the subsequent analytic and infinite-family
investigations, the algebraic/transcendental alternative, and the later
18-page arithmetic-transfer paper. Its purpose is to preserve the actual
mathematical content across these stages, with consistent hypotheses and
an explicit distinction between auxiliary theorems, published inputs,
finite certificates, conditional implications, and unfinished targets.
The constructions involve logarithmic elimination, positive moments,
primitive lattice norms, orthogonal polynomials, denominator arithmetic,
differential Galois theory, local root counts, probability, and arithmetic
Gevrey series. They yield concrete identities and obstructions for the
stated classes of constructions. They do not prove the irrationality or
transcendence of Euler's constant. A source-coverage ledger and the original
research packages accompany this edition so that corrections and retained
historical claims can be traced.
\end{abstract}
\noindent\textbf{Editorial status.} Research manuscript; no assertion of
external peer review or of novelty for every auxiliary result. The word
\emph{complete} refers to the identified source corpus and its gamma-related
mathematics, not to a completed transcendence proof or to all mathematics in
the two large companion manuscripts.
\clearpage
\tableofcontents
\clearpage

\section{Scope, status, and how the parts fit together}
\label{master:scope}

The common target is the following statement:
\begin{equation}\label{master:target}
 P(\gamma)\ne0\qquad
 \text{for every }0\ne P\in\Z[X],
 \qquad
 \gamma=\lim_{n\to\infty}(H_n-\log n).
\end{equation}
An exact identity, a positive integral, or a finite rational certificate
does not by itself establish this universally quantified assertion. The
purpose of integrating the research is to display which implications have
been proved and which additional implication would still be needed.

\paragraph{The source corpus.}
This edition uses 59 identified research files and packages, including the
complete source of the 35-page paper of 8 September and the 18-page
arithmetic-transfer revision of 10 September. The companion sources [001]
and [002] are the versions preserved in the earlier complete paper archive.
Their gamma-relevant foundations are integrated here; their complete
original texts are preserved in that archive. The earlier document recalled
as 27 pages has not been uniquely identified in the available corpus, so no
claim is made to have compared that particular pagination. The source
manifest records exactly what was available for this edition.

\paragraph{The integration rule.}
The earlier 35-page and 18-page mathematical bodies are retained, with
their labels and citations renamed to avoid collisions. Intervening notes
are synthesized by subject, with missing early results added separately.
Duplicate arguments are cross-referenced or shortened; statements corrected
in a later source are presented in their corrected form. The preserved
source packages contain the original wording and finite data. When a broad
companion assertion is not re-proved here, it is identified as a source
input and is not silently used to close the arithmetic problem.

\begin{center}\small
\begin{tabularx}{\linewidth}{@{}p{.14\linewidth}X@{}}
\toprule
Part & Mathematical role \\
\midrule
I & Gamma-related foundations from [001], [002], and the earliest certificates.\\
II & The complete positive-moment, Sondow, primitive-elimination, and finite-certificate development.\\
III & Additional early results absent from the 35-page synthesis.\\
IV & Later analytic variations, filters, residues, and comparisons with other methods.\\
V & Polynomial infinite families, exact arithmetic targets, and conditional transcendence criteria.\\
VI & Algebraic/transcendental alternatives, local counts, finite camouflage, and logical limits.\\
VII & Simultaneous approximation, tensor reduction, statistical cancellation, and Euler-quotient denominators.\\
VIII & Unified open implications, corrections, source coverage, and reproducibility.\\
\bottomrule
\end{tabularx}
\end{center}

\paragraph{Five different statuses.}
An \emph{auxiliary theorem} has the displayed hypotheses and a mathematical
proof. A \emph{published input} is attributed to its external source. A
\emph{finite certificate} concerns explicitly bounded data and can be
verified by exact arithmetic. A \emph{conditional implication} proves
``hypothesis implies conclusion'' without asserting its hypothesis for
$\gamma$. An \emph{open target} is a statement not proved by the construction.
Historical observations or approximate experimental exponents are not
silently promoted from the last two statuses to the first.

\paragraph{Notation is local when necessary.}
Throughout, $H_n=\sum_{j=1}^n1/j$, $H_0=0$, and $d_n=\lcm(1,\ldots,n)$.
When a symbol such as $D_N$, $P_N$, or $E$ is re-used in a different
construction, its definition is local to that section. The entire function
\[
 \Ein(z)=\sum_{n\ge1}\frac{(-1)^{n+1}z^n}{n\,n!}
\]
is kept distinct from a plus-sign polyexponential series. On the positive
real branch $\Ein(x)=\gamma+\log x+E_1(x)$, where
$E_1(x)=\int_x^\infty e^{-t}\,dt/t$. Thus
$\Ein(1)=\gamma+\delta/e$ for $\delta=eE_1(1)$.
No independence theorem for $\Ein(1)$ is interpreted as an independence
theorem for $\gamma$ without an additional argument.

\paragraph{What counts as progress toward the target.}
Three gaps recur: nonzero coefficient vectors need not have nonzero values;
fast analytic decay can disappear after taking primitive arithmetic
heights; and an identity of functions need not transfer its arithmetic
information to a chosen constant. The parts below address specific forms
of these gaps. None is a universal impossibility theorem for other methods.

\clearpage
\part{Foundations from the companion manuscripts}\label{master:foundations}
% BEGIN INTEGRATED PART I
\section{Analytic and arithmetic foundations from the companion papers}
\label{fd:sec:foundations}

This section integrates the gamma-related inputs of the companion manuscripts
\cite{pm:KEIO001,pm:KEIO002}.  Their broader results on independent spectra, inverse
monodromy, and algebraic models are retained in the source archive; they are
not additional proofs of the arithmetic nature of Euler's constant.
The order-one independence statement needed here is supplied with its
argument.  Statements involving all higher polyexponential orders are
identified separately as companion inputs.

We use
\begin{equation}\label{fd:eq:definitions}
 \gamma=\lim_{n\to\infty}(H_n-\log n),\qquad
 \Ein(z)=\sum_{n\ge1}\frac{(-1)^{n-1}z^n}{n\,n!},\qquad
 \delta=\int_0^\infty\frac{e^{-t}}{1+t}\,dt,
\end{equation}
where $H_n=\sum_{j=1}^n1/j$ and $H_0=0$.  On the principal slit plane,
\begin{equation}\label{fd:eq:connection}
 E_1(z)=\Ein(z)-\gamma-\log z,\qquad
 \Ein'(z)=\frac{1-e^{-z}}z,\qquad
 \eta:=\Ein(1)=\gamma+\delta/e.
\end{equation}
The normalization agrees with \cite[\S6.2]{fd:DLMF}.
The entire function $\Ein$ is an $E$-function: its normalized Taylor
coefficients are $(-1)^{n-1}/n$, with exponentially bounded common
denominators.  The function $E_1$ has a logarithmic singularity at zero
and is not an $E$-function.  In particular an $E$-function value theorem
must not be applied to $E_1(1)$ merely because \eqref{fd:eq:connection}
expresses it through an $E$-value and $\gamma$.

\subsection{The finite-point independence input and its exact scope}
Put
\[
 K_{\exp}=\Qbar(e^\alpha:\alpha\in\Qbar).
\]
Every coefficient of a polynomial over this field involves only finitely
many exponential values.

\begin{lemma}[Laurent derivative obstruction]\label{fd:lem:laurent}
Let $\nu_1,\ldots,\nu_d\in\Qbar$ be linearly independent over $\Q$,
and endow
\[
 R=\Qbar(z)[X_1^{\pm1},\ldots,X_d^{\pm1}],\qquad K=\operatorname{Frac}(R)
\]
with $D X_j=-\nu_jX_j$.  If finitely supported constants $c_u\in\Qbar$
satisfy
\[
 Dg=z^{-1}\sum_{u\in\Z^d}c_uX^u,\qquad g\in K,
\]
then every $c_u$ is zero.
\end{lemma}
\begin{proof}
First suppose an irreducible nonunit $P\in R$ divides $DP$.  Writing
$DP=AP$, comparison of Newton polytopes shows that $A\in\Qbar(z)$;
if $DP=0$, the following coefficient argument applies directly.
For two nonzero coefficients $a_u,a_v$ of $P$,
\[
 (a_u/a_v)'/(a_u/a_v)=(u-v)\cdot\nu.
\]
A rational function has no nonzero constant logarithmic derivative.
Thus $u=v$, so $P$ is a Laurent monomial times a unit, a contradiction.
If a reduced denominator of $g$ had a nonunit factor $P$, differentiating
would increase its pole order unless $P\mid DP$.  Hence $g\in R$.
Writing $g=\sum a_u(z)X^u$ gives
\[
 a_u'-(u\cdot\nu)a_u=c_u/z.
\]
For $u=0$, rational derivatives have zero residue.  For $u\ne0$, a finite
pole of $a_u$ creates a higher-order pole on the left; if it has no pole
at zero, that side is regular there.  A nonzero $c_u/z$ is therefore
impossible.  The same argument with zero right side identifies the
constant field of $K$ as $\Qbar$.
\end{proof}

\begin{theorem}[Multipoint $\Ein$ input]\label{fd:thm:multipoint}
For distinct $\lambda_1,\ldots,\lambda_s\in\Qbar^\times$, let
\[
d=\dim_{\Q}\operatorname{span}_{\Q}\{\lambda_1,\ldots,\lambda_s\}.
\]
Then
\[
 \operatorname{trdeg}_{\Qbar}
 \Qbar(e^{\lambda_i},\Ein(\lambda_i):1\le i\le s)=d+s.
\]
The only algebraic relations among the exponential and $\Ein$ coordinates
are the toric relations among the exponentials.  In particular the family
$\{\Ein(\lambda):\lambda\in\Qbar^\times\}$ is algebraically independent
over $K_{\exp}$, in the finite-subset sense.
\end{theorem}
\begin{proof}
Choose a $\Z$-basis $\nu_1,\ldots,\nu_d$ of the group generated by the
$\lambda_i$, and write $\lambda_i=m_i\cdot\nu$.  Distinct exponential
characters $e^{-u\cdot\nu z}$ are linearly independent over $\Qbar(z)$,
so $X_j=e^{-\nu_jz}$ realizes the field of Lemma~\ref{fd:lem:laurent}.
Put $Y_i(z)=\Ein(\lambda_i z)$.  Then
\[
 DY_i=(1-X^{m_i})/z.
\]
The field generated by these functions embeds in $\Qbar((z))$, so it has
no new differential constants.  The primitive Kolchin--Ostrowski theorem
\cite{fd:Kolchin} would turn any algebraic dependence of the $Y_i$ over
$K$ into an identity
\[
 Dg=\frac{\sum_i c_i-\sum_i c_iX^{m_i}}z
\]
with some $c_i\in\Qbar$ not all zero.  Lemma~\ref{fd:lem:laurent} excludes
this, since the $m_i$ are distinct and nonzero.

The functional transcendence degree is unchanged on extending constants
from $\Qbar$ to $\C$: any relation involves finitely many monomials
whose Taylor coefficients are algebraic, and scalar extension preserves
the rank of their coefficient matrix.  For specialization include $1$,
the required exponential characters, and
all $Y_i$ in a finite first-order system.  Its coefficients lie in
$\Qbar(z)$ and have no finite pole except possibly zero.  All entries
are $E$-functions.  The Siegel--Shidlovskii theorem, in the refined form
of Beukers \cite{pm:Beukers2006}, therefore transfers the functional
transcendence degree at $z=1$.  Algebraic independence of the $Y_i$
over the exponential torus makes its binomial character relations the
full relation ideal.  Finally, to enlarge the exponential field to
$K_{\exp}$, add to the finite set of arguments all algebraic exponents
occurring in a hypothetical relation's coefficients.  The finite result
again excludes that relation.
\end{proof}

\begin{corollary}[Full-rank linear images]\label{fd:cor:linear}
If $Y=(\Ein(\lambda_1),\ldots,\Ein(\lambda_s))^t$ and
$B\in M_{r,s}(\Qbar)$ has row rank $r$, the entries of $BY$ are
algebraically independent over $K_{\exp}$.  For arbitrary $B$, their
relation ideal is generated by the linear relations in the left kernel
of $B$.
\end{corollary}
\begin{proof}
Theorem~\ref{fd:thm:multipoint} identifies $Y$ with independent
indeterminates for the purpose of polynomial evaluation.  Invertible
linear changes of coordinates reduce $B$ to a coordinate projection.
\end{proof}

\begin{remark}[Higher orders retained as companion input]\label{fd:rem:polyexp}
For $r\ge1$ put
\[
 \Ein_r(z)=\sum_{n\ge1}\frac{(-1)^{n-1}z^n}{n^r n!},\qquad
 z\Ein_r'(z)=\Ein_{r-1}(z),\quad \Ein_0(z)=1-e^{-z}.
\]
Each fixed-order function is an $E$-function.  The all-order statement
in \cite[``Multipoint one-level polyexponentials of arbitrary weight'']{pm:KEIO001}
is
\[
 \operatorname{trdeg}_{\Qbar}
 \Qbar(e^{\lambda_i},\Ein_r(\lambda_i):1\le i\le s, 1\le r\le R)
 =d+Rs.
\]
Its iterated-primitive proof is retained in the companion source.  It is
not a new theorem of this revision, and none of the gamma arithmetic
conclusions below requires an unproved extension from order one to
unbounded order.  In particular, finite-point independence does not give
a uniform estimate as $R$, the arguments, or their heights increase.
\end{remark}

\section{Moving levels, zero traces, and exact arithmetic detectors}
\label{fd:sec:moving}

\begin{proposition}[Balanced dyadic levels]\label{fd:prop:dyadic}
For $N\ge1$ define
\[
 Y_N=(N+1)\Ein(2^N)-N\Ein(2^{N+1}).
\]
Then the $Y_N$ are algebraically independent over $K_{\exp}$ in the
finite-subset sense, and
\begin{equation}\label{fd:eq:epsilon}
 Y_N=\gamma+\varepsilon_N,\qquad
 \varepsilon_N=(N+1)E_1(2^N)-NE_1(2^{N+1})>0,
 \qquad
 \varepsilon_N\sim\frac{N+1}{2^N}e^{-2^N}.
\end{equation}
Moreover $Y_N$ strictly decreases to $\gamma$.
\end{proposition}
\begin{proof}
The rows $(N+1)e_N-Ne_{N+1}$ have full rank for every finite selection:
inspect the smallest occurring coordinate and proceed inductively.
Corollary~\ref{fd:cor:linear} gives independence.  The connection formula
cancels the logarithms exactly.  Positivity follows from the strict
decrease of $E_1$ on the positive axis.  Put $g(t)=E_1(2^t)$; then
$g''(t)=(\log2)^2 2^t e^{-2^t}>0$, and
\[
 \varepsilon_N-\varepsilon_{N+1}
 =(N+1)(g(N)-2g(N+1)+g(N+2))>0.
\]
Integration by parts gives $E_1(x)=e^{-x}x^{-1}(1+O(x^{-1}))$, proving
the asymptotic and the limit.
\end{proof}

\begin{proposition}[Critical slope is root multiplicity]\label{fd:prop:slope}
For every nonzero $P\in\C[X]$,
\begin{equation}\label{fd:eq:slope}
 \lim_{N\to\infty}\frac{-\log|P(Y_N)|}{2^N}
 =\operatorname{ord}_{X=\gamma}P.
\end{equation}
Consequently $\gamma$ is algebraic if and only if there exists
$0\ne P\in\Q[X]$ for which the limit equals $1$.
\end{proposition}
\begin{proof}
Write $P(X)=(X-\gamma)^mQ(X)$ with $Q(\gamma)\ne0$.  Equation
\eqref{fd:eq:epsilon} gives $-\log\varepsilon_N=2^N+O(N+\log(N+1))$;
the contribution of $\log|Q(Y_N)|$ stays bounded.  This proves the limit.
An irreducible polynomial over $\Q$ is separable, so the minimal
polynomial of an algebraic $\gamma$ has multiplicity one.
\end{proof}

This is an exact detector, not an estimate showing that its positive
alternative cannot occur.  It also explains why a lower bound uniform in
moving arguments cannot be inferred from a fixed-function, fixed-point
$E$-value theorem.

For $c\ne0$ let $Z_c$ be the zero multiset of $\Ein(z)-c$.  Define
\[
 P_m(X)=-m[z^m]\log(1-X\Ein(z)),\qquad m\ge2.
\]

\begin{proposition}[Trace recovery and finite-stage closure]\label{fd:prop:trace}
The inverse power sums converge absolutely and satisfy
\[
 T_m(c):=\sum_{\rho\in Z_c}\rho^{-m}=P_m(c^{-1}).
\]
The polynomials $P_m\in\Q[X]$ are monic of degree $m$.  In particular
\[
 P_2=X^2-X/2,\qquad
 P_4=X^4-X^3+25X^2/72-X/24,
\]
and
\begin{equation}\label{fd:eq:traceclosure}
 c^{-1}=144T_4(c)-144T_2(c)^2-14T_2(c),\qquad
 K[T_m(c):m\ge2]=K[c^{-1}]
\end{equation}
for every characteristic-zero subfield $K\subset\C$.
\end{proposition}
\begin{proof}
The growth estimate $|\Ein(z)|\le |z|e^{|z|}$ gives order at most one;
its negative-axis growth gives order one.  Genus-one Hadamard
factorization of $\Ein(z)-c$ gives
\[
 [z^m]\log(1-\Ein(z)/c)=-\frac1m\sum_{\rho\in Z_c}\rho^{-m}
 \quad(m\ge2).
\]
This proves the trace formula, including multiplicities.  Expanding the
formal logarithm shows monicity and the displayed polynomials.
The identity $144P_4-144P_2^2-14P_2=X$ proves the ring equality.
\end{proof}

One trace chosen at each independent $Y_N$ remains algebraically
independent, because a nonconstant polynomial image is an algebraic
extension in the reverse direction.  In contrast, all orders at one level
generate only $K(Y_N)$: increasing trace order at a fixed level adds no
new transcendence degree.  Their limits are $S_m=P_m(\gamma^{-1})$.
For $S_2=\gamma^{-2}-(2\gamma)^{-1}$, set
\[
 L_N(s)=2sY_N^2+Y_N-2.
\]
The identity
\[
 L_N(S_2)=(Y_N-\gamma)\{2S_2(Y_N+\gamma)+1\}
\]
shows slope one, whereas for fixed $s\ne S_2$, $L_N(s)$ tends to a
nonzero constant and has slope zero.  Proving that no algebraic $s$
can have the first behavior would prove transcendence, but is not
provided by the construction.

\subsection{A common-shift alternative with an independent anchor}
\begin{proposition}[Common-shift defect]\label{fd:prop:shift}
Let $z_0,z_1,\ldots,z_r$ be algebraically independent over a field
$K\subset\C$, let $c,e_0\in K^\times$, and let $g\in\C$.  Put
$T_j=c(z_j-g)$ and $\Delta=e_0(z_0-g)$.  Then
\[
 \operatorname{trdeg}_K K(T_1,\ldots,T_r)\ge r-1.
\]
If the $T_j$ are algebraically dependent, $g$ and $\Delta$ are
algebraically independent over $K$.
\end{proposition}
\begin{proof}
The $r-1$ differences $T_j-T_1=c(z_j-z_1)$ are independent.  Under
dependence put $L=K(T_1,\ldots,T_r)$, of transcendence degree $r-1$.
The equality
$L(g)=K(z_1,\ldots,z_r,g)$ forces its transcendence degree to equal $r$.
Thus $g$ is transcendental over $L$ and algebraic over
$K(z_1,\ldots,z_r)$.  Independence of the anchor $z_0$ makes it
transcendental over the algebraic closure of this last field, hence over
$K(g)$.  Therefore $g$ and $e_0(z_0-g)$ are independent.
\end{proof}

Apply this with $K=K_{\exp}$, $g=\gamma$, $z_0=\eta$, and $c=e_0=e$.
The remaining coordinates may be any finite selection of $Y_N$ and
$\Ein(-n)$, whose joint independence follows from disjoint supports and
Corollary~\ref{fd:cor:linear}.  Thus either all the corresponding residuals
are independent, or a finite relation proves algebraic independence of
$\gamma$ and $\delta$.  This alternative does not construct that relation;
expanding the family alone cannot rule out its independent branch.

\section{Reciprocal witnesses and the distinction between a level and its carrier}
\label{fd:sec:reciprocal}
Let $J(z)=z+z^{-1}$.  With $C_0(w)=2$, $C_1(w)=w$ and
$C_{n+1}=wC_n-C_{n-1}$, define
\[
 H(w)=\sum_{n\ge1}\frac{(-1)^{n-1}C_n(w)}{n\,n!}.
\]
Both roots of $z^2-wz+1$ remain bounded on compact $w$-sets, so this
series is locally uniformly convergent.  It gives an entire function with
\begin{equation}\label{fd:eq:reciprocal}
 H(J(z))=\Ein(z)+\Ein(z^{-1}),\qquad
 H'(w)=2e^{-w/2}\frac{\sinh(\sqrt{w^2-4}/2)}{\sqrt{w^2-4}}.
\end{equation}
The apparent singularities at $w=\pm2$ are removable.

\begin{proposition}[Arithmetic reciprocal witness]\label{fd:prop:reciprocal}
Distinct algebraic $w_1,\ldots,w_s$ give algebraically independent values
$H(w_i)$ over $K_{\exp}$.  Hence an algebraic solution of
$H(w)=2\gamma$ would make $\gamma$ transcendental.
The real equation has a unique solution $w_*\in(-2,2)$.  If $w_*$ is
algebraic, then $\gamma$ and $\delta$ are algebraically independent over
$K_{\exp}$.
\end{proposition}
\begin{proof}
The reciprocal pairs above distinct algebraic $w_i$ are disjoint sets of
nonzero algebraic numbers.  Their sums of $\Ein$ values are nonzero linear
forms on disjoint independent supports; repeated roots at $\pm2$ have
coefficient two.  Corollary~\ref{fd:cor:linear} applies.
The derivative in \eqref{fd:eq:reciprocal} is positive on $\R$: for
$|w|<2$ its trigonometric numerator has argument in $[0,1]$, and for
$|w|\ge2$ the hyperbolic expression is positive.  Its endpoint limits
show $H(\R)=\R$.  Since $\Ein(-1)<0<\gamma<\Ein(1)$, the Euler root
lies in $(-2,2)$.
If it is algebraic, its roots $a,a^{-1}$ are algebraic, distinct, and
different from $1$.  The three numbers
$\Ein(a),\Ein(a^{-1}),\eta$ are independent over $K_{\exp}$.
The invertible transformation
\[
 \gamma=\tfrac12(\Ein(a)+\Ein(a^{-1})),\quad
 \delta=e(\eta-\gamma),\quad
 \Ein(a)-\Ein(a^{-1})
\]
proves the last assertion.
\end{proof}

The identity $H(w_*)=2\gamma$ does not show that $w_*$ is algebraic.
Conversely, transcendence of $w_*$ alone does not decide the nature of
$\gamma$.  The rational differential equation for $H$ is unchanged by
adding a constant; a differential carrier which forgets its level does
not supply the missing arithmetic witness.

\section{Stokes data, Frobenius splitting, and the comparison gap}
\label{fd:sec:framing}
Set $F(x)=-\Ein(x)$ and
\[
 \mathcal L=(\partial+1)\circ\partial\circ x\circ\partial
 =x\partial^3+(x+2)\partial^2+\partial.
\]
The bases $(1,\log x,E_1(x))$ and $(1,F(x),\log x)$ are related by
\begin{equation}\label{fd:eq:Cgamma}
 C_\gamma=\begin{pmatrix}1&0&0\\-\gamma&-1&-1\\0&1&0\end{pmatrix}.
\end{equation}
Their Wronskian is nonzero, since
$W(1,\log x,E_1(x))=e^{-x}/x^2$.

\begin{proposition}[The lateral transition is gamma-free]\label{fd:prop:stokes}
Across the negative real axis the principal lateral values satisfy
\[
 \begin{pmatrix}1\\\log_+x\\E_{1,+}(x)\end{pmatrix}
 =\begin{pmatrix}1&0&0\\2\pi i&1&0\\-2\pi i&0&1\end{pmatrix}
 \begin{pmatrix}1\\\log_-x\\E_{1,-}(x)\end{pmatrix}.
\]
The row $(-g,-1,-1)$ is fixed by this transition for every $g\in\C$.
Adding an algebraic constant $t$ to $F$ preserves its $E$-function
status and changes the entry $-\gamma$ in \eqref{fd:eq:Cgamma} to
$t-\gamma$.  The unframed lateral transition, its trace, determinant,
and exterior characteristic data do not detect this entry.
\end{proposition}
\begin{proof}
The logarithm jumps by $2\pi i$, whereas \eqref{fd:eq:connection} makes
$E_1$ jump by $-2\pi i$.  The two contributions cancel in the displayed
row.  The change of $F$ adds $t$ times the constant basis vector, which
is precisely the stated shear of the connection matrix.
\end{proof}

Here a distinguished asymptotic coefficient such as $-\gamma$ may be
called a framed Stokes constant, as in the arithmetic theory of
$E$-operators \cite{pm:FischlerRivoal2016}.  It must be distinguished from a strict
Stokes factor or an unframed conjugacy invariant.  An arithmetic
containment theorem for possible coefficients is not a separation
theorem asserting that a particular coefficient is transcendental.

\begin{proposition}[Bare Frobenius splits, filtration can retain an extension]
\label{fd:prop:frobenius}
Let $K$ be complete non-archimedean, let $\sigma$ be an isometric
automorphism, and let $q\in K^\sigma$ satisfy $|q|\ne1$.  For every
$c\in K$ there is a unique $t$ solving $qt-\sigma(t)=c$, and
\[
 \begin{pmatrix}1&-t\\0&1\end{pmatrix}
 \begin{pmatrix}1&c\\0&q\end{pmatrix}
 \begin{pmatrix}1&\sigma(t)\\0&1\end{pmatrix}
 =\begin{pmatrix}1&0\\0&q\end{pmatrix}.
\]
If a filtration line is $K(e_1+h e_0)$, its displacement from the
Frobenius splitting is $\kappa=h-t$, and it is invariant under changing
$e_1$ to $e_1+u e_0$.  A splitting preserving this line exists exactly
when $\kappa=0$.
\end{proposition}
\begin{proof}
For $|q|<1$ use $t=-\sum_{n\ge0}q^n\sigma^{-n-1}(c)$; for $|q|>1$
use $t=\sum_{n\ge0}q^{-n-1}\sigma^n(c)$.  Both series converge and
satisfy the equation.  A difference of two solutions obeys
$qu=\sigma(u)$, hence is zero by the norm.  Matrix multiplication gives
the splitting.  Under the basis change, $c$ becomes
$c+\sigma(u)-qu$, while $h$ and $t$ both become their old values minus
$u$.  Thus their difference is unchanged.
\end{proof}

For example, over $\Q_p$, $\varphi=\operatorname{diag}(1,p)$ and
$\operatorname{Fil}^1=\Q_p(e_1+a e_0)$, with the usual two Hodge weights
$0,1$, are weakly admissible for every $a\in\Q_p$.  The total Hodge and
Newton numbers are one; on the two Frobenius eigenlines the Hodge number
is at most the Newton number.  The filtered extension splits exactly at
$a=0$, but its Frobenius traces, including tensor and symmetric-power
traces, do not depend on $a$.  The example has no assumed identification
of $a$ with any $p$-adic Euler constant.

The companion construction of the Euler exponential motive uses the
external input of Fres\'an--Jossen \cite[\S12.8]{fd:FJ}: a non-split
extension of $\Q(-1)$ by $\Q(0)$ has adapted comparison matrix
\[
 P_\gamma=\begin{pmatrix}1&\gamma\\0&2\pi i\end{pmatrix}.
\]
This motivic statement is cited, not reproved here.  With columns in the
de Rham basis and rows in the Betti basis,
\[
 \operatorname{comp}(xd_0+d_1)=(x+\gamma)b_0+2\pi i b_1.
\]
Consequently transverse lines defined over $\Qbar$ are matched exactly
when $x+\gamma=2\pi i u$ for some $x,u\in\Qbar$.

\begin{proposition}[Conditional matched-line implication]\label{fd:prop:matched}
If every such matched pair is induced by a subobject splitting the
Euler extension, then $1,\pi,\gamma$ are linearly independent over
$\Qbar$.  The version restricted to the line $\Qbar b_1$ already implies
transcendence of $\gamma$.
\end{proposition}
\begin{proof}
A relation $\gamma\in\Qbar+\pi\Qbar$ supplies the matched pair above.
Its hypothetical lift splits the non-split extension, a contradiction.
For the restricted version an algebraic $\gamma$ supplies $x=-\gamma$
and $u=0$.
\end{proof}

The lifting premise remains unproved.  Existence of realizations,
non-splitting of the motive, uniqueness of a bare Frobenius splitting,
and a numerical relation in the comparison matrix are different
statements.  Neither the matrix calculation nor known non-splitting
permits one to delete the algebraic-$\gamma$ alternative.

\section{The earlier interior-zero and Stieltjes attacks}
\label{fd:sec:early}

The first exploration \cite{fd:early} introduced signed-coefficient
kernels while retaining positivity of their values.  This escapes a
restriction on coefficientwise-positive perturbations, but not the
height cost of removing their logarithmic coordinates.
Put
\[
 W(u)=\frac1{\log u}+\frac1{1-u},\quad
 \omega(t)=\frac1{t\{\log^2(t^{-1}-1)+\pi^2\}},\quad
 D(u,t)=1-(1-u)t.
\]
The Markov representation used in the Pr\'evost family is
$W(u)=\int_0^1\omega(t)D(u,t)^{-1}dt$; see \cite{fd:Prevost} and
\cite[``Fixed-prime Rodrigues twists'']{pm:KEIO002}.

\begin{proposition}[Positive Rodrigues transducer]\label{fd:prop:rodrigues}
Let $G\in\Z[u]$ be nonnegative on $[0,1]$ with $G(1)>0$, let
$n=\ell+\deg G$ and $0\le m\le\ell$.  Put
\[
 P(u)=\frac{(-1)^\ell}{\ell!}\frac{d^\ell}{du^\ell}
       \{u^\ell(1-u)^\ell G(u)\}=\sum_{j=0}^n c_ju^j,
 \qquad J=\int_0^1u^{\ell-m}P(u)W(u)\,du.
\]
Then $P\in\Z[u]$, $P(1)=G(1)$, and
\begin{align}
 J&=G(1)\gamma-\sum_jc_jH_j+
                    \sum_jc_j\log(\ell-m+j+1),\label{fd:eq:rodlinear}\\
 (-1)^mJ&=\int_0^1\!\int_0^1
 \frac{u^\ell(1-u)^\ell G(u)(1-t)^{\ell-m}t^m}
 {D(u,t)^{\ell+1}}\omega(t)\,dt\,du>0.\label{fd:eq:rodpositive}
\end{align}
\end{proposition}
\begin{proof}
Dividing a derivative of a monomial by $\ell!$ gives an integer binomial
coefficient.  Integration by parts gives
$\int_0^1u^jP(u)du=0$ for $0\le j<\ell$.
The identity
\[
 \int_0^1u^rW(u)du=\gamma+\log(r+1)-H_r
\]
then yields \eqref{fd:eq:rodlinear}, since
$\sum_jc_j(H_{j+\ell-m}-H_j)=0$ by these moment equalities.
For the positive form differentiate $u^{\ell-m}/D(u,t)$ exactly $\ell$
times:
\[
 \frac{d^\ell}{du^\ell}\frac{u^{\ell-m}}{D(u,t)}
 =(-1)^m\ell!\frac{(1-t)^{\ell-m}t^m}{D(u,t)^{\ell+1}}.
\]
A further integration by parts proves \eqref{fd:eq:rodpositive}; the
integrand is nonzero on an open rectangle.
\end{proof}

For $n=400k$, $\ell=392k$, $m=114k$, and
$G(u)=(5u-2)^{4k}(2u-1)^{4k}$, the original note gives
\[
 J_k=3^{4k}\gamma-A_k+L_k>0,
 \qquad\lim_{k\to\infty}J_k^{1/(400k)}=\rho.
\]
The constant is specified algebraically by
\[
 \rho^{200}=\max_{0\le u\le1}
 \left(\frac{57}{196}\right)^{57}
 \left(\frac{139}{196}\right)^{139}
 u^{139}(1-u)^{196}(5u-2)^2(2u-1)^2.
\]
Indeed maximizing in $t$ first gives, for $\alpha=m/\ell$,
\[
 \max_t\frac{u(1-u)(1-t)^{1-\alpha}t^\alpha}{D(u,t)}
 =\alpha^\alpha(1-\alpha)^{1-\alpha}u^{1-\alpha}(1-u).
\]
Laplace upper and lower bounds on compact neighborhoods of maxima give
the root limit.  The three interior maxima are roots of
$3390u^3-4463u^2+1939u-278$ on the intervals separated by $2/5$ and
$1/2$.  The source's rational interval certificate proves
\[
 -\log\rho>\log(3+2\sqrt2),\qquad
 -\log\rho=1.7666899507\ldots.
\]
The exact interval certificate is archived; this revision does not label
its decimal evaluation as a proof.  The surviving $L_k$ in
\eqref{fd:eq:rodlinear} is a sum of prime logarithms.  The improvement is
therefore an analytic decay improvement before logarithm elimination.
It does not by itself furnish small integer forms in $1,\gamma$.

The companion Pr\'evost ledger has optimal raw rate
$\log(3+2\sqrt2)=1.7627471740\ldots$.  Since
$\log\operatorname{lcm}(1,\ldots,n)=n+o(n)$ and every fixed finite set
of primes contributes only $O(\log n)$ to this least common multiple,
its favorable height-one Ridout comparison misses the threshold by
$2-\log(3+2\sqrt2)=0.2372528259\ldots$.  This is a ledger for the
specified template with logarithms still present, not a universal
impossibility theorem for their linear combinations.  The first exact
seven-column logarithm cancellation in \cite{fd:early} produced
\[
 -8586246065985805791983028
 +15272776765356434827382880\,\gamma,
\]
whose coefficients and value are large.  Its exact cancellation is
preserved as a finite certificate, with no optimality or limiting claim.

\subsection{A positive Pad\'e error and its moving-scale denominator cost}
The Binet measure and its moments are
\begin{equation}\label{fd:eq:binetmeasure}
 d\mu(u)=\frac{du}{e^{2\pi\sqrt u}-1},\quad
 \mu_j=\frac{|B_{2j+2}|}{2j+2},\quad
 S(s)=\int_0^\infty\frac{d\mu(u)}{s+u}.
\end{equation}
Binet's formula \cite[\S5.9]{fd:DLMF} gives
$\gamma=H_n-\log n-(2n)^{-1}+S(n^2)$.
For the monic orthogonal polynomial $p_m$ and the Gaussian rational
approximant $R_m$, polynomial division and orthogonality give
\begin{equation}\label{fd:eq:padeerror}
 E_m(s):=S(s)-R_m(s)
 =\frac1{p_m(-s)^2}\int_0^\infty\frac{p_m(u)^2}{s+u}\,d\mu(u)>0.
\end{equation}
All zeros of $p_m$ are positive, so $E_m$ decreases on $s>0$.
For $n=2^N$ put
\[
 r_{m,N}=(N+1)\{H_n-(2n)^{-1}+R_m(n^2)\}
 -N\{H_{2n}-(4n)^{-1}+R_m(4n^2)\}.
\]
Then $r_{m,N}\in\Q$ and
\[
 \gamma-r_{m,N}=(N+1)E_m(n^2)-NE_m(4n^2)\ge E_m(n^2)>0.
\]

\begin{proposition}[Slow-degree Stieltjes barrier]\label{fd:prop:slow}
If $Q_{m,N}$ is the reduced denominator of $r_{m,N}$, there is an
absolute constant $C>0$ such that, for $m,N\ge1$ and $n=2^N$,
\[
 \log\{Q_{m,N}(\gamma-r_{m,N})\}
 \ge\vartheta(2n)-\vartheta(n)
 -C\{m^2\log(m+1)+(m+1)\log n\},
\]
where $\vartheta(x)=\sum_{p\le x}\log p$.  Hence
$Q_{m,N}(\gamma-r_{m,N})\to\infty$ whenever the expression in braces
is $o(n)$, in particular if $m=o(\sqrt{n/\log n})$.
\end{proposition}
\begin{proof}
Von Staudt--Clausen gives a common denominator $D_m$ of the first $2m+1$
Binet moments with $\log D_m=O(m)$.  Their integerized logarithmic
heights are $O(m\log(m+1))$.  Cramer's rule on the $m$ by $m$ moment
matrix gives integer polynomials $U_m,V_m$ such that
$R_m(s)=U_m(s)/(D_mV_m(s))$ and coefficient logarithmic heights
$O(m^2\log(m+1))$.  The same bound controls $p_m$.
For $n<p\le2n$, unless $p$ divides
$D_mV_m(n^2)V_m(4n^2)$, the harmonic term $-NH_{2n}$ contributes the
unique nonintegral $p$-part, $-N/p$; $p>N$ excludes cancellation.
The total logarithm of these exceptional divisors is
$O(m^2\log(m+1)+m\log n)$, proving the required lower bound on $\log Q$.
For the error, the density of $\mu$ is bounded below on $[1,2]$ and
\[
 \int_1^2p_m(u)^2du\ge
 \frac1{(2m+1)\binom{2m}{m}^2}\ge\frac1{(2m+1)16^m}.
\]
Insert this in \eqref{fd:eq:padeerror} and use the coefficient bound on
$p_m(-n^2)$.  Combining the two inequalities proves the result.
The asymptotic consequence follows from the prime number theorem.
\end{proof}

\section{Positive rational kernels: the corrected infinite-family target}
\label{fd:sec:positivehistory}
Let $h_m=H_{m-1}+(2m)^{-1}$ and $L_M=\operatorname{lcm}(1,\ldots,M)$.
For real coefficients $c_m$, $2\le m\le M$, set
\[
 F_c(u)=-\sum_{m=2}^M\frac{c_m}{u+m^2}.
\]
The exact constraints
\begin{equation}\label{fd:eq:positiveconstraints}
 \sum c_m=-1,\qquad \sum c_mv_p(m)=0\quad\hbox{for every prime }p
\end{equation}
yield
\begin{equation}\label{fd:eq:positiveidentity}
 \int_0^\infty F_c(u)d\mu(u)=\gamma-r_c,
 \qquad r_c=-\sum c_mh_m.
\end{equation}
If $F_c\ge0$ it is not identically zero, since $uF_c(u)\to1$;
its integral is therefore strictly positive.

\begin{proposition}[Target denominator, not residue denominator]
\label{fd:prop:residues}
Suppose \eqref{fd:eq:positiveconstraints} and $F_c\ge0$ hold, and impose
$r_c=P/(2L_M)$ with $P\in\Z$.  Then
\[
 2L_M\int F_c,d\mu=2L_M\gamma-P>0
\]
regardless of the denominators of the individual $c_m$.  If the last
quantity is less than $1$ for unbounded $M$, then $\gamma$ is irrational.
\end{proposition}
\begin{proof}
The displayed equality is exact.  If $\gamma=a/b$, then $b\mid L_M$
for every $M\ge b$, making its right side an integer strictly between
zero and one.  This is impossible.
\end{proof}

This is the correction introduced by \cite{fd:positivev2}: integer
residues are a sufficient design constraint, not a necessary one.  The
fixed-order Smith lattice of \cite{fd:multiscale} is therefore not an
obligatory gate.  Its homogeneous spacing calculations neither solve the
affine target condition nor supply an inner radius for the positive cone.
An upper bound on available coefficient width must not be used as a
lower bound ensuring successful rounding.

\begin{proposition}[Compact feasible cone and exact intersection condition]
\label{fd:prop:compact}
For $M\ge4$, the set $\mathcal C_M$ of real $c$ satisfying
\eqref{fd:eq:positiveconstraints} and $F_c\ge0$ is nonempty and compact.
Its minimum energy
\[
 E_M^*=\min_{c\in\mathcal C_M}\int F_c\,d\mu
\]
is attained and is positive.  Put $P_M=\lfloor2L_M\gamma\rfloor$.
For every sufficiently large $M$, a feasible real kernel with
$r_c=P_M/(2L_M)$ exists if and only if
\begin{equation}\label{fd:eq:intersection}
 E_M^*\le\frac{\{2L_M\gamma\}}{2L_M}.
\end{equation}
\end{proposition}
\begin{proof}
The kernel $2/(u+4)-1/(u+16)$ is feasible and has $r_c=13/24$.
The set is closed.  If it were unbounded, divide an unbounded sequence
by its norm and pass to a unit limit $d$.  Its kernel $F_d$ is
nonnegative, has $\sum d_m=0$, and is log-balanced.  Because its leading
$1/u$ term vanishes, it is integrable against Lebesgue measure, and
\[
 \int_0^\infty F_d(u)du=2\sum_m d_m\log m=0.
\]
Nonnegativity forces $F_d=0$, hence $d=0$ by uniqueness of partial
fractions, a contradiction.  The objective is continuous and positive
at every feasible point, so its attained minimum is positive.
Convexity makes the range of $r_c$ an interval whose upper endpoint is
$\gamma-E_M^*$.  Since $P_M/(2L_M)\to\gamma>13/24$, the target lies in
that interval exactly when \eqref{fd:eq:intersection} holds.
\end{proof}

The proof does not assert that a boundary real solution has rational
coefficients.  This distinction does not affect
Proposition~\ref{fd:prop:residues}, whose arithmetic target is the fixed
rational value $r_c$.  It does affect claims about rational certificates.

The exact finite construction of \cite{fd:positivev2} uses
$\mathcal S=\{2,3,4,6,8,9,12,16,18,24\}$ and
\[
 F(u)=\frac{u^6(au^3+bu^2+cu+d)}{a\prod_{m\in\mathcal S}(u+m^2)},
\]
where
\[
 \begin{split}
 a&=5967500403865,\qquad b=40760598624375361,\\
 c&=13281147167266501260,\qquad d=161940004068861643104.
 \end{split}
\]
All four integers are positive, and exact partial fractions give
\[
 \int F\,d\mu=\gamma-\frac{3090544783}{5354228880},
 \qquad 0<10708457760\gamma-6181089566<1.
\]
The final inequality is certified by independent rational Binet and
logarithm bounds in the archived verifier.  Later selected stages
$M=24,32,40,48,64,80,128$ have exact primal--dual certificates for
optimization in the smaller cone of numerators with nonnegative monomial
coefficients \cite{fd:uniformv3}.  These are finite optimizations, not
proofs for every intervening $M$ or for an unbounded family.

\subsection{Two explicit infinite families and their denominator limitations}
For integers $a\ge2$, $k\ge1$,
\[
 \Phi_{a,k}(u)=\frac{k+1}{u+a^{2k}}-\frac{k}{u+a^{2k+2}}>0.
\]
Its residues satisfy the required logarithmic balance, and
\[
 \int\Phi_{a,k}\,d\mu
 =\frac{k+1-k/a^2}{12a^{2k}}+O_a(ka^{-4k}).
\]
At the natural top node $M=a^{k+1}$, this is of order $k/M^2$, so
multiplication by $2L_M$ makes it diverge.  The multiscale butterfly
identities in \cite{fd:multiscale} construct this same two-pole family;
they do not change its error scale.

A genuinely many-pole infinite family from \cite{fd:uniformv3} is as
follows.  Fix $q\ge2$, put $a=q^2$, and for $s\ge2$ define
\[
 C_s=\sum_{k=1}^{s-1}\frac1{a^k-1},\quad
 t_s=(a^s-a)(s+C_s),\quad
 F_s(u)=\frac{u^{s-2}(u+t_s)}{\prod_{j=1}^s(u+a^j)}.
\]

\begin{proposition}[Geometric positive family and reduced-denominator barrier]
\label{fd:prop:geometric}
Writing $F_s(u)=\sum_{j=1}^sR_j/(u+a^j)$ gives
\[
 \sum R_j=1,\qquad \sum jR_j=0,
 \qquad \gamma-r_s=\int F_s\,d\mu>0,
 \quad r_s=\sum_{j=1}^sR_jh_{q^j}\in\Q.
\]
If $r_s=P_s/Q_s$ is reduced and $M=q^s$, then
\[
 \log(\gamma-r_s)=-(\log q)s^2+O_q(s\log(s+1)),\qquad
 \log Q_s\ge(1-1/q)M+o_q(M).
\]
In particular $Q_s(\gamma-r_s)\to\infty$.
\end{proposition}
\begin{proof}
If $A_j$ are the residues of $u^{s-1}/\prod(u+a^j)$, their polynomial is
\[
 W_s(z)=\sum A_jz^j=z\prod_{k=1}^{s-1}\frac{a^kz-1}{a^k-1}.
\]
It is characterized by $W_s(1)=1$ and the zeros $a^{-k}$.
Thus $W_s'(1)=s+C_s$ and
$a^{-1}W_s'(a^{-1})=1/(a^s-a)$.  Substitution of
$R_j=A_j(1-t_s/a^j)$ proves both sums.  Positivity is immediate.

A lower bound for the integral on $[1,2]$ is
$\exp\{-(\log q)s^2+O_q(s)\}$.
For an upper bound use $\prod(u+a^j)\ge a^{s(s+1)/2}$ and the moments
\eqref{fd:eq:binetmeasure}, whose logarithms are $O(s\log(s+1))$.
These bounds prove the stated leading error scale.
The residue formula shows that a common denominator of the $R_j$, and
the numerator of the nonzero $R_s$, have logarithmic size $O_q(s^2)$.
For clarity a power of $a$ times
$\prod_{k=1}^{s-1}(a^k-1)^3$ with exponent $O(s^2)$ on $a$ clears all
these residues.  Outside divisors of these integers, a prime
$M/q<p\le M$ does not occur in the earlier $h_{q^j}$, while
\[
 p h_M\equiv H_{\lfloor M/p\rfloor}\pmod p.
\]
Here the index is in $\{1,\ldots,q-1\}$; the finitely many possible
nonzero rational harmonic numbers are units at all sufficiently large
$p$.  Thus each such unexceptional prime divides the reduced $Q_s$.
The excluded primes have total logarithm $O_q(s^2)$, and the prime
number theorem gives the denominator lower bound.
\end{proof}

The geometric denominator barrier applies to this free rational target
$r_s$, not to all prescribed targets $P/(2L_M)$.  Given any feasible
rational kernel with $r_1>P/(2L_M)>13/24$, convex interpolation with
the basic two-pole kernel realizes the prescribed rational target with
positive rational coefficients.  What remains is precisely the
unbounded intersection in \eqref{fd:eq:intersection}; small unconstrained
energy does not bound the fractional part on its right away from zero.

\begin{remark}[Historical corrections retained]\label{fd:rem:history}
The earlier multiscale note also records local fixed-order Smith laws,
a two-by-two failure of total positivity for its atom kernel, and a
cyclotomic denominator obstruction for a fixed-base flattening ansatz.
Their detailed certificates and proof architecture remain in the archive;
the stronger global rounding conclusion suggested there is not assumed.
The later $v2$ relaxation supersedes the claim that integral residues are
necessary.  The $v3$ computations also refute the rule that maximal
moment cancellation is always optimal: at $M=80$ it is not optimal in
the specified coefficientwise cone, and at $M=100$ that particular
candidate leaves the cone.  Neither observation proves nonexistence of
other positive kernels at those levels.
\end{remark}

% END INTEGRATED PART I

\clearpage
\part{Positive moments and primitive logarithmic elimination}\label{master:moments}
% BEGIN INTEGRATED PART II
\section{Positive moments: framework and scope}\label{pm:sec:introduction}

Euler's constant
\[
 \gamma=\lim_{n\to\infty}\left(\sum_{k=1}^n\frac1k-\log n\right)
\]
appears both in convergent special-function identities and as a connection
constant between analytic solutions and asymptotic expansions.  The companion
manuscripts \cite{pm:KEIO001,pm:KEIO002} develop these two viewpoints through
exponential-integral periods, positive moment systems, and moving
polyexponential levels.  This part concentrates on the arithmetic
step which those viewpoints leave unresolved: producing useful integer
forms in $1$ and $\gamma$ after all auxiliary logarithms have been removed.

Our starting point is Sondow's positive double integral
\cite{pm:Sondow2003}.  We combine its exact logarithmic expansion with rational
linear programming, polynomial positivity on a compact interval, and the
geometry of the primitive target lattice.  These tools have distinct roles.
Prime valuations describe which logarithms can cancel; positivity certifies
that a resulting form does not vanish; rational optimization accounts for
the common factor removed when the target is made primitive.  Correlated
moment inequalities improve the finite optimization without introducing
new transcendental variables.

The principal conclusions are as follows.  All unrestricted assertions below
are proved in the indicated sections; numerical exponents refer to finite
certificates, not asymptotic extrapolations.
\begin{enumerate}[leftmargin=*,itemsep=3pt]
\item The prime support of Sondow's logarithmic integer is exactly the primes
at most $2n$.  This gives the exact rank of the prefix logarithm matrix,
explicit prime-adaptive multiplicative relations, and an obstruction to
every fixed-order recurrence with algebraic rational-function coefficients.
\item For a primitive target $(q,p)$, a rational quotient-norm problem is
exactly equivalent to integer logarithm elimination followed by removal of
the target common factor.  Positive rational changes of column normalization
do not change this primitive optimization problem.
\item Bernstein coefficients give sufficient certificates of nonvanishing
which are strictly less restrictive than first-coordinate dominance.
Two-sided rational bounds for the integrals lead to an exact interval model
for $\gamma$, with rational primal and dual certificates.
\item Positive differences of moving order have a fully explicit integral
representation and an exact exponential decay rate.  Their best rate after
the displayed integer normalization is
\[
 \log16-2-\log(1-e^{-2})=0.9180021801\ldots.
\]
These forms still contain prime logarithms.  Their bounds can nevertheless
be used as additional constraints in the logarithm-elimination problem.
\item The finite certificates include one primitive pair satisfying
$0<p-q\gamma<q^{-103/100}$.  Adding fifteen correlated moment constraints
contracts the computed intervals by factors of approximately $11.32$ and
$9.33$ at levels $32$ and $40$, respectively.
\item A separate rational moment construction admits a graded Todd
normalization, with an explicit improvement in its denominator upper bound.
The mixed exponential-integral viewpoint supplies exact reduction
statements, but its arithmetic specialization step remains open.
\end{enumerate}

\paragraph{Boundary of the positive-moment results.}
It proves neither the irrationality nor the transcendence of $\gamma$.
In particular, no infinite family with an unbounded primitive denominator
and a fixed positive Roth margin is established.  This distinction is part
of the mathematical statement, not a qualification to be removed by further
finite precision.  Section~\ref{pm:sec:infinite-target} formulates the exact
missing implications.  The work also makes no claim that the constructions
exhaust the approaches in the literature, or that the sufficient criteria
used here are necessary for every possible proof of transcendence.

\paragraph{Relation to earlier work.}
The integral and its logarithmic decomposition are due to Sondow; the
rational-function interpretation is developed by Sondow and Zlobin
\cite{pm:Sondow2003,pm:SondowZlobin2004}.  Rational approximants associated with
multiple orthogonal polynomials provide a different established approach;
see Wolfs and Van Assche \cite{pm:WolfsVanAssche2025}.  The present draft
organizes the subsequent arguments and certificates developed after the
companion manuscripts.  Those manuscripts are unpublished sources, not
external validation.  Results requiring only the classical integral and
elementary finite-dimensional arguments are proved here independently of
their broader algebraic-independence claims.  Historical priority for the
particular combinations of methods has not been established.

\paragraph{Notation and status.}
All logarithms are natural.  Put $H_0=0$ and
$H_n=\sum_{k=1}^n1/k$, and write
$d_j=\lcm(1,\ldots,j)$ for $j\ge1$.  The word \emph{primitive} means
$\gcd(q,p)=1$, with $q>0$ when $p/q$ is viewed as an approximant.
A \emph{finite certificate} consists of exact integer or rational data
and identities or inequalities which can be checked without evaluating
$\gamma$.  Rounded decimal summaries are not used in these checks.

\section{The integral, its integer coefficients, and nonvanishing}
\label{pm:sec:setup}

For $n\ge1$ define
\begin{equation}\label{pm:eq:sondow-integral}
 I_n=\int_0^1\!\int_0^1
 \frac{[x(1-x)y(1-y)]^n}{(1-xy)(-\log(xy))}\,dx\,dy.
\end{equation}
Let $t=x(1-x)y(1-y)$ and
$h(x,y)=((1-xy)(-\log(xy)))^{-1}$.  Then $0<t\le1/16$ in the
open square and $h>0$.  The singular weight $h\,dx\,dy$ alone has
infinite mass, so finite-measure arguments below always include a factor
$t$ before taking a pushforward.

The classical evaluation \cite{pm:Sondow2003,pm:SondowZlobin2004} is
\begin{equation}\label{pm:eq:sondow-decomp}
 I_n=C_n\gamma-A_n+L_n,\qquad
 C_n=\binom{2n}{n},\qquad
 A_n=\sum_{r=0}^n\binom nr^2H_{n+r},
\end{equation}
where
\begin{equation}\label{pm:eq:sondow-log-positive}
 L_n=2\sum_{k=1}^n\left(
 \sum_{r=0}^{\min(k-1,n-k)}
 \binom nr^2(H_{n-r}-H_r)\right)\log(n+k).
\end{equation}
Equivalently,
\begin{equation}\label{pm:eq:sondow-log-factorial}
 L_n=2\sum_{r=0}^n\binom nr^2(H_r-H_{n-r})\log((n+r)!).
\end{equation}
For completeness, the equivalence of the two logarithmic expressions is
purely finite.  Set $c_r=2\binom nr^2(H_r-H_{n-r})$.  Then
$c_{n-r}=-c_r$ and $\sum c_r=0$.  On expanding each factorial logarithm,
the terms $\log j$ for $j\le n$ disappear, while the coefficient of
$\log(n+k)$ is $\sum_{r=k}^n c_r$.  Pairing indices $r$ and $n-r$
reduces this last sum to the coefficient in
\eqref{pm:eq:sondow-log-positive}.
The integral evaluation itself is a cited classical input, not a new
identity claimed here.

Every denominator occurring in $A_n$ and in the coefficients of $L_n$
divides $d_{2n}$.  Consequently
\begin{equation}\label{pm:eq:integer-form}
 e_n:=d_{2n}I_n
 =a_n\gamma-b_n+\sum_{p\le2n}\beta_{p,n}\log p,
 \quad a_n=d_{2n}C_n,\quad b_n=d_{2n}A_n,
\end{equation}
with integer $a_n,b_n,\beta_{p,n}$; throughout a prime index is denoted
by $p$ only when no target numerator is in use.  The integer
\begin{equation}\label{pm:eq:sondow-integer}
 S_n=\exp(d_{2n}L_n)=\prod_{p\le2n}p^{\beta_{p,n}}
\end{equation}
is well-defined and positive.

For a finite set of indices $\mathcal J$, let
$B=(\beta_{p,n})_{p\le2\max\mathcal J,\ n\in\mathcal J}$, where the row index runs over all
primes at most $2\max\mathcal J$.  Write $a=(a_n)_{n\in\mathcal J}$,
$b=(b_n)_{n\in\mathcal J}$ and
\[
 \Phi(w)=(a\cdot w,b\cdot w).
\]
An integer vector $u$ in $\ker B$ gives the exact identity
\begin{equation}\label{pm:eq:eliminated-form}
 \sum_{n\in\mathcal J}u_ne_n=Q\gamma-P,
 \qquad (Q,P)=\Phi(u),\qquad
 \prod_{n\in\mathcal J}S_n^{u_n}=1.
\end{equation}
Unique prime factorization proves the equivalence between the multiplicative
identity and $Bu=0$.

\begin{lemma}[Two-sided rational integral bounds]\label{pm:lem:integral-bounds}
For every $n\ge1$,
\begin{equation}\label{pm:eq:two-sided-bounds}
 \ell_n:=\frac{d_{2n}}{4(2n+1)^2C_n^2}
 <e_n<
 h_n:=\frac{d_{2n}}{n^2C_n^2}.
\end{equation}
In particular, all integrals in \eqref{pm:eq:sondow-integral} are finite.
\end{lemma}
\begin{proof}
For $0<u<1$, one has $-\log u<1/u$ and $-\log u\ge1-u$.
It follows that $h(x,y)>xy$ and
\[
 h(x,y)\le\frac1{(1-xy)^2}
 \le\frac1{(1-x)(1-y)}.
\]
The comparisons are strict on a subset of positive measure.  Thus
\[
 B(n+2,n+1)^2<I_n<B(n+1,n)^2,
\]
where $B(\ ,\ )$ denotes Euler's beta function only in this display.
The two beta values equal $[2(2n+1)C_n]^{-1}$ and $(nC_n)^{-1}$,
respectively.  Multiplication by $d_{2n}$ proves the assertion.
\end{proof}

\begin{example}[The first two indices]
There is no singular starting case for the Sondow family:
\[
 I_1=2\gamma-\frac52+2\log2,\qquad
 I_2=6\gamma-\frac{131}{12}+3\log12.
\]
In particular $S_1=16$.  The bounds and integer normalization above
already hold at $n=1$.  The order $N$ of the separate moment matrix in
Section~\ref{pm:ht:section} is a different parameter.
\end{example}

Three logically distinct conditions must be maintained:
\[
 u\ne0,\qquad \Phi(u)\ne(0,0),\qquad Q\gamma-P\ne0.
\]
The first does not imply the second.  The second does not imply the third
without an additional argument.  Positive integral representations and
strict rational interval certificates will supply that argument in the
finite examples below.


% ---- structure.tex ----
\section{Prime support and canonical logarithm elimination}
\label{pm:str:sec-primes}

The logarithms in Sondow's evaluation have a particularly rigid support.
This section determines that support and the rank of every prefix of the
prime-exponent matrix.  Prefixes and moving windows have different rank
questions; we distinguish them throughout.

For $n\geq1$, define
\begin{equation}\label{pm:str:def-c}
 c_{n,k}=2\sum_{i=0}^{\min(k-1,n-k)}
       \binom ni^2(H_{n-i}-H_i)\qquad(1\leq k\leq n).
\end{equation}
The evaluation of the integral in \cite{pm:Sondow2003,pm:SondowZlobin2004}
gives $L_n=\sum_{k=1}^n c_{n,k}\log(n+k)$.  Put
\begin{equation}\label{pm:str:def-S}
 S_n=\prod_{k=1}^n(n+k)^{d_{2n}c_{n,k}},
 \qquad \beta_{p,n}=v_p(S_n).
\end{equation}
In particular, $d_{2n}L_n=\log S_n=\sum_p\beta_{p,n}\log p$.

\begin{lemma}\label{pm:str:positive-exponents}
For $n\geq1$, each $c_{n,k}$ is strictly positive and
$d_{2n}c_{n,k}$ is a positive integer.
Consequently $S_n$ is a positive integer greater than one.
\end{lemma}

\begin{proof}
In \eqref{pm:str:def-c}, the summation limit is at most
$\lfloor(n-1)/2\rfloor$.  Thus every harmonic difference in that
sum is positive.  The denominators of all those harmonic numbers divide
$d_{2n}$, proving the integrality assertion.
\end{proof}

\begin{theorem}[Exact support and prefix rank]\label{pm:str:prime-rank}
For every prime $p$ and integer $n\geq1$,
\begin{equation}\label{pm:str:prime-support}
 \beta_{p,n}>0\quad\Longleftrightarrow\quad p\leq2n.
\end{equation}
Let $\pi(x)$ denote the number of primes at most $x$, and define
\[
 n(2)=1,\qquad n(p)=\frac{p+1}{2}\quad(p\text{ odd prime}).
\]
For every $M\geq1$, the prefix matrix
\[
 B_{[1,M]}=(\beta_{p,n})_{p\leq2M,\ 1\leq n\leq M}
\]
has rank $\pi(2M)$.  Its columns with indices $n(p)$ form a basis of
its column space over $\mathbb Q$, and
\begin{equation}\label{pm:str:prefix-nullity}
 \dim_{\mathbb Q}\ker B_{[1,M]}=M-\pi(2M).
\end{equation}
\end{theorem}

\begin{proof}
All bases in the product \eqref{pm:str:def-S} lie in $(n,2n]$, so no
prime greater than $2n$ occurs.  If $n<p\leq2n$, the base $p$ itself
occurs.  If $p\leq n$, then the interval of $n$ consecutive integers
$n+1,\ldots,2n$ contains a multiple of $p$.  Strict positivity of every
exponent proves \eqref{pm:str:prime-support} in both cases.

Order the primes at most $2M$ increasingly.  Their indices $n(p)$
are distinct and at most $M$.  In the square submatrix formed by these
columns, the entry in a row below the diagonal is zero: a larger prime
has not yet appeared at that column.  Every diagonal entry is positive
by \eqref{pm:str:prime-support}.  The determinant is the product of the
diagonal entries and is nonzero.  The number of rows gives the reverse
rank inequality, and rank--nullity proves \eqref{pm:str:prefix-nullity}.
\end{proof}

\begin{corollary}[Canonical prime-adaptive relations]\label{pm:str:canonical}
If $n$ is not one of the indices $n(p)$, there is a unique rational
expansion
\begin{equation}\label{pm:str:canonical-rational}
 (\beta_{p,n})_{p\leq2n}
 =\sum_{\substack{q\leq2n\\q\text{ prime}}}\rho_{q,n}
       (\beta_{p,n(q)})_{p\leq2n}.
\end{equation}
If $D$ is any common positive denominator of the $\rho_{q,n}$, then
\begin{equation}\label{pm:str:canonical-product}
 S_n^D\prod_{\substack{q\leq2n\\q\text{ prime}}}
                  S_{n(q)}^{-D\rho_{q,n}}=1.
\end{equation}
There are infinitely many such relations, constructed without the value
of $\gamma$.
\end{corollary}

\begin{proof}
The square triangular submatrix in Theorem~\ref{pm:str:prime-rank} is
invertible, so back substitution gives the unique coefficients in
\eqref{pm:str:canonical-rational}.  Clearing denominators makes the
valuation at every prime of the left side of
\eqref{pm:str:canonical-product} zero.  Unique factorization proves that
identity.  For example, every $n>2$ with $n\equiv2\pmod3$ is not an
entry index, since $2n-1>3$ is divisible by $3$.  There are infinitely
many such $n$.
\end{proof}

The support of these relations can extend back to level $1$.
The construction supplies neither a short coefficient vector in a
moving window nor a small primitive linear form in $1,\gamma$.
Thus the existence of exact multiplicative identities and their
usefulness for Diophantine approximation are separate questions.

For completeness, the two target coordinates already survive
independently in a small prefix.  This finite fact has an all-prefix
consequence, but no corresponding all-window consequence.

\begin{proposition}[Rank of the prefix target]\label{pm:str:prefix-target}
For every $M\geq8$, the map
\[
 \Phi_M:\ker_{\mathbb Q}B_{[1,M]}\longrightarrow\mathbb Q^2,
 \qquad u\longmapsto\left(\sum_{n=1}^M a_nu_n,
                            \sum_{n=1}^M b_nu_n\right)
\]
has rank two.  Its kernel has dimension $M-\pi(2M)-2$.
\end{proposition}

\begin{proof}
The following two integer vectors, displayed by coordinates, give an
explicit certificate at $M=8$:
\begin{center}
\begin{tabular}{r|r|r}
$n$ & $u_n$ & $v_n$\\\hline
1 & $-3923225775$ & $353653621155957622500$\\
2 & $-158150125$ & $-24892511625147813075$\\
3 & $392805$ & $-633334105016310555$\\
4 & $-200868$ & $-518196421884141$\\
5 & $35150$ & $0$\\
6 & $0$ & $57656027697450$\\
7 & $0$ & $-309266976225$\\
8 & $0$ & $5037311350$
\end{tabular}
\end{center}
Substitution in \eqref{pm:str:def-c}--\eqref{pm:str:def-S}, with prime rows
$2,3,5,7,11,13$, gives $B_{[1,8]}u=B_{[1,8]}v=0$ and
\begin{align*}
 \Phi_8(u)&=(-16097728500,-8969752951),\\
 \Phi_8(v)&=(-27120467275073944200,-44055302128555509412).
\end{align*}
The determinant of these two target vectors is the nonzero integer
\[
 465926401277865247962411307800.
\]
These equalities involve only integer and rational arithmetic and are
also checked by the accompanying verifier.  Appending zero coordinates
preserves the two vectors and their target images in every larger
prefix.  Rank--nullity and \eqref{pm:str:prefix-nullity} finish the proof.
\end{proof}

\begin{remark}\label{pm:str:prefix-scope}
Theorems about $B_{[1,M]}$ cannot be applied to $B_{[N,M]}$ merely by
replacing $1$ by $N$.  Removing the early entry columns removes the
triangular certificate.  In particular, no full prime-row rank or
rank-two target assertion for every window $[N,2N]$ is claimed here.
\end{remark}

\section{A fixed-order recurrence obstruction}
\label{pm:str:sec-recurrence}

We use P-recursiveness in the following sense: a complex sequence $f_n$
is P-recursive over a field $K\subseteq\mathbb C$ if it satisfies a
nonzero recurrence $\sum_{j=0}^r p_j(n)f_{n+j}=0$, with
$p_j\in K[n]$, for all sufficiently large integers $n$.  Its initial
values need not belong to $K$.

\begin{theorem}\label{pm:str:log-nonprec}
Neither $(L_n)_{n\geq1}$ nor $(\log S_n)_{n\geq1}$ is P-recursive
over $\overline{\mathbb Q}$.  Equivalently, for either sequence $f_n$,
an identity
\begin{equation}\label{pm:str:forbidden-rec}
 \sum_{j=0}^r c_j(n)f_{n+j}=0,
 \qquad c_j\in\overline{\mathbb Q}(n),
\end{equation}
valid for every sufficiently large integer $n$, forces all $c_j$ to be
identically zero.  Their ordinary generating series are not D-finite
over $\overline{\mathbb Q}(z)$.
\end{theorem}

\begin{proof}
First take coefficients in $\mathbb Q(n)$.  Distinct prime logarithms
are linearly independent over $\mathbb Q$: clear denominators in a
putative relation, exponentiate, and apply unique factorization.
Suppose \eqref{pm:str:forbidden-rec} holds, and discard its trailing zero
coefficients so that $c_r\ne0$.  For each sufficiently large odd prime
$p$, put $n=(p+1)/2-r$.  The final term has index $n+r=(p+1)/2$,
so its logarithmic expansion contains $\log p$ with a positive
coefficient.  Every preceding term has index at most $(p-1)/2$ and
therefore contains no $\log p$.  Linear independence forces $c_r(n)=0$.
This happens at infinitely many distinct integers, contradicting that
$c_r$ is a nonzero rational function.  The argument applies to both
sequences: division by $d_{2n}$ changes coefficients but not support.

For algebraic coefficients, use Baker's linear-independence theorem
\cite{pm:Baker1966}: logarithms of algebraic numbers that are linearly
independent over $\mathbb Q$ remain linearly independent over
$\overline{\mathbb Q}$.  At every integer outside the finitely many
poles of the $c_j$, their values are algebraic.  The same isolated
coefficient of $\log p$ again forces $c_r(n)=0$ infinitely often.

Finally, comparison of coefficients in a nonzero linear differential
equation with polynomial coefficients for the ordinary generating
series produces a nonzero polynomial-coefficient recurrence on a tail.
Such a recurrence has just been excluded.
\end{proof}

The isolation argument also constrains any individual relation.

\begin{corollary}\label{pm:str:last-gap}
Suppose $m_1<\cdots<m_s$, $u_j\in\mathbb Q\setminus\{0\}$, and
$\prod_{j=1}^s S_{m_j}^{u_j}=1$.  Then $s\geq2$, and the interval
$(2m_{s-1},2m_s]$ contains no prime.
\end{corollary}

\begin{proof}
One factor cannot give a relation because $S_m>1$.  A prime in the
stated interval would have positive valuation in the last factor and
zero valuation in all other factors.  After clearing denominators of
the $u_j$, its valuation in the product would be nonzero.
\end{proof}

The non-P-recursiveness is also present before the lcm normalization
and before the logarithmic part is separated from the integral.

\begin{lemma}\label{pm:str:A-holonomic}
The rational sequence $A_n=\sum_{k=0}^n\binom nk^2H_{n+k}$ is
P-recursive over $\mathbb Q$.
\end{lemma}

\begin{proof}
Introduce the parameter polynomial
\begin{equation}\label{pm:str:F-parameter}
 F_n(t)=\sum_{k=0}^n\binom nk^2\frac{(1+t)_{n+k}}{(n+k)!},
 \qquad (a)_m=\prod_{j=0}^{m-1}(a+j).
\end{equation}
Vandermonde's identity and differentiation of the finite product give
\begin{equation}\label{pm:str:F-specialize}
 F_n(0)=C_n,\qquad F_n'(0)=A_n.
\end{equation}
For generic $t$, the summand is the proper hypergeometric term
\[
 T(n,k;t)=\frac{\Gamma(n+1)^2\Gamma(n+k+1+t)}
 {\Gamma(k+1)^2\Gamma(n-k+1)^2
  \Gamma(n+k+1)\Gamma(1+t)}.
\]
Its two discrete shift ratios are
\begin{align*}
 \frac{T(n+1,k;t)}{T(n,k;t)}
 &=\frac{(n+1)^2(n+k+1+t)}{(n+1-k)^2(n+k+1)},\\
 \frac{T(n,k+1;t)}{T(n,k;t)}
 &=\frac{(n-k)^2(n+k+1+t)}{(k+1)^2(n+k+1)}.
\end{align*}
The reciprocal Gamma factors give its terminating support $0\leq k\leq n$;
outside this interval the term vanishes for generic $t$.  The
proper-hypergeometric summation theorem, over the parameter field
$\mathbb Q(t)$, therefore supplies a nonzero recurrence for the finite
sum \cite[Sections 2.1 and 3.3]{pm:WZ1992}.  Clearing coefficient
denominators, write that recurrence as
\begin{equation}\label{pm:str:parameter-operator}
 P(t)F(t)=0,\qquad
 P(t)=\sum_{j=0}^r p_j(n,t)E^j,
 \qquad p_j\in\mathbb Q[n,t],
\end{equation}
where $E$ shifts $n$ to $n+1$.

Remove the largest common power of $t$ from the coefficients of $P$.
Then $P(0)$ is a nonzero recurrence operator.  Since every $F_n(t)$
is a polynomial, \eqref{pm:str:parameter-operator} extends to $t=0$ and
may be differentiated there.  By \eqref{pm:str:F-specialize},
\begin{equation}\label{pm:str:A-inhomogeneous}
 P(0)A=-P'(0)C.
\end{equation}
Each ratio $C_{n+j}/C_n$ is rational in $n$, so the right side is
$C_nR(n)$ for some $R\in\mathbb Q(n)$.  If $R=0$, $P(0)$ already
annihilates $A$.  Otherwise $C_nR(n)$ is annihilated, on a nonsingular
tail, by
\[
 Q=E-\frac{2(2n+1)R(n+1)}{(n+1)R(n)}.
\]
The product $QP(0)$ is nonzero: in the ring of recurrence operators
over $\mathbb Q(n)$, the leading coefficient of a product is the
product of a nonzero coefficient and a shift of a nonzero coefficient.
This product annihilates $A$.  Clearing its rational denominators
completes the proof.
\end{proof}

\begin{theorem}\label{pm:str:I-nonprec}
The unscaled Sondow sequence $(I_n)_{n\geq1}$ is not P-recursive
over $\overline{\mathbb Q}$, and its ordinary generating series is
not D-finite over $\overline{\mathbb Q}(z)$.
\end{theorem}

\begin{proof}
The sequence $\gamma C_n$ obeys
\[
 (n+1)\gamma C_{n+1}-2(2n+1)\gamma C_n=0
\]
regardless of the arithmetic nature of $\gamma$.  If $I_n$ were
P-recursive over $\overline{\mathbb Q}$, closure under addition and
Lemma~\ref{pm:str:A-holonomic} would make
$L_n=I_n-\gamma C_n+A_n$ P-recursive over that field.  This contradicts
Theorem~\ref{pm:str:log-nonprec}.  The closure property does not require
algebraic initial values: the shifts of each sequence span a
finite-dimensional space over $\overline{\mathbb Q}(n)$, and shifts
of their sum lie in the sum of those spaces.  The generating-series
assertion follows by coefficient comparison.
\end{proof}

\begin{remark}\label{pm:str:recurrence-scope}
Theorems~\ref{pm:str:log-nonprec} and \ref{pm:str:I-nonprec} exclude a
homogeneous recurrence of fixed order, with coefficients in
$\overline{\mathbb Q}(n)$, holding at every sufficiently large
integer.  They do not exclude growing-order transformations,
prime-adaptive coefficients, exceptional subsequences, inhomogeneous
relations, or other integral families.  In particular, they are
consistent with Corollary~\ref{pm:str:canonical} and do not settle the
arithmetic nature of $\gamma$.
\end{remark}

\section{Normalization, short relations, and survival of target directions}
\label{pm:str:sec-escape}

For a window $N\leq n\leq M$, write $B=B_{[N,M]}$ and
\[
 \Phi(u)=(Q,P)=(a\mathbin{\cdot}u,b\mathbin{\cdot}u).
\]
When $Bu=0$, the logarithms cancel and
\begin{equation}\label{pm:str:target-value}
 \sum_{n=N}^M u_ne_n=Q\gamma-P.
\end{equation}
The assertions $u\ne0$, $\Phi(u)\ne0$, and $Q\gamma-P\ne0$ are
different.  The first two are statements about rational linear maps;
the last is a statement about a value at $\gamma$.

\begin{proposition}[Dependence of crude width bounds on normalization]
\label{pm:str:normalization}
Fix $w>0$, put $M=\lfloor(1+w)N\rfloor$, and let
$c=\log16-2>0$.  As $N\to\infty$,
\begin{align}
 \max_{N\leq n\leq M}d_{2n}I_n
 &\leq\exp(-cN+o(N)),\label{pm:str:individual-normalization}\\
 \max_{N\leq n\leq M}d_{2M}I_n
 &\leq\exp(-(c-2w)N+o(N)).\label{pm:str:common-normalization}
\end{align}
Thus the restriction $w<c/2$ is sufficient for exponential decay in
the second crude bound, but is not a restriction in the first.
\end{proposition}

\begin{proof}
Write $T=x(1-x)y(1-y)$ and
$h=((1-xy)(-\log(xy)))^{-1}$ for the integral kernel.  For
$0<x,y<1$, the inequality $-\log(xy)>1-xy$ gives
\[
 Th<\frac{xy(1-x)(1-y)}{(1-xy)^2}\leq xy.
\]
The last inequality follows from $1-xy\geq1-x$ and
$1-xy\geq1-y$.  Therefore $I_1<1/4$.  Since $T\leq1/16$,
\begin{equation}\label{pm:str:crude-I}
 0<I_n\leq16^{-(n-1)}I_1<4\,16^{-n}\qquad(n\geq1).
\end{equation}
The prime number theorem implies $\log d_m=m+o(m)$.
On the fixed-width window this yields $\log d_{2n}=2n+o(N)$
uniformly.  Combining with \eqref{pm:str:crude-I} proves
\eqref{pm:str:individual-normalization}.  Using $d_{2M}$ for every
column instead gives $2M-N\log16=-(c-2w)N+O(1)$, proving
\eqref{pm:str:common-normalization}.
\end{proof}

The two presentations are related by positive rational column
rescaling.  To see exactly what is preserved, let $C=[B;a;b]$,
let $R$ be a positive rational diagonal matrix, and let
$\epsilon_i>0$ be error bounds for the columns of $C$.
For a fixed target $v$, the substitution $x=Rw$ gives a bijection
\[
 \{w:(CR)w=(0,v)\}\longleftrightarrow
 \{x:Cx=(0,v)\},
 \qquad
 \sum_i(\epsilon_iR_{ii})|w_i|=\sum_i\epsilon_i|x_i|.
\]
Consequently the rational weighted optimization for that target is
unchanged.  An unweighted bound on integer coefficients, or on their
relative sizes, is not preserved by this substitution.  The width
penalty in \eqref{pm:str:common-normalization} must therefore not be
combined with a coefficient bound from a different presentation.

\begin{lemma}[An elementary kernel bound]\label{pm:str:pigeonhole}
Let an integer matrix have $R$ rows and $s$ columns, with entries of
absolute value at most $H\geq1$, and suppose $s>R$.  It has a nonzero
integer kernel vector $u$ with
\begin{equation}\label{pm:str:kernel-K}
 \|u\|_\infty\leq
 K:=\left\lfloor(2sH+1)^{R/(s-R)}\right\rfloor.
\end{equation}
No row-rank hypothesis is needed.
\end{lemma}

\begin{proof}
The $(K+1)^s$ points of $\{0,\ldots,K\}^s$ have at most
$(2sHK+1)^R$ images.  Since
\[
 (K+1)^{s-R}>(2sH+1)^R,
 \qquad 2sHK+1\leq(2sH+1)(K+1),
\]
two distinct grid points have the same image.  Their difference is
the required vector.
\end{proof}

\begin{theorem}[Small relations can have zero target]\label{pm:str:small-zero}
Fix $w>0$ and $M=\lfloor(1+w)N\rfloor$.  For every sufficiently
large $N$ there is a nonzero integer vector $u$, supported in $[N,M]$,
such that
\begin{equation}\label{pm:str:zero-equations}
 Bu=0,\qquad \Phi(u)=0,
\end{equation}
and
\begin{align}
 \log\|u\|_\infty
 &\leq\left(\frac{2(2+\log4)(1+w)^2}{w}+o(1)\right)
                 \frac{N}{\log N},\label{pm:str:short-zero-height}\\
 \sum_{n=N}^M|u_n|e_n
 &\leq\exp(-cN+o(N)).\label{pm:str:short-zero-error}
\end{align}
The signed sum in \eqref{pm:str:target-value} is exactly zero for these
vectors.  The same size estimates, without the requirement
$\Phi(u)=0$, also follow directly from the kernel bound.
\end{theorem}

\begin{proof}
The binomial identity $\sum_i\binom ni^2=C_n\leq4^n$ gives
$c_{n,k}\leq2H_n4^n$.  Therefore
\[
 0\leq\beta_{p,n}\leq
 2n\,d_{2n}H_n4^n\log_2(2n)
 \leq4n^3d_{2n}4^n.
\]
Here $H_n\leq n$ and $\log_2(2n)\leq2n$ for $n\geq1$.
The last bound also dominates $a_n$ and $b_n$, since
$A_n\leq H_{2n}C_n\leq2n4^n$.  Thus the integer matrix $[B;a;b]$
has at most $R=\pi(2M)+2$ rows, $s=M-N+1$ columns, and entries
bounded by
\[
 H=4M^3d_{2M}4^M.
\]
As $N\to\infty$, the prime number theorem gives
\[
 R=\left(2(1+w)+o(1)\right)\frac{N}{\log N},\quad
 s=wN+O(1),\quad
 \log H=(2+\log4)(1+w)N+o(N).
\]
In particular $s>R$ for large $N$.  Apply
Lemma~\ref{pm:str:pigeonhole} to obtain \eqref{pm:str:zero-equations} and
\eqref{pm:str:short-zero-height}.  The number of terms and their
coefficient bound contribute only $\exp(o(N))$, so
Proposition~\ref{pm:str:normalization} proves
\eqref{pm:str:short-zero-error}.  Omitting the last two rows makes the
same argument apply to $B$ alone.
\end{proof}

The coefficient $(1+w)^2/w=w+2+1/w$ in this particular estimate
is minimized at $w=1$.  No finite onset, optimal window for primitive
approximation, or necessary coefficient lower bound follows from
this asymptotic observation.

\begin{lemma}[Span of short relations]\label{pm:str:short-span}
For an $R$-by-$s$ integer matrix $A$ with entries bounded by $H$, let
\[
 V_K=\operatorname{span}_{\mathbb Q}
       \{u\in\mathbb Z^s:Au=0,\ \|u\|_\infty\leq K\}.
\]
For every integer $K\geq1$,
\begin{equation}\label{pm:str:span-bound}
 \dim V_K\geq
 \frac{s\log(K+1)-R\log(2sHK+1)}{\log(2K+1)}.
\end{equation}
\end{lemma}

\begin{proof}
A fiber of $A$ on $\{0,\ldots,K\}^s$ contains at least
$(K+1)^s/(2sHK+1)^R$ points.  Subtract one point of that fiber.
The resulting distinct vectors lie in $V_K\cap[-K,K]^s$.
If $d=\dim V_K$, some projection onto $d$ coordinate positions
is injective on $V_K$: choose $d$ independent rows of a matrix whose
columns are a basis of $V_K$.  There are at most $(2K+1)^d$ projected
integer vectors.  Taking logarithms gives \eqref{pm:str:span-bound}.
\end{proof}

\begin{corollary}\label{pm:str:many-zero}
With the window and notation of Theorem~\ref{pm:str:small-zero}, there
are at least $s-O(N/\sqrt{\log N})$ linearly independent integer
vectors satisfying \eqref{pm:str:zero-equations}, each with
\[
 \|u\|_\infty\leq
   \left\lceil\exp\left(\frac{N}{\sqrt{\log N}}\right)\right\rceil,
 \qquad
 \sum_{n=N}^M|u_n|e_n\leq\exp(-cN+o(N)).
\]
\end{corollary}

\begin{proof}
Apply Lemma~\ref{pm:str:short-span} to $[B;a;b]$ with the displayed
$K$.  In its right side, $R=O(N/\log N)$, $\log H=O(N)$, and
$\log K\sim N/\sqrt{\log N}$.  Hence its deficit from $s$ is
$O(N/\sqrt{\log N})$.  Choose an independent subset from the
vectors spanning $V_K$.  Proposition~\ref{pm:str:normalization}
provides their stated absolute-sum bounds.
\end{proof}

Even a subspace of nearly full dimension generated by short relations
can therefore consist entirely of zero-target relations.  The next
bound addresses the target map explicitly.

\begin{theorem}[A quantitative target-survival bound]\label{pm:str:escape}
Let $B\in\mathbb Z^{r\times s}$ have rank $r$ and entries of
absolute value at most $H\geq1$.  Suppose a rational-linear map
$\Phi:\ker_{\mathbb Q}B\to\mathbb Q^t$ has rank $t$.
Let $1\leq j\leq t$ and $K\geq1$ be integers.  If
\begin{equation}\label{pm:str:escape-condition}
 (K+1)^{t-j+1}>(2sH+1)^r,
\end{equation}
then the images under $\Phi$ of the integer vectors
$u\in\ker B$ with $\|u\|_\infty\leq K$ span dimension at least $j$.
\end{theorem}

\begin{proof}
Suppose those images lie in a subspace $T\subseteq\mathbb Q^t$
of dimension at most $j-1$.  Two points in the same $B$-fiber of
$\{0,\ldots,K\}^s$ differ by one of the specified short vectors.
Each such fiber consequently lies in a translate of
$\ker B\cap\Phi^{-1}(T)$, of dimension at most
$s-r-t+j-1$.  An injective coordinate projection on that linear
space bounds the size of the fiber by
$(K+1)^{s-r-t+j-1}$.  The number of $B$-images is at most
$(2sHK+1)^r$.  Thus
\[
 (K+1)^s\leq(2sHK+1)^r(K+1)^{s-r-t+j-1}
 \leq(2sH+1)^r(K+1)^{s-t+j-1},
\]
contrary to \eqref{pm:str:escape-condition}.
\end{proof}

For a rank-two target, this guarantees one surviving direction when
$K+1>(2sH+1)^{r/2}$, and two independent directions when
$K+1>(2sH+1)^r$.  If its rank assumptions hold for a Sondow window,
the coefficient estimates used above yield a sufficient bound of
order $\log K=O(N^2/\log N)$.  That sufficient bound is too large
to combine with the exponential error estimate; it is not a lower
bound on the size of actual Sondow target vectors.

\begin{proposition}[Limits of rank-and-height information]
\label{pm:str:chain-examples}
There are integer matrices of rank $r$, coefficient bound $H$, and
rank-two target maps for which every nonzero target requires an
integer lift of sup norm at least $H^{r/2}$ when $r$ is even.
There are also such examples in which obtaining two independent
target images requires a lift of sup norm at least $H^r$.
Arbitrarily many zero-target kernel coordinates can be appended
without changing these conclusions.
\end{proposition}

\begin{proof}
Let $H\geq2$ be an integer and first put $r=2m$.  Use the equations
\[
 x_i=Hx_{i+1},\qquad y_i=Hy_{i+1}\qquad(0\leq i<m),
 \qquad \Phi=(x_m,y_m).
\]
They have rank $2m$, and an integer solution has
$x_i=H^{m-i}a$, $y_i=H^{m-i}b$, with $a,b\in\mathbb Z$.
Any nonzero target has norm at least $H^m$ in its lift; this bound
is attained with $(a,b)=(1,0)$ or $(0,1)$.

For the second assertion use a single chain
$x_i=Hx_{i+1}$, $0\leq i<r$, an unconstrained coordinate $y$,
and $\Phi=(x_r,y)$.  The $y$ direction has a lift of norm one.
Two independent images, however, require a vector with $x_r\ne0$,
whose lift has norm at least $H^r$.  This is also attained.
In both constructions further unconstrained coordinates may be
added and mapped to zero.
\end{proof}

These examples concern general integer matrices, not the particular
Sondow matrix.  They show why additional information about the prime
exponents and the primitive target coordinates is required: an
abundance of short relations, by itself, does not guarantee short
relations that survive in the target.


% ---- convex.tex ----
\section{Primitive targets and an exact quotient norm}
\label{pm:cvx:sec:quotient}

Fix a finite block $\mathcal I$ of $s$ indices.  Retain the integer
coefficient identity from the preceding section in the vector form
\begin{equation}\label{pm:cvx:eq:vector-identity}
 e=a\gamma-b+B^T\lambda_\gamma,
 \qquad (\lambda_\gamma)_p=\log p,
 \qquad \Phi w=(a^Tw,b^Tw).
\end{equation}
Here $B$ has $r$ rows, indexed by the relevant primes, and $a,b,e$ are
column vectors of length $s$.  Put
\[
 \mathcal C=\begin{pmatrix}B\\a^T\\b^T\end{pmatrix}.
\]
Throughout this section we assume
\begin{equation}\label{pm:cvx:eq:rank-assumption}
 \rank\mathcal C=r+2.
\end{equation}
This assumption implies that
$\Phi:\ker B\longrightarrow\R^2$ is surjective.  We do not assert
\eqref{pm:cvx:eq:rank-assumption} for every moving block; it is checked
exactly for each block used in the finite computations below.

Let $\varepsilon_n\in\Q_{>0}$ be proved upper bounds for $e_n$.
The usual triangle inequality assigns to an integer relation $Bu=0$ the
bound $\sum_n\varepsilon_n|u_n|$.  The size of this expression before
dividing the target by its greatest common divisor can be misleading.
The following quotient construction performs the division exactly.

\begin{theorem}[Primitive quotient optimization]\label{pm:cvx:thm:primitive}
For $v\in\R^2$, define
\begin{equation}\label{pm:cvx:eq:F}
 F(v)=\min_{Bw=0,\ \Phi w=v}\sum_{n\in\mathcal I}
             \varepsilon_n|w_n|.
\end{equation}
Then $F$ is a norm on $\R^2$ whose unit ball is a centrally symmetric
rational polygon.  If $v=(q,p)\in\Z^2$ is primitive, then
\begin{equation}\label{pm:cvx:eq:primitive-equality}
 F(v)=
 \min_{\substack{u\in\Z^s,\ Bu=0,\ \Phi u=gv\\g\in\Z_{>0}}}
       \frac{\sum_n\varepsilon_n|u_n|}{g}.
\end{equation}
Both minima are attained.  For a rational target, a rational optimal lift in
\eqref{pm:cvx:eq:F} can be chosen with at most $r+2$ nonzero coordinates.
Moreover,
\begin{equation}\label{pm:cvx:eq:F-error}
 |q\gamma-p|\le F(q,p).
\end{equation}
\end{theorem}

\begin{proof}
Surjectivity gives a feasible lift for every target.  Positivity of all
weights makes the intersections of a fiber with sublevel sets of the
objective compact.  Consequently the minimum exists.  For a rational
target this is a rational linear program, after writing
$w=w^+-w^-$ with $w^+,w^-\ge0$, and it has a rational optimum.

Every integer lift on the right of
\eqref{pm:cvx:eq:primitive-equality} gives a feasible rational lift $u/g$
on the left.  Conversely, let $w$ be a rational optimum and choose a
positive common denominator $g$ of its coordinates.  Then $u=gw$ is
integral, $Bu=0$, and $\Phi u=g(q,p)$.  Since $\gcd(q,p)=1$,
\begin{equation}\label{pm:cvx:eq:gcd-exact}
 \gcd(a^Tu,b^Tu)=g.
\end{equation}
Thus the denominator used to clear the lift is exactly the greatest
common divisor removed from the target.  This proves the equality and
attainment in \eqref{pm:cvx:eq:primitive-equality}.

The quotient of the weighted $\ell^1$ norm is symmetric, homogeneous,
and subadditive.  If $F(v)=0$, an attained minimum has zero lift, hence
$v=0$.  The unit ball is the image under $\Phi$ of the intersection
of $\ker B$ with the rational weighted cross-polytope
$\sum\varepsilon_n|w_n|\le1$.  Surjectivity makes this image
two-dimensional; it is a centrally symmetric rational polygon.

For the support bound, choose a rational optimum.  If its support has
more than $r+2$ entries, there is a nonzero rational vector in
$\ker\mathcal C$ supported on those entries.  In a sufficiently small
interval about the optimum, the objective is affine along this vector.
Its slope must be zero, or one direction would improve the objective.
Move to an endpoint of the interval on which the coordinate signs are
unchanged.  At least one coordinate vanishes, and the objective remains
optimal.  Repetition proves the assertion.

Finally, \eqref{pm:cvx:eq:vector-identity} gives
$e^Tw=q\gamma-p$ whenever $Bw=0$ and $\Phi w=(q,p)$.
Taking the triangle inequality and then the minimum proves
\eqref{pm:cvx:eq:F-error}.
\end{proof}

The lattice $\Phi(\ker B\cap\Z^s)$ can have a nontrivial index in
$\Z^2$.  The theorem does not identify it with $\Z^2$.
Instead, it allows a multiple $gv$ of a primitive target and then
divides by its actual target gcd.  An obstruction formulated solely in
terms of the unnormalized lift cannot be carried over to $F$ without
accounting for \eqref{pm:cvx:eq:gcd-exact}.

\begin{corollary}[Invariance under positive column normalization]
\label{pm:cvx:cor:normalization}
Let $D=\operatorname{diag}(t_n)$, where $t_n\in\Q_{>0}$.
Replace $(B,a^T,b^T)$ by $(BD,a^TD,b^TD)$ and replace
$\varepsilon_n$ by $t_n\varepsilon_n$.  The function $F$ is unchanged.
\end{corollary}
\begin{proof}
The substitution $x=Dw$ identifies both feasible sets and their
objectives.  All data remain rational, so denominator clearing remains
available even if the normalized presentation is not integral.
\end{proof}

In particular, the multipliers $d_{2n}$ may be included in the columns
or removed from them together with the corresponding error weights.
Raw coefficient norms and raw dominance tests need not be invariant
under this operation.  The optimized primitive bound is invariant.

\begin{proposition}[Dual certificate]\label{pm:cvx:prop:dual}
For every $(q,p)\in\R^2$,
\begin{equation}\label{pm:cvx:eq:F-dual}
 F(q,p)=\max\left\{\eta_1q+\eta_2p:
 \left|B^T\lambda+a\eta_1+b\eta_2\right|\le\varepsilon\right\},
\end{equation}
where inequalities and absolute values are coordinatewise.
For a rational target, rational primal and dual optima exist.
\end{proposition}
\begin{proof}
The linear program with matrix $[\mathcal C,-\mathcal C]$ and
nonnegative variables $(w^+,w^-)$ has costs
$(\varepsilon,\varepsilon)$.  Its dual is precisely
\eqref{pm:cvx:eq:F-dual}.  The primal is feasible and has finite attained
minimum by Theorem~\ref{pm:cvx:thm:primitive}, so finite-dimensional linear
programming duality applies.
\end{proof}

In particular a finite optimum can be checked without trusting the
optimization algorithm.  It is enough to supply a rational feasible
lift $w$ and rational $(\lambda,\eta_1,\eta_2)$ satisfying the dual
inequalities and equality of the two objectives.  Equivalently, if
$z\in\operatorname{row}(\mathcal C)$ satisfies
$|z_n|\le\varepsilon_n$ and
$z_nw_n=\varepsilon_n|w_n|$ for every $n$, then $z$ certifies the
global optimality of $w$ on its fiber.  These are rational equalities
and inequalities.

\section{Nonvanishing from polynomial positivity}
\label{pm:cvx:sec:positivity}

The conditions $u\ne0$ and $\Phi u\ne0$ do not by themselves imply
$e^Tu\ne0$.  A separate analytic sign certificate is possible for
many Sondow combinations.  Put
\[
 T(x,y)=x(1-x)y(1-y),\qquad
 h(x,y)=\frac{1}{(1-xy)(-\log(xy))}.
\]
Define a measure on $[0,1/16]$ by
\begin{equation}\label{pm:cvx:eq:measure}
 \nu(E)=\int_{\{(x,y)\in(0,1)^2:T(x,y)\in E\}}
                  T(x,y)h(x,y)\,dx\,dy.
\end{equation}
Its mass is $I_1<1/4$, as follows also from
Lemma~\ref{pm:lem:integral-bounds}.  Every nonempty open
subinterval of $(0,1/16)$ has positive $\nu$-measure: the continuous
map $T$ takes every such value in the open square, and its inverse
image of the subinterval is nonempty and open.  Moreover,
\begin{equation}\label{pm:cvx:eq:moments}
 I_n=\int_0^{1/16}t^{n-1}\,d\nu(t),\qquad n\ge1.
\end{equation}
The factor $T$ in \eqref{pm:cvx:eq:measure} ensures that the measure is
finite; no finite zeroth integral $I_0$ is used.

For a consecutive block $\mathcal I=\{N,\ldots,N+m\}$ with $N\ge1$
and a real lift $w=(w_0,\ldots,w_m)$, define
\begin{equation}\label{pm:cvx:eq:Rw}
 R_w(X)=\sum_{j=0}^m w_jd_{2(N+j)}16^{-j}X^j.
\end{equation}
The notation $w_j$ here refers to the coordinate at index $N+j$.
Equation~\eqref{pm:cvx:eq:moments} gives the exact identity
\begin{equation}\label{pm:cvx:eq:positive-integral}
 \sum_{j=0}^m w_jd_{2(N+j)}I_{N+j}
 =\int_0^{1/16}t^{N-1}R_w(16t)\,d\nu(t).
\end{equation}

\begin{proposition}[A gamma-independent sign certificate]
\label{pm:cvx:prop:polynomial-sign}
Suppose $Bw=0$, $\Phi w=(q,p)$, and $R_w$ is nonnegative on
$[0,1]$ and not identically zero.  Then $q\gamma-p>0$.
\end{proposition}
\begin{proof}
A nonzero nonnegative polynomial is positive on a nonempty open
subinterval of $(0,1)$.  Full support of $\nu$, together with
$t^{N-1}>0$ for $t>0$, makes the right side of
\eqref{pm:cvx:eq:positive-integral} strictly positive.  Its left side is
$q\gamma-p$ by prime-log cancellation.
\end{proof}

\subsection{Exact Bernstein coefficients}

For $R(X)=\sum_{j=0}^m r_jX^j$, the degree-$m$ Bernstein expansion is
\begin{equation}\label{pm:cvx:eq:bernstein}
 R(X)=\sum_{k=0}^m b_k\binom mkX^k(1-X)^{m-k},\qquad
 b_k=\sum_{j=0}^k r_j\frac{\binom kj}{\binom mj}.
\end{equation}
Indeed, the identity
\[
 X^j=\sum_{k=j}^m\frac{\binom kj}{\binom mj}
                     \binom mkX^k(1-X)^{m-k}
\]
follows from $\binom mk\binom kj=\binom mj\binom{m-j}{k-j}$
and the binomial theorem.  Thus rational monomial coefficients give
rational Bernstein coefficients by an explicit triangular transform.
If all $b_k\ge0$ and at least one is positive, then $R(X)>0$ for
$0<X<1$.  If all $b_k>0$, the partition-of-unity identity gives
$R(X)\ge\min_kb_k>0$ throughout $[0,1]$.

The elementary sufficient condition
\begin{equation}\label{pm:cvx:eq:dominance}
 r_0>\sum_{j=1}^m|r_j|
\end{equation}
implies positivity of every $b_k$, since each binomial ratio in
\eqref{pm:cvx:eq:bernstein} belongs to $[0,1]$.  Its converse fails.
Consequently Bernstein positivity supplies additional sign
certificates beyond first-coefficient dominance; the finite Sondow
example below exhibits this difference within an actual log-free
kernel.

\begin{proposition}[Eventual degree elevation]
\label{pm:cvx:prop:elevation}
Let $R\in\Q[X]$ have degree at most $m$ and be strictly positive
on $[0,1]$.  For all sufficiently large integers $M\ge m$, every
coefficient in its degree-$M$ Bernstein expansion is positive.
More precisely, put
\[
 C_R=\frac12\sum_{j=2}^m |r_j|j(j-1).
\]
For all $0\le k\le M$,
\begin{equation}\label{pm:cvx:eq:elevation-bound}
 \left|b_k^{(M)}-R(k/M)\right|
 \le\frac{C_R}{M-m+1},\qquad
 b_k^{(M)}=\sum_{j=0}^{\min(k,m)}
                    r_j\frac{\binom kj}{\binom Mj}.
\end{equation}
In particular, if $R\ge\delta>0$, it suffices that
$C_R/(M-m+1)<\delta$.
\end{proposition}
\begin{proof}
The contributions $j=0,1$ agree exactly with evaluation at $k/M$.
For $j\ge2$ and $k\ge j$, compare
\[
 \frac{\binom kj}{\binom Mj}
 =\prod_{i=0}^{j-1}\frac{k-i}{M-i}
 \quad\hbox{with}\quad
 (k/M)^j.
\]
All factors belong to $[0,1]$, and
\[
 \left|\frac{k-i}{M-i}-\frac{k}{M}\right|
 =\frac{i(M-k)}{M(M-i)}\le\frac{i}{M-m+1}.
\]
Telescoping the difference between the two products gives the bound
$j(j-1)/(2(M-m+1))$.  If $k<j$, the binomial ratio is zero and
\[
 (k/M)^j\le (j-1)/M
 \le \frac{j(j-1)}{2(M-m+1)}.
\]
Summing these bounds with weights $|r_j|$ proves
\eqref{pm:cvx:eq:elevation-bound}.  Compactness gives
$\delta=\min_{[0,1]}R>0$, proving eventual positivity.
\end{proof}

The proposition gives a rational semidecision procedure for strict
positivity of a fixed polynomial.  It makes no bound uniform over
all lifts or all levels.  A nonnegative polynomial with an interior
zero cannot have a nonzero expansion with all Bernstein coefficients
nonnegative, because every Bernstein basis function is positive in
the interior.  Such polynomials require a different exact
nonnegativity certificate.

For an explicit lower bound, on $[1/4,3/4]^2$ one has
$T\ge9/256$ and $h>1/3$, using $-\log(xy)\le\log16<3$.
The square has area $1/4$, so
\begin{equation}\label{pm:cvx:eq:rough-lower}
 I_N>\frac1{12}(9/256)^N.
\end{equation}
If $s_0\in\{1,-1\}$ orients $R_w$ so that all its Bernstein
coefficients are positive, and $b_*$ is their minimum, then
\begin{equation}\label{pm:cvx:eq:sign-lower}
 s_0(q\gamma-p)>\frac{b_*}{12}(9/256)^N>0.
\end{equation}
Everything on the right is rational for a rational lift.

\subsection{Positivity inside the primitive optimization}

Let $\mathcal B$ denote the rational linear map from $w$ to the
Bernstein coefficients of $R_w$ at a fixed degree.  Define
\begin{equation}\label{pm:cvx:eq:F-positive}
 F^+(v)=\min_{\substack{Bw=0,\ \Phi w=v\\\mathcal Bw\ge0}}
                              \sum_j\varepsilon_j|w_j|,
\end{equation}
with value $+\infty$ if infeasible.  This is an asymmetric,
extended-valued optimization, not a norm.  For a nonzero rational
target, every feasible lift has a nonzero polynomial: the monomial
coefficient map in \eqref{pm:cvx:eq:Rw} is invertible.  Thus
\begin{equation}\label{pm:cvx:eq:F-positive-error}
 0<q\gamma-p\le F^+(q,p)
\end{equation}
whenever the minimum is finite.  Compactness of weighted $\ell^1$
sublevel sets and rational linear programming give an attained
rational optimum.  For a primitive integer target one has
\begin{equation}\label{pm:cvx:eq:positive-gcd}
 F^+(v)=
 \min_{\substack{u\in\Z^s,\ Bu=0,\ \Phi u=gv,\ g\in\Z_{>0}
                 \\\mathcal Bu\ge0}}
       \frac{\sum_j\varepsilon_j|u_j|}{g}.
\end{equation}
The proof is exactly denominator clearing as in
Theorem~\ref{pm:cvx:thm:primitive}: the cone inequalities are
homogeneous and survive multiplication or division by $g$.
This shows that an explicit sign condition can be imposed before
optimization without losing the exact gcd cancellation.  It does
not establish feasibility or decay at arbitrarily large levels.

\section{Two-sided bounds and an affine feasible interval}
\label{pm:cvx:sec:interval}

We use the strict bounds $\ell_n<e_n<h_n$ proved in
Lemma~\ref{pm:lem:integral-bounds}.

Retain a fixed block satisfying \eqref{pm:cvx:eq:rank-assumption} and
write $w^+_n=\max(w_n,0)$, $w^-_n=\max(-w_n,0)$.  The two-sided
information suggests the support function
\begin{equation}\label{pm:cvx:eq:U}
 U(v)=\min_{Bw=0,\ \Phi w=v}
                  \{h^Tw^+-\ell^Tw^-\}.
\end{equation}

\begin{theorem}[Asymmetric primitive optimization]
\label{pm:cvx:thm:asymmetric}
For every target the minimum in \eqref{pm:cvx:eq:U} is attained; a
rational optimum exists for every rational target.  For $v=(q,p)\ne0$,
\begin{equation}\label{pm:cvx:eq:U-interval}
 -U(-v)<q\gamma-p<U(v).
\end{equation}
For a primitive integer target,
\begin{equation}\label{pm:cvx:eq:U-gcd}
 U(v)=\min_{\substack{u\in\Z^s,\ Bu=0,\ \Phi u=gv
                         \\g\in\Z_{>0}}}
           \frac{h^Tu^+-\ell^Tu^-}{g}.
\end{equation}
Its dual is
\begin{equation}\label{pm:cvx:eq:U-dual}
 U(q,p)=\max\left\{\eta_1q+\eta_2p:
       \ell\le B^T\lambda+a\eta_1+b\eta_2\le h\right\}.
\end{equation}
\end{theorem}
\begin{proof}
For a lift in the fiber of $v$, subtracting its exact value from the
objective gives
\[
 h^Tw^+-\ell^Tw^- -(q\gamma-p)
   =(h-e)^Tw^++(e-\ell)^Tw^-.
\]
All coefficients on the right are strictly positive.  On this fixed
finite block their minimum is positive, so the objective has compact
sublevel sets on each fiber.  This proves attainment and, for every
nonzero lift, strict domination of $q\gamma-p$.  A nonzero target
has a nonzero optimal lift.  Applying the same argument to $-v$
proves \eqref{pm:cvx:eq:U-interval}.  Rational linear programming gives
a rational optimum for rational $v$.  Clearing its denominators and
using \eqref{pm:cvx:eq:gcd-exact} proves \eqref{pm:cvx:eq:U-gcd}.

To compute the dual, use costs $(h,-\ell)$ for nonnegative variables
$(w^+,w^-)$.  The dual inequalities are
$B^T\lambda+a\eta_1+b\eta_2\le h$ and its negative at most
$-\ell$, which gives \eqref{pm:cvx:eq:U-dual}.  Strong duality applies
because the primal optimum is finite and attained.
\end{proof}

If the only lower information is $e_n\ge0$, the origin belongs to
the dual feasible set and $U(v)\ge0$ in every direction.  Such a
support function alone cannot give a strictly negative upper bound
or a strictly positive lower bound.  The quantitative lower bounds
in \eqref{pm:eq:two-sided-bounds} contain additional sign information.

\subsection{Fixing the constant coordinate}

Equation~\eqref{pm:cvx:eq:vector-identity} specifies the coefficient
of $b$ exactly: it is $-1$.  Preserving this coordinate gives the
following relaxation of the gamma identities:
\begin{equation}\label{pm:cvx:eq:G}
 G=\left\{\theta\in\R:
   \exists y\in\R^r\quad
   \ell\le a\theta-b+B^Ty\le h\right\}.
\end{equation}
The variables $y$ are arbitrary real numbers.  In particular, this
model does not impose that $y$ consists of actual prime logarithms.

\begin{theorem}[Exact interval of the finite boxed model]
\label{pm:cvx:thm:interval}
The feasible set $G$ is a compact rational interval
$[\alpha,\beta]$ with
\begin{equation}\label{pm:cvx:eq:gamma-interior}
 \alpha<\gamma<\beta.
\end{equation}
Its exact endpoints are
\begin{align}
 \beta&=\min_{Bw=0,\ a^Tw=1}
                  \{b^Tw+h^Tw^+-\ell^Tw^-\},
                       \label{pm:cvx:eq:beta-LP}\\
 -\alpha&=\min_{Bw=0,\ a^Tw=-1}
                  \{b^Tw+h^Tw^+-\ell^Tw^-\}.
                       \label{pm:cvx:eq:alpha-LP}
\end{align}
Both linear programs have rational primal and dual optima.  For every
real target $(q,p)$, the best worst-case absolute error permitted by
the model \eqref{pm:cvx:eq:G} is
\begin{equation}\label{pm:cvx:eq:J}
 J(q,p):=\max_{\theta\in G}|q\theta-p|
       =\max\{|q\alpha-p|,|q\beta-p|\}.
\end{equation}
For $q>0$, this gives the strict signed bounds
\begin{equation}\label{pm:cvx:eq:sliced-error}
 q\alpha-p<q\gamma-p<q\beta-p,
\end{equation}
which are at least as strong as \eqref{pm:cvx:eq:U-interval}.
\end{theorem}
\begin{proof}
The point $(\theta,y)=(\gamma,\lambda_\gamma)$ satisfies every
inequality strictly.  It therefore has a neighborhood within the
feasible set.  The map $(y,\theta)\mapsto B^Ty+a\theta$ is
injective because $[B;a^T]$ has full row rank.  Its inverse on its
image is continuous, so the inverse image of the bounded box
$b+[\ell,h]$ is bounded.  It is a nonempty compact rational
polytope.  Its projection onto the $\theta$-axis is a compact
rational interval, and the neighborhood of the actual point places
$\gamma$ strictly between its endpoints.

For \eqref{pm:cvx:eq:beta-LP}, split $w$ into nonnegative parts.  The
costs become $(b+h,-b-\ell)$ and the constraint matrix is
$[B;a^T]$.  The dual maximizes $\theta$ under
\[
 \ell\le B^Ty+a\theta-b\le h,
\]
which is exactly the maximum in \eqref{pm:cvx:eq:G}.  Changing the
target coordinate $1$ to $-1$ changes the dual objective to
$-\theta$ and proves \eqref{pm:cvx:eq:alpha-LP}.  Feasibility,
boundedness, and rationality give the asserted attained optima.

The interval formula \eqref{pm:cvx:eq:J} is immediate from linearity.
The points $(y,\theta,-1)$ form a slice of the feasible set in
\eqref{pm:cvx:eq:U-dual}.  Taking the maximum of $q\theta-p$ on this
slice cannot exceed $U(q,p)$; taking the minimum cannot be less than
$-U(-q,-p)$.  This proves the comparison of signed bounds.  Strictness
in \eqref{pm:cvx:eq:sliced-error} follows from
\eqref{pm:cvx:eq:gamma-interior} when $q>0$.
\end{proof}

The qualification ``of the finite boxed model'' is essential.
Endpoints are feasible in the relaxation, hence no smaller interval
is implied solely by those linear inequalities.  Additional valid
constraints, for example the correlated positive-difference
inequalities developed elsewhere in this part, can shrink this
interval.  The improvement over $F$ also does not contradict the
optimality of the fixed-target quotient norm: an endpoint certificate
allows $b^Tw$ to vary and subsequently adds an exactly known rational
constant when estimating a prescribed $(q,p)$.

For completeness, let $c=(\alpha+\beta)/2$ and
$\rho=(\beta-\alpha)/2>0$.  The elementary identity
$\max(|u+v|,|u-v|)=|u|+|v|$ yields
\begin{equation}\label{pm:cvx:eq:parallelogram}
 J(q,p)=\rho|q|+|cq-p|.
\end{equation}
Thus $J$ is a norm whose unit ball is a rational parallelogram.
The change $(q,p)\mapsto(q,cq-p)$ has determinant $-1$, so its
area is $2/\rho=4/(\beta-\alpha)$.  This explicit geometry does
not supply denominator growth or an infinite approximation theorem:
the placement of the parallelogram relative to $\Z^2$ remains an
arithmetic question.

\section{Exact finite certificates at two levels}
\label{pm:cvx:sec:finite}

The computations use the blocks $\{N,\ldots,2N\}$ for $N=32,40$.
All Sondow coefficients, bounds, linear programs, and comparisons
are rational or integral.  The targets are retained from the
preceding quotient computation.  That construction chose continued
fractions of a rational center computed from a rational right inverse
of $\mathcal C$ and a weighted-median minimization; it did not take
continued fractions of a supplied decimal value of $\gamma$.
Nothing about this finite selection establishes a uniform assertion
as $N$ varies.

Write $(q_N,p_N)$ for the first retained target at level $N$.  They are
\begin{equation}\label{pm:cvx:eq:target32}
 \begin{aligned}
 q_{32}&=75465755513411273280863840027,\\
 p_{32}&=43560016245970207457736606578,
 \end{aligned}
\end{equation}
\begin{equation}\label{pm:cvx:eq:target40}
 \begin{aligned}
 q_{40}&=1466031611818694840452495676451550525,\\
 p_{40}&=846216411582593862687645010331522899.
 \end{aligned}
\end{equation}
Their numerator--denominator gcds are $1$.  In the following table,
$\alpha_N,\beta_N$ are the exact rational endpoints of
Theorem~\ref{pm:cvx:thm:interval}, $\mathsf L_N=q_N\alpha_N-p_N$, and
$\mathsf R_N=q_N\beta_N-p_N$.  Decimal logarithms are descriptive rounded
values; they are not used to prove any displayed rational inequality.

\begin{table}[htbp]
\centering
\small
\begin{tabular}{c r r r c}
\hline
$N$ & $\log(\beta_N-\alpha_N)$ & $\log|\mathsf L_N|$ &
       $\log|\mathsf R_N|$ & signs of $(\mathsf L_N,\mathsf R_N)$\\
\hline
32 & $-132.84450590$ & $-68.20232457$ & $-66.20516987$ & $(+,+)$\\
40 & $-174.74952280$ & $-86.07078916$ & $-86.07530188$ & $(-,-)$\\
\hline
\end{tabular}
\caption{Exact endpoint certificates summarized on the target-form scale.
The large rational endpoints and their primal--dual witnesses are stored
in the accompanying certificate file.}
\label{pm:cvx:tab:endpoints}
\end{table}

The corresponding finite changes of error bound are shown in
Table~\ref{pm:cvx:tab:improvement}.  Here the symmetric bound uses
$\varepsilon=h$ in \eqref{pm:cvx:eq:F}, the two-sided bound uses
\eqref{pm:cvx:eq:U-interval}, and the projected bound uses
\eqref{pm:cvx:eq:J}.

\begin{table}[htbp]
\centering
\begin{tabular}{c r r r}
\hline
$N$ & symmetric log bound & two-sided log bound & projected log bound\\
\hline
32 & $-65.730623$ & $-65.782655$ & $-66.205170$\\
40 & $-84.632505$ & $-84.633619$ & $-86.070789$\\
\hline
\end{tabular}
\caption{Three certified upper bounds for the same two primitive targets.}
\label{pm:cvx:tab:improvement}
\end{table}

There is also a direct polynomial sign certificate for each of these
targets, using the optimal lift for the symmetric bound.  For $N=32$
the polynomial fails \eqref{pm:cvx:eq:dominance}, but all its $33$
degree-$32$ Bernstein coefficients are positive.  For $N=40$ the
negative of the lift polynomial has all $41$ degree-$40$ Bernstein
coefficients positive.  Formula~\eqref{pm:cvx:eq:sign-lower} gives
positive rational lower bounds whose logarithms are approximately
$-85.084754$ and $-107.596244$, respectively.  These establish the
signs directly from the integral representation at the same level.
The affine interval bounds in Table~\ref{pm:cvx:tab:endpoints} are
stronger for these targets.

\begin{proposition}[A finite exponent certificate]
\label{pm:cvx:prop:finite-exponent}
For the integers $q=q_{40}$ and $p=p_{40}$ in
\eqref{pm:cvx:eq:target40}, the exact certificates prove
\begin{equation}\label{pm:cvx:eq:103}
 0<p-q\gamma<q^{-103/100}.
\end{equation}
\end{proposition}
\begin{proof}
The certificate supplies exact rational endpoints
$\alpha_{40},\beta_{40}$ satisfying
\[
 q\beta_{40}-p<0,\qquad
 E:=p-q\alpha_{40}>0,\qquad E^{100}q^{103}<1.
\]
Theorem~\ref{pm:cvx:thm:interval} therefore gives
$0<p-q\gamma<E<q^{-103/100}$.  All three certificate comparisons
are rational or integral; no decimal evaluation of $\gamma$ enters
the proof.
\end{proof}

To make the role of computation precise, the supplied verifier
reconstructs $B,a,b,\ell,h$ from the coefficient formulas.  For each
linear program it checks a rational lift, the exact target equations,
a rational dual vector, every dual inequality, and equality of the
objectives.  In the endpoint programs the constraint matrix is
$[B;a^T]$ and the extra objective offset is $b^Tw$.
For an arbitrary rational offset $o$ and target $t$, a certificate
for the problem
\[
 \min_{Aw=t}\{o^Tw+h^Tw^+-\ell^Tw^-\}
\]
consists of $w,y$ satisfying
\[
 Aw=t,\qquad \ell\le A^Ty-o\le h,\qquad
 o^Tw+h^Tw^+-\ell^Tw^-=t^Ty.
\]
Weak duality proves optimality directly from these relations.
The verifier separately checks denominator clearing and
\eqref{pm:cvx:eq:gcd-exact}, the signs in
Table~\ref{pm:cvx:tab:endpoints}, and $E^{100}q^{103}<1$.
The positivity verifier also reconstructs and inverts the Bernstein
transform rather than sampling a grid of real points.

The principal data files are
\nolinkurl{gamma_push_quotient_sharp_certificate.json},
\nolinkurl{gamma_sondow_intervals_certificate.json}, and
\nolinkurl{gamma_sondow_positivity_certificate.json}.
The interval certificate can be checked without reoptimizing by
running
\begin{verbatim}
python3 gamma_sondow_intervals.py --verify
\end{verbatim}
the corresponding positivity checks are implemented in
\nolinkurl{gamma_sondow_positivity.py}.  The accompanying sources
specify the dependencies and preserve the exact integers and rational
witnesses omitted from the printed tables.

Proposition~\ref{pm:cvx:prop:finite-exponent} concerns one pair, not an
infinite sequence.  In particular, the exponent $103/100$ in the
integer linear form, or $203/100$ after division by $q$, supplies no
irrationality or transcendence conclusion from this finite
certificate.  The results of these sections are the exact
normalization of primitive bounds, explicit nonvanishing criteria,
the optimal interval in a specified finite relaxation, and verified
finite applications of those constructions.


% ---- damping.tex ----
\section{Positive differences of moving order}\label{pm:sec:damping}

The absolute value of each coefficient in an alternating combination is
often an unnecessarily expensive way to estimate its error.  A common
positive moment representation can retain the correlation between the
integrals.  Here that observation gives an explicit construction at every
pair of indices, together with its exact exponential rate.

\begin{lemma}[A finite moment carrier]\label{pm:dm:measure}
Let $\sigma$ be the pushforward of $t h(x,y)\,dx\,dy$ under $s=16t$.
It is a finite positive measure on $[0,1]$, gives positive mass to every
nonempty open subinterval of $(0,1)$, and satisfies
\begin{equation}\label{pm:dm:moment}
 \sigma([0,1])=I_1<\frac14,\qquad
 I_r=16^{-(r-1)}\int_0^1s^{r-1}\,d\sigma(s)\quad(r\ge1).
\end{equation}
\end{lemma}
\begin{proof}
The upper comparison in Lemma~\ref{pm:lem:integral-bounds} gives
$t h(x,y)\le xy$, whose integral is $1/4$.  The strict mass bound follows.
For any interior value of $s$, an open subinterval around it has a nonempty
open preimage inside the square; on that preimage the density is positive.
This proves full interior support.  The moment identity is the definition
of pushforward applied to $(16t)^{r-1}$.
\end{proof}

\begin{theorem}[All-index positive differences and rational bounds]
\label{pm:dm:positive}
For integers $n\ge1$ and $m\ge0$, put
\[
 D_{n,m}=\sum_{k=0}^m(-16)^k\binom mk I_{n+1+k}.
\]
Then
\begin{equation}\label{pm:dm:identity}
 D_{n,m}=16^{-n}\int_0^1s^n(1-s)^m\,d\sigma(s)>0
\end{equation}
and, with $0^0=1$,
\begin{align}
 D_{n,m}&<\frac{16^{-n}}4
       \frac{n^nm^m}{(n+m)^{n+m}},\label{pm:dm:upper}\\
 D_{n,m}&>\frac1{1024}
       \left(\frac9{256}\right)^{n+1}
       \left(\frac{31}{256}\right)^m.\label{pm:dm:lower}
\end{align}
\end{theorem}
\begin{proof}
The binomial theorem applied inside \eqref{pm:dm:moment} proves
\eqref{pm:dm:identity}; the integrand is positive in the interior.  Its maximum
is $n^nm^m/(n+m)^{n+m}$, at $s=n/(n+m)$ when $m>0$; the case $m=0$
is immediate.  The mass bound in Lemma~\ref{pm:dm:measure} proves the upper
bound.  For the lower bound, return to the square and write
\[
 D_{n,m}=\int_{(0,1)^2}t^{n+1}(1-16t)^m h(x,y)\,dx\,dy.
\]
On $[1/4,3/8]^2$ one has $t\ge9/256$,
$1-16t\ge31/256$, and $h>xy\ge1/16$.  Its area is $1/64$,
so restriction to this square proves \eqref{pm:dm:lower}.
\end{proof}

The bound obtained instead from $I_r<4\,16^{-r}$ by taking absolute
values term by term is $(16^{-n}/4)2^m$.  Its ratio to
\eqref{pm:dm:upper} is
\[
 2^m\frac{(n+m)^{n+m}}{n^nm^m},
\]
equal to $8^n$ when $m=n$.  This is a comparison of certified upper bounds,
not a claim that the actual errors change under a change of representation.

\begin{theorem}[Exact moving-order rate]\label{pm:dm:rate}
Suppose $m_n\ge0$ are integers and $m_n/n\to\lambda\in[0,\infty)$.
Define
\[
 \mathcal H(\lambda)=(1+\lambda)\log(1+\lambda)
                       -\lambda\log\lambda,\qquad \mathcal H(0)=0.
\]
Then
\begin{equation}\label{pm:dm:rate-raw}
 \lim_{n\to\infty}\frac{\log D_{n,m_n}}n
 =-\log16-\mathcal H(\lambda).
\end{equation}
For $E_{n,m}=d_{2(n+m+1)}D_{n,m}$, one consequently has
\begin{equation}\label{pm:dm:rate-int}
 \lim_{n\to\infty}\frac{\log E_{n,m_n}}n=-\kappa(\lambda),
 \quad \kappa(\lambda)=\log16-2-2\lambda+\mathcal H(\lambda).
\end{equation}
The maximum over $\lambda\ge0$ is
\begin{equation}\label{pm:dm:optimum}
 \kappa_{\max}=\log16-2-\log(1-e^{-2})=0.9180021801\ldots,
 \qquad \lambda_*=(e^2-1)^{-1}.
\end{equation}
\end{theorem}
\begin{proof}
The upper limit in \eqref{pm:dm:rate-raw} follows from \eqref{pm:dm:upper}.
If $\lambda>0$, the function
$\log s+\lambda\log(1-s)$ has its unique maximum at
$s_*=(1+\lambda)^{-1}$.  On every closed interval contained in $(0,1)$,
the exponents with $m_n/n$ converge uniformly to this function.  Choose a
small interval around $s_*$ and restrict \eqref{pm:dm:identity} to it.
Its $\sigma$-mass is positive.  The logarithm of that fixed mass, divided
by $n$, tends to zero, which gives the matching lower limit as the interval
shrinks.  When $\lambda=0$, use a fixed interval arbitrarily close to $1$;
on each such interval the term $(m_n/n)\log(1-s)$ tends uniformly to zero.
This gives the same conclusion.

The prime number theorem in the form $\log d_j=j+o(j)$ proves
\eqref{pm:dm:rate-int}.  Finally
\[
 \kappa'(\lambda)=-2+\log(1+1/\lambda),\qquad
 \kappa''(\lambda)=-\frac1{\lambda(1+\lambda)}<0.
\]
The unique critical point is $\lambda_*$, and substitution gives
\eqref{pm:dm:optimum}.  The endpoint $\lambda=0$ has rate $\log16-2$.
\end{proof}

A construction with rational index rules is already available:
$m_n=\lfloor n/6\rfloor$ gives $\lambda=1/6$ and a rate larger than
$\log16-2$.  There is no need to know the decimals of $\gamma$ to choose
these orders.  Nevertheless $E_{n,m}$ is an integer linear form in
$1,\gamma$, and prime logarithms, not yet in $1,\gamma$ alone.

\begin{proposition}[Growth of the displayed gamma coefficient]
\label{pm:dm:coefficient}
Let
\[
 \mathcal A_{n,m}=\sum_{k=0}^m(-16)^k\binom mk C_{n+1+k}.
\]
For every sequence in Theorem~\ref{pm:dm:rate}, this coefficient is eventually
nonzero with sign $(-1)^{m_n}$ and
\begin{equation}\label{pm:dm:coefficient-rate}
 \lim_{n\to\infty}\frac{\log|\mathcal A_{n,m_n}|}n
 =\log4+\lambda\log63.
\end{equation}
Thus the gamma coefficient in $E_{n,m_n}$ has exponential rate
$2(1+\lambda)+\log4+\lambda\log63$.
\end{proposition}
\begin{proof}
The arcsine moment formula for the central binomial coefficients gives
\[
 \mathcal A_{n,m}=\int_0^1
 (4y)^{n+1}(1-64y)^m\frac{dy}{\pi\sqrt{y(1-y)}}.
\]
On $y>1/64$ the integrand has constant sign $(-1)^m$.  Its logarithmic
absolute value divided by $n$ has supremum tending to
$\log4+\lambda\log63$, approached as $y\uparrow1$.  Restriction to a
fixed interval near $1$ gives the matching lower exponential bound, just
as in the proof of Theorem~\ref{pm:dm:rate}.  The absolute contribution of
$0<y<1/64$ is at most $16^{-n-1}$, and hence cannot cancel the first
contribution for large $n$.  This proves the sign and the limit; the
integer normalization contributes $2(1+\lambda)$ to the rate.
\end{proof}

This coefficient estimate is not a statement about the final primitive
height after another elimination.  The analytic decay rate, the
coefficient growth rate, and the primitive height must remain separate.

\subsection{Correlated constraints in a fixed window}

Consider a block $\mathcal J=\{N,\ldots,M\}$ and choose $i\ge N$,
$m\ge0$ with $i+m\le M$ and $i\ge2$.  Put $D=d_{2M}$ and define
the integer vector
\begin{equation}\label{pm:dm:cut-vector}
 u_{i+k}=\frac{D}{d_{2(i+k)}}(-16)^k\binom mk\quad(0\le k\le m),
 \qquad u_j=0\ \text{otherwise}.
\end{equation}
Then $u\cdot e=D D_{i-1,m}$.  The rational lower and upper bounds in
Theorem~\ref{pm:dm:positive}, multiplied by $D$, give valid inequalities
on this linear combination.  Append them to the individual interval
constraints in the same variables $\theta$ and $y$ used for the gamma
projection.  Equivalently, append the column $(Bu,a\cdot u)$ and the
constant offset $b\cdot u$ in the dual optimization.  No new prime-logarithm
variables are introduced.

The feasible set can only decrease, so its projected gamma interval can
only contract.  Strict contraction is a finite statement requiring a
certificate.  For $\mathcal J=\{N,\ldots,2N\}$, the stored computations
take starts
\[
 i\in\{N,N+\lfloor N/4\rfloor,N+\lfloor N/2\rfloor\}
\]
and orders $m\in\{1,2,4,8,16\}$ at $N=32,40$.  All fifteen pairs
are contained in their respective windows.  The exact primal-dual
certificates give Table~\ref{pm:tab:correlated}; its decimals are only
summaries of rational endpoint comparisons.

\begin{table}[htbp]
\centering
\caption{Finite contraction after fifteen correlated constraints.
Widths are positive; all logarithms are natural.}\label{pm:tab:correlated}
\begin{tabular}{rrrr}
\toprule
$N$ & $\log$(box width) & $\log$(correlated width) & width ratio\\
\midrule
32 & $-132.8445058965$ & $-135.2714119062$ & $11.32379212$\\
40 & $-174.7495228030$ & $-176.9827640464$ & $9.33005811$\\
\bottomrule
\end{tabular}
\end{table}

Both new endpoints lie strictly inside the old interval.  Since additional
constraints were imposed, this improvement is compatible with optimality
of the earlier box relaxation.  These two contractions do not prove
convergence of moving-window widths, much less a persistent improvement in
the Diophantine exponent.


% ---- heights.tex ----
\section{Rational moment approximants and graded denominator bounds}
\label{pm:ht:section}

This section concerns a rational approximation family distinct from the
logarithm-elimination construction.  Its analytic error has a fixed
polynomial scale, whereas the remaining arithmetic question concerns
the denominator after complete cancellation.  The positive-moment
representation and the Pad\'e construction are classical mechanisms;
the statements below make the normalization and the remaining
denominator question explicit.

\subsection{A positive rational moment problem}

Write
\begin{equation}\label{pm:ht:moments}
 \mathscr B(z)=\frac{z}{-\log(1-z)}=\sum_{m\geq0}\mathfrak b_mz^m,
 \qquad c_j=-\frac{\mathfrak b_{j+1}}{j+1}
 =\frac{1}{(j+1)(j+1)!}
   \int_0^1x\prod_{k=1}^j(k-x)\,dx.
\end{equation}
In particular, every $c_j$ is a positive rational number,
$c_0=1/2$, and $c_1=1/24$.  The coefficient identity follows by
integrating $(1-z)^x$ over $0<x<1$.

For $0<t<1$, put
\begin{equation}\label{pm:ht:measure}
 \rho(t)=\frac{1}{t\{\log^2((1-t)/t)+\pi^2\}},\qquad
 v(t)=\int_t^1\frac{\rho(s)}s\,ds,\qquad d\nu(t)=v(t)\,dt.
\end{equation}
Schr\"oder's classical representation of the Gregory coefficients gives
$-\mathfrak b_{j+1}=\int_0^1t^j\rho(t)\,dt$; a contour derivation and its history
are given in \cite{pm:Blagouchine2017}.  One can check the density directly:
the function $(1-\mathscr B(z))/z$ is analytic off $[1,\infty)$, tends to zero
at infinity, and its upper-bank jump at $z=1/t$ is
$2\pi i/(\log^2((1-t)/t)+\pi^2)$.  Cauchy's formula around the slit
therefore yields the stated moment representation.  Tonelli's theorem
now gives
\begin{equation}\label{pm:ht:gamma-markov}
 c_j=\int_0^1t^j\,d\nu(t),\qquad
 \gamma=\sum_{j\geq0}c_j
       =\int_0^1\frac{d\nu(t)}{1-t}.
\end{equation}
For completeness, the middle identity follows without an arithmetic
assumption on $\gamma$: integrating $(1-\mathscr B(z))/z$ over $(0,1)$ and
substituting $t=1-z$ gives
\[
 \sum_{j\geq0}c_j
 =\int_0^1\left(\frac1{1-t}+\frac1{\log t}\right)dt
 =\lim_{m\to\infty}(H_m-\log(m+1))=\gamma.
\]
The last equality follows by multiplying the integrand by $1-t^m$,
using the integral formulas for $H_m$ and $\log(m+1)$, and applying
monotone convergence.  All displayed positive measures have positive
density on every open subinterval of $(0,1)$.

Define
\begin{equation}\label{pm:ht:gram}
 K_N=(c_{i+j}-c_{i+j+1})_{0\leq i,j<N},\qquad
 \boldsymbol c_N=(c_0,\ldots,c_{N-1})^{\mathsf T},\qquad
 \widehat\gamma_N=\boldsymbol c_N^{\mathsf T}K_N^{-1}\boldsymbol c_N.
\end{equation}

\begin{theorem}[Positive remainder and nonsingularity]\label{pm:ht:positive}
For every $N\geq1$, $K_N$ is positive definite, $\widehat\gamma_N\in\Q$, and
\begin{equation}\label{pm:ht:variational}
 R_N:=\gamma-\widehat\gamma_N
 =\min_{\substack{\deg P\leq N\\P(1)=1}}
     \int_0^1\frac{P(t)^2}{1-t}\,d\nu(t)>0.
\end{equation}
The sequence $\widehat\gamma_N$ is strictly increasing and converges to $\gamma$.
Moreover,
\begin{equation}\label{pm:ht:error-sharp}
 R_N=\Theta(N^{-2}\log^{-2}N),\qquad
 0<R_N<\frac1{9N(N+1)}\quad(N\geq1).
\end{equation}
There is no exceptional singular case at $N=1$: $K_1=(11/24)$ and
$\widehat\gamma_1=6/11$.
\end{theorem}

\begin{proof}
For a nonzero polynomial $r$ of degree below $N$,
$\int(1-t)r(t)^2d\nu(t)>0$, proving positive definiteness.
Every admissible $P$ is uniquely $1-(1-t)r$.  If $\boldsymbol r$ is
the coefficient vector of $r$, the objective in
\eqref{pm:ht:variational} equals
\[
 \gamma-2\boldsymbol r^{\mathsf T}\boldsymbol c_N
       +\boldsymbol r^{\mathsf T}K_N\boldsymbol r.
\]
Its unique minimum is attained at
$\boldsymbol r=K_N^{-1}\boldsymbol c_N$, which proves the identity
and rationality.  The minimum is positive because $P(1)=1$ and the
measure has full support.  The minimizing polynomial is orthogonal
to every polynomial of degree below $N$ for $\nu$, by the first
variation.  If two successive minima were equal, uniqueness would
give a nonzero polynomial of degree at most $N$ orthogonal to itself
for $\nu$.  Thus the decrease of $R_N$ is strict.

Set $w(t)=v(t)/(1-t)$.  From \eqref{pm:ht:measure},
\begin{equation}\label{pm:ht:weight-bounds}
 0<w(t)<\frac1{\pi^2t},\qquad
 w(t)\geq\frac1{2\{\log^2(2/(1-t))+\pi^2\}}\quad(t\geq1/2).
\end{equation}
The first bound discards the logarithmic square in the denominator
and integrates $s^{-2}$.  For the second, restrict the integral to
$t\leq s\leq(1+t)/2$.  On that interval,
$s^{-2}\geq1$ and
$|\log((1-s)/s)|\leq\log(2/(1-t))$.

Let $J_j$ be the shifted Jacobi polynomials normalized by
$J_j(1)=1$ for the weight $t$ on $[0,1]$.  Rodrigues' formula and
integration by parts give
\[
 J_j(t)=\frac1{j!t}\frac{d^j}{dt^j}
        \{t^{j+1}(t-1)^j\},\qquad
 \int_0^1tJ_j(t)J_k(t)\,dt=\frac{\mathbf1_{j=k}}{2j+2}.
\]
Consequently
$r=\sum_{j=0}^{N-1}(2j+2)J_j/[N(N+1)]$ satisfies $r(1)=1$ and
$\int tr^2=1/[N(N+1)]$.  The trial $P=tr$ in
\eqref{pm:ht:variational} proves
$R_N<1/[\pi^2N(N+1)]<1/[9N(N+1)]$, and in particular convergence.

Here are details of the sharp order, since it determines the
denominator target.  If $t\geq1-\epsilon$, $0<\epsilon<1/2$, then
\[
 w(t)\leq\frac1{t\log^2((1-\epsilon)/\epsilon)}.
\]
Let $r_d$ be the normalized endpoint reproducing kernel of degree
$d$ for Lebesgue measure on $[0,1]$.  Thus $r_d(1)=1$ and
$\int_0^1r_d^2=1/(d+1)^2$.  With $k+d=N$ and $k\geq1$, use
$P=t^kr_d$ and split the integral at $1-\epsilon$ to get
\begin{equation}\label{pm:ht:upper-sharp}
 R_N\leq\frac1{(d+1)^2}
 \left\{\frac1{\log^2((1-\epsilon)/\epsilon)}
       +\frac{(1-\epsilon)^{2k-1}}{\pi^2}\right\}.
\end{equation}
The choices $k=\lfloor N/\log N\rfloor$, $d=N-k$,
$\epsilon=(\log N)^3/N$ for sufficiently large $N$ show
$\limsup N^2\log^2N\,R_N\leq1$.

For the lower bound take $I=[1/2,1-N^{-2}]$, $N\geq2$, of length
$L\geq1/4$.  The second estimate in \eqref{pm:ht:weight-bounds} reduces
the problem to the endpoint evaluation norm for polynomials on $I$.
After mapping $I$ to $[-1,1]$, the evaluation point is
$x_0=1+2/(N^2L)\leq1+8/N^2$.  The ordinary Legendre polynomials
satisfy $|P_j(x_0)|\leq\exp(4j/N)\leq e^4$ for $j\leq N$,
by their integral representation and
$\operatorname{arcosh}(1+y)\leq\sqrt{2y}$.
The reproducing kernel is therefore at most
$L^{-1}\sum_{j=0}^N(2j+1)P_j(x_0)^2\leq9e^8N^2$.
Cauchy--Schwarz for evaluation implies
\begin{equation}\label{pm:ht:lower-sharp}
 R_N\geq\frac1{18e^8N^2\{\log^2(2N^2)+\pi^2\}}\quad(N\geq2).
\end{equation}
Together with \eqref{pm:ht:upper-sharp}, this proves the asserted order.
\end{proof}

\subsection{Degree-dependent integral normalization}

Write $\widehat\gamma_N=p_N/q_N$ in lowest terms with $q_N>0$.
The following integers are the Todd numbers:
\begin{equation}\label{pm:ht:todd}
 T_m=\prod_{p\ {\rm prime}}p^{\lfloor m/(p-1)\rfloor},\qquad
 \mathsf T_N=\operatorname{diag}(T_1,\ldots,T_N),\qquad
 \Delta_N=\operatorname{lcm}(1,\ldots,2N)
       \prod_{\substack{3\leq p\leq2N+1\\p\ {\rm prime}}}p.
\end{equation}
Their classical topological interpretation is discussed in
\cite{pm:BuchstaberVeselov2024}.  Only the explicit product and the
elementary valuation bounds below are used here.

\begin{lemma}[Coefficient denominators]\label{pm:ht:coefficient}
For every prime $p$ and $m\geq0$,
\[
 v_p(\mathfrak b_m)\geq-\left\lfloor\frac{m}{p-1}\right\rfloor,
 \qquad
 v_p(c_j)\geq-\left\lfloor\frac{j+1}{p-1}\right\rfloor-v_p(j+1).
\]
\end{lemma}
\begin{proof}
Expand the reciprocal of $F(z)=\sum_{m\geq0}z^m/(m+1)$.
Every coefficient of $\mathscr B=F^{-1}$ is a finite sum over compositions
$m=m_1+\cdots+m_k$ of products $(-1)^k\prod_i(m_i+1)^{-1}$.
If $p^a\mid m_i+1$, then
$(p-1)a\leq p^a-1\leq m_i$.  Summing these bounds proves the
first inequality.  The second follows from \eqref{pm:ht:moments}.
\end{proof}

\begin{theorem}[Graded integral Gram matrix]\label{pm:ht:integral}
For every $N\geq1$, the matrix and vector
\[
 \mathsf A_N=\Delta_N\mathsf T_NK_N\mathsf T_N\in M_N(\Z),\qquad
 \boldsymbol v_N=\Delta_N\mathsf T_N\boldsymbol c_N\in\Z^N
\]
are integral, $\mathsf A_N$ is positive definite, and
\begin{equation}\label{pm:ht:primitive}
 \widehat\gamma_N=\frac1{\Delta_N}\boldsymbol v_N^{\mathsf T}\mathsf A_N^{-1}\boldsymbol v_N,
 \quad
 q_N=\frac{\mathcal Z_N}{\gcd(\mathcal Y_N,\mathcal Z_N)},
 \quad
 \begin{cases}
 \mathcal Y_N=\boldsymbol v_N^{\mathsf T}\operatorname{adj}(\mathsf A_N)\boldsymbol v_N,\\
 \mathcal Z_N=\Delta_N\det \mathsf A_N>0.
 \end{cases}
\end{equation}
\end{theorem}
\begin{proof}
Put $a=i+1$, $b=j+1$ for a matrix entry.  The relevant moments are
$c_{a+b-2}$ and $c_{a+b-1}$.  For $h=p-1$ and
$m\in\{a+b-1,a+b\}$,
\[
 \left\lfloor\frac m h\right\rfloor
 -\left\lfloor\frac a h\right\rfloor
 -\left\lfloor\frac b h\right\rfloor\leq1.
\]
For odd $p$, the possible carry is cleared by the prime product in
$\Delta_N$, and $v_p(m)$ is cleared by its least common multiple.
For $p=2$ the floor difference is at most zero, so that the least
common multiple suffices.  No prime larger than $2N+1$ occurs.
For $T_ac_{a-1}$ the only remaining denominator exponent is at most
$v_p(a)$, also cleared by $\Delta_N$.  This proves integrality.
Positive definiteness follows by congruence.  Direct inversion
gives the scalar identity, and reduction of the fraction $\mathcal Y_N/\mathcal Z_N$
gives the exact denominator formula.
\end{proof}

\begin{theorem}[Primitive denominator upper bound]\label{pm:ht:height}
For every $N\geq1$,
\begin{equation}\label{pm:ht:finite-height}
 q_N\leq \Delta_N^{N+1}\left(\prod_{j=1}^NT_j\right)^2
       \left(\frac{11}{24}\right)^N4^{-(N-1)^2}.
\end{equation}
In particular,
\begin{equation}\label{pm:ht:asymptotic-height}
 \log q_N\leq N^2\log N+
       \left(\frac52-\log4\right)N^2+o(N^2).
\end{equation}
The coefficient in this upper bound is not an asserted asymptotic
constant for the actual reduced denominator.
\end{theorem}
\begin{proof}
Theorem~\ref{pm:ht:integral} gives
$q_N\leq \mathcal Z_N=\Delta_N^{N+1}(\prod T_j)^2\det K_N$.
The matrix $K_N$ is the Gram matrix of monomials for the measure
$(1-t)d\nu(t)$, of total mass $c_0-c_1=11/24$.
Replace the monomials by the monic shifted Chebyshev polynomials
$1$ and $2^{1-2j}\mathcal T_j(2t-1)$, $1\leq j<N$.
This triangular change has determinant one and preserves the Gram
determinant.  Hadamard's inequality and their uniform norms give
\[
 \det K_N\leq\left(\frac{11}{24}\right)^N
                  \prod_{j=1}^{N-1}2^{2-4j}
       =\left(\frac{11}{24}\right)^N4^{-(N-1)^2}.
\]
This proves \eqref{pm:ht:finite-height}.

We give the arithmetic estimate needed for the second term.
Legendre's formula, with $s_p(m)$ the sum of the base-$p$ digits,
implies
\[
 v_p(T_m/m!)=
 \left\lfloor\frac{s_p(m)}{p-1}\right\rfloor\geq0.
\]
The primes $p\leq\sqrt m+1$ contribute
$O(\sqrt m\log(m+1))$ to $\log(T_m/m!)$, since the displayed
exponent is at most $1+\lfloor\log_pm\rfloor$.
If $p>\sqrt m+1$, then $p^2>m$, and that exponent equals
$\lfloor m/(p-1)\rfloor-\lfloor m/p\rfloor\in\{0,1\}$.
When it is one there is an integer $k$ with
$m/p<k\leq m/(p-1)$, hence $p=\lfloor m/k\rfloor+1$ and
$k\leq\sqrt m+1$.  There are at most $O(\sqrt m)$ such primes.
Consequently
\[
 \log T_m=\log(m!)+O(\sqrt m\log(m+1)),\qquad
 2\sum_{j=1}^N\log T_j
 =N^2\log N-\frac32N^2+O(N^{3/2}\log(N+1)).
\]
Finally the prime number theorem gives
$\log \Delta_N=\psi(2N)+\vartheta(2N+1)-\log2=4N+o(N)$.
Taking logarithms in \eqref{pm:ht:finite-height} proves
\eqref{pm:ht:asymptotic-height}.
\end{proof}

\begin{proposition}[An exact forced common factor]\label{pm:ht:gcd}
Put $T=T_{2N}$ and $W=\Delta_NT$.  Define the uniform normalization
$\mathsf A_N^{\rm u}=WK_N$, $\boldsymbol v_N^{\rm u}=W\boldsymbol c_N$,
and $\mathcal Y_N^{\rm u},\mathcal Z_N^{\rm u}$ by \eqref{pm:ht:primitive}, with $W$
in place of $\Delta_N$.  Then
\[
 \mathcal R_N=\frac{T_{2N}^{N+1}}{(\prod_{j=1}^NT_j)^2}\in\Z_{>0},
 \qquad \mathcal Y_N^{\rm u}=\mathcal R_N\mathcal Y_N,
 \qquad \mathcal Z_N^{\rm u}=\mathcal R_N\mathcal Z_N,
\]
and therefore
\[
 \gcd(\mathcal Y_N^{\rm u},\mathcal Z_N^{\rm u})
 =\mathcal R_N\gcd(\mathcal Y_N,\mathcal Z_N),\qquad
 \log\mathcal R_N=N^2\log N+
       \left(\log4-\frac12\right)N^2+o(N^2).
\]
\end{proposition}
\begin{proof}
Integrality of the uniform normalization follows from
Lemma~\ref{pm:ht:coefficient}.  For each prime $p$, the exponent of
$p$ in $\mathcal R_N$ is nonnegative because
\[
 (N+1)\left\lfloor\frac{2N}{p-1}\right\rfloor
 -2\sum_{j=1}^N\left\lfloor\frac j{p-1}\right\rfloor\geq0.
\]
The determinant scaling proves the equality for $\mathcal Z_N$;
both scalar ratios equal $\widehat\gamma_N$, which proves the equality
for $\mathcal Y_N$ and the gcd identity.  The estimate for $T_m$ in the preceding proof gives
the stated asymptotic expansion.
\end{proof}

\subsection{What a determinant does not determine}

Let $D$ be any common denominator of $c_0,\ldots,c_{2N-1}$, and put
$M=D^2K_N$, $\boldsymbol u=D\boldsymbol c_N$.  The finite group
$G=\Z^N/M\Z^N$ has a nondegenerate pairing
\[
 \lambda([x],[y])=x^{\mathsf T}M^{-1}y\pmod\Z.
\]
Its well-definedness and nondegeneracy follow immediately by
changing representatives and testing against the coordinate vectors.
The exact denominator invariant is
\begin{equation}\label{pm:ht:self-pairing}
 q_N=\operatorname{ord}_{\Q/\Z}\lambda([\boldsymbol u],[\boldsymbol u]),
\end{equation}
since $\boldsymbol u^{\mathsf T}M^{-1}\boldsymbol u=\widehat\gamma_N$.
Neither $|G|$ nor the order of $[\boldsymbol u]$ alone determines
this self-pairing.  For example, $M=5I_2$ and
$\boldsymbol u=(1,2)^{\mathsf T}$ give a class of order five with
integral self-pairing.  Real positive definiteness therefore does
not exclude denominator cancellation.

There is nevertheless a uniform local statement at prime indices.
Write $\Z_p$ for the $p$-adic integers.
\begin{proposition}[Prime entry after primitive reduction]\label{pm:ht:prime}
For every prime $p\geq5$,
\[
 \widehat\gamma_{p-1}\in\Z_p,\qquad p\widehat\gamma_p\equiv4\pmod p,
 \qquad p\nmid q_{p-1},\qquad v_p(q_p)=1.
\]
\end{proposition}
\begin{proof}
Up to degree $2p$, write
\[
 F(z)=\mathscr B(z)^{-1}=F_0(z)+\frac{z^{p-1}}p+\frac{z^{2p-1}}{2p},
 \qquad F_0\in\Z_p[[z]].
\]
Since $3p-3>2p$, reciprocal expansion gives, modulo $z^{2p+1}$,
\[
 \mathscr B=F_0^{-1}-\left(\frac{z^{p-1}}p+
      \frac{z^{2p-1}}{2p}\right)F_0^{-2}
       +\frac{z^{2p-2}}{p^2}F_0^{-3}.
\]
The coefficients $[z]F_0^{-2}=-1$ and
$[z^2]F_0^{-3}=1/2$ imply
\[
 p^2c_{p-1}\equiv-1,\qquad
 p^3c_{2p-1}\equiv-\frac14\pmod p.
\]
The same expansion shows that $c_j$ is integral for $j\leq p-3$,
$v_p(c_{p-2})\geq-1$, and $v_p(c_j)\geq-2$ for $j\leq2p-2$.
Partition
\[
 K_p=\begin{pmatrix}A&v\\v^{\mathsf T}&d\end{pmatrix},
 \qquad\boldsymbol c_p=\binom{b}{\beta},\qquad A=K_{p-1}.
\]
The reduction of $p^2A$ vanishes strictly before its anti-diagonal
and has anti-diagonal entries one.  Thus $A^{-1}\in p^2M_{p-1}(\Z_p)$.
Here $b\in p^{-1}\Z_p^{p-1}$, $v\in p^{-2}\Z_p^{p-1}$,
$\beta=c_{p-1}$, and $d=c_{2p-2}-c_{2p-1}$.  Hence
\[
 b^{\mathsf T}A^{-1}b\in\Z_p,\quad
 p^2(\beta-v^{\mathsf T}A^{-1}b)\equiv-1,\quad
 p^3(d-v^{\mathsf T}A^{-1}v)\equiv\tfrac14\pmod p.
\]
The Schur-complement formula for $\widehat\gamma_p$ proves the result.
\end{proof}

\begin{remark}[The remaining arithmetic requirement]\label{pm:ht:scope}
Theorem~\ref{pm:ht:positive} gives the exact equivalence
\[
 q_N\gamma-p_N\longrightarrow0
 \quad\Longleftrightarrow\quad
 q_N=o(N^2\log^2N),
\]
also along any specified infinite subsequence.  No such denominator
estimate is proved here.  The upper bound
\eqref{pm:ht:asymptotic-height} is far above the required scale and
does not show that the integer forms tend to zero.  Conversely,
Proposition~\ref{pm:ht:prime} gives only a local lower constraint at
different indices; multiplying its prime factors as though they
belonged to a single $q_N$ would be invalid.

There is an elementary collective lower bound.  If
$M_N=\max_{N\leq j\leq2N}q_j$, strict monotonicity gives $N$
successive rational gaps, each at least $M_N^{-2}$.  Since their
sum is smaller than $R_N=O(N^{-2}\log^{-2}N)$,
\[
 \max_{N\leq j\leq2N}q_j=\Omega(N^{3/2}\log N).
\]
This does not exclude exceptional sparse subsequences.  The results
of this section concern rational approximants and their exact
normalization, and establish neither irrationality nor transcendence
of $\gamma$.
\end{remark}


% ---- bridges.tex ----
\section{Companion constructions and the specialization problem}
\label{pm:br:section}

The companion manuscripts \cite{pm:KEIO001,pm:KEIO002} develop constructions
based on exponential integrals, moving levels, and arithmetic Gevrey
series.  This section records the parts needed to interpret the present
arithmetic approximation results.  The claims proved below concern
functions and exact identities; their specialization at Euler's constant
is identified separately.  In particular, no assertion concerning joint
algebraic independence at arbitrarily many sampling points is needed.

\subsection{Conventions and a balanced approach to Euler's constant}

On the positive real axis put
\[
 \operatorname{Ein}(x)=\sum_{k\geq1}\frac{(-1)^{k-1}x^k}{k\,k!},
 \qquad E_1(x)=\int_x^\infty\frac{e^{-t}}t\,dt,
 \qquad \delta=\int_0^\infty\frac{e^{-t}}{1+t}\,dt.
\]
The connection identity, with the principal logarithm, is
\begin{equation}\label{pm:br:connection}
 \operatorname{Ein}(x)=\gamma+\log x+E_1(x),
 \qquad \delta=eE_1(1).
\end{equation}
The derivative of the difference in the first identity is zero, and its
limit at zero is \(\gamma\); the second follows by translating the
integral.  Thus \(\operatorname{Ein}(1)=\gamma+\delta/e\).

The dyadic level in \cite{pm:KEIO001,pm:KEIO002} is
\begin{equation}\label{pm:br:dyadic}
 Y_N=(N+1)\operatorname{Ein}(2^N)
          -N\operatorname{Ein}(2^{N+1}),\qquad N\geq1.
\end{equation}
Its definition uses values of an explicitly specified entire function,
without using a decimal expansion of \(\gamma\).

\begin{proposition}[Dyadic error and polynomial detector]
\label{pm:br:dyadic-detector}
The numbers \(Y_N\) decrease strictly to \(\gamma\).  More precisely,
\[
 \varepsilon_N:=Y_N-\gamma
 =(N+1)E_1(2^N)-NE_1(2^{N+1})>0,
 \qquad
 \varepsilon_N\sim\frac{N+1}{2^N}e^{-2^N}.
\]
For every nonzero \(P\in\mathbb C[X]\), the following limit exists:
\begin{equation}\label{pm:br:slope}
 \lim_{N\to\infty}\frac{-\log|P(Y_N)|}{2^N}
       =\operatorname{ord}_{X=\gamma}P.
\end{equation}
The expression on the left is defined for all sufficiently large \(N\).
Consequently \(\gamma\) is algebraic if and only if there exists a
nonzero \(P\in\mathbb Q[X]\) for which the limit in
\eqref{pm:br:slope} equals one.
\end{proposition}

\begin{proof}
Substitution of \eqref{pm:br:connection} into \eqref{pm:br:dyadic} cancels the
logarithms.  Positivity follows because \(E_1\) is positive and strictly
decreasing.  For \(g(t)=E_1(2^t)\) one has
\[
 g''(t)=(\log2)^2 2^t e^{-2^t}>0,
 \qquad
 \varepsilon_N-\varepsilon_{N+1}
 =(N+1)\{g(N)-2g(N+1)+g(N+2)\}>0.
\]
Integration by parts gives
\(E_1(x)=e^{-x}x^{-1}(1+O(x^{-1}))\), proving the asymptotic formula
and \(-\log\varepsilon_N\sim2^N\).
Write \(P(X)=(X-\gamma)^mR(X)\), where \(R(\gamma)\ne0\).
Then \(P(Y_N)=\varepsilon_N^mR(Y_N)\), which proves
\eqref{pm:br:slope}.  If \(\gamma\) is algebraic, its minimal polynomial
has a simple root there, since the characteristic is zero.  Conversely,
a positive value of the limit implies \(P(\gamma)=0\).
\end{proof}

This is an exact classification of a detector, not a lower bound excluding
its critical slope.  The same proof works whenever a specified sequence
approaches a limit at a known rate.  It does not turn the transcendental
approximants \(Y_N\) into rational approximants of controlled height.
The specialized quadratic detector of \cite{pm:KEIO002} is covered by
\eqref{pm:br:slope}; excluding its critical value still requires arithmetic
information not furnished by the rapid convergence alone.

\subsection{A normalized mixed family and the exact zero-removal bridge}

An \(E\)-function is a holonomic series
\(\sum a_nz^n/n!\) with algebraic coefficients whose conjugates and
common denominators grow at most exponentially.  An arithmetic Gevrey
series of order one is a series \(\sum n!b_nw^n\) with
\(\sum b_nw^n\) a \(G\)-function.  Let \(\mathcal M\) denote the
mixed space of formal sums
\[
 F(z)+\widehat f(1/z),\qquad \widehat f(0)=0,
\]
with these two respective types.  We use direction-zero sums when that
direction is admissible.  Summation is injective on the mixed space,
and differentiation and multiplication by \(z\) preserve it
\cite[Definition~1 and Lemma~1]{pm:FischlerRivoal2024}.

Consider the normalized inverse series
\begin{equation}\label{pm:br:normalized-family}
 \widehat D(z)=\sum_{n\geq0}(-1)^n n!z^{-n-1},\qquad
 \widehat W(z)=e^z\operatorname{Ein}(z)-\widehat D(z),\qquad U(z)=e^z.
\end{equation}
Its Borel transform is \((1+\xi)^{-1}\), so that
\begin{equation}\label{pm:br:laplace}
 D_0(z)=\int_0^\infty\frac{e^{-zt}}{1+t}\,dt=e^zE_1(z),
 \qquad W_0(z)=e^z(\gamma+\log z).
\end{equation}
Both the formal family and its sectorial sum satisfy
\begin{equation}\label{pm:br:system}
 U'=U,\qquad \widehat W'=\widehat W+U/z.
\end{equation}
The inverse tail used in the specific bridge of \cite{pm:KEIO001} is
\(z\widehat D(z)\), rather than \(\widehat D(z)\).  Its sectorial
formula therefore has the extra term \((1-z)e^zE_1(z)\).  The two
normalizations have the same value at \(z=1\), but their functional
identities must not be interchanged.

\begin{lemma}\label{pm:br:linear-values}
The numbers \(1,e,e\operatorname{Ein}(1)\) are linearly independent
over \(\overline{\mathbb Q}\).
\end{lemma}

\begin{proof}
The \(E\)-functions \(1,e^z,e^z\operatorname{Ein}(z)\) form a
differential system over \(\mathbb Q(z)\) with no finite singularity
except zero.  They are linearly independent over \(\mathbb C(z)\).
Indeed, a relation involving the third would imply
\(\operatorname{Ein}(z)=A(z)+B(z)e^{-z}\) for rational functions
\(A,B\).  Differentiation and comparison of the independent functions
\(1,e^{-z}\) give \(A'(z)=1/z\), impossible for a rational function.
The remaining two functions are also independent.  The linear
specialization theorem for \(E\)-functions at the ordinary algebraic
point \(1\) now gives the assertion \cite{pm:Beukers2006}.
\end{proof}

\begin{proposition}[Exact status of the restricted bridge]
\label{pm:br:zero-removal}
For every \(\beta\in\overline{\mathbb Q}\), put
\(\widehat R_\beta=\widehat W-\beta U\).  Then
\[
 (R_\beta)_0(1)=e(\gamma-\beta),\qquad
 \widehat R_\beta\notin(z-1)\mathcal M.
\]
Hence the implication
\begin{equation}\label{pm:br:restricted-zr}
 (R_\beta)_0(1)=0\quad\Longrightarrow\quad
                  \widehat R_\beta\in(z-1)\mathcal M
\end{equation}
holds if and only if \(\gamma\ne\beta\).  Requiring
\eqref{pm:br:restricted-zr} for every algebraic \(\beta\) is equivalent
to the transcendence of \(\gamma\).
\end{proposition}

\begin{proof}
The value identity follows from \eqref{pm:br:laplace}.  If
\(\widehat R_\beta=(z-1)(F(z)+\widehat f(1/z))\), comparison of the
nonnegative and nonpositive powers in
\[
 e^z(\operatorname{Ein}(z)-\beta)-(z-1)F(z)
       =\widehat D(z)+(z-1)\widehat f(1/z)
\]
forces both sides to be an algebraic constant \(c\).  Evaluation at one
gives \(e\operatorname{Ein}(1)-\beta e=c\), contradicting
Lemma~\ref{pm:br:linear-values}.  The consequent of
\eqref{pm:br:restricted-zr} is therefore always false, so the implication
holds exactly when its antecedent is false.
\end{proof}

This recovers the logical content of the bridge theorem in
\cite{pm:KEIO001} with the normalization used in \cite{pm:KEIO002}.  It
does not prove the required zero-removal assertion.  Analytic removability
of a zero at one is distinct from membership of the resulting quotient
in an arithmetic class.

Let \(\mathbf E\) be the ring of values of \(E\)-functions at algebraic
points, and \(\mathbf D\) the ring of directionally summed Gevrey-one
values with the convention of \cite{pm:FischlerRivoal2024}.  The standard
conditional conclusion can be stated without an uncertain closure claim.

\begin{corollary}[Conditional intersection consequence]
\label{pm:br:intersection}
If \(\gamma\) is algebraic, then
\[
 \delta\in(\mathbf E\cap\mathbf D)\setminus\overline{\mathbb Q}.
\]
In particular, \(\mathbf E\cap\mathbf D=\overline{\mathbb Q}\)
implies that \(\gamma\) is transcendental.
\end{corollary}

\begin{proof}
Under the stated assumption, \(e^z(\operatorname{Ein}(z)-\gamma)\)
is an \(E\)-function and its value at one is \(\delta\).
The equality \(D_0(1)=\delta\) gives membership in \(\mathbf D\).
If \(\delta\) were algebraic, the identity
\(e\operatorname{Ein}(1)-\gamma e-\delta=0\) would contradict
Lemma~\ref{pm:br:linear-values}.
\end{proof}

This implication is already contained in
\cite[Theorem~2]{pm:FischlerRivoal2024}; it is included to fix the
normalization.  In particular, writing \(E_1(1)=e^{-1}\delta\) does
not by itself prove \(E_1(1)\in\mathbf D\), since the needed closure
under multiplication by \(e^{-1}\) has not been established here.
Nor should \(\gamma\) be identified with an unframed Stokes jump
multiplier: \(W_0\) changes by \(2\pi i e^z\) around zero, while adding
\(c e^z\) changes its constant term without changing that jump.
The arithmetic study of such connection and Stokes coefficients requires
specified normalizations \cite{pm:FischlerRivoal2016}.

\subsection{Mixed Gamma jets and the cost of polynomial targets}

The all-order family in \cite{pm:KEIO002} permits a precise classification
at the function level.  For \(r\geq1\) set
\[
 \operatorname{Ein}_r(z)=\sum_{n\geq1}\frac{(-1)^{n-1}z^n}{n^r n!},
 \qquad
 \widehat D_r(z)=\sum_{n\geq0}(-1)^n
                  \genfrac{[}{]}{0pt}{}{n+1}{r}z^{-n-1},
\]
where brackets denote unsigned Stirling numbers of the first kind.  Put
\(\widehat W_0=e^z\) and
\(\widehat W_r=e^z\operatorname{Ein}_r(z)-\widehat D_r(z)\).
These are mixed functions.  For example, the normalized Borel transform
of the inverse tail is the \(G\)-function
\[
 \mathcal B\widehat D_r(\xi)
  =\frac{(-1)^{r-1}}{(r-1)!}
         \frac{\log^{r-1}(1+\xi)}{1+\xi},
\]
which has no positive-ray singularity.  Write \(W_r\) for the
direction-zero sum, and \(\mathcal M_0\) for such sectorial mixed sums.

\begin{lemma}[Generating function and arithmetic normal form]
\label{pm:br:normal-form}
Coefficientwise near \(t=0\),
\begin{equation}\label{pm:br:jet-generating}
 \sum_{r\geq0}W_r(z)t^r=e^z\Gamma(1-t)z^t.
\end{equation}
Moreover, as analytic functions on a common sector,
\begin{equation}\label{pm:br:mixed-intersection}
 \mathcal M_0\cap e^z\mathbb C[\log z]
        =\bigoplus_{r\geq0}\overline{\mathbb Q}\,W_r.
\end{equation}
All sums on the right are finite.
\end{lemma}

\begin{proof}
The incomplete Gamma decomposition gives
\[
 \Gamma(1-t)z^t
   =1+\sum_{r\geq1}\operatorname{Ein}_r(z)t^r
                         -t z^t\Gamma(-t,z).
\]
The expansion
\(e^zz^t\Gamma(-t,z)\sim
\sum_{n\geq0}(-1)^n(t+1)_nz^{-n-1}\), together with
\([t^{r-1}](t+1)_n=\genfrac{[}{]}{0pt}{}{n+1}{r}\), proves
\eqref{pm:br:jet-generating}.  Equivalently, one can sum the displayed
Borel transforms before integrating for sufficiently small \(|t|\).

Set \(T=z(\partial_z-1)\).  Then \(TW_0=0\) and
\(TW_r=W_{r-1}\).  To compute its mixed kernel, write
\(\Psi=F+\sum_{n\geq1}a_nz^{-n}\) and suppose \(T\Psi_0=0\).
Injectivity of summation gives \(T\Psi=0\) formally.  Its constant
coefficient is \(-a_1\), and its negative coefficients give
\(-na_n-a_{n+1}=0\).  Thus all \(a_n\) vanish and
\(F'=F\), whence \(F=c e^z\) with \(c\) algebraic.  Applying \(T\)
lowers the degree of a logarithmic polynomial.  Induction on that degree,
followed by subtraction of the corresponding \(W_r\), proves
\eqref{pm:br:mixed-intersection}.  Uniqueness follows from their distinct
logarithmic degrees, whose leading coefficients are \(1/r!\).
\end{proof}

Put \(X(z)=\gamma+\log z\).  The Gamma expansion gives
\[
 e^{-z}W_r(z)=[t^r]\exp\left(X(z)t+
                           \sum_{k\geq2}\frac{\zeta(k)}k t^k\right).
\]
The next proposition identifies the exact extra coefficients introduced
by polynomial targets.

\begin{proposition}[Integral coefficient conversion]
\label{pm:br:polynomial-conversion}
Let \(P(X)=\sum_{k=0}^d a_kX^k\in\mathbb Z[X]\),
\(H=\max_k|a_k|\), and define the formal polynomials
\[
 c_n(Z)=[t^n]\exp\left(-\sum_{k\geq2}\frac{Z_k}k t^k\right),
 \qquad B_r(Z)=\sum_{k=r}^d a_k k!c_{k-r}(Z).
\]
Then \(B_r\in\mathbb Z[Z_2,\ldots,Z_d]\), with weighted degree at
most \(d-r\) when \(\operatorname{wt}(Z_k)=k\), and
\[
 \|B_r\|_1\leq H(d+1)d!,\qquad
 e^zP(X(z))=\sum_{r=0}^d B_r(\zeta(2),\ldots,\zeta(d))W_r(z).
\]
Here \(\|\cdot\|_1\) is the sum of absolute polynomial coefficients.
For every \(P\in\overline{\mathbb Q}[X]\), including zero,
\begin{equation}\label{pm:br:degree-obstruction}
 e^zP(\gamma+\log z)\in\mathcal M_0
                  \quad\Longleftrightarrow\quad \deg P\leq1,
\end{equation}
where zero is included on the right by convention.
\end{proposition}

\begin{proof}
The signed cycle expansion is
\[
 n!c_n(Z)=
 \sum_{\sum k m_k=n}(-1)^{\sum m_k}
       \frac{n!}{\prod_{k\geq2}k^{m_k}m_k!}
                          \prod_{k\geq2}Z_k^{m_k}.
\]
Each unsigned coefficient counts permutations with prescribed cycle
lengths and no fixed points.  It is an integer, and their sum is the
derangement number, at most \(n!\).  Hence the stated integrality,
weighted degree, and bound follow by summing over \(k\).
Multiplying the two formal generating functions with opposite zeta
exponents gives
\(X^k=k!\sum_{r=0}^k c_{k-r}(\zeta)e^{-z}W_r\), proving the identity.

For the last assertion, degree at most one follows from
\(W_0=e^z\), \(W_1=e^zX\).  For degree \(d\geq2\), the unique
expansion has coefficient
\[
 B_{d-2}=(d-2)!\left(a_{d-2}-\binom d2 a_d\zeta(2)\right).
\]
This is transcendental since \(a_d\ne0\) and
\(\zeta(2)=\pi^2/6\) is transcendental.  The normal form
\eqref{pm:br:mixed-intersection} therefore excludes mixed membership.
\end{proof}

For instance, \(e^zX^2=2W_2-\zeta(2)W_0\).
The coefficient bound is for integer polynomials in \emph{formal}
period variables; it is not a height bound for algebraic numbers after
substituting \(\zeta(k)\).  At \(z=1\) the conversion represents
\(eP(\gamma)\), but no nonvanishing theorem for that value has been
proved here.  This obstruction also does not rule out the degree-one
bridge: under any hypothetical algebraic equality \(\gamma=\beta\),
the function \(W_1-\beta W_0\) is already mixed, irrespective of the
degree of \(\beta\).  The unresolved step remains arithmetic
specialization, rather than the existence of a functional encoding.


% ---- targets.tex ----
\section{What an infinite construction must establish}
\label{pm:sec:infinite-target}

The finite arithmetic data above support exact theorems about particular
forms.  Passing to an infinite sequence requires a separate statement
about the primitive targets and their nonzero values.  The following
criteria make that distinction explicit.

\begin{proposition}[A nonzero small-form criterion]\label{pm:tg:irrational}
If there are integers $q_j,p_j$ with
\[
 0<|q_j\gamma-p_j|\longrightarrow0,
\]
then $\gamma$ is irrational.
\end{proposition}
\begin{proof}
If $\gamma=a/b\in\Q$ with $b>0$, every nonzero value of
$q_j\gamma-p_j$ has absolute value at least $1/b$.
\end{proof}

An alternative is to construct two linearly independent integer targets
at each of infinitely many levels, with both error bounds tending to
zero.  If $\gamma$ were rational, both errors would eventually be zero;
their targets would then lie on the same one-dimensional rational line,
contradicting independence.  This version avoids certifying the sign of
each individual form, but still requires control of the second independent
target, not only one short kernel vector.

\begin{theorem}[Fixed positive margin and transcendence]\label{pm:tg:roth}
Suppose there are a fixed $\varepsilon>0$ and primitive integer pairs
$(q_j,p_j)$, with $q_j>0$ and $q_j\to\infty$, such that
\begin{equation}\label{pm:tg:margin}
 0<|q_j\gamma-p_j|<q_j^{-1-\varepsilon}.
\end{equation}
Then $\gamma$ is transcendental.
\end{theorem}
\begin{proof}
Proposition~\ref{pm:tg:irrational} first excludes rationality.  Division by
$q_j$ gives
\[
 0<\left|\gamma-\frac{p_j}{q_j}\right|<q_j^{-2-\varepsilon}.
\]
Roth's theorem \cite{pm:Roth1955} permits only finitely many such reduced
rational approximations to an algebraic irrational number for each fixed
$\varepsilon>0$.  The unbounded denominators give a contradiction.
\end{proof}

The theorem requires only an infinite subsequence of levels.  The constant
$0.03$ in the finite certificate is not special: every fixed positive
margin would suffice.  Conversely, a sequence of exponents decreasing to
$1$ in the linear form, or a fixed excellent exponent at finitely many
levels, does not supply the hypothesis of Theorem~\ref{pm:tg:roth}.
The criterion is sufficient and is not necessary for transcendence.

\subsection{An exact interval formulation of the missing estimate}

Suppose a gamma-free exact procedure produces rational intervals
$(\alpha_j,\beta_j)$ containing $\gamma$.  Put
\[
 m_j=\frac{\alpha_j+\beta_j}2,\qquad
 r_j=\frac{\beta_j-\alpha_j}2.
\]
For a target $(q,p)$ with $q>0$, every point of the interval gives
\begin{equation}\label{pm:tg:interval-budget}
 |q\gamma-p|\le |qm_j-p|+q r_j.
\end{equation}
If $p/q\notin(\alpha_j,\beta_j)$, the strict inclusion of $\gamma$
also proves $q\gamma-p\ne0$.  Therefore one concrete sufficient target
for the present method is to prove, at infinitely many levels, that
primitive pairs with $q_j\to\infty$ satisfy
\begin{equation}\label{pm:tg:remaining}
 \frac{p_j}{q_j}\notin(\alpha_j,\beta_j),\qquad
 |q_jm_j-p_j|+q_jr_j<q_j^{-1-\varepsilon}
\end{equation}
for a fixed $\varepsilon>0$.  The first condition may instead be replaced
by any independent positive-integral certificate.

Since the terms in the right side of \eqref{pm:tg:interval-budget} are
nonnegative, meeting this particular budget necessarily requires
\[
 r_j<q_j^{-2-\varepsilon},\qquad
 |m_j-p_j/q_j|<q_j^{-2-\varepsilon}.
\]
Thus small interval width alone is insufficient: the arithmetic position
of its midpoint and the primitive denominator chosen near it matter at
the same time.  These necessary inequalities refer to this interval
budget, not to every possible proof of transcendence.

\begin{problem}[Uniform primitive control]\label{pm:tg:problem}
For the Sondow interval models, possibly augmented by positive moving
differences or stronger moment constraints, construct an infinite
sequence of primitive targets satisfying \eqref{pm:tg:remaining}, or prove
another nonvanishing estimate and the bound \eqref{pm:tg:margin}.
\end{problem}

No solution of Problem~\ref{pm:tg:problem} is given here.  Neither a
large-dimensional kernel, a scalar denominator bound, nor numerical
improvement in a fixed window proves it.  In particular, an exact relation
in $\ker B\cap\ker\Phi$ can be arbitrarily inexpensive while carrying
no information about $\gamma$.

\subsection{Scope of the methodological tests}

The construction rule, effective target dimension, error-to-height ratio,
and nonvanishing certificate are useful tests for an explicit small-form
method.  They should not be promoted to universal restrictions on all
transcendence arguments.  A gamma-independent construction rule prevents
an unsupported numerical selection from replacing a proof, but access to
a rational enclosing interval is not by itself circular.  Additional
variables may legitimately occur if they are eventually eliminated or
treated by a theorem whose hypotheses control their joint arithmetic.
The distinction is whether that elimination or arithmetic theorem has
actually been proved.

Likewise, Dirichlet's approximation theorem is an existence guarantee,
not a universal upper bound on how good approximations can be for arbitrary
real numbers.  It cannot supply an unconditional ceiling for an exponent
of an unknown transcendental constant.  Nor does a large coefficient in
one representation of an algebraic approximant prove a lower bound for
its minimal-polynomial height.  These logical distinctions exclude several
premature impossibility conclusions as well as premature positive proofs.

\section{Reproducibility and the theorem boundary}\label{pm:sec:reproduction}

The original package \texttt{KEIO\_Gamma\_Paper\_v1\_Complete.zip},
retained in the cumulative archive, contains the exact certificates and
Python verifiers for the finite statements of this part.  The
main finite claims can be checked without evaluating $\gamma$, by running
the following commands from its unpacked \texttt{certificates} directory:
\begin{verbatim}
python3 gamma_sondow_intervals.py --verify
python3 gamma_sondow_damping.py --verify
python3 gamma_sondow_relations.py
python3 gamma_sondow_positivity.py
python3 gamma_sondow_threshold_verify.py
\end{verbatim}
The first two check rational primal-dual equalities and inequalities,
the endpoint identities, exact target common factors, and exponent
comparisons by integer powers.  The next two verify prime-logarithm
cancellation and polynomial-positivity certificates.  The last checks
the elementary finite constants accompanying the general kernel estimates.
The all-index theorems rest on the written proofs, not on these finite
checks.  Package versions and execution results are recorded with the
archive; some JSON entries exceed Python's default decimal-digit limit,
which the supplied scripts adjust explicitly.



The results organize an explicit interaction between analytic positivity,
prime arithmetic, and primitive lattice optimization.  The remaining
question is quantitative and arithmetic: whether the surviving primitive
targets can be controlled along an unbounded sequence while the certified
error retains a fixed improvement over the Roth threshold.


% ---- bibliography.tex ----

% END INTEGRATED PART II

\clearpage
\part{Additional early constructions and arithmetic audits}\label{master:early}
% BEGIN INTEGRATED PART III
\section{Early constructions not contained in the first gamma paper}
\label{ec:scope}

The first gamma paper concentrated on Sondow windows, positive moments,
primitive quotient optimization and arithmetic transfer.  The research
records of 6 September also contain several different constructions.  This
part retains their principal mathematical content.  A statement called a
source calculation below is retained with its exact proof-file reference;
it is not silently promoted to a new theorem of this edition.  The complete
source files, including their finite certificates and internal reviews,
are included in the accompanying source archive.

Throughout this part, $H_0=0$, $H_n=\sum_{j=1}^n j^{-1}$ and
$L_n=\operatorname{lcm}(1,\ldots,n)$, with $L_0=1$.  We use
$\delta=\int_0^\infty e^{-t}(1+t)^{-1}\,dt$ and
$\Ein(z)=\sum_{k\ge1}(-1)^{k+1}z^k/(k\,k!)$.
Thus $\Ein(1)=\gamma+e^{-1}\delta$.  Integral positivity, a nonzero
coefficient vector, a nonzero determinant and a nonzero value at $\gamma$
are four different statements.

\section{A full-rank logarithmic Laguerre compression}
\label{ec:laguerre}

For $N\ge1$ define the rational symmetric matrix
\[
 M_N=(m_{ij})_{0\le i,j<N},\qquad
 m_{ii}=H_i,\qquad m_{ij}=-\frac1{|i-j|}\quad(i\ne j).
\]
The standard Laguerre polynomials are normalized by
$\int_0^\infty e^{-t}L_i(t)L_j(t)\,dt=\delta_{ij}$.

\begin{proposition}[The logarithmic compression]
\label{ec:laguerre-identity}
For every $i,j\ge0$,
\[
 \int_0^\infty e^{-t}L_i(t)L_j(t)\log t\,dt
 =m_{ij}-\gamma\delta_{ij}.
\]
Every $M_N$ has simple real eigenvalues, and the eigenvalues of
$M_N$ and $M_{N+1}$ strictly interlace.
\end{proposition}
\begin{proof}
Differentiate the generalized Laguerre orthogonality relation
\[
 \int_0^\infty t^s e^{-t}L_i^{(s)}(t)L_j^{(s)}(t)\,dt
 =\frac{\Gamma(i+s+1)}{i!}\delta_{ij}
\]
at $s=0$, using
$\partial_sL_n^{(s)}|_{s=0}=\sum_{k<n}L_k/(n-k)$.
The diagonal gives $H_i-\gamma$, and an off-diagonal entry gives
$-1/|i-j|$.

Let $x$ be an eigenvalue and let $p$ be the corresponding polynomial of
degree at most $N-1$.  Set $a=e^{x-\gamma}$.  Then
\[
 \int_0^\infty e^{-t}\log(t/a)p(t)q(t)\,dt=0
 \qquad(\deg q<N).
\]
If $\deg p\le N-2$, the admissible choice $q=(t-a)p$ gives a strictly
positive integral, a contradiction.  Hence every eigenvector has degree
$N-1$.  An eigenspace of dimension at least two would contain a nonzero
combination of smaller degree, so the eigenvalues are simple.
If two consecutive compressions shared an eigenvalue, the equality case
of symmetric-matrix interlacing would give an eigenvector of the larger
compression supported in the smaller space.  Its degree would be at most
$N-1$, contradicting the same argument for the size $N+1$ compression.
Cauchy's interlacing theorem now gives strict interlacing.
\end{proof}

Let $P_N$ be the positive-leading primitive integer multiple of
$\det(XI_N-M_N)$.  For example,
\[
 P_1=X,\quad P_2=X^2-X-1,\quad
 P_3=4X^3-10X^2-3X+11.
\]
Clearing matrix entries by $L_{N-1}$ and bounding the row sums gives
\[
 \log H(P_N)=O(N^2).
\]
For $S_N=P_N^2+P_{N+1}^2$, strict interlacing implies
$S_N(x)>0$ for every real $x$, while
$\deg S_N=2N+2$ and $\log H(S_N)=O(N^2)$.

\begin{corollary}[An unfulfilled determinant criterion]
\label{ec:laguerre-criterion}
If
\[
 \limsup_{N\to\infty}\frac{-\log S_N(\gamma)}{N^2}=+\infty,
\]
then $\gamma$ is transcendental.
\end{corollary}
\begin{proof}
For a fixed algebraic number $\alpha$, the nonzero resultant of its
primitive minimal polynomial with $S_N$ gives
$|S_N(\alpha)|\ge\exp(-C_\alpha N^2)$.  Nonvanishing is automatic here
because $S_N$ is positive on the real line.  This contradicts the stated
limsup.
\end{proof}
The small-value hypothesis is not proved.  The finite tests in
\texttt{gamma\_spectral\_attempt.md} instead give large integer values
in their tested range.  This construction changes the rank-one occurrence
of $\gamma$ into a full-rank shift; it does not solve the height comparison.

\section{Exact logarithmic denominators and the limits of a fixed transform}
\label{ec:logden}

\begin{theorem}[Prime residues and cumulative denominators]
\label{ec:logden-theorem}
Let $K$ be a number field, $a\in K$ and
$P(z,Y)\in K[z,Y]$ have $Y$-degree $d\ge1$.  Write
\[
 P(z,a+\log(1-z))=\sum_{n\ge0}c_nz^n.
\]
If $\Delta_N$ is the least positive integer such that
$\Delta_Nc_0,\ldots,\Delta_Nc_N$ are algebraic integers, then
\[
 \log\Delta_N=H_dN+o(N).
\]
More explicitly, write the first nonzero coefficient of
$T_r(z)=\partial_Y^rP(z,a+\log(1-z))/r!$ as
$\beta_rz^{k_r}$.  For every sufficiently large prime $p$ outside a fixed
finite exceptional set, and every prime ideal $\mathfrak p$ above $p$,
\[
 p^rc_{rp+k_r}\equiv(-1)^r\beta_r\pmod{\mathfrak p}
 \qquad(1\le r\le d).
\]
In particular the local denominator exponent is exactly $r$.
\end{theorem}
\begin{proof}
The logarithm is transcendental over $K(z)$, so each $T_r$ is nonzero.
Remove the finitely many primes at which the coefficients or the
$\beta_r$ are nonintegral or at which a $\beta_r$ is not a unit.
Choose $p>k_r$ with $rp+k_r<p^2$.  In a contribution to the coefficient
of degree $rp+k_r$, at most $r$ logarithmic factors can have a denominator
divisible by $p$.  After multiplying by $p^r$, only contributions with
exactly $r$ such factors survive modulo $\mathfrak p$.  Each must be
$-z^p/p$: replacing one by a term of degree $2p$ would exceed the degree
budget.  Their remaining coefficient is exactly $\beta_r$ by Taylor's
formula in $Y$.  This proves the congruence and the exact valuation.

Set $K_0=\max k_r$.  The congruences force
\[
 \log\Delta_N\ge\sum_{r=1}^d\vartheta((N-K_0)/r)-O(1),
 \qquad \vartheta(x)=\sum_{p\le x}\log p.
\]
For $p>\sqrt N$ a term of degree at most $N$ contains at most
$\min(d,\lfloor N/p\rfloor)$ factors with a $p$ in their denominator.
For $p\le\sqrt N$ the bound $d\lfloor\log_pN\rfloor$ suffices, apart
from a fixed denominator contribution.  These small primes contribute
$o(N)$.  Thus
\[
 \log\Delta_N\le\sum_{r=1}^d\vartheta(N/r)+o(N).
\]
The prime number theorem proves the result.
\end{proof}

\begin{corollary}[A delimited failure of global boundedness]
\label{ec:global-bound}
If $P(z,Y)\in\Q(z)[Y]$ has positive $Y$-degree and $a$ is algebraic,
then a Taylor germ of $P(z,a+\log(1-z))$, after removal of a fixed pole
at zero, is not globally bounded.  Here globally bounded means that a
fixed nonzero algebraic change $z\mapsto cz$, followed by multiplication
by a fixed nonzero algebraic constant, makes every coefficient integral.
\end{corollary}
\begin{proof}
Multiply by a fixed polynomial to clear the rational functions in $z$.
Global boundedness would be preserved.  Fixed algebraic scalings affect
only finitely many finite places, whereas the theorem forces a denominator
at every sufficiently large good prime.
\end{proof}

For the harmonic bridge
\[
 B_a(z)=\frac{-a-\log(1-z)}{1-z},
\]
every fixed differential polynomial with coefficients in $\Q(z)$ is a
polynomial in $a+\log(1-z)$ with rational-function coefficients.  Such a
transform can become globally bounded only after losing that variable,
and therefore losing the shift $a$.  For example,
\[
 ((1-z)\partial_z-1)B_a=\frac1{1-z}.
\]
The exact constant $H_d$ is asserted for polynomial coefficients in $z$;
rational coefficients may add geometric denominators at finitely many
primes.  Growing-degree transforms and other coupled functions are outside
this obstruction.

\section{The Vacca generator and a rational Markov realization}
\label{ec:vacca}

Put $\ell(n)=1+\lfloor\log_2n\rfloor$ and
\[
 A(z)=\sum_{n\ge1}\frac{\ell(n)}{2n(2n+1)}z^n,
 \qquad f(z)=\sum_{j\ge0}z^{2^j}.
\]
The classical Vacca--Sondow identity gives $A(1)=\gamma$, and direct
coefficient comparison gives
\[
 2\theta(2\theta+1)A(z)=\frac{f(z)}{1-z},\qquad
 f(z)=z+f(z^2),\qquad \theta=z\frac{d}{dz}.
\]

\begin{proposition}[Loss of the Mahler class under this integration]
\label{ec:not-mahler}
The series $A(z)$ is not a rational-coefficient $k$-Mahler function for
any integer $k\ge2$.
\end{proposition}
\begin{proof}
Suppose a rational series $C(z)$ satisfies
\[
 P_0(z)C(z)+\sum_{i=1}^qP_i(z)C(z^{k^i})=0,
 \qquad P_i\in\Z[z],\quad P_0\ne0.
\]
Let $cz^v$ be the first nonzero term of $P_0$.  For all sufficiently
large $n$, comparison of the coefficient of $z^{n+v}$ expresses the
$n$th coefficient of $C$ as $1/c$ times an integer combination of
earlier coefficients.  Consequently only finitely many primes occur in
the coefficient denominators.
For every odd prime $p$, the coefficient of $A$ at $n=(p-1)/2$ has
$p$-valuation $-1$, because $0<\ell(n)<p$.  This is a contradiction.
\end{proof}

Set $a_n=\ell(n)/(2n(2n+1))$.  A birth-and-death chain on the positive
integers is specified by
\[
 P_{n,n+1}=\frac{a_{n+1}}{4a_n},\qquad
 P_{n,n-1}=\frac14\quad(n\ge2),
\]
and the remaining row mass is assigned to $P_{n,n}$.  Since
$a_{n+1}/a_n<2$, these are nonnegative rational transition probabilities.
Detailed balance gives the stationary probability
$\pi_n=a_n/\gamma$.  The chain is irreducible and positive recurrent, so
the mean positive return time to state $1$ is
\[
 \mathbb E_1\tau_1^+=\frac1{\pi_1}=6\gamma.
\]
Its spectral gap is zero.  Indeed, for the indicator of the tail
$\{n\ge N\}$, the Dirichlet energy is
$O(\log N/N^2)$ and the variance is comparable to $\log N/N$.
Thus the Rayleigh quotient tends to zero.  All transition probabilities
are rational; no arithmetic theorem about this infinite-chain return time
has been established here.

\section{The full logarithm-free floor-integral space}
\label{ec:floor}

For $m\ge2$ put $Q_m(u)=\prod_{j=1}^m(u+j)$ and let
\[
 V_m=\{P/Q_m^2:P\in\Q[u],\ \deg P\le2m-3\}.
\]
Write a member $g\in V_m$ in partial fractions as
\[
 g(u)=\sum_{j=1}^m\left(\frac{\alpha_j}{u+j}
                         +\frac{\beta_j}{(u+j)^2}\right).
\]
The decay at infinity is equivalent to
$\sum\alpha_j=0$ and $\sum\beta_j=\sum j\alpha_j$.  For
\[
 A=\sum\beta_j,\qquad B=\sum\beta_jH_{j-1},\qquad
 \lambda_p=\sum\alpha_jv_p((j-1)!),
\]
integration over successive unit intervals yields
\[
 \mathcal T(g):=\int_0^\infty(\lfloor u\rfloor+1)g(u)\,du
 =A\gamma-B+\sum_{p\le m-1}\lambda_p\log p.
\]
Absolute convergence follows from $g(u)=O(u^{-3})$.  The two decay
identities cancel the divergent terms and the common $\log(2\pi)$ term
in the partial-fraction calculation.

\begin{theorem}[Exact dimension and target freedom]
\label{ec:floor-dimension}
Let $W_m\subset V_m$ be the subspace where every $\lambda_p$ vanishes.
Then
\[
 \dim W_m=2m-2-\pi(m-1).
\]
The target map $W_m\to\Q^2$, $g\mapsto(A,B)$ is onto, and its kernel
has dimension $2m-4-\pi(m-1)$.
\end{theorem}
\begin{proof}
Put $T_j=\sum_{r>j}\alpha_r$ for $1\le j<m$.  These coordinates are
free and
\[
 \sum_r\alpha_r\log((r-1)!)=\sum_{j=1}^{m-1}T_j\log j,
 \qquad A=\sum_{j=1}^{m-1}T_j.
\]
The prime-exponent matrix of the numbers $1,\ldots,m-1$ has full row
rank: its columns at prime indices form an identity matrix.  It imposes
$\pi(m-1)$ independent conditions.  Choosing only $T_1$ nonzero gives
arbitrary $A$ without introducing logarithms.  With that $A$ fixed,
opposite changes of $\beta_1$ and $\beta_2$ vary $B$ arbitrarily, since
$H_0=0$ and $H_1=1$.  The dimension statements follow.
\end{proof}

This theorem concerns the whole numerator space, not the more restricted
Sondow numerator.  Its surjectivity is a reason to keep the primitive-height
question explicit: the existence of a target pair is no smallness theorem.

\begin{theorem}[Completeness for strict rational linear inequalities]
\label{ec:positive-complete}
For $A,B\in\Q$, the inequality $A\gamma-B>0$ holds if and only if some
$m$ and some nonzero $g=P/Q_m^2\in W_m$ satisfy
$P(u)\ge0$ for $u\ge0$ and have target pair $(A,B)$.
\end{theorem}
\begin{proof}
One direction is immediate from the positive integral.  For the converse,
use
\[
 \mathcal K[R](u)=\int_0^1x^u(-\log x)(1-x)R(x)\,dx.
\]
For a polynomial $R$ this is a rational function with double poles at
negative integers.  Direct calculation and Tonelli's theorem give
\[
 \mathcal K[x^{j-1}]=\frac1{(u+j)^2}-\frac1{(u+j+1)^2},
 \qquad \mathcal T(\mathcal K[R])=\int_0^1R(x)\,dx.
\]
Let
\[
 F(x)=\frac1{1-x}+\frac1{\log x}
     =\sum_{n\ge0}a_n(1-x)^n,
\]
where the positive rational Gregory coefficients satisfy
\[
 a_n=\frac1{(n+1)!}\int_0^1t\prod_{j=1}^n(j-t)\,dt,
 \qquad \sum_{n\ge0}a_n=1.
\]
Then $\mathcal K[F]=((u+1)^2(u+2))^{-1}$ has integral $\gamma$.
Set $L_N(x)=\sum_{n<N}a_n(1-x)^n$, $S_N=\int_0^1L_N$,
$E_N=\mathcal K[F-L_N]$, and $b_N=\mathcal K[(1-x)^N]$.
Both $E_N$ and $b_N$ are positive rational functions.  The ratio
$E_N(u)/b_N(u)$ is the average of the decreasing function
$J_N(x)=\sum_{k\ge0}a_{N+k}(1-x)^k$ under a density proportional to
$x^u(-\log x)(1-x)^{N+1}$.  Its derivative is
$\operatorname{Cov}(\log x,J_N(x))\le0$.  Consequently
\[
 E_N^+(u):=\delta_Nb_N(u)-E_N(u)\ge0,
 \qquad \delta_N=E_N(0)/b_N(0)\in\Q.
\]
Here $E_N(0)$ is rational because both terms defining $E_N$ are rational
functions.  Also $0\le J_N\le\sum_{n\ge N}a_n\to0$, so
$\delta_N\to0$.  Therefore
\[
 S_N<\gamma<U_N:=S_N+\frac{\delta_N}{N+1},
 \qquad S_N,U_N\longrightarrow\gamma.
\]
Let $h_1=(u+1)^{-2}-(u+2)^{-2}$, so $h_1>0$ and
$\mathcal T(h_1)=1$.  If $A>0$, choose $N$ with $AS_N>B$ and put
$g=AE_N+(AS_N-B)h_1$.  If $A<0$, choose $N$ with $AU_N>B$ and put
$g=(-A)E_N^++(AU_N-B)h_1$.  If $A=0$, take $g=-Bh_1$.
These are nonnegative rational functions with the required integral.
Putting their finitely many poles over one $Q_m^2$ gives the asserted
numerator.  All their logarithmic coefficients vanish by construction.
\end{proof}

Completeness here is a terminating certificate search for a true strict
inequality.  It cannot reject equality.  The quantitative interval bound,
proved in \texttt{gamma\_bridge\_cones.md}, is
\[
 U_N-S_N<\frac{N+2}{(N+1)^2H_N^2},\qquad
 U_N-S_N\sim\frac1{N\log^2N}.
\]
This is an explicit family of feasible intervals, not the optimal rate of
the full positive cone.  The first upper endpoint is $U_1=53/88$.

\section{Gamma filters: positivity, integer sampling and algebraic elimination}
\label{ec:filters}

For $N\ge2$ and $q\ge1$ define
\[
 \mathcal E_{N,q}=\gamma+
 q\sum_{j=1}^N\frac{(-1)^{j+1}}j\binom Nj
       \log\Gamma(1+j/q).
\]
The binomial finite-difference identities remove the powers of degrees
$2,\ldots,N$ in the local Gamma expansion and retain its linear term.
More usefully, the Gamma integral and the beta identity give the positive
remainder
\begin{equation}\label{ec:beta-tail}
 \mathcal E_{N,q}
 =q\sum_{m\ge1}\int_{mq}^\infty
       \frac{N!\,du}{u^{N+2}\prod_{j=1}^N(1+j/u)}>0.
\end{equation}
One way to verify this formula is to differentiate in $1/q$, apply
$\sum_{j=0}^N(-1)^j\binom Nj/(u+j)=N!/[u(u+1)\cdots(u+N)]$,
and match the vanishing limit at $1/q=0$; the absolutely convergent
Gamma integral justifies the interchange.

\begin{theorem}[All-parameter error and a unit gamma coefficient]
\label{ec:filter-error}
With $S=N(N+1)/2$ and
$A_{N,q}=N!\zeta(N+1)/((N+1)q^N)$, one has
\[
 e^{-S/q}A_{N,q}\le\mathcal E_{N,q}\le A_{N,q}.
\]
For $q=L_N$, every Gamma logarithm has an integer coefficient and
\[
 \mathcal E_{N,L_N}\sim\frac{N!}{(N+1)L_N^N},\qquad
 \frac{\log\mathcal E_{N,L_N}}{N^2}\longrightarrow-1.
\]
\end{theorem}
\begin{proof}
In \eqref{ec:beta-tail}, $u\ge q$ and
$1\le\prod_{j=1}^N(1+j/u)\le e^{S/q}$.  Integrating the comparison
power $N!u^{-N-2}$ and summing over $m$ gives the two bounds.
The prime number theorem gives $\log L_N=N+o(N)$ and
$N^2/L_N\to0$; Stirling's formula proves the limit.  Integer coefficients
follow from $j\mid L_N$.
\end{proof}

\begin{proposition}[A necessary sampling denominator]
\label{ec:lcm-necessity}
Suppose $q,M$ are positive integers, $c_i,n_i\in\Z$, and
\[
 \sum_i c_i(n_i/q)=M,\qquad
 \sum_i c_i(n_i/q)^r=0\quad(2\le r\le N).
\]
Then $L_N\mid Mq$.
\end{proposition}
\begin{proof}
The integer-valued polynomial $\binom{x}{r}$ has zero constant term
and linear coefficient $(-1)^{r-1}/r$.  Hence the assumptions give
\[
 \sum_i c_i\binom{n_i}{r}=(-1)^{r-1}Mq/r\in\Z.
\]
Thus every $1\le r\le N$ divides $Mq$.
\end{proof}
This applies to unequal nodes and arbitrarily many nodes; at coefficient
$M=1$, the denominator $L_N$ is minimal within this exact moment filter.
It is not a universal denominator theorem for all gamma approximants.

An earlier odd-cumulant filter has the distinct form
\[
 A_r=\sum_{i=1}^r
 \frac{2r(-1)^{i+1}(r!)^2}{i(r-i)!(r+i)!}
 \log\frac{\Gamma(1-i/(2r))}{\Gamma(1+i/(2r))}.
\]
The source calculation in \texttt{gamma60\_probability.md}, using the
partial-fraction expansion at the nodes $i^2/(4r^2)$, proves
\[
 0<(-1)^{r+1}(A_r-\gamma)
 <\frac{(r!)^3}{r(3r)!}\sim\frac{2\pi}{\sqrt3}27^{-r}.
\]
This is an alternating family of upper and lower bounds; monotonicity of
the two subfamilies is not asserted.  The surviving Gamma logarithms
cannot be assumed to be logarithms of algebraic numbers.

\subsection{The full coefficient denominator and uniform phase law}
\label{ec:phase}
For integer $q\ge1$, the least positive integer making every coefficient
$q(-1)^{j+1}\binom Nj/j$ integral is exactly
\[
 D_{N,q}=\frac{L_N}{\gcd(L_N,q)}.
\]
To see the lower bound, for a prime $p\le N$ let $p^e$ be its largest
power at most $N$, and write $N=bp^e+c$, $1\le b<p$, $0\le c<p^e$.
The congruence $(1+X)^N\equiv(1+X^{p^e})^b(1+X)^c\pmod p$
gives $\binom N{p^e}\equiv b\pmod p$.  The coefficient at $j=p^e$
therefore attains the required denominator exponent.  The upper bound
is immediate.  In particular $D_{N,N}=L_N/N$, not $L_N$.

Put $\lambda_N(u)=2/u+\sum_{j=1}^N1/(u+j)$ and
$\epsilon_N(q)=1/(q\lambda_N(q))$.  The stronger uniform estimate
from \texttt{gamma\_bridge\_phase.md} is
\[
 \frac{B(q,N+1)}{\lambda_N(q)}\le\mathcal E_{N,q}
 \le\frac{1+1/N}{1-\epsilon_N(q)}
       \frac{B(q,N+1)}{\lambda_N(q)},
\]
where $B(u,N+1)=N!/[u(u+1)\cdots(u+N)]$.  Here is the short argument
retaining its uniform scope.  For $f_N(u)=B(u,N+1)/u$,
$-f_N'/f_N=\lambda_N$, and
$-\lambda_N'\le\lambda_N/u$.  Integration by parts gives
\[
 \frac{f_N(v)}{\lambda_N(v)}\le\int_v^\infty f_N(u)\,du
 \le\frac{f_N(v)}{(1-1/(v\lambda_N(v)))\lambda_N(v)}.
\]
Also $u\lambda_N(u)$ increases,
$m\lambda_N(mq)\ge\lambda_N(q)$, and
\[
 \sum_{m\ge2}\frac{B(mq,N+1)}{B(q,N+1)}
 \le\sum_{m\ge2}\frac{B(m,N+1)}{B(1,N+1)}=\frac1N.
\]
Apply these estimates to the beta-tail formula.
Since $\epsilon_N(q)\le1/(1+H_{N+1})$, the equivalent
$\mathcal E_{N,q}\sim B(q,N+1)/\lambda_N(q)$ is uniform over
all $q\ge1$.  In particular
\[
 \mathcal E_{N,q}\sim\frac{\Gamma(q)}{N^q\log N}
 \quad(q\text{ fixed}),
\]
with constant $\Gamma(q)$, not $\Gamma(q+1)$.  If $q/N\to\rho>0$,
then
\[
 -N^{-1}\log\mathcal E_{N,q}\to
 (1+\rho)\log(1+\rho)-\rho\log\rho.
\]
The exact error comparison in Theorem~\ref{ec:filter-error} is more
convenient at $q=L_N$; this uniform phase formula covers the smaller
sampling regimes as well.

There is also a useful formal classical-Gamma obstruction.  At conductor
$q$, let $e_r$ denote $\log\Gamma(r/q)$ modulo logarithms of algebraic
numbers and $\log\pi$, with $e_q=0$, and impose the stated reflection
and multiplication relations.  The integer functional
\[
 \Psi_q(e_r)=2r-q\quad(1\le r<q)
\]
annihilates them: both reflection and each multiplication relation are
checked by summing an arithmetic progression.  For the filter vector
$V_{N,q}=q\sum_{j=1}^N(-1)^{j+1}\binom Nj e_{j\bmod q}/j$ it gives
\[
 \Psi_q(V_{N,q})=q(2-qH_N)\quad(q>N),
\]
and $N(2-NH_N-(-1)^{N+1})$ when $q=N$.  These are nonzero for
$q\ge N$, $q\ge3$.  Thus an individual filter does not lose every
Gamma symbol under the specified classical relations.  The statement
concerns that formal quotient and does not assert completeness of all
relations between actual Gamma values.  Combinations of different
filters, including the following algebraic-valued construction, are not
excluded.

\subsection{Removing the Gamma and pi coordinates}
\label{ec:alpha-family}
There is a different, algebraic-valued family.  For $N\ge2$, let
\[
 r_i=\frac18+\frac{i}{24N},\quad 0\le i<N,\qquad
 I_k(r)=r^k-(1/2-r)^k-(2r)^k,
\]
and solve
\[
 \sum_iw_i=0,\qquad \sum_iw_iI_k(r_i)=-1\quad(2\le k\le N).
\]
Define the positive algebraic number
\[
 \alpha_N=\prod_{i=0}^{N-1}
 \left(\frac{2^{1-2r_i}\cos(\pi r_i)}{1-2r_i}\right)^{w_i}.
\]
The weights are rational; the powers denote their positive real branches.
The source theorem in \texttt{gamma\_next\_relations.md} states and proves
that the system is nonsingular for every $N\ge2$ and
\begin{equation}\label{ec:alpha-bound}
 |\gamma+\log\alpha_N|
 \le150(N+1)^2(90/139)^N.
\end{equation}
For clarity, the all-order mechanism is recorded here.  The operator
\[
 (Tu)(r)=u(r)-u(1/2-r)-u(2r)
\]
maps polynomials of degree at most $N$ onto the codimension-one space
$p(0)=p(1/4)$.  Its kernel consists of multiples of $r-1/2$:
the coefficient of the highest term of degree $k\ge2$ is
$1-(-1)^k-2^k\ne0$, and the degree-one calculation gives that kernel.
For $W_N(z)=\prod_i(z-r_i)$, the source's elementary node-product bound
$|W_N(1/4)/W_N(0)|\le5/7$ proves the constrained interpolation is
nonsingular.  Differentiating $Tu=E$ twice gives
\[
 u''(x)=\tfrac14u''(x/2)-\tfrac14u''((1-x)/2)
        -\tfrac14E''(x/2).
\]
This is a contraction in the second-derivative norm.  Reflection and
duplication for $\Gamma$ then reduce the error to at most
$\tfrac14\|(f-p_N)''\|_{[0,1/2]}$, where
$f(r)=\log(2^{1-2r}\cos(\pi r)/(1-2r))$ and $p_N$ is the constrained
interpolant.  Cauchy's formula on the circle with center $7/48$ and
radius $3/5$, with the source bound $|f|<7$, yields
\eqref{ec:alpha-bound}.  The full contour-distance and constant audit is
retained in that exact proof file.

Under the assumption that $\gamma$ is algebraic,
$\gamma+\log\alpha_N$ is nonzero for every $N$: otherwise the
Hermite--Lindemann theorem would make the positive algebraic number
$\alpha_N=e^{-\gamma}$ transcendental.  This resolves nonvanishing under
the contradiction hypothesis, not the degree-height comparison.  The
source's Cijsouw-based sufficient condition
$D_N^2(D_N+\log H_N)=o(N)$ on an infinite subsequence, where $D_N,H_N$
are the actual degree and primitive polynomial height of $\alpha_N$,
remains unproved.  Large displayed powers do not by themselves determine
those primitive invariants.

\subsection{The E-function replacement and what the constant three quarters measures}
\label{ec:e-replacement}
For rational $a>0$ put
\[
 E_a(z)=\sum_{k\ge0}\frac{a}{a+k}\frac{(-z)^k}{k!},\qquad
 G(a,X)=\frac{X^a}{a}E_a(X)=\int_0^Xt^{a-1}e^{-t}\,dt.
\]
The functions $E_a$ are rational-coefficient E-functions, with
$zE_a'+aE_a=ae^{-z}$.  Take
\[
 q=L_N,\qquad d_j=q(-1)^{j+1}\binom Nj/j,\qquad a_j=1+j/q,
\]
and let $X$ be the least
power of two at least $8(N+1)(\lceil\log_2q\rceil+1)$.  Replacing
$\Gamma(a_j)$ by $G(a_j,X)$ gives a remainder
$\widetilde\varepsilon_N$ satisfying
\[
 \tfrac34\mathcal E_{N,q}<\widetilde\varepsilon_N
 <\tfrac54\mathcal E_{N,q}.
\]
Indeed $0<\log\Gamma(a)-\log G(a,X)\le4(X+1)e^{-X}$ for
$1\le a\le2$, $X\ge4$, and
$\sum|d_j|\le q2^N$; comparison with the lower bound of
Theorem~\ref{ec:filter-error} gives the displayed one-quarter budget.

The identities $\sum d_ja_j=qH_N+1\in\Z$ give a rational prefactor
\[
 C_N=X^{qH_N+1}\prod_j a_j^{-d_j}>0
\]
and E-functions
\[
 P_N(z)=\prod_{j\ \mathrm{odd}}E_{a_j}(Xz)^{d_j},\qquad
 Q_N(z)=\prod_{j\ \mathrm{even}}E_{a_j}(Xz)^{-d_j}
\]
such that
\[
 e^{\widetilde\varepsilon_N}
 =C_Ne^\gamma P_N(1)/Q_N(1).
\]
If $\gamma$ is algebraic, $C_Ne^{\gamma z}P_N(z)-Q_N(z)$ is an
E-function with a positive small value at $1$.  Its defining family,
orders and heights still vary with $N$.

The proof files \texttt{gamma\_next\_rescaling.md} and
\texttt{gamma\_next\_efunction.md} provide two precise cost statements:
\[
 \frac{\log h(C_N)}N\longrightarrow1+\log2,
 \qquad
 \frac{\log\operatorname{ord}P_N}{N^2},\
 \frac{\log\operatorname{ord}Q_N}{N^2}\longrightarrow\frac34.
\]
Here $h$ is the logarithmic height of the displayed reduced rational
prefactor, and $\operatorname{ord}$ is the minimal order of a scalar
linear differential equation over $\C(z)$.  Thus the $3/4$ is a constant
for the logarithm of a differential order, not a transcendence exponent
of $\gamma$ or a denominator exponent for a rational gamma approximation.
The minimum-order input, after the ramification $z=w^q$, is algebraic
independence of the primitives
$\int_0^w v^{j-1}e^{-Xv^q}\,dv$ and their translation orbits.  It gives
\[
 \operatorname{ord}\prod_{j\in J}E_{1+j/q}(Xz)^{n_j}
 =\prod_{j\in J}(n_j+1),\qquad
 J\subseteq\{1,\ldots,q-1\},\quad n_j\ge0.
\]
The exceptional endpoint $j=q$, when present at the first small indices,
is treated separately; it does not affect the asymptotic statement.
These costs do not establish a contradiction with a value lower bound.

\section{Primitive exterior elimination in the Euler--Gompertz construction}
\label{ec:exterior}

Let $s=e^{-1}$ and $E(z)=\sum_{k\ge1}z^k/(k\,k!)$, so
$E(-1)=-\gamma-s\delta$.  For integer rows $r_i=(A_i,B_i,C_i)$, set
\[
 R_i=A_is+B_iE(-1)+C_i
     =s(A_i-B_i\delta)+C_i-B_i\gamma.
\]
Define
\[
 U_{ij}=A_iC_j-A_jC_i,\quad
 V_{ij}=A_iB_j-A_jB_i,\quad
 W_{ij}=C_iB_j-C_jB_i,
\]
and $L_{ij}=U_{ij}-V_{ij}\gamma+W_{ij}\delta$.

\begin{proposition}[The surviving primitive pair]
\label{ec:shared-row}
For $D=\det(r_1,r_2,r_3)$,
\[
 W_{23}L_{12}-W_{12}L_{23}=D(B_2\gamma-C_2).
\]
If $D\ne0$ and $(B_2,C_2)\ne(0,0)$, removal of the full coefficient
content leaves exactly $(q,p)=(B_2/g,C_2/g)$, up to common sign, where
$g=\gcd(|B_2|,|C_2|)$.
\end{proposition}
\begin{proof}
Expansion makes the $\delta$ coefficient zero, the $\gamma$ coefficient
$DB_2$, and the constant coefficient $-DC_2$.  The common content is
therefore exactly $|D|g$.
\end{proof}
The cofactor matrix satisfies
\[
 \operatorname{cof}(\operatorname{cof}M)=\det(M)M
\]
for invertible $3$ by $3$ matrices.  This follows from
$\operatorname{cof}M=\det(M)M^{-T}$.  Consequently repeated full
cofactor formation, followed by primitive normalization, has projective
period two.  Moving the two outside rows while retaining the common row
does not change its surviving rational approximation.  This corrects a
crude product-of-heights estimate; it does not prove the remaining form
small or nonzero.

\subsection{The actual Rivoal rows and an all-degree determinant}
\label{ec:rivoal}
The published Hermite--Pad\'e input is
\[
 R_n(z)=\frac{z^{3n+1}}{n!^2}
 \int_0^1\!\int_0^1e^{zuv}u^{2n}(1-u)^n v^{2n}(1-v)^n\,du\,dv
 =A_n(z)e^z+B_n(z)E(z)+C_n(z),
\]
with $\deg A_n\le n-1$ and $\deg B_n,\deg C_n\le n$.
Writing
\[
 f_{n,k}=\frac{(2n+k)!}{k!^2(n-k)!^2},\qquad
 h_{n,k}=H_{2n+k}-2H_k+2H_{n-k},
\]
integration by parts, or the partial-fraction formulas in
\texttt{gamma\_cont\_rivoal.md}, gives
\[
 B_n(z)=\sum_{k=0}^n(-1)^kf_{n,k}z^{n-k},\qquad
 C_n(z)=\sum_{k=0}^n(-1)^kf_{n,k}h_{n,k}z^{n-k}.
\]
In particular $B_n(-1)=(-1)^n\sum_k f_{n,k}\ne0$; the primitive
shared-row output is the rational number $C_n(-1)/B_n(-1)$.

For completeness, define
\[
 D_n(z)=\det\begin{pmatrix}
 A_n&B_n&C_n\\A_{n+1}&B_{n+1}&C_{n+1}\\A_{n+2}&B_{n+2}&C_{n+2}
 \end{pmatrix}(z).
\]
Put $m=n+1$ and
\[
 \beta_t=(-1)^t(3t)!/(t!^2),\qquad
 d_m=\frac{27m^2+18m+1}{3(m+1)(3m+1)(3m+2)},
\]
\[
 K_n=\frac{(2n)!^2}{(3n+1)!^2},\quad
 M_{0,n}=-\beta_m\beta_{m+1}d_m>0,
\]
\[
 t_n=\frac{(2m-1)^2}{(3m-1)^2}
      -\frac{4(9m^2+2m-2)}{3(27m^2+18m+1)}.
\]

\begin{proposition}[Consecutive-row determinant]
\label{ec:rivoal-det}
For all $n\ge0$,
\[
 D_n(z)=K_nM_{0,n}z^{3n+1}(1+t_nz),\qquad
 (-1)^{n+1}D_n(-1)>0.
\]
\end{proposition}
\begin{proof}
The degree bounds give $\deg D_n\le3n+2$.  Replacing the final column
by the remainders shows $z^{3n+1}\mid D_n$.  Only the first two
Taylor coefficients of $R_n$ and of
$M=A_mB_{m+1}-A_{m+1}B_m$ are needed.  Beta integration gives
\[
 R_n=K_nz^{3n+1}
 \left(1+\frac{(2n+1)^2}{(3n+2)^2}z+O(z^2)\right).
\]
Put $a_t=2H_t-H_{3t}$.  The displayed coefficient formulas and the
vanishing of the first coefficients of $R_t$ give, for $t\ge1$,
\[
 B_t(0)=\beta_t,\quad B_t'(0)=-t\beta_t/3,
 \quad A_t(0)=\beta_ta_t,
\]
\[
 A_t'(0)=\beta_t\left[-(t/3+1)a_t+2t/3-4/9\right].
\]
Since $a_m-a_{m+1}=-d_m$, one obtains
\[
 M(0)=M_{0,n},\qquad
 M'(0)/M(0)=-\frac{4(9m^2+2m-2)}{3(27m^2+18m+1)}.
\]
Multiplication proves the formula.  The subtracted term in $t_n$ is
positive and $(2m-1)/(3m-1)<2/3$, hence $t_n<4/9<1$.
\end{proof}
This determinant nonvanishing holds in every degree.  Nonvanishing of the
evaluated primitive form was certified only at the finitely tested indices
in the source; these two scopes must not be conflated.

\section{Gregory arithmetic: exact cancellation before integration}
\label{ec:gregory}

Let
\[
 F(z)=\frac{-\log(1-z)}z=\sum_{n\ge0}\frac{z^n}{n+1},\quad
 B(z)=F(z)^{-1}=\sum_{n\ge0}b_nz^n,
\]
and let $g_n=-b_n>0$ for $n\ge1$.  Then
\[
 G(z)=\sum_{n\ge1}\frac{g_n}{n}z^n,\qquad
 G'(z)=\frac{1-B(z)}z,\qquad G(1)=\gamma.
\]
The moments used in the first gamma paper are
$c_n=g_{n+1}/(n+1)$.

\begin{theorem}[Todd denominators]
\label{ec:todd}
For every $N\ge0$, the least integer clearing $b_0,\ldots,b_N$ is
\[
 T_N=\prod_{p\le N+1}p^{\lfloor N/(p-1)\rfloor}.
\]
For every prime $p$ and every $k\ge1$,
\[
 v_p(b_n)\ge-\lfloor n/(p-1)\rfloor,\qquad
 p^kb_{k(p-1)}\equiv(-1)^k\pmod p.
\]
Moreover $\log T_N=N\log N+O(N)$.
\end{theorem}
\begin{proof}
For $j\ge1$, $(p-1)v_p(j+1)\le j$, with equality only when $j=p-1$.
Expand $B=(1+(F-1))^{-1}$ geometrically.  Each contribution to $b_n$
is a signed product of $1/(j_i+1)$ for an ordered composition
$\sum j_i=n$.  Its denominator exponent is at most
$\lfloor n/(p-1)\rfloor$.  At $n=k(p-1)$ only the composition with
$k$ entries all equal to $p-1$ attains exponent $k$.  This proves the
congruence and the exact cumulative denominator.  Finally
\[
 \log T_N=\sum_{p\le N+1}\lfloor N/(p-1)\rfloor\log p
 =N\log N+O(N)
\]
by the weighted prime estimate
$\sum_{p\le x}\log p/(p-1)=\log x+O(1)$.
\end{proof}

\begin{theorem}[Reciprocal Hankel cancellation]
\label{ec:reciprocal-hankel}
For every $N\ge1$,
\[
 \det(b_{i+j})_{0\le i,j<N}
 =\frac{(-1)^{N-1}}
 {\displaystyle\prod_{j=0}^{N-2}(2j+3)\binom{2j+2}{j}^{2}}.
\]
The empty product is $1$.  The absolute determinant has $N^2$-root
limit $1/4$; its exact reduced denominator has $N^2$-root limit $4$.
\end{theorem}
\begin{proof}
For a general reciprocal pair $F=\sum f_nz^n$, $B=1/F$, $f_0=1$,
let $T=(b_{i-j})_{0\le j\le i<N-1}$ and
$H=(f_{i+j+2})_{0\le i,j<N-1}$.  Coefficient comparison gives
\[
 (b_{i+j+2}-b_{i+1}b_{j+1})=-THT^{\mathsf T}.
\]
It follows either by multiplying the bivariate generating function
\[
 \frac{(F(x)-1)/x-(F(y)-1)/y}{x-y}
\]
by $-B(x)B(y)$, or directly by convolution.  Taking the Schur complement
of the initial entry $b_0=1$ gives the determinant $(-1)^{N-1}\det H$.
Here $H_{ij}=1/(i+j+3)$; the Cauchy determinant formula gives the displayed
product.  Stirling's formula proves the two limits.
\end{proof}
Thus coefficientwise factorial-scale denominators can collapse in a full
Hankel determinant.  The exact integer after clearing this determinant
is nevertheless always $\pm1$: the capacity and denominator exponents
balance.  This identity applies before the division by $n$ defining $G$,
not automatically to the gamma Pad\'e output.

\subsection{Formal-group Frobenius and an exact Lyapunov equation}
\label{ec:lyapunov}
For $\psi_p(z)=1-(1-z)^p$, the multiplicative formal-group logarithm gives
\[
 B(\psi_p(z))=\frac{\psi_p(z)}{pz}B(z),\qquad
 \psi_p(z)\equiv z^p\pmod p.
\]
This is not the exact Mahler substitution $z\mapsto z^p$ over $\Q$.
Theorem~\ref{ec:todd} also gives the exact $p$-adic radius
\[
 R_p(B)=R_p(G)=p^{-1/(p-1)}.
\]
Division of the $n$th coefficient by $n$ changes its $n$th-root growth
by a factor tending to one.  The coordinate $x=-\log(1-z)$ gives
\[
 B(1-e^{-x})=\frac{1-e^{-x}}x,\qquad
 G(1-e^{-x})=\Ein(x)+\log\frac{1-e^{-x}}x.
\]
The first transformed function is an E-function.  In the second formula,
$\gamma$ still occurs at the limit $x\to+\infty$, not as a regular
algebraic-argument E-value.

Let
\[
 T=(b_{i-j})_{0\le j\le i<N},\quad
 J=(1/(i+j+2))_{0\le i,j<N},\quad
 M=(c_{i+j})_{0\le i,j<N},
\]
and $\Lambda=\operatorname{diag}(1/2,3/2,\ldots,N-1/2)$.

\begin{proposition}[The integrated moment matrix as a Gramian]
\label{ec:lyapunov-theorem}
The matrix $M$ is the unique solution of
\[
 \Lambda M+M\Lambda=TJT^{\mathsf T},
 \qquad
 M=\int_0^\infty e^{-t\Lambda}TJT^{\mathsf T}e^{-t\Lambda}\,dt.
\]
If $D=\operatorname{diag}(0,1,\ldots,N-1)$, $U$ is lower triangular
with every entry on and below the diagonal equal to $1$, and
$X=T^{-1}MT^{-\mathsf T}$, then
\[
 \mathcal A X+X\mathcal A^{\mathsf T}=J,
 \qquad \mathcal A=\Lambda+I-UT.
\]
Both $M$ and $X$ are positive definite.
\end{proposition}
\begin{proof}
The reciprocal-series identity
\[
 (b_{i+j+1})=-T(f_{i+j+1})T^{\mathsf T}
\]
and $(i+j+1)c_{i+j}=-b_{i+j+1}$ give the first equation.
The Lyapunov operator has positive eigenvalues $i+j+1$ on matrix units,
so its inverse is the stated convergent integral.  Positivity follows
because $J$ is positive definite and $T$ is invertible.
Differentiating $B=z/[-\log(1-z)]$ gives
\[
 zB'=B-B^2/(1-z).
\]
Using truncated Toeplitz multiplication and
$[D,T(B)]=T(zB')$ gives
$T^{-1}DT=D+I-UT$.  Congruence of the first equation proves the second.
\end{proof}
The inversion of the factors $i+j+1$ is the exact new arithmetic cost
introduced by integration.  Positive stability alone does not control
the primitive denominator of a Schur-complement value.

For the elementary partial sums $S_N=\sum_{n=1}^Ng_n/n=A_N/Q_N$ in
lowest terms, the same valuation calculation proves
\[
 v_2(Q_{2^k})=2^k+k.
\]
Indeed, at $N=2^k$ the last summand has valuation $-(N+k)$ and every
preceding summand has larger valuation.  The coefficient integral gives
$g_{N+1}\ge1/(6N(N+1))$, hence
\[
 Q_N(\gamma-S_N)\ge\frac{2^N}{6(N+1)^2}\longrightarrow\infty
 \qquad(N=2^k).
\]
This is a failure theorem for these dyadic truncations, not for every
subsequence or for the Pad\'e family.

\section{Moving factorial parameters and the exact target they approach}
\label{ec:moving}

For integers $n\ge1$, $m\ge0$, define
\[
 f_{n,m,k}=\frac{(m+k)!}{k!^2(n-k)!^2},\quad
 h_{n,m,k}=H_{m+k}-2H_k+2H_{n-k},\quad
 r_{n,m}=\frac{\sum_{k=0}^nf_{n,m,k}h_{n,m,k}}
                 {\sum_{k=0}^nf_{n,m,k}}.
\]
Up to common factors, $m=0$, $m=n$ and $m=2n$ correspond respectively
to the Laguerre, Aptekarev and surviving Rivoal families.

\begin{theorem}[Complete convergence classification]
\label{ec:moving-theorem}
For every sequence of nonnegative integers $m_n$,
\[
 r_{n,m_n}\to\gamma\quad\Longleftrightarrow\quad m_n=o(n^2).
\]
If $m_n/n^2\to a\in(0,\infty)$, then
\[
 r_{n,m_n}\longrightarrow
 \gamma+\frac{2K_0(2/\sqrt a)}{I_0(2/\sqrt a)}>\gamma.
\]
If $m_n/n^2\to\infty$, then
\[
 r_{n,m_n}=\log(m_n/n^2)-\gamma
 +O(1/n+n^2/m_n)\longrightarrow+\infty.
\]
\end{theorem}
\begin{proof}
Put $j=n-k$ and normalize by the endpoint term:
\[
 w_j=\frac{(n^{\underline j})^2}
           {(m+n)^{\underline j}(j!)^2},\qquad
 g_j=H_{m+n-j}-2H_{n-j}+2H_j.
\]
Then $r_{n,m}$ is the weighted average of $g_j$ and
$0\le g_j-g_0\le4H_j$.
At quadratic scale $m/n^2\to a>0$, for each fixed $j$,
$w_j\to a^{-j}/(j!)^2$ and $g_j\to\log a-\gamma+2H_j$.
The bounds $w_j\le(2/a)^j/(j!)^2$ and the displayed harmonic inequality
permit dominated convergence.  The classical series
\[
 K_0(2\sqrt x)=-(\tfrac12\log x+\gamma)I_0(2\sqrt x)
               +\sum_{j\ge0}\frac{H_jx^j}{(j!)^2}
\]
gives the stated limit.  Both Bessel functions are positive at positive
real arguments.  If $m/n^2\to\infty$, then
$w_j\le\varepsilon^j/(j!)^2$ with $\varepsilon=n^2/(m+1)\to0$,
so $r_{n,m}=H_{m+n}-2H_n+O(\varepsilon)$, proving the last formula.

Suppose now that $m=o(n^2)$ and set $\rho=n/\sqrt{m+n}\to\infty$.
Compare the normalized $w_j$ with the Bessel distribution
\[
 \mathbb P_\rho(J=j)=\frac{\rho^{2j}}{(j!)^2I_0(2\rho)}.
\]
Their unnormalized ratio is
\[
 R_j=\prod_{s<j}\frac{(1-s/n)^2}{1-s/(m+n)}\le1.
\]
For $j\le2\rho$ it is at least $e^{-8}$ for large $n$.  Since
$\mathbb E_\rho J^2=\rho^2$, this region has probability at least
$3/4$, and the actual distribution $\mathbb Q$ satisfies
$\mathbb Q(J=j)\le(4e^8/3)\mathbb P_\rho(J=j)$.
The Bessel moments
\[
 \mathbb E_\rho J=\rho I_1(2\rho)/I_0(2\rho),\qquad
 \mathbb E_\rho J^4=\rho^4+2\rho^3I_1(2\rho)/I_0(2\rho)+\rho^2
\]
imply $J/\rho\to1$ in probability under both distributions.
Moreover
\[
 \mathbb E_\rho (J+1)/\rho\le2,\qquad
 \mathbb E_\rho\rho/(J+1)=I_1(2\rho)/I_0(2\rho)\le1.
\]
The inequality $(\log y)^2\le C(y+y^{-1})$ gives uniform
integrability, and therefore
$\mathbb E_{\mathbb Q}H_J=\log\rho+\gamma+o(1)$.
For $j\le n/2$, $H_n-H_{n-j}\le2j/n$; for the remaining tail use
$H_n$ and the fourth-moment bound.  Thus both endpoint harmonic shifts
have expectation $o(1)$, and
\[
 r_{n,m}=H_{m+n}-2H_n+2\log\rho+2\gamma+o(1)=\gamma+o(1).
\]
Finally, if $m_n/n^2$ does not tend to zero, some subsequence has a
further ratio limit in $(0,\infty]$.  The already proved alternatives
exclude convergence to $\gamma$ on that subsequence.
\end{proof}

The manifest integer denominator
$Q^0_{n,m}=\sum_k\binom nk^2(m+k)!/m!$ satisfies
\[
 \frac{(m+n)!}{m!}\le Q^0_{n,m}
 \le\frac{(m+n)!}{m!}I_0(2n/\sqrt{m+n}).
\]
This is not a lower bound on the denominator of the reduced fraction
$r_{n,m}$.  The finite source test of $5\le n\le100$, $0\le m\le2n$
contains $10{,}176$ pairs and certifies $|Q\gamma-P|>1$ for every tested
primitive pair.  It does not prove failure at all indices.  The quadratic
movement changes the limiting target; below that scale the primitive
arithmetic problem remains.

\section{Sign variation, normal indices and growing prime windows}
\label{ec:normal-indices}

An additional local result preceding the first paper is not implied by
its first-prime-entry theorem.  Its real component has a useful abstract
form.

\begin{theorem}[A quadratic sign change leaves no three-index gap]
\label{ec:normal-theorem}
Let $\omega$ be a finite positive measure with infinite support on a
compact interval, and let $s_2$ be a real monic quadratic polynomial.
Write $h_j=\int t^js_2(t)\,d\omega(t)$ and
$D_r=\det(h_{i+j})_{0\le i,j<r}$, with $D_0=1$.
Then $(D_r,D_{r+1},D_{r+2})\ne(0,0,0)$ for every $r\ge0$.
Each Hankel matrix has nullity at most two.
\end{theorem}
\begin{proof}
If $D_n\ne0$, let $P_n$ be the monic polynomial orthogonal to degrees
less than $n$ for the signed functional
$\mathcal L(f)=\int fs_2\,d\omega$.  The first $\kappa$ with
$\mathcal L(t^\kappa P_n)\ne0$ is at most $n+2$, because otherwise
\[
 \mathcal L(s_2P_n^2)=\int s_2^2P_n^2\,d\omega>0
\]
would vanish.  Write $\kappa=n+k-1$, $1\le k\le3$.  In the monic
basis $1,\ldots,t^{n-1},P_n,tP_n,\ldots,t^{k-1}P_n$, the cross block
vanishes and the last block has nonzero constant anti-diagonal and zeros on the side $i+j<k-1$.  Therefore
\[
 D_{n+k}=D_n(-1)^{k(k-1)/2}\mathcal L(t^\kappa P_n)^k\ne0,
\]
while the intervening determinants vanish.  Start with $D_0=1$ and
iterate.  For nullity, a kernel of dimension at least three would contain
a nonzero polynomial $f$ of degree at most $r-3$.  Testing against
$s_2f$ contradicts $\int s_2^2f^2\,d\omega>0$.
\end{proof}
The gap bound is sharp in this class: with $d\omega=dt$ on $[0,1]$ and
$s_2=t^2-t+1/6$, one has $D_1=D_2=0$ and $D_3=-1/180^3$.

For the gamma Pad\'e local model, define from the Gregory series $B$
\[
 \lambda_s=\frac{[z^{s+2}]B^2}{s+1},\quad
 \mu_s=-\frac{[z^{s+3}]B^3}{s+1},\quad
 d_s=c_s-8\lambda_s+16\mu_s,\quad h_s=d_s-d_{s+1}.
\]
The source \texttt{gamma\_push\_local.md} identifies this $h_s$ with
the preceding signed-moment form.  The identification is explicit:
with $x=\log(t/(1-t))$ put
\[
 \rho(t)=\frac{x^4+16x^3+(2\pi^2+48)x^2+16\pi^2x+
                    \pi^4-16\pi^2}{t(x^2+\pi^2)^3},\quad
 v(t)=\int_t^1\rho(s)\,\frac{ds}s.
\]
Then $h_j=\int_0^1t^j(1-t)v(t)\,dt$.  The quartic numerator has two
real roots: after $y=x+4$, its derivative is four times
$y^3+(\pi^2-24)y+32$, which has negative discriminant and one real
root, and its minimum is negative.  Consequently $\rho$ has signs
$+,-,+$.  The derivative $v'=-\rho/t$, the endpoint positivity and
$\int v=d_0=-1/6$ show that $v$ has precisely two simple zeros
$\alpha<\beta$ with the same sign pattern.  Therefore
$d\omega=(1-t)v(t)\,dt/((t-\alpha)(t-\beta))$ is positive and finite,
which verifies the sign-change hypothesis.

The arithmetic transfer uses the fixed-offset congruence from
\texttt{gamma\_cont\_pade.md}: if $p>2r+3$ and $D_r$ is a $p$-adic
unit, then $p\ell_{p+r}\equiv4\pmod p$, hence the reduced denominator
$q_{p+r}$ has valuation $v_p(q_{p+r})=1$.  The rational moments above
give an explicit increasing bound $\mathcal B_r$ on the reduced numerator
of $D_r$ with
\[
 \log\mathcal B_r=6r^2\log r+O(r^2).
\]
Indeed, through degree $M=2r+2$, $M!L_{M+1}$ clears $B$, its cube
clears $B^3$, and an extra $L_M$ clears the divisions by $s+1$.
Expanding the determinant gives the asserted leading growth.  The exact
bound retained in the source is
\[
 \mathcal B_r=r!(C_MQ_M)^r,\quad
 Q_M=(M!L_{M+1})^3L_M,\quad
 C_M=2\{1+8(M+1)(M+3)\}.
\]
Combining this with Theorem~\ref{ec:normal-theorem} proves: for every
$c<1/\sqrt3$, every sufficiently large prime $p$, and simultaneously
every integer
\[
 1\le r\le c\sqrt{\log p/\log\log p},
\]
at least one of $q_{p+r},q_{p+r+1},q_{p+r+2}$ has $p$-valuation one.
The growing window is too narrow to imply a uniform prime-count lower
bound.  In particular this is not an exponential lower bound on every
$q_N$.  The previously suggested stronger assertion that every prime
$N/2<p\le N$ survives is false at $(N,p)=(51,43)$.

\section{A raw quotient index and why primitive cancellation remains separate}
\label{ec:quotient-index}

Let $B\in\Z^{r\times s}$ have full row rank, and suppose the matrix
$C$ formed by adjoining the two target rows $a,b$ has rank $r+2$.
Let $\Lambda=\ker_{\Z}B$ and
$\Omega=\{(a\cdot u,b\cdot u):u\in\Lambda\}\subset\Z^2$.
For a full-row-rank integer matrix $M$, write $\Delta(M)$ for the gcd
of its maximal minors.

\begin{proposition}[Exact target-lattice index]
\label{ec:index-theorem}
The lattice $\Omega$ has rank two and
\[
 [\Z^2:\Omega]=\frac{\Delta(C)}{\Delta(B)}.
\]
Equivalently, for a saturated integer-kernel basis $u_1,\ldots,u_k$,
the index is the gcd of all determinants
$\det(\Phi(u_i),\Phi(u_j))$, where $\Phi(u)=(a\cdot u,b\cdot u)$.
\end{proposition}
\begin{proof}
An integer column Hermite reduction gives a unimodular $V$ with
$BV=[H\ 0]$ and $H$ invertible.  The last columns of $V$ form a basis
of the entire integer kernel.  In this basis
\[
 CV=\begin{pmatrix}H&0\\ *&F\end{pmatrix},
\]
where the columns of $F$ generate $\Omega$.  Every nonzero maximal
minor contains all first $r$ columns, so
$\Delta(C)=|\det H|\Delta(F)$, while
$\Delta(B)=|\det H|$.  The rank-two image index is $\Delta(F)$.
\end{proof}

For a Sondow window, write the positive integer-normalized integrals as
$e_n=a_n\gamma-b_n+\sum_p\beta_{p,n}\log p$ and choose rational
upper bounds $e_n\le E_n$.  Set $S_a=\sum|a_n|$ and $S_e=\sum E_n$.
If $u,v\in\Lambda$ have independent target images and
$D=[\Z^2:\Omega]$, then
\[
 \|u\|_\infty\|v\|_\infty\ge\frac{D}{2S_aS_e}.
\]
Indeed, for $L_u=Q_u\gamma-P_u$, the nonzero target determinant is a
multiple of $D$, whereas
\[
 |Q_uP_v-Q_vP_u|=|Q_vL_u-Q_uL_v|
 \le2\|u\|_\infty\|v\|_\infty S_aS_e.
\]
This holds without any assumption on $\gamma$.  It bounds raw
representatives of two independent targets.  Dividing target pairs by
their gcd changes that cost: the lattice $M\Z^2$ has index $M^2$, but
its two coordinate generators reduce to primitive pairs of determinant
one.  The primitive quotient norm in the first paper is precisely the
later correction needed to account for this operation.

The source \texttt{gamma\_cont\_quotient.md} checks a proposed prime
divisibility law at finitely many windows $N=32,40,48,64$.  The
all-$N$ divisibility law remains a conjecture.  Neither the raw index nor
those finite checks give a no-go theorem for primitive output forms.

\section{Retained auxiliary routes and their exact status}
\label{ec:retained-register}

The following items are part of the cumulative record, but their long
special-purpose calculations are retained in the source archive rather
than duplicated in the main proof development.

\begin{enumerate}
\item \textbf{Restricted positive cones and square certificates.}
\texttt{gamma60\_operators.md} solves the full nonnegative numerator
cone at $m=3$, rather than only a chosen finite sample, and gives its
exact algebraic endpoints and a dual certificate.  The resulting bound
$|q\gamma-p|>q/200$ belongs to that fixed cone.  The larger-pole
completeness theorem above supersedes an interpretation of this finite
obstruction as a global one.  The distinct Laguerre square-certificate
calculation in \texttt{gamma60\_arithmetic.md} obtains
$|a-b\gamma|\ge |b|/(\sqrt{80}N^{3/2})$ in its specified square
family, with tensor extensions.  It does not cover arbitrary signed
integrals or the whole floor-integral space.

\item \textbf{Growing-degree exponential transforms.}
\texttt{gamma60\_geometry.md} constructs increasing-degree transforms
outside the fixed-degree logarithmic denominator obstruction of
Theorem~\ref{ec:logden-theorem}.  Their limit coefficients retain
$e^\gamma$.  This is a genuine scope distinction; no arithmetic
classification of that remaining constant is derived from the transform.

\item \textbf{Positive tensors and finite support.}
\texttt{gamma\_next\_polynomial.md} proves finite positive tensor
certificates for strict polynomial signs and constructs positive rational
atomic measures realizing alternative values in a fixed-pole feasible
interval.  Consequently increasing tensor degree at fixed poles does
not isolate $\gamma$.  The particular $m=3$ alternative is $4/7$.
\texttt{gamma\_frontier\_arithmetic\_cones.md} also constructs
totally real algebraic alternatives of arbitrary degree under finitely
many moment conditions.  Their positivity is a real condition, not an
integrality theorem for all divided powers.

\item \textbf{Divided powers.}
The same arithmetic-cones note proves that an algebraic number $x$ for
which every $\binom xn$ is an algebraic integer must be an integer.
No such factorial divisibility is established for $\gamma$ or for a
fixed multiple of it.  This implication is a source theorem with
unverified target hypothesis, not an irrationality proof.

\item \textbf{Geometric nodes and actual reduced denominators.}
\texttt{gamma\_frontier\_scaling.md} gives a Gaussian-binomial
Gamma filter with positive error
$\log R_N=-\tfrac12(\log B)N^3+O_B(N)$, but with a displayed
rational prefactor of logarithmic height $\Theta_B(q\log q)$,
$\log q\sim\tfrac12(\log B)N^2$.
\texttt{gamma\_frontier\_spectral.md} constructs a distinct purely
rational family with
$e^{-n}/6<\gamma-r_n<5e^{-n}$ for $n\ge4$ and proves for its
actual reduced denominator $d_n$ that
$\log d_n=\Theta(n^4\log n)$.  Thus $d_n(\gamma-r_n)\to+\infty$
for that family, even after cancellation.

\item \textbf{Mahler means and fixed kernels.}
\texttt{gamma\_frontier\_mahler.md} isolates a Fourier mean
$1/2-\gamma/\log b$; the arithmetic status of that mean does not
follow from a Mahler special-value theorem.  The fixed-kernel analysis
in \texttt{gamma\_next\_renormalization.md} uses
\[
 W(x)=\kappa(1-x)+(1-x^2)A(x),\qquad \kappa=W(-1)/2,
\]
for $W(1)=0$, to distinguish retention and loss of the gamma coefficient.
The denominator exponents proved there concern its specified truncations
and do not apply to every generating function for $\gamma$.

\item \textbf{Information and diffusion.}
\texttt{gamma60\_physics.md} records the classical normalized identity
$D_{\rm KL}(\operatorname{Exp}(1)\Vert\operatorname{Gamma}(2,1))=\gamma$.
The integer moments $n!$ determine the exponential distribution, but the
associated Ser--Sondow reconstruction has error asymptotic to
$1/(N\log N)$ and a prime-log coefficient denominator with logarithm
at least $(1/2+o(1))N$.  Uniqueness of the distribution supplies no
small integer form.  For the diffusion generator
$L=x\partial_x^2+(1-x)\partial_x$, the formal identity
$L\log x=-1$ cannot be integrated against stationarity without the
nonzero boundary term.  The source computes that term; the purported
contradiction from dropping it is withdrawn.

\item \textbf{Multiscale quadrature and spectral realizations.}
\texttt{gamma60\_analysis.md} combines two Gamma quadratures and
proves an additional factor $N^{-3}$ in its error, with a retained
algebraic-logarithmic target.  It does not remove those coordinates.
\texttt{gamma\_next\_spectral.md} constructs transcendental finite
approximants in $\Qbar(\pi)$ converging to $\gamma$ with error
$\Theta(n^{-2})$.  Transcendence of each approximant is not preserved
under a limit.  The earlier spectral note also realizes any prescribed
positive total trace by rapidly decreasing, finitely algebraically
independent positive eigenvalues.  Those qualitative spectral properties
alone therefore cannot decide the arithmetic of a total trace.

\item \textbf{Mixed Gamma jets and local discriminant pairings.}
The first gamma paper incorporates the integral coefficient conversion
and graded denominator bound from the continuation notes.  The additional
closed multiplication law over a larger period ring is retained in
\texttt{gamma\_cont\_mixed.md}; enlarging the coefficient ring does
not prove a zero-lifting theorem.  The local Pad\'e pairing
$[v]^TM^{-1}[v]\pmod\Z$ can be isotropic even for a positive definite
real Gram matrix.  Its order, not the order of $[v]$ or the size of the
discriminant group, is the reduced output denominator.  The source
example $(N,p)=(7,5)$ exhibits this distinction.
\end{enumerate}

The source register is part of completeness: it identifies which claims
have been re-proved here, which are inputs with retained detailed proofs,
which are finite computations and which remain conditional.  No item in
this part proves that $\gamma$ is irrational or transcendental.

% END INTEGRATED PART III

\clearpage
\part{Analytic variations and imported methods}\label{master:analytic}
% BEGIN INTEGRATED PART IV
\section{Automatic series, periodic identities, and normalization}
\label{ar:automatic}
The following investigations concern explicit representations of Euler's
constant and the information preserved by changes of representation.  A
functional transcendence theorem is not, without an additional specialization
argument, a theorem about a particular value.  The counterexamples below are
kept alongside the positive constructions because they delimit precisely
which information is insufficient.

\subsection{A Mahler derivative whose primitive leaves every Mahler class}
For an integer $B\ge2$, put
\[
 \epsilon_B(n)=B\mathbf1_{B\mid n}-1,\qquad
 f_B(z)=\sum_{n\ge1}\epsilon_B(n)
       \frac{\lfloor\log_Bn\rfloor}{n}z^n,\qquad
 G_B(z)=\sum_{k\ge1}z^{B^k}.
\]
The generalized Vacca identity \cite{ar:SondowVacca} and direct differentiation
give
\begin{equation}\label{ar:vacca}
 f_B(1)=\gamma,\qquad
 zf_B'(z)=G_B(z)\left(\frac B{1-z^B}-\frac1{1-z}\right),\qquad
 G_B(z)-G_B(z^B)=z^B.
\end{equation}
Here the boundary series converges.  Indeed, its sum through $B^M-1$ is
\[
 M H_{B^{M-1}-1}-(M-1)H_{B^M-1},
\]
which tends to $\gamma$; a partial last block is $O_B(MB^{-M})$.
In particular $f_B'$ is $B$-Mahler.  The primitive has a different
arithmetic property.

\begin{lemma}[Finite denominator support for a Mahler series]
\label{ar:mahlerfinite}
If $H(z)=\sum h_nz^n\in\Qbar[[z]]$ satisfies
$\sum_{i=0}^r a_i(z)H(z^{b^i})=0$, where $b\ge2$, $a_i\in\Qbar(z)$,
and $a_0\ne0$, then all $h_n$ belong to the ring of $S$-integers of one
number field for one finite set $S$ of finite places.
\end{lemma}
\begin{proof}
Clear rational-function denominators, and write
$d=\operatorname{ord}_0a_0$, $\alpha=[z^d]a_0\ne0$.  The coefficient of
$z^{n+d}$ solves for $h_n$ by division by $\alpha$.  All other indices in
$a_0H$ are smaller than $n$, and those in $a_iH(z^{b^i})$, $i>0$, are of
the form $(n+d-k)/b^i<n$ for sufficiently large $n$.  Thus a finite initial
segment and the finitely many coefficients of the $a_i$, together with
$\alpha^{-1}$, generate all subsequent coefficients by ring operations.
\end{proof}

\begin{proposition}[Persistent new-prime denominators]
\label{ar:vaccaprime}
The series $f_B$ is not Mahler in any integer base.  More generally, let
$c\in\Q^\times$, $R\in z\Q[[z]]$ be algebraic, and $r=R'(0)$.  If
$\exp(cf_B+R)=\sum e_nz^n$, then at every sufficiently large prime $p$,
\[
 pe_p\equiv-c\lfloor\log_Bp\rfloor-r\pmod p,
 \qquad v_p(e_p)=-1.
\]
Consequently this exponential is neither algebraic nor Mahler in any
integer base.  Multiplication by an algebraic unit in $\Q[[z]]$, or a
fixed nonzero rational rescaling of the variable, preserves the conclusion.
\end{proposition}
\begin{proof}
For sufficiently large $p\nmid B$, $p[z^p]f_B=-\lfloor\log_Bp\rfloor$
is nonzero modulo $p$, while all earlier coefficients are $p$-integral.
Eisenstein's denominator theorem makes all coefficients of $R$ simultaneously
$p$-integral outside a fixed finite set.  In the coefficient of degree $p$
of $\exp h$, $h=cf_B+R$, only $h_p$ and $h_1^p/p!$ can have a $p$-denominator.
Wilson's and Fermat's congruences give the displayed residue.  After a fixed
common denominator is cleared, its negative is a nonzero integer of size
$O(\log p)<p$ for all sufficiently large $p$.  Lemma~\ref{ar:mahlerfinite}
and Eisenstein's theorem apply.  For an algebraic unit $U$, only
$U(0)e_p$ contributes a nonintegral summand to $[z^p]U\exp h$; the rescaling
claim follows because a fixed nonzero rational number is a unit at almost
all primes.
\end{proof}

The same conclusion holds for a nontrivial finite combination
$F=\sum_{d\in\mathcal D}c_df_{b^d}$, including $\exp(F+R)$.
For $b^n<p<b^{n+1}$ the detecting residue is
$-A(n)$, respectively $-A(n)-R'(0)$, where
$A(n)=\sum_dc_d\lfloor n/d\rfloor$.  This function is not eventually
constant: otherwise its periodic difference
$\sum_{d\mid n}c_d$ would vanish identically, and triangular induction
would give every $c_d=0$.  Bertrand's theorem supplies a prime in
$(b^n,2b^n)$ for each sufficiently large detecting index.  This proves
persistence for powers of one fixed base, not for arbitrary independent
bases.  Nevertheless
\begin{equation}\label{ar:boundarycounter}
 (f_b-f_{b^2})(1)=0,
 \qquad \exp\bigl((f_b-f_{b^2})(1)\bigr)=1.
\end{equation}
Thus the function-level obstructions just proved cannot themselves decide
the arithmetic of the boundary value.

\subsection{Finite Gamma data and global rigidity}
\begin{proposition}[Finite-jet normalization counterexamples]
\label{ar:finitejet}
Fix finitely many positive rational evaluation points, finitely many
rational test points for Gamma multiplication formulas, and a finite jet
order $R\ge1$.  For every real $\eta$ there is a meromorphic function
$\widetilde\Gamma$ satisfying the Gamma recurrence and reflection formula
globally, with the same poles, residues, zeros and positive integer values
as $\Gamma$, the prescribed rational function values, and the multiplication
identities through jet order $R$, such that
$\widetilde\Gamma'(1)=\eta$ and $\widetilde\Gamma(x)>0$ for $x>0$.
It is differentially transcendental over $\C(z)$.
\end{proposition}
\begin{proof}
Choose $N$ clearing the denominators of all test arguments, including their
multiplication translates.  Put $M=\lceil R/2\rceil$ and
\[
 \phi(z)=\frac1{2\pi N}\sum_{k=0}^M
 \frac{\binom{2k}{k}}{4^k(2k+1)}\sin^{2k+1}(2\pi Nz),\qquad
 \widetilde\Gamma(z)=\Gamma(z)e^{(\eta+\gamma)\phi(z)}.
\]
The truncated arcsine identity gives
$\phi(t+h)=h+O(h^{2M+3})$ at every test argument $t$.
Moreover $\phi$ is entire, $1$-periodic, zero at the integers, and
$\phi(1-z)=-\phi(z)$.  These facts prove recurrence, reflection, divisor,
residue, value and derivative assertions.  In a multiplication identity
the added exponents have difference
$\sum_{j=0}^{m-1}\phi(z+j/m)-\phi(mz)$; its linear jets cancel and all
jets through $R$ vanish.  Positivity follows on the positive axis.
The difference-Galois form of H\"older's theorem applies to every nonzero
meromorphic solution of $F(z+1)=zF(z)$ \cite{ar:HardouinSinger}.
This is an existence counterexample, explicitly using $\gamma$ in its
definition, not a proposed gamma-independent approximation algorithm.
\end{proof}

This freedom disappears if the finite multiplication tests are replaced by
a global duplication identity and $\widetilde\Gamma/\Gamma$ is entire and
nowhere zero.  In that case
\[
 \widetilde\Gamma(z)=e^{4\pi ikz}\Gamma(z)\quad(k\in\Z),
\]
and positivity selects $k=0$.  To see this, the periodic logarithmic
derivative $h$ of the quotient satisfies
$2h(2z)=h(z)+h(z+1/2)$.  Its Fourier coefficients obey $a_{2n}=a_n$;
exponential decay on horizontal lines forces $a_n=0$ for $n\ne0$.
Periodicity, the value at $1$, and duplication then give the displayed
exponential.  This is analytic uniqueness, not arithmetic specialization.
An even simpler counterexample to unrestricted difference-value principles
is
\[
 F(z)=\psi(z)+\gamma=
 \sum_{k\ge1}\left(\frac1k-\frac1{k+z-1}\right),\qquad
 F(z+1)-F(z)=z^{-1},\qquad F(n)=H_{n-1}\in\Q.
\]
It is differentially transcendental yet has an entire algebraic shift orbit
of rational values.  Shift, reflection and multiplication identities for
$\psi$ are invariant under adding a constant \cite{ar:DLMFGamma}.

\subsection{Periodic constants and cyclotomic averaging}
If $f$ is algebraic-valued and $q$-periodic, let
$m=q^{-1}\sum_{a=1}^qf(a)$ and
$\widehat f(r)=q^{-1}\sum_{a=1}^qf(a)\zeta_q^{-ra}$.  Fourier expansion
and Dirichlet convergence give
\begin{equation}\label{ar:periodic}
 \lim_{x\to\infty}\left(\sum_{n\le x}\frac{f(n)}n-m\log x\right)
 =m\gamma-\sum_{r=1}^{q-1}\widehat f(r)\log(1-\zeta_q^r).
\end{equation}
Thus the finite periodic family adds algebraic logarithmic directions to
one unresolved $\gamma$ coordinate.  For example,
\[
 \gamma(\Omega)=\prod_{p\in\Omega}(1-p^{-1})
 \left(\gamma+\sum_{p\in\Omega}\frac{\log p}{p-1}\right).
\]
Baker's theorem controls a nonzero logarithmic part but does not by itself
exclude algebraicity of the remaining $\gamma$ coordinate.

For $L(q)=\sum_{n\ge1}d(n)q^n$, an integer $p>1$ and prime $m$, the
root-of-unity filter gives
\begin{equation}\label{ar:unweighted}
 \frac1m\sum_{j=0}^{m-1}L(\zeta_m^jp^{-1/m})
 =2L(p^{-1})-L(p^{-m}).
\end{equation}
Indeed only indices divisible by $m$ survive, and
$d(mn)=2d(n)-\mathbf1_{m\mid n}d(n/m)$.  Weighting instead by
$1+X+\cdots+X^{m-1}$, with $m$ now arbitrary, gives
\[
 W_m=(1-p^{-1})\sum_{\ell\ge1}p^{-\ell}D(m\ell),\quad
 D(x)=\sum_{n\le x}d(n).
\]
With $D(x)=x\log x+(2\gamma-1)x+\Delta(x)$, set
\[
 C_p=\frac{(p-1)^2}{p}\sum_{\ell\ge1}\frac{\ell\log\ell}{p^\ell},
 \quad T_m=\frac12\left(\frac{p-1}{m}W_m-\log m+1-C_p\right).
\]
Then the exact identity
\begin{equation}\label{ar:weighted}
 T_m-\gamma=\sum_{\ell\ge1}
 \frac{(p-1)^2\ell}{p^{\ell+1}}\frac{\Delta(m\ell)}{2m\ell}
\end{equation}
has positive weights summing to one.  It retains $\gamma$, but introduces
the logarithmic moment $C_p$ and is not a rational approximation.  Infinite
analytic averages in these formulas are not field traces of known algebraic
numbers; only their finite polynomial truncations have that interpretation.

A separate truncation makes the height cost explicit.  For prime $M\ge101$,
$q=1-1/M$, $K=\lceil4M\log M\rceil$, the reduced denominator of
$\sum_{n\le K}d(n)q^n=A_M/D_M$ is $D_M=M^K$: its numerator modulo $M$
is $(-1)^Kd(K)\ne0$, since $d(K)\le2\sqrt K<M$.
The normalized logarithmic approximant
$G_M=-\log M-\log(q)A_M/D_M$ satisfies
$\gamma-G_M=(4M)^{-1}+O(M^{-2})$, by Euler--Maclaurin and the geometric
tail estimate.  Its primitive integer logarithmic form is
\[
 D_M\gamma-(A_M-D_M)\log M+A_M\log(M-1),
\]
of size asymptotic to $M^K/(4M)$.  The approximation is neither rational nor
a small linear form in $1,\gamma$.

\section{Divisor sums, exact dyadic denominators, and logarithmic jets}
\label{ar:divisor}
Throughout this section $\Delta(x)=D(x)-x\log x-(2\gamma-1)x$.
The elementary hyperbola estimate $|\Delta(x)|\le8\sqrt x$ suffices for
all explicitly stated elementary error bounds.

\begin{theorem}[A gamma-independent reduced dyadic family]
\label{ar:dyadic}
For odd $j\ge3$, define
\[
 P_j=2^{j+1}+4jD(2^{j-2})-(j-2)D(2^j),\quad
 r_j=P_j/2^{j+2},\quad
 \mathcal J=\{j:\lfloor\sqrt{2^j}\rfloor\text{ is odd}\}.
\]
The set $\mathcal J$ is infinite.  For $j\in\mathcal J$, $r_j$ is reduced,
its denominator is exactly $q_j=2^{j+2}$, these fractions are distinct, and
\begin{equation}\label{ar:dyadicerror}
 |r_j-\gamma|\le(6j-4)2^{-j/2},\qquad
 P_j-q_j\gamma=4j\Delta(2^{j-2})-(j-2)\Delta(2^j).
\end{equation}
At most one fraction equals $\gamma$.  Consecutive selected indices satisfy
$j'\le2j+3$, and $q_{j'}\le2q_j^2$.
\end{theorem}
\begin{proof}
For $g(n)=\tfrac12(D(n)/n-\log n+1)$,
$g(n)=\gamma+\Delta(n)/(2n)$.  The combination
$((k+h)g(b^k)-kg(b^{k+h}))/h$ cancels the logarithm and preserves the
coefficient of $\gamma$.  Taking $b=2$, $k=j-2$, $h=2$ proves the formula
and error estimate.  Since $d(n)$ is odd exactly at squares,
$D(n)\equiv\lfloor\sqrt n\rfloor\pmod2$, so $P_j$ is odd on
$\mathcal J$.  Writing $j=2h+1$, selection means that the $h$th binary
digit of $\sqrt2$ is $1$; there are infinitely many such digits.
If the next $1$ occurs at $H>h$ and $A=\lfloor2^h\sqrt2\rfloor$, then
\[
 (3\cdot2^{2h})^{-1}<\sqrt2-A/2^h<2^{1-H}.
\]
The lower bound follows by rationalizing and using the positive integer
$2\cdot2^{2h}-A^2$.  Hence $H\le2h+2$.  Distinct reduced denominators
imply distinct fractions and therefore eventual nonvanishing of their errors.
\end{proof}

The error bound for the integer form in \eqref{ar:dyadicerror} grows like
$j2^{j/2}$.  Even an individual bound $\Delta(x)=O(x^\theta)$ yields only
$O(j2^{\theta j})$ for that integer form.  This is a limitation of the
available upper bound, not a lower bound on actual cancellation.
Ridout's theorem \cite{ar:Ridout} gives a precise sufficient condition:
if, for some fixed $\varepsilon>0$, infinitely many selected $j$ satisfy
\begin{equation}\label{ar:ridouttarget}
 0<|4j\Delta(2^{j-2})-(j-2)\Delta(2^j)|<q_j^{-\varepsilon},
\end{equation}
then $\gamma$ is transcendental.  Rational $\gamma$ is excluded by a fixed
integer-lattice lower bound; algebraic irrational $\gamma$ is excluded by
Ridout after division by $q_j$.  Condition \eqref{ar:ridouttarget} remains
unproved and is not asserted equivalent to transcendence.

\subsection{An interior value and its divided logarithmic jet}
Put $A_h=P_{2h+1}$, $C_h=A_{h+1}-4A_h$, and
\[
 F(z)=\sum_{h\ge1}A_hz^h,\qquad G(z)=\sum_{h\ge1}C_hz^h.
\]
The bound $A_h=8\gamma4^h+O(h2^h)$ and a telescoping sum imply
\begin{equation}\label{ar:interiorvalue}
 G\text{ is analytic for }|z|<\tfrac12,\qquad
 \gamma=1+G(1/4)/32.
\end{equation}
A finite partial sum is exactly the old dyadic approximation, so the
identity does not accelerate it.

\begin{proposition}[Functional transcendence and a stronger countermodel]
\label{ar:functioncounter}
The series $F,G$ are transcendental over $\Q(z)$ and have no finite
linear $2^m$-Mahler relation for any $m\ge1$.  Nevertheless, the initial
value $A_1=32$, coefficient parities, divisibility
$C_h\in(2h+1)\Z$, and the displayed value identity do not determine the
arithmetic of its value: they admit models with any prescribed real value
$\alpha$ in place of $\gamma$.
\end{proposition}
\begin{proof}
Let $b_h=\lfloor2^h\sqrt2\rfloor\bmod2$.  Then
$A_h\equiv b_h$ and $C_h\equiv b_{h+1}\pmod2$.
Adamczewski--Bugeaud \cite{ar:AB} implies that these binary digits are not
$2$-automatic, while Christol's theorem \cite{ar:Christol} makes the
corresponding reduced series transcendental over $\mathbb F_2(z)$.
A primitive integer polynomial relation would reduce to a nonzero relation
modulo $2$.  A primitive $2^m$-Mahler relation reduces, by Frobenius, to
$\sum_i\overline Q_i(z)\overline F(z)^{2^{mi}}=0$, also a nonzero
polynomial relation.  The same proof applies to $G$.

For the countermodel start with $\widetilde A_1=32$ and let
$e_h=\widetilde A_h-8\alpha4^h$.  Choose
$v_h\in b_{h+1}+2\Z$ nearest to $-4e_h/(2h+1)$, and put
$\widetilde A_{h+1}=4\widetilde A_h+(2h+1)v_h$.
Then $|e_h|\le2h-1$ for $h\ge2$, and $|v_h|\le4$ for $h\ge2$.
The same telescoping argument gives
$\alpha=1+\widetilde G(1/4)/32$.
The countermodel series has polynomially bounded integer coefficients and
the same nonautomatic parity.  Its radius is $1$, it is not rational,
and P\'olya--Carlson gives the unit circle as natural boundary
\cite{ar:PolyaCarlson}.  In particular, even this additional analytic
property does not preclude an algebraic prescribed value.
\end{proof}

The exact division $C_h=(2h+1)V_h$ has further content.  Set
\[
 B_h=D(2^{2h-1}),\qquad V_h=8B_{h+1}-B_{h+2}-16B_h,
 \qquad V(z)=\sum_{h\ge1}V_hz^h.
\]
Then
\begin{equation}\label{ar:jet}
 V(1/4)=8-16\log2,\qquad
 V'(1/4)=64\gamma+32\log2-80.
\end{equation}
To prove this, for $J(z)=\sum B_hz^h$ use
\[
 J(z)=\frac{4\log2\,z}{(1-4z)^2}
 +\frac{4(\gamma-(1+\log2)/2)z}{1-4z}+E(z),
 \qquad
 V(z)=-\frac{(1-4z)^2}{z^2}J(z)+\frac3z-4,
\]
where $E$ is analytic on $|z|<1/2$.  Its contribution to both the value
and first derivative vanishes.  For general $b\ge2$, the analogous values
are
\begin{align*}
 X_b&=D(b^3)-b^2D(b)-2b^3\log b,\\
 Y_b&=b^5(2\gamma-1+\log b)-b^4D(b),
 \qquad \Qbar(X_b,Y_b)=\Qbar(\gamma,\log b).
\end{align*}
Every $X_b$ is transcendental by Hermite--Lindemann; at most one $Y_b$ is
algebraic because $Y_b/b^5-Y_c/c^5$ is a rational translate of $\log(b/c)$.
Neither statement excludes algebraicity of $\gamma$.

Even the fixed value in \eqref{ar:jet} admits an arbitrary derivative
parameter.  For $h\ge3$ round
\[
 4^h\bigl(h\log2+\alpha-(1+\log2)/2\bigr)
\]
down to the nearest integer of the parity of $B_h$, retaining $B_1=3$,
$B_2=20$.  The error lies in $(-2,0]$, so the resulting $\widetilde V_h$
are bounded, with $-16<\widetilde V_h<34$ for $h\ge3$.
The preceding double-zero calculation preserves $V(1/4)$ but replaces
$\gamma$ by $\alpha$ in $V'(1/4)$.  This countermodel preserves the
finite-difference structure and all indicated parities, not the actual
identity $B_h=D(2^{2h-1})$ or the full multiplicativity of $d$.

\subsection{Local Euler factors and Riesz smoothing}
For a prime $p$, interpreting $d(x)=0$ at nonintegral $x$, one has
\[
 d(n)-2d(n/p)+d(n/p^2)=d(n)\mathbf1_{p\nmid n}.
\]
The local identity $(1-X)^2(1-X)^{-2}=1$ proves it.  Accordingly
$O_k=D(2^k)-2D(2^{k-1})+D(2^{k-2})$ is the odd-part divisor sum, and
\[
 R_j=\frac{2^{j-1}+4(j+2)O_{j-2}-jO_j}{2^j}\longrightarrow\gamma.
\]
For odd $j=2h+1\ge5$, the numerator is odd whenever $b_h\ne b_{h-1}$.
This gives infinitely many distinct reduced denominators $2^j$, with
\[
 |R_j-\gamma|\le2(2+\sqrt2)^2(3j+4)2^{-j/2}.
\]
Indeed the parity of $O_j$ counts odd squares, and the main-term expansion
proves convergence and the bound.  Removing finitely many Euler factors
changes constants, not the exponent guaranteed by individual error bounds.
For a fixed filter $\sum_ic_iD(\lambda_ix)$ the main term is exactly
\[
 x\{a\log x+(2\gamma-1)a+b\},\qquad
 a=\sum_ic_i\lambda_i,\quad b=\sum_ic_i\lambda_i\log\lambda_i.
\]
Thus cancelling $x\log x$ by a fixed filter also cancels its $\gamma$
coefficient.  A shift polynomial at geometric nodes $b_0^h$ cancels both
main terms exactly when it has a double root at $b_0$.

For comparison, let
\[
 M(x)=\sum_{n\le x}d(n)(x-n),\qquad
 A(x)=M(x)/x^2+3/4-1/(4x)-\tfrac12\log x.
\]
The integrated Voronoi estimate gives $A(x)-\gamma=O(x^{-5/4})$
\cite{ar:Ivic}.  For odd $j\ge5$ and $x=2^j$,
\[
 s_j=\frac{jA(x/4)-(j-2)A(x)}2
 =\frac{16jM(x/4)-(j-2)M(x)+3x^2/2-(3j+2)x/4}{2x^2}.
\]
If $\lfloor\sqrt{2^j}\rfloor\bmod4\in\{1,2\}$, the numerator is odd.
Counting square contributions to $M(x)$ modulo $2$ proves this; the
condition occurs infinitely often because it detects changes in consecutive
binary digits of $\sqrt2$.  Hence there is an infinite reduced family with
$q_j=2^{2j+1}$ and
\[
 |s_j-\gamma|=O((\log q_j)q_j^{-5/8}).
\]
The stronger real-variable smoothing has squared the denominator.  Even
the classical unsmoothed bound $\Delta(x)=O(x^{1/3}\log x)$ gives exponent
$2/3$ in the original family, whereas the Ridout sufficient threshold is
$1+\varepsilon$.  These are guaranteed upper-bound exponents, not proved
optimal errors.

\subsection{The first singular local layer of the moment denominator}
\label{ar:localmoment}
For the previously defined moment approximation
$\ell_N=\boldsymbol c_N^TK_N^{-1}\boldsymbol c_N$, retain its local
Schur representation at $N=p+r$, $p>2r+3$:
\[
 \ell_{p+r}=\beta+z^TT^{-1}z,\quad
 T\in\operatorname{Sym}_r(\Z_p),\quad z\in\Z_p^r,\quad
 p\beta\equiv4\pmod p.
\]
Here $p\beta$ is integral; the displayed representation is the input from
the earlier moment calculation, not an assertion for arbitrary matrices.
If $Z$ lifts a basis of $\ker(T\bmod p)$ and
\[
 \kappa=p^{-1}Z^TTZ\bmod p,\qquad h=Z^Tz\bmod p,
\]
then, provided $\kappa$ is nonsingular,
\begin{equation}\label{ar:schurlayer}
 p\ell_{p+r}\equiv4+h^T\kappa^{-1}h\pmod p.
\end{equation}
To prove this, choose a complementary basis and write
$T=\left(\begin{smallmatrix}A&pC\\pC^T&pD\end{smallmatrix}\right)$,
where $A$ is invertible and $D\bmod p=\kappa$.  Its Schur complement is
$p(D-pC^TA^{-1}C)$, so the only negative valuation layer is $p^{-1}$ and
has the residue in \eqref{ar:schurlayer}.  The formula is invariant under
changing lifts and bases.  A nonzero residue means valuation $1$ in the
reduced denominator; a zero residue means valuation $0$.  In dimension
one, $(-\kappa/p)=-1$ excludes cancellation.  The verified example
$(N,p,\kappa,h)=(51,43,19,28)$ has residue $0$, unlike
$(15,13,10,11)$, $(49,47,3,39)$ and $(24,19,6,15)$, whose residues are
$7,41,13$.  If $\kappa$ is itself singular, a deeper layer is required.

\section{Exact filter optimization and continuous divisor correlations}
\label{ar:filters}
\subsection{What individual error intervals can certify}
\begin{proposition}[Two-support minimax principle]
\label{ar:minimax}
For $N<K$ and positive $\epsilon_k$,
\[
 \min_{\sum c_k=1,\ \sum kc_k=0}\sum_{k=N}^K\epsilon_k|c_k|
 =\min_{N\le i<j\le K}\frac{j\epsilon_i+i\epsilon_j}{j-i}.
\]
A minimizer has $c_i=j/(j-i)$, $c_j=-i/(j-i)$ and all other coefficients
zero.  If $\epsilon_k=C2^{-ak}$, $a>0$, the minimum is
$\Theta_a(CK2^{-aK})$, uniformly in the left endpoint $N<K$.
\end{proposition}
\begin{proof}
Write $c_k=x_k-y_k$, $x_k,y_k\ge0$.  An optimal basic feasible solution of
this positive-cost linear program with two equations has at most two
positive variables.  One variable cannot meet both equations; two distinct
indices with opposite signs give the stated solution.  For the asymptotic
bound the pair $K-1,K$ is sufficient above.  Below, put $s=j-i$; when
$j\ge K/2$, the first cost term is at least
$(CK/2)2^{-aK}2^{as}/s$, bounded below by a positive constant times
$CK2^{-aK}$.  When $j<K/2$, the exponential bound is stronger.
\end{proof}

There is an exact optimum for arbitrary variable length in the elementary
divisor bound.  For $4\le L\le k$, impose
$\sum_{\ell=0}^Lw_\ell=1$, $\sum\ell w_\ell=k$.  Then
\begin{equation}\label{ar:exactminimax}
 \min\sum_{\ell=0}^L|w_\ell|2^{\ell/2}=5k/4-1,
 \qquad w_0=1-k/4,\quad w_4=k/4.
\end{equation}
The dual certificate is $|5\ell/4-1|\le2^{\ell/2}$, with equality only
at $\ell=0,4$.  Its signs also prove uniqueness.  This yields
\[
 R_{k,4}=\frac{2^{k+2}+(4-k)D(2^k)+16kD(2^{k-4})}{2^{k+3}},
 \quad |R_{k,4}-\gamma|\le(5k-4)2^{-k/2}.
\]
The denominator is exactly $2^{k+3}$ on the selected odd indices of
Theorem~\ref{ar:dyadic}.  Optimizing the error is different from optimizing
the integerized error: among two-point formulas with displayed denominator
$2m2^k$, that latter elementary budget is minimized at $m=1$, not $m=4$.
Actual primitive-height optimization over all filters is not solved here.

Allowing continuous nodes $Xe^{-s_i}$ with $0\le s_i\le\log X$ leads to
the analogous exact relaxation.  If $u=W(e^{-1})>0$ and
$\log X\ge(1+u)/\beta$, then
\[
 \min_{\sum w_i=1,\ \sum w_is_i=\log X}\sum_i|w_i|e^{\beta s_i}
 =\frac\beta u\log X-1.
\]
The line $(\beta/u)s-1$ is bounded in absolute value by $e^{\beta s}$
and tangent at $s_*=(1+u)/\beta$; equality uses $s=0,s_*$.
The optimal continuous nodes need not be rational.  These optimizations
concern independent error intervals, leaving actual correlations untouched.

\subsection{Covariance, uniform growing filters, and exact prediction}
Let
\[
 X_r(x)=2^{-r/2}\Delta(4^rx),\quad
 t=2^{-3/2},\quad c=\frac{2(1-t)}{1+t},\quad
 \rho_h=t^h(1+ch),\quad S=\sum_{n\ge1}d(n)^2n^{-3/2}.
\]
For fixed $r,s$, the continuous covariance is
\begin{equation}\label{ar:covariance}
 \lim_{T\to\infty}T^{-3/2}\int_1^TX_r(x)X_s(x)\,dx
 =\frac S{6\pi^2}\rho_{|r-s|}.
\end{equation}
Its normalized spectral density is
\begin{equation}\label{ar:density}
 f(\theta)=\frac{(1-t)^3}{1+t}
 \frac{1+4t+t^2+2t\cos\theta}{(1-2t\cos\theta+t^2)^2},
 \quad f_{\min}=\left(\frac{1-t}{1+t}\right)^3>0.
\end{equation}
These identities use the actual divisor error, not an assumed stochastic
model.

\begin{theorem}[Uniform continuous mean for growing dyadic filters]
\label{ar:uniform}
Fix $0<\eta<1/2$.  Uniformly for all distinct nonnegative integers
$r_1,\ldots,r_m$, all real vectors $w$, and $m\le T^{1/2-\eta}$,
\[
 T^{-3/2}\int_T^{2T}\left(\sum_{i=1}^mw_iX_{r_i}(x)\right)^2dx
 =C_*\sum_{i,j}w_iw_j\rho_{|r_i-r_j|}
 +O_\eta(T^{-\eta/8})\|w\|_2^2,
 \quad C_*=(2^{3/2}-1)S/(6\pi^2).
\]
There is no upper bound on the $r_i$.  They and $w$ may depend on $T$ but
are held fixed during integration in $x$.
\end{theorem}
\begin{proof}
Use the truncated Voronoi formula \cite{ar:LauTsang}
\[
 V_M(u)=\frac{u^{1/4}}{\pi\sqrt2}\sum_{n\le M}
 d(n)n^{-3/4}\cos(4\pi\sqrt{nu}-\pi/4).
\]
Its standard truncation, followed by integration of distinct frequencies,
gives, for $2\le M\le Y$,
\[
 Y^{-3/2}\int_Y^{2Y}|\Delta(u)-V_M(u)|^2du
 \ll_\varepsilon M^{-1/2+2\varepsilon}
       +Y^{-1/2+2\varepsilon}\log(2Y).
\]
For completeness, first truncate at $\lfloor Y\rfloor$; the pointwise
remainder contributes $O_\varepsilon(Y^{1+2\varepsilon})$ to the square
integral.  The diagonal tail is
$O(Y^{3/2}\sum_{n>M}d(n)^2n^{-3/2})$.  The off-diagonal integral is
$O(Y/|\sqrt a-\sqrt b|)$; splitting $a<b$ at $a=b/2$ bounds its sum by
$O_\varepsilon(Y^{1+2\varepsilon}\log(2Y))$.
Rescaling $u=4^rx$ preserves this normalized estimate uniformly for $r\ge0$.

Use a Hilbert space with equal frequencies $2^r\sqrt n$ identified.  The
vector $v_r$ of coefficients $d(n)n^{-3/4}$ has inner products
\[
 \langle v_r,v_s\rangle
 =4^{-3h/4}\sum_{n\ge1}d(n)d(4^hn)n^{-3/2}=S\rho_h,
 \qquad h=|r-s|.
\]
The last equality follows by writing $n=2^\nu m$, $m$ odd, and summing
$(\nu+1)(\nu+2h+1)t^\nu$.  Summing the Poisson kernel and its derivative
gives \eqref{ar:density}.  Parseval bounds the limiting quadratic form above
and below by fixed positive multiples of $\|w\|_2^2$.

For $a,b\le M$, unequal frequencies satisfy
$|2^r\sqrt a-2^s\sqrt b|\ge(3\sqrt M)^{-1}$: after assuming $r\le s$,
either $4^{s-r}b\le4M$ and rationalization applies, or the gap exceeds
$\sqrt M$.  Group equal frequencies before squaring.  If the distinct
frequencies are ordered, their successive gaps give row sums bounded by
$O(\sqrt M\log(2mM))$ for the reciprocal-gap matrix, using
$2|A_iA_j|\le A_i^2+A_j^2$.  Therefore the nonconstant terms contribute
\[
 O(T^{-1/2}\sqrt M\log(2mM))
 \left\|\sum_iw_iv_{r_i,M}\right\|^2.
\]
The Hilbert-space tail has norm
$O_\varepsilon(\sqrt mM^{-1/4+\varepsilon})\|w\|_2$.
The replacement of the actual functions by their truncated sums has
normalized norm at most
\[
 O_\varepsilon\left(\sqrt mM^{-1/4+\varepsilon}
 +\sqrt mT^{-1/4+\varepsilon}\sqrt{\log(2T)}\right)\|w\|_2.
\]
Take $M=\lfloor T^{1-\eta}\rfloor$, $\varepsilon=\eta/16$.  All these
perturbations are $O_\eta(T^{-\eta/8})\|w\|_2$, and their quadratic effects
have the stated bound.  The constant term of the cosine-square integral
is exactly $(2^{3/2}-1)/(6\pi^2)$ times the Hilbert-space squared norm.
\end{proof}

The exact prediction problem also has a closed solution.  Set
\[
 A=1+4t+t^2,\quad \kappa=\frac{A+\sqrt{A^2-4t^2}}2,
 \quad s=t/\kappa,\quad \sigma^2=\frac{(1-t)^3}{1+t}\kappa.
\]
Then $f(\theta)=\sigma^2|(1+se^{i\theta})/(1-te^{i\theta})^2|^2$.
For every polynomial $W$ with $W(0)=1$, Parseval gives
\[
 \frac1{2\pi}\int_{-\pi}^{\pi}f(\theta)|W(e^{i\theta})|^2d\theta
 \ge\sigma^2.
\]
The infimum is attained among infinite absolutely summable filters by
$W_\infty(z)=(1-tz)^2/(1+sz)$; its degree-$H$ Taylor truncations have
excess $O(s^{2H})$.  Here $\sigma^2\simeq0.4967594891$.
Under the different constraint $\sum w_i=1$, the exact finite optimum is
$({\bf1}^TR_m^{-1}{\bf1})^{-1}=\Theta(1/m)$, where
$R_m=(\rho_{|i-j|})$.  None of these continuous mean statements supplies
a lower bound at any selected dyadic point.

\subsection{Exact resonance removal and its normalization}
Let $h_\ell=\tfrac12(D(4^\ell)/4^\ell+1)$ and
$P(Z)=\sum_jc_jZ^j\in\Z[Z]$, $q=P(1)\ne0$.
The combination $q^{-1}\sum_jc_jh_{k+j}$ cancels its logarithm precisely
when $P'(1)=-kP(1)$.  In a common finite Voronoi truncation, its square-index
resonance vanishes precisely when $P(2^{-3/2})=0$, or
$8Z^2-1\mid P$.
\begin{proposition}\label{ar:resonance}
These two conditions force $49\mid q$.  The least positive value $49$ is
attained by
\[
 P_k(Z)=(8Z^2-1)\{7-(7k+16)(Z-1)\}.
\]
More generally, for primitive $F\in\Z[Z]$, $A=F(1)\ne0$,
$B=F'(1)$, the least positive $|P(1)|$ under $F\mid P$ and
$P'(1)=-kP(1)$ is $A^2/\gcd(|A|,|B|)$.
\end{proposition}
\begin{proof}
Write $P=FQ$ by Gauss's lemma.  Then
$AQ'(1)=-(kA+B)Q(1)$, so $A/\gcd(A,B)$ divides $Q(1)$.
Conversely choose $Q$ linear with that constant value at $1$ and the
required integer derivative.  For $F=8Z^2-1$ use $A=7$, $B=16$.
\end{proof}
The concrete approximation has elementary error at most
$(126k+372)2^{-k}/49$.  The nonsquare frequencies remain.  The number
$49$ is a filter normalization, not necessarily a surviving factor of the
reduced rational denominator.

\section{Reciprocity, special-function conversion, and arithmetic transfer tests}
\label{ar:reciprocity}
\subsection{A positive exponentially convergent reciprocity formula}
Define $S(y)=\sum_{n\ge1}d(n)K_0(2\pi ny)$.  Tonelli's theorem applied
to the positive Bessel integral gives
$S(y)=\int_0^\infty L(e^{-2\pi y\cosh u})\,du$.
The classical Ramanujan--Koshliakov relation \cite{ar:DixitKumar} is
\[
 \sqrt y\{\gamma-\log(4\pi)+\log y+4S(y)\}
 =\frac{\gamma-\log(4\pi)-\log y+4S(1/y)}{\sqrt y}.
\]
One derivation writes
$4S(y)=(2\pi i)^{-1}\int_{(\sigma)}\Omega(s)^2y^{-s}\,ds$, where
$\Omega(s)=\pi^{-s/2}\Gamma(s/2)\zeta(s)$, and uses
$\Omega(s)=\Omega(1-s)$.  The residues at $1,0$ are, respectively,
$(c-\log y)/y$ and $-c-\log y$, $c=\gamma-\log(4\pi)$.
Differentiation at the self-dual point gives
\begin{equation}\label{ar:besselsum}
 \gamma=\log(4\pi)-2+
 \sum_{n\ge1}d(n)\{16\pi nK_1(2\pi n)-4K_0(2\pi n)\}.
\end{equation}
Every summand is positive.  In fact, with $x=2\pi n$,
\[
 8xK_1(x)-4K_0(x)
 =4\int_0^\infty(2x\cosh u-1)e^{-x\cosh u}\,du>0.
\]
From $K_1(x)\le e^{-x}(\sqrt{\pi/(2x)}+x^{-1})$ and
$d(n)\le2\sqrt n$, the tail after $N$ is at most
\begin{equation}\label{ar:besseltail}
 16(\pi+1)\rho^{N+1}
 \left(\frac{N+1}{1-\rho}+\frac\rho{(1-\rho)^2}\right),
 \qquad \rho=e^{-2\pi}.
\end{equation}
Thus this route proves one-sided nonvanishing and exponential analytic
convergence.  Its partial sums contain Bessel values and $\log(4\pi)$;
there is no associated fixed integer lattice.  The identity
\[
 K_0(z)=\sum_{k\ge1}\frac{H_k(z^2/4)^k}{(k!)^2}
       -\{\gamma+\log(z/2)\}I_0(z)
\]
reintroduces $\gamma$ and does not split \eqref{ar:besselsum} into two
convergent sums, since the separate $I_0$ contribution grows.  The dual
arguments $2\pi y$ and $2\pi/y$ cannot both be algebraic because their
product is $4\pi^2$.

For the parameter family
$F_a(y)=4\sum_{n\ge1}n^{-a}\sigma_{2a}(n)K_a(2\pi ny)$, the
Ramanujan--Guinand relation reads
\[
 F_a(y)-y^{-1}F_a(1/y)=
 \Omega(1+2a)(y^{-1-a}-y^a)+
 \Omega(2a)(y^{-1+a}-y^{-a}),
\]
interpreted by removable limits at $a=0,\pm1/2$.  At $a=1/2$ the
expansions
$\Omega(1+t)=t^{-1}+c/2+O(t)$ and
$\Omega(t)=-t^{-1}+c/2+O(t)$ cancel the simple-pole coefficient carrying
$\gamma$.  Since $K_{1/2}(x)=\sqrt{\pi/(2x)}e^{-x}$, the resulting
identity is the Dedekind eta transformation
$\log\eta(i/y)=\tfrac12\log y+\log\eta(iy)$, not a rational formula
for $\gamma$.  Recovering the order derivative gives instead
\begin{equation}\label{ar:besseldegree}
 \left.\partial_\nu K_\nu(x)\right|_{\nu=1/2}
 =\sqrt{\frac\pi{2x}}e^xE_1(2x)
\end{equation}
\cite{ar:DLMFBessel}.  The differentiated full reciprocity relation is
preserved in the companion source note; it contains additionally
$\zeta'(2)/\zeta(2)$ and an $E_1$ series.  Its termwise $E_1$ coefficient
cannot be removed by a fixed finite same-scale half-integer order
combination.  This last assertion concerns the displayed functions, not
linear independence of their specialized values.

\subsection{Fixed rational maps amplify existing denominator costs}
\begin{proposition}[Fixed-map height and error]
\label{ar:rationalmap}
Let $R=A/B\in\Q(X)$ be nonconstant with coprime integer polynomials,
$D=\max(\deg A,\deg B)$.  Suppose $B(\alpha)\ne0$ and $R$ has a zero
of order $m\ge1$ at $\alpha\in\R$.  If $x_n=a_n/q_n\to\alpha$ are
reduced, $q_n\to\infty$, $x_n\ne\alpha$, then
\[
 \operatorname{den}(R(x_n))\asymp q_n^D,\qquad
 |\operatorname{num}(R(x_n))|\asymp q_n^D|x_n-\alpha|^m.
\]
If $\alpha$ is algebraic of degree $d$, then $D\ge dm$.
In particular $q_n|x_n-\alpha|\to\infty$ implies divergence of the
reduced numerators after every such fixed transformation.
\end{proposition}
\begin{proof}
Homogenize $A,B$ to degree $D$.  Their common divisor at a primitive pair
$(a,q)$ is uniformly bounded: Bezout identities in the original coordinate
bound it by $Cq^L$, and identities for the reciprocal polynomials bound it
by $C'a^{L'}$.  Coprimality of $a,q$ gives a fixed bound $|CC'|$.
Moreover $B_h(a_n,q_n)=q_n^DB(x_n)\sim q_n^DB(\alpha)$.
Taylor expansion gives $|R(x_n)|\asymp|x_n-\alpha|^m$.  If $\alpha$ is
algebraic, its minimal polynomial divides $A$ with multiplicity at least
$m$.  Finally $q_n^D|x_n-\alpha|^m=
q_n^{D-m}(q_n|x_n-\alpha|)^m$.
\end{proof}
For a polynomial $P$ with leading coefficient $A_d$, every denominator
prime $p\nmid A_d$ survives with valuation multiplied by $d$; uniformly,
$q^d/|A_d|^d\le\operatorname{den}(P(a/q))\le q^d$ when the value is
nonzero.  A hypothetical minimal polynomial has a simple root at $\gamma$,
so $P(x_n)$ gains only one power of the approximation error, not $d$.
Algebraic norms have the same local denominator amplification.  These
statements apply after an endpoint is formed, not to increasing-degree
maps or simultaneous coefficient design before integration.

Positivity also cannot be transported through arbitrary conjugations.
For $\alpha=\sqrt2$,
$\int_0^1t^n(1+\alpha t)dt=(n+1)^{-1}+\alpha(n+2)^{-1}$ is a positive
Hausdorff sequence, whereas its conjugated second moment is
$1/3-\sqrt2/4<0$.  An argument transporting the real positive model to
all algebraic conjugates therefore needs an additional arithmetic theorem.

\subsection{Prime defects of the polyexponential trace tower}
Let
\[
 F_s(z)=\sum_{n\ge1}\frac{(-1)^{n+1}z^n}{n^s n!},\qquad
 P_{s,n}(X)=-n[z^n]\log(1-XF_s(z)),\quad s\ge1.
\]
For every prime $p$,
\begin{equation}\label{ar:traceprime}
 p^sP_{s,p}(X)\in\Z_{(p)}[X],\qquad
 p^sP_{s,p}(X)\equiv-X\pmod p.
\end{equation}
Indeed expand $p\sum_{j=1}^pX^j[z^p]F_s^j/j$.  Only $j=1$ survives
multiplication by $p^s$ and reduction; its residue is $-X$ by Wilson's
theorem.  Thus at every nonzero algebraic level $c$, the trace
$P_{s,p}(c^{-1})$ has valuation $-s$ at every sufficiently large good prime
(normalizing $v_{\mathfrak p}(p)=1$).  Fixed multipliers and rescalings
cannot make the ordinary trace sequence globally bounded.

More strongly, let a globally bounded algebraic-coefficient series have
first nonconstant term $R(w)=a_0+a_rw^r+\cdots$, $a_r\ne0$.
If $R(F_s(z))=\sum b_nz^n$, then at all sufficiently large good primes
\begin{equation}\label{ar:compositionprime}
 p^{s+1}b_{p+r-1}\equiv-ra_r\pmod{\mathfrak p},\qquad
 v_{\mathfrak p}(b_{p+r-1})=-(s+1).
\end{equation}
The only composition involving an index at least $p$ is one of the $r$
permutations of $(p,1,\ldots,1)$ in $F_s^r$; all larger outer powers use
indices below $p$.  This proves \eqref{ar:compositionprime}.
Removing the first $R\ge1$ outer logarithmic terms merely moves the defect:
for $U_{s,n}^{[R]}=n\sum_{j=R+1}^nX^j[z^n]F_s^j/j$ and $p>R+1$,
\[
 p^{s+1}U_{s,p+R}^{[R]}(X)\equiv-RX^{R+1}\pmod p.
\]
These are obstructions to global boundedness and ordinary integral
realizations, not impossibility theorems for capacity criteria that allow
explicit denominator costs.

Order of operations matters for Borel conversion.  Multiplication of the
coefficients of $F_s$ by $n!$ gives $-\operatorname{Li}_s(-z)$, a
$G$-function.  But the ordinary trace generating function at $s=1$ is
\[
 H_X(z)=\sum_{n\ge1}P_{1,n}(X)z^n
       =\frac{Xz\Ein'(z)}{1-X\Ein(z)},\qquad X\ne0,
\]
and is not $D$-finite.  The entire order-one function $\Ein-X^{-1}$ has
infinitely many zeros: otherwise Hadamard factorization would make it
$e^{az+b}P(z)$, whose derivative has finitely many zeros, contrary to
$\Ein'(z)=(1-e^{-z})/z$.  Its logarithmic derivative has a pole at every
such zero.  A $D$-finite function has only finitely many finite
singularities.  Multiplying or dividing a coefficient sequence by $n!$
preserves $P$-recursiveness in both directions \cite{ar:Stanley}; therefore
neither factorial transform of this full trace series is $D$-finite.

\subsection{Finite algebraic specializations and filtered local data}
\begin{proposition}[No exceptional Euler shift at a finite algebraic point]
\label{ar:Einshift}
For every $c\in\C$, $F_c=\Ein-c$ has minimal inhomogeneous differential
order $2$, with $zF_c''+(z+1)F_c'=1$, and minimal homogeneous order $3$.
For every algebraic $c$ and nonzero algebraic $\alpha$, the numbers
$e^{-\alpha},F_c(\alpha)$ are algebraically independent over $\Qbar$.
Thus the algebraic-value set of $F_c$ among algebraic arguments is exactly
$\{0\}$; it has an algebraic zero only if $c=0$, at zero.
\end{proposition}
\begin{proof}
The derivative proves the equation.  Functional independence of $e^{-z}$
and $\Ein(z)$ follows along the positive axis: in a polynomial relation
$\sum_je^{-jz}P_j(z,\Ein(z))=0$, select the smallest nonzero $j$.
Since $\Ein(x)=\log x+\gamma+O(e^{-x}/x)$, that term has a nonzero
leading power-logarithmic asymptotic, while every subsequent exponential
is smaller.  Translation by $c$ preserves independence.  Hence
$1,F_c,F_c'$ are linearly independent over $\C(z)$, giving the asserted
minimal orders by the displayed equation and its derivative.
For algebraic $c$, the vector $(1,e^{-z},F_c)$ consists of $E$-functions
satisfying a system with only finite pole $0$.  Siegel--Shidlovskii
specialization at $\alpha\ne0$ preserves transcendence degree two
\cite{pm:Beukers2006}.  At zero $F_c(0)=-c$.
\end{proof}
This theorem applies to every algebraic shift, including a hypothetical
algebraic $\gamma$; it does not distinguish the finite part at infinity.
The exact Laplace formula
\[
 s\int_0^\infty e^{-sx}\Ein(bx)\,dx=\log(1+b/s),\qquad b,s>0,
\]
shows why a simple regularization loses that coordinate: the contribution
$\gamma$ cancels the universal correction
$s\int_0^\infty e^{-sx}\log x\,dx=-\log s-\gamma$.

A local filtration retains an extension parameter but does not force it
nonzero.  For
\[
 \phi e_0=e_0,\quad \phi e_1=ce_0+pe_1,\qquad
 \operatorname{Fil}^1D=\Q_p(e_1+ae_0),\quad
 \operatorname{Fil}^0D=D,\quad \operatorname{Fil}^2D=0,
\]
every $(a,c)\in\Q_p^2$ is weakly admissible.  The invariant
$\kappa=a-c/(p-1)$ is zero exactly for the filtered split extension.
Indeed the eigenlines are $\Q_pe_0$ and
$\Q_p(e_1+c e_0/(p-1))$, with Newton numbers $0,1$ and Hodge numbers
$0$ and either $0$ or $1$; total numbers are both $1$.
Under $e_1'=e_1+te_0$, one has $c'=c+(1-p)t$, $a'=a-t$, leaving
$\kappa$ unchanged.  Arbitrary local choices have not been identified with
a canonical Euler comparison realization.  The historical source's
suggested rapid-decay realization is retained as a research proposal, not
as an established relation-preserving map.

\section{Residue-first positivity, convex tensor tests, and two-pole escape}
\label{ar:residue}
Write $T(R)=\int_1^\infty\lfloor t\rfloor R(t)\,dt$ and
\[
 r_j(t)=(t+j)^{-2}-[t(t+1)]^{-1},\qquad H_0=0.
\]
For rational $A_0,\ldots,A_M$ and $q=\sum_jA_j$, the exact log-free
identity is
\begin{equation}\label{ar:residueidentity}
 R_A=\sum_jA_jr_j=O(t^{-3}),\qquad
 T(R_A)=q\gamma-\sum_jA_jH_j.
\end{equation}
To prove it, expand $\lfloor t\rfloor=\sum_{m\ge1}\mathbf1_{t\ge m}$
and integrate first.  The first $L$ primitives sum to
$\sum_jA_j(H_{L+j}-H_j)-q\log(L+1)$ and tend to the stated limit.
Absolute convergence follows from $O(t^{-3})$.

For $q=1$, put $z=1/t$ and
\[
 K_A(z)=\frac1{1+z}-\sum_{j=1}^M\frac{A_jj(2+jz)}{(1+jz)^2},
 \qquad D_M(z)=\prod_{j=1}^M(1+jz)^2.
\]
Then $R_A=z^3K_A$, and $N_A=D_MK_A$ is a polynomial of degree at most
$2M-1$.  Since $D_M>0$ on $[0,1]$, positivity is exactly a polynomial
interval question; nonnegative Bernstein coefficients are a sufficient,
not necessary, certificate.  For example
\[
 A=(17/48,5/6,-3/16),\quad
 \sum_jA_jH_j=53/96,\quad
 N_A(z)=\tfrac1{12}(1-z)^3+\tfrac32z^3>0,
\]
proves $\gamma>53/96$ with signed coefficients.  The upper vector
$(13/44,10/11,-9/44)$ gives $\gamma<53/88$.
The error is exactly
$\int_0^1z\lfloor1/z\rfloor N_A(z)/D_M(z)\,dz$, which admits rational
upper bounds by interval subdivision.

The archived finite optimization proposed $16$ bounds and certified $14$;
the degree-$M=10$ upper and $M=12$ lower proposals failed exact sign checks.
For $M=2,4,6,8$, the certified lower endpoints had denominator bit lengths
$7,14,53,85$ and integer errors exceeding $2,86,10^{13},10^{22}$,
respectively.  These are feasibility certificates, not global optimum
certificates.  No asymptotic impossibility is inferred from them.

\subsection{Why the positive-semidefinite coefficient cone is narrower}
Suppose a symmetric matrix $X\succeq0$ has $X\mathbf1=0$, trace $Q>0$,
and all logarithms in its floor-integral transform cancel.  Its remaining
value is $Q\gamma-A$, $A=\sum_jH_jX_{jj}$.  Then
\begin{equation}\label{ar:psdgap}
 \gamma-A/Q>2^{-25}.
\end{equation}
Indeed normalize $Q=1$ and write $X_{jk}=\langle v_j,v_k\rangle$,
$h(t)=\sum_jv_j/(t+j)$.  The integral is $\int_1^\infty\lfloor t\rfloor
\|h(t)\|^2dt$.  Since $A<\gamma<3/5$, $\|v_0\|^2>2/5$.
Cauchy--Schwarz and $\zeta(2)-1<13/20$ imply
\[
 \|h(1)\|>\sqrt{2/5}-\sqrt{39/100}>1/128,
 \qquad \|h'(t)\|<2.
\]
Thus $\|h(t)\|>1/256$ on $[1,1+1/512]$, giving \eqref{ar:psdgap}.
This excludes convergence in that PSD subclass; signed residue kernels
are outside its hypotheses.

The natural Lagrange-quadrature subclass is classical.  With
$G_n=\int_0^1\binom xn\,dx$ and
$C_n=\int_0^1x(x+1)\cdots(x+n-1)\,dx/n!$, $C_0=1$, it gives
\[
 \ell_M=\sum_{n=1}^M\frac{|G_n|}{n}\uparrow\gamma,\qquad
 u_M=1-\sum_{n=1}^M\frac{C_n}{n(n+1)}\downarrow\gamma,
 \qquad u_M-\ell_M=\frac{C_M}{M+1}.
\]
These are Fontana--Mascheroni/Cauchy formulas \cite{ar:Blagouchine}; their
rediscovery is not a new irrationality result.

\subsection{The complete acceptance set of every fixed residue basis}
Let $h=r_0-r_1>0$, $\nu(dt)=\lfloor t\rfloor h(t)dt$ and
$v_j=r_j/h$.  Then $\nu$ is a full-support probability measure on the
compactification $T=[1,\infty]$, and
\[
 v_0-v_1=1,\quad v_0(t)=\frac{t+1}{2t+1},\quad
 v_j(\infty)=\tfrac12-j,\quad
 \int v_j\,d\nu=\gamma-H_j.
\]
For finite $M\ge1$, define
\[
 C_M=\operatorname{conv}\{(v_0(t),\ldots,v_M(t)):t\in T\},\quad
 I_M=\{\alpha:(\alpha-H_j)_{j=0}^M\in C_M\}.
\]

\begin{theorem}[All tensor degrees inherit the linear interval]
\label{ar:convextensor}
Each $I_M$ is a compact nondegenerate interval contained in $[1/2,2/3]$,
with $\gamma$ in its interior.  Every nonnegative rational-coefficient
residue tensor
\[
 K_C(\boldsymbol t)=\sum_{j_1,\ldots,j_d}C_{j_1\cdots j_d}
                  \prod_{i=1}^d r_{j_i}(t_i)
\]
has polynomial integral
$P_C(\gamma)$, where
$P_C(X)=\sum C_{j_1\cdots j_d}\prod_i(X-H_{j_i})$.
It satisfies $P_C(\alpha)\ge0$ for every $\alpha\in I_M$, and strictly
$>0$ on the interior if the kernel is nonzero.  Conversely every
$\alpha\notin I_M$ fails a rational-coefficient positive linear test.
Thus the acceptance sets for all tensor degrees and for all linear tests
are exactly the same set $I_M$.
\end{theorem}
\begin{proof}
The rational functions $r_0,\ldots,r_M$ are linearly independent: their
unique double poles first force each coefficient for $j\ge2$ to vanish,
and then the double poles at $0,-1$ force the remaining two.  Consequently
the affine span of $C_M$ is exactly $x_0-x_1=1$.
The full-support barycenter $y(\gamma)=\int v\,d\nu$ lies in its relative
interior.  Otherwise a proper supporting affine functional would be
nonnegative on the curve and have integral zero, forcing it to vanish on
the entire affine span.  Intersecting with the line $y(\alpha)$ proves
nondegeneracy and interiority of $I_M$.

Every interior point has a full-support representing measure: mix $\nu$
with a measure representing a farther point on the same interval.  For
an arbitrary point of $I_M$ a representing probability measure suffices.
Integrating the normalized tensor $K_C/\prod_i h(t_i)$ against its product
measure gives $P_C(\alpha)$; full support makes a nonzero nonnegative
continuous kernel strictly positive in integral.

For a point outside $C_M$, strict convex separation gives an affine
functional nonnegative on $C_M$ and negative at the point.  Use
$v_0-v_1=1$ to absorb its constant.  Adding a sufficiently small positive
multiple of that constant function makes positivity uniform, while
negativity at the point persists.  Rational approximation of the finite
coefficient vector then produces the required rational linear test.
\end{proof}

For $M=1$, the theorem can be sharpened to an exact coefficient criterion:
a polynomial $P$ of degree at most $d$ has a nonnegative tensor
representation precisely when its degree-$d$ Bernstein coefficients on
$[1/2,2/3]$ are nonnegative.  Symmetrization makes the multiaffine lift
unique, and its vertex values are those coefficients.  In particular the
lift of $(X-a)^2$, $a\in(1/2,2/3)$, is $(x_1-a)(x_2-a)$, which changes
sign despite its nonnegative diagonal.

For $M=2$ one obtains explicitly
\[
 I_2=[\beta,53/88],\quad
 8\beta^3-27\beta^2+18\beta-3=0,\quad \beta\in(1/2,2/3).
\]
To compute it, use $x=v_0$ and
$v_2=x(x-1)(5x-1)/(3x-1)^2$, a strictly concave function on this interval.
Intersect its graph and endpoint chord with $y=x-3/2$.
The intervals are nested and the classical quadrature bounds show
$\bigcap_MI_M=\{\gamma\}$, an analytic identification with no arithmetic
classification.  If $I_M=[\ell,u]$ and
$\epsilon_M=\min(\gamma-\ell,u-\gamma)>0$, any positive nonzero integer
linear form represented in this fixed basis obeys
\[
 q\gamma-p\ge\epsilon_M|q|\quad(q\ne0),\qquad
 -p\ge1\quad(q=0).
\]
There is no claim of a uniform positive $\epsilon_M$ as $M\to\infty$.

\subsection{Two fixed poles with increasing order}
Setting tensor variables equal changes the integral.  Direct partial
fractions give
\[
 T(r_0^2)=4\gamma-1-\zeta(2)+\tfrac13\zeta(3),\quad
 T(r_0r_1)=6\gamma-\tfrac72,
\]
\[
 T(r_1^2)=4\gamma-\tfrac{13}3+\zeta(2)+\tfrac13\zeta(3).
\]
Cancelling the displayed $\zeta(2),\zeta(3)$ coefficients in their span
leaves only the cross product.  This is formal coefficient cancellation,
not an assumption about independence of the actual constants.

For $k\ge2$, put $d_k(t)=t^{-k}-(t+1)^{-k}$.  Telescoping primitives give
$T(d_k)=1/(k-1)$.  Consequently
\begin{equation}\label{ar:twopole}
 R_{q,B}=qr_0-\sum_{k=2}^LB_kd_k,
 \qquad T(R_{q,B})=q\gamma-\sum_{k=2}^L\frac{B_k}{k-1}.
\end{equation}
This is the complete formally zeta-free space of rational functions with
poles only at $0,-1$, orders at most $L$, and decay $O(t^{-3})$.
Indeed writing its coefficients as $a_k/t^k+b_k/(t+1)^k$ gives
$a_1+b_1=0$, $a_2+b_2=b_1=q$ and
\[
 T(R)=q\gamma-b_2+
 \sum_{k=3}^L\frac{(a_k+b_k)\zeta(k-1)-b_k}{k-1}.
\]
The formal zeta coefficients vanish exactly when $a_k+b_k=0$ for $k\ge3$,
which is precisely \eqref{ar:twopole}.  The null directions
$(k-1)d_k-(\ell-1)d_\ell$ have zero integral and necessarily change sign
unless zero.  They can alter shape, but do not supply positive zero integrals.

For $u=1/t$, positivity is equivalent to nonnegativity on $[0,1]$ of
\begin{equation}\label{ar:twopoly}
 N_{q,B}(u)=q(1+u)^{L-1}-\sum_{k=2}^LB_ku^{k-3}
       \{(1+u)^L-(1+u)^{L-k}\}.
\end{equation}
The apparent pole at $k=2$ cancels.  Its degree is at most $2L-3$ for
$L\ge3$ and at most one for $L=2$.

The arithmetic-first sublattice $B_k=(k-1)c_k$, $c_k\in\Z$, has
$p=\sum c_k\in\Z$.  It produced finitely many positive forms:
$(L,q,p)=(3,14,8),(5,14,8),(8,208,120),(10,532,307)$.
The first has the explicit certificate
\[
 14r_0-7d_2-2d_3=
 \frac{u^4(1-u)(1+2u)}{(1+u)^3}\ge0.
\]
For $L=8$ the coefficients are
$(c_2,\ldots,c_8)=(104,17,2,-8,10,-7,2)$; the exact Bernstein
certificate for \eqref{ar:twopoly} is nonnegative and nonzero.
After reducing by $8$ it gives $0<26\gamma-15<0.007608$.
These are finite certificates, with duplicate rational endpoints among
runs, not an infinite small-integer-form sequence.

\begin{theorem}[Explicit analytic escape with no arithmetic rate]
\label{ar:twosequence}
The two-pole space has gamma-independent rational endpoints
$a_m^-<\gamma<a_m^+$ for every $m\ge1$, using pole order $L=m^2+2$, such
that $a_m^+-a_m^-\le3/(8m)$.
\end{theorem}
\begin{proof}
Let $f(t)=\psi_1(t)-1/t=\sum_{n\ge0}r_0(t+n)$,
$F(u)=f(1/u)$, and $g(u)=F'(u)/u$, with $g(0)=1$.
Writing $v=u/(1+u)$ gives $F(u)-F(v)=r_0(t)$.
For $u>0$,
\[
 g(u)=2\sum_{n\ge0}(1+nu)^{-3}-u^{-1}
 =\int_0^\infty s^2e^{-s}
   \left(\frac1{1-e^{-us}}-\frac1{us}\right)ds.
\]
If $h(x)=(1-e^{-x})^{-1}-x^{-1}$, then
$0<h'(x)=x^{-2}-(4\sinh^2(x/2))^{-1}\le1/12$.
The upper bound follows by differentiating the partial fraction expansion
of $\coth$: each term
$2(4\pi^2a^2-x^2)/(4\pi^2a^2+x^2)^2$ is at most
$2/(4\pi^2a^2)$.  Thus $0\le g'\le1/2$.  The same telescoping integral gives
\[\gamma=\int_0^1(1-u)g(u)\,du.\]

Put $n=m^2$, $N=m^3$, $\delta_m=1/(4m)$ and $u_j=j/n$.
For $j>0$, define the rational quantities
\[
 s_j=2\sum_{k=0}^N(1+ku_j)^{-3}-u_j^{-1},\quad
 \ell_j=s_j+\frac1{u_j(1+(N+1)u_j)^2},\quad
 v_j=s_j+\frac1{u_j(1+Nu_j)^2},
\]
with $\ell_0=v_0=1$.  The integral test gives
$\ell_j\le g(u_j)\le v_j$ and
$v_j-\ell_j\le2/(1+m)^3\le\delta_m$.
Let $b_{n,j}(u)=\binom nj u^j(1-u)^{n-j}$ and set
\[
 p_- =\sum_j\ell_jb_{n,j}-\delta_m,\qquad
 p_+ =\sum_jv_jb_{n,j}+\delta_m.
\]
The Lipschitz bound and binomial variance give
$|g(u)-\sum_jg(j/n)b_{n,j}(u)|\le1/(4m)$, so
$p_-\le g\le p_+$ and $p_+-p_-\le3\delta_m$.
For $P_\pm(u)=\int_0^u sp_\pm(s)ds=\sum B_{\pm,k}u^k$, the
kernels $r_0-[P_\pm(u)-P_\pm(v)]$ have the required opposite signs.
Their rational endpoints are explicitly
\[
 a_m^- =\sum_{j=0}^n\frac{n-j+1}{(n+1)(n+2)}\ell_j-\delta_m/2,
 \quad
 a_m^+ =\sum_{j=0}^n\frac{n-j+1}{(n+1)(n+2)}v_j+\delta_m/2.
\]
These follow from
$\sum_kB_k/(k-1)=\int_0^1(1-u)p(u)du$.
The interval width is at most $3\delta_m/2$.
Strictness holds because $r_0$ has nonzero simple-pole residues and cannot
be a finite combination of the $d_k$.
\end{proof}
The denominator bit lengths for $m=1,2,3$ in this explicit rule are
$(7,5),(89,79),(642,627)$ for the two endpoints.  No estimate proves that
the denominator times the error tends to zero.  Maximal cancellation at
infinity in the same space instead returns the classical divergent
trigamma expansion
\[
 f(t)\sim\frac1{2t^2}+\sum_{n\ge1}\frac{\mathsf B_{2n}}{t^{2n+1}},
 \quad |\mathsf B_{2n}|=\frac{2(2n)!\zeta(2n)}{(2\pi)^{2n}}.
\]
Its first signed examples give $23/40<\gamma<7/12$; adding more terms at
fixed $t=1$ eventually worsens that particular rule.  This is not an
obstruction to every global positivity optimization.

\subsection{The exact Todd cost of integer-residue flatness}
\begin{proposition}[Integer quadrature flatness]
\label{ar:toddflat}
Let $A_j\in\Z$, $q=\sum_jA_j$, and suppose
$R_A(t)=O(t^{-L-3})$.  Then
\[
 T_L\mid q,\qquad T_L=\prod_{p\text{ prime}}p^{\lfloor L/(p-1)\rfloor}.
\]
For $M\ge L$ this is necessary and sufficient for the existence of integer
residues with that flatness, already supported on $0,\ldots,L$.
The condition $\sum A_jH_j\in\Z$ is separate.
\end{proposition}
\begin{proof}
Writing $\mu_r=\sum_jA_jj^r$, expansion gives coefficients
$(-1)^r((r+1)\mu_r-q)t^{-r-2}$; hence flatness means
$\mu_r=q/(r+1)$ for $0\le r\le L$.  In the integer-valued binomial basis
this is $\sum_jA_j\binom jk=qG_k$, where
$G_k=\int_0^1\binom xk dx$.
The matrix $(\binom jk)_{0\le k,j\le L}$ is unimodular triangular, so
integer residues exist exactly when all $qG_k$ are integral.

For clarity the required denominator identity is proved here.  The inverse
of $\sum_{j\ge0}(-1)^jz^j/(j+1)$ is $\sum G_nz^n$.
For each prime $p$, its coefficient recurrence gives
$v_p(G_n)\ge-\lfloor n/(p-1)\rfloor$ by induction, since
$v_p(j+1)\le\lfloor j/(p-1)\rfloor$.
When $n=m(p-1)$, the recurrence term with $j=p-1$ has valuation $-m$;
all other terms have valuation at least $-m+1$.
For $j=k(p-1)>p-1$ this uses $v_p(j+1)\le k-1$, because
$p^k-1>k(p-1)$; the remaining indices lose a floor unit.
Thus $v_p(G_{m(p-1)})=-m$, proving
$\operatorname{lcm}_{k\le L}\operatorname{den}(G_k)=T_L$.
\end{proof}
These are the classical Todd/Hirzebruch denominators
\cite{pm:BuchstaberVeselov2024}.  Legendre's formula gives $L!\mid T_L$.
For minimal support $M=L$, the unique normalized rule is the classical
Lagrange rule with endpoint $\ell_L=\sum_{k=1}^L|G_k|/k$.
Keeping integer residues and $q>0$ gives the actual lower bound
\[
 q\gamma-p\ge\frac{q}{16(L+1)^2H_L^2}
 \ge\frac{L!}{16(L+1)^2H_L^2}\longrightarrow\infty.
\]
Indeed on $0\le x\le(2H_L)^{-1}$,
$\prod_{j=1}^L(1-x/j)\ge1-xH_L\ge1/2$, yielding
$|G_{L+1}|\ge1/[16(L+1)H_L^2]$; use the next positive series term.
Subsequent primitive reduction can destroy the integer-residue hypothesis,
so this last assertion is not a lower bound for every reduced endpoint.
Nor is high flatness necessary for every small positive kernel.

\section{Ten further routes: completed algebra and the remaining value problem}
\label{ar:tenroutes}
This section preserves the unique results of the ten-route audit.  Route
R02, classification of $G$-function solutions to a rationally forced
first-order exponential equation, is incorporated in the later Euler
zero-division analysis.  It is not counted a second time here.  The other
routes address different transformations and do not combine automatically
into one arithmetic proof.

\subsection{R01: signed factorial-tail acceleration}
Put $X=2^N$, $Y_N=(N+1)\Ein(X)-N\Ein(2X)$, and
$A_m(x)=\sum_{k<m}(-1)^kk!x^{-k-1}$.  Then
\[
 Z_N=Y_N-(N+1)e^{-X}A_X(X)>\gamma
\]
eventually, with
\begin{equation}\label{ar:accelerated}
 Z_N-\gamma\sim(N+1)\sqrt{\frac\pi{2X}}e^{-2X}.
\end{equation}
Indeed integration by parts gives
$E_1(x)=e^{-x}A_m(x)+(-1)^mr_m(x)$,
$r_m(x)=m!\int_x^\infty e^{-t}t^{-m-1}dt>0$.
Since $X$ is even,
$Z_N-\gamma=(N+1)r_X(X)-NE_1(2X)$.
The substitution $t=X+s$ and Stirling's formula give
$r_X(X)\sim\sqrt{\pi/(2X)}e^{-2X}$, while
$E_1(2X)\sim e^{-2X}/(2X)$ is smaller.  If $m/X\to\theta>0$, the
exponential rate is $1+\theta-\theta\log\theta$, uniquely maximized at
$\theta=1$.

The primitive coefficient height in the specified coordinates
$(\Ein(X),\Ein(2X),e^{-X})$ is exactly
\begin{equation}\label{ar:accelheight}
 H_N=(N+1)2^{(N-1)X+N+1-v_2(N+1)}.
\end{equation}
For the consecutive summands of $A_X(X)$ the valuation differences are
$v_2(k+1)-N<0$, so its last summand uniquely determines its denominator,
$2^{(N-1)X+N+1}$.  Multiplication by $N+1$ gives the displayed clearing
factor; consecutive integers $N,N+1$ and the odd reduced final numerator
make the resulting vector primitive.  Since $0<A_X(X)<1/X$, its first
coefficient is largest.  Hence the analytic height ratio tends to
$2/\log X=0$.  These coordinates are $E$-values, so clearing their
coefficients does not give an integer linear form in $1,\gamma$.
For fixed nonzero $P\in\C[T]$, the new slope is
\[
 \lim_N\frac{-\log|P(Z_N)|}{2^N}
 =2\operatorname{ord}_{T=\gamma}P,
\]
by Taylor expansion.  This doubles an analytic slope without proving that
no rational polynomial vanishes at $\gamma$.

\subsection{R03: a quantitative all-prime budget}
The classical Euler--Gompertz Pad\'e sequences \cite{ar:MZ} satisfy
\[
 q_n=\sum_{k=0}^nk!\binom nk^2,\quad
 u_{n+1}=2(n+1)u_n-n^2u_{n-1},
 \quad(q_0,q_1)=(1,2),\quad(s_0,s_1)=(0,1),
\]
\[
 q_ns_{n+1}-q_{n+1}s_n=(n!)^2,\quad
 q_n\delta-s_n=n!\int_0^\infty\frac{t^ne^{-t}}{(1+t)^{n+1}}dt>0,
\]
\[
 n!\le q_n\le n!e^{2\sqrt n},\quad 0\le s_n<q_n,\qquad
 q_n\delta_p-s_n\in(n!)^2\Z_p,
 \quad \delta_p=\sum_{k\ge0}(-1)^kk!.
\]
The height bound follows from
$q_n/n!=\sum_{j=0}^n\binom nj/j!\le\sum_{j\ge0}n^j/(j!)^2
\le e^{2\sqrt n}$.

\begin{proposition}[Unmatched factorial mass]\label{ar:unmatched}
For reduced $\alpha=a/b\in\Q$, $b>0$, set
$\mathcal B_\alpha=\{p:\delta_p\ne\alpha\}$.  For every $n$ except
at most one,
\[
 \sum_{\substack{p\le n\\p\in\mathcal B_\alpha}}v_p(n!)\log p
 \ge\tfrac12\log(n!)-\sqrt n-\tfrac12\log(|a|+b).
\]
Consequently, for sufficiently large $n$,
\[
 \sum_{\substack{p\le n\\p\in\mathcal B_\alpha}}\frac{\log p}{p-1}
 \ge\tfrac12\log n-\tfrac12-\frac1{\sqrt n}
             -\frac{\log(|a|+b)}{2n}.
\]
\end{proposition}
\begin{proof}
The determinant identity makes $s_n/q_n$ strictly increasing.  Thus
$L_n=aq_n-bs_n$ vanishes at most once.  At a matched prime, $p\nmid b$
and $v_p(L_n)\ge2v_p(n!)$.  The product of those prime powers divides
the nonzero integer $L_n$, while
$|L_n|\le(|a|+b)n!e^{2\sqrt n}$.  Taking logarithms proves the first
bound.  Use $v_p(n!)\le n/(p-1)$ and $\log(n!)\ge n\log n-n$ for the second.
\end{proof}
This quantitative consequence of the classical Pad\'e method is stronger
than mere divergence of the unmatched prime weight; it does not assert
natural density or identify the real value with finite-place values.
For $P\in\Z[T]$ of degree $d$ and height $H$, the integer
$q_n^dP(s_n/q_n)$ vanishes at most $d$ times.  At a prime satisfying
$P(\delta_p)=0$ it lies in $(n!)^2\Z_p$, by subtracting homogeneous
polynomial evaluations, even if $q_n$ is not a $p$-adic unit.
Its archimedean upper bound is $(d+1)Hq_n^d$.  The resulting budget,
after division by $(n!)^2$, tends to zero for $d=1$ but is not small for
$d\ge2$; already for $d=2$ it is at least $3H(n+1)^2$.
This diagnoses those estimates, not all possible refinements.

\subsection{R04: minimal universal Bell--Gould filters}
Let $B_0=1$, $B_{n+1}=\sum_{k=0}^n\binom nkB_k$, and
$G_0=0,G_1=1$, with the same recurrence for $G_{n+1}$ when $n\ge1$.
A common monic filter modulo $M$ is $C\in\Z[X]$ with
$C(E)B_n=C(E)G_n=0\pmod M$ for every $n\ge0$, where $E$ is forward
shift.  Write $\mu(M)$ for its least degree.

\begin{proposition}\label{ar:bellgould}
Over $\mathbb F_p[X]$, the common annihilator ideal is
$((X^p-X)(X^p-X-1))$.  Thus $\mu(p)=2p$, and
\[
 \mu(M)=\max_{p^r\parallel M}\mu(p^r).
\]
If a known monic common recurrence modulo $p^r$ has degree $D$, a proposed
filter is universal exactly when it kills the first $D$ shifts of both
sequences.
\end{proposition}
\begin{proof}
The Bell--Gould recurrence identities, retaining the initial boundary term,
give $T(E)B=0$, $T(E)G=a$ over $\mathbb F_p$, where
\[
 T=X^p-X-1,\qquad
 a_n=\sum_{j=0}^{p-1}(-1)^{p-1-j}(p-1-j)!j^n,
\]
with $0^0=1$; these follow by applying the binomial recurrence
\cite{ar:Gould}.  All weights are nonzero.  Vandermonde independence of
$j^n$, beginning at $n=0$, makes the annihilator of $a$ equal to
$(X^p-X)$.  The polynomial $T$ is irreducible because a root has
$\theta^{p^k}=\theta+k$, a Frobenius orbit of length $p$.  Since $B$ is
nonzero its annihilator is $(T)$.  Every common filter thus contains both
coprime factors, whose product works.
Reduction proves the lower CRT bound; equalize local degrees by multiplying
by powers of $X$ and apply coefficientwise CRT for the upper bound.
Finally shift operators commute, so the filtered sequence retains the
known degree-$D$ monic recurrence, and its first $D$ zero values propagate.
\end{proof}
The finite exact systems retained in the source certificates give
\[
\begin{array}{c|rrrrrr}
M&8&64&81&125&12!&20!\\\hline
\text{previous degree}&10&20&21&30&32&56\\
\mu(M)&6&12&15&20&22&38
\end{array}
\]
For example $X^6+X^5-3X^4+3X^3-X^2+3X-2$ works modulo $8$ and gives
$28\delta-17$ after the integral normalization.  These finite certified
minima are not a formula for every $\mu(p^r)$.
If $p_N$ is the largest prime at most $N$, reduction gives
$\mu(N!)\ge2p_N$.  Combined with the previously constructed degree-$3N$
upper bound and the prime number theorem, this gives
$2\le\liminf\mu(N!)/N\le\limsup\mu(N!)/N\le3$.
No existence of that limit or small real values follows.  The $19$ retained
forms at $N=2,\ldots,20$ were certified nonzero, but the final value was
approximately $1.25\times10^{16}$.

\subsection{R05: a rational spectral edge with an actual localization cost}
For $0\le i,j<N$, let
\[
 A_N(i,j)=(i+j)!,\quad
 B_N(i,j)=(i+j)!\{H_{i+j}-i-j-1\},\quad
 \lambda_N=\max_{v\ne0}\frac{v^TB_Nv}{v^TA_Nv}.
\]
Both matrices are integral and $A_N$ is positive definite.

\begin{theorem}\label{ar:spectral}
One has $\lambda_N<\gamma-1$, $\lambda_N\le\lambda_{N+1}$ and
$\lambda_N\to\gamma-1$.  Equivalently,
$cA_N-B_N\succeq0$ for every $N$ exactly when $c\ge\gamma-1$; every
finite pencil is positive definite at the threshold.  In the standard
Laguerre basis the same eigenvalues are those of
\[
 T_N(i,j)=\begin{cases}
 H_i-2i-1,&i=j,\\
 \min(i,j),&|i-j|=1,\\
 -1/|i-j|,&|i-j|\ge2.
 \end{cases}
\]
The infinite matrix on finitely supported sequences is essentially
self-adjoint, with purely absolutely continuous spectrum
$(-\infty,\gamma-1]$ of multiplicity two.  Moreover, for
$\gamma_N=1+\lambda_N$,
\begin{equation}\label{ar:spectralgap}
 \gamma-\gamma_N\ge\frac1{64N^2(1+1/(4N))}.
\end{equation}
\end{theorem}
\begin{proof}
Gamma integrals and their derivatives identify the Rayleigh quotient with
the average of $f(x)=\gamma+\log x-x$ against $p(x)^2e^{-x}dx$.
Its unique maximum is $\gamma-1$ at $x=1$; a nonzero polynomial cannot
concentrate its mass at that single point.  Nested spaces give monotonicity.
Polynomials are dense in $L^2((1+|f|)e^{-x}dx)$: for a function
orthogonal to them, the measure $g(1+|f|)e^{-x}dx$ has a Laplace transform
analytic on $\Re z<1/2$ with all derivatives zero at zero; Fourier
uniqueness gives the zero measure.  Approximation near $x=1$ proves the
limit.

The Laguerre generating function yields
\[
 \int_0^\infty x^\alpha e^{-x}
 \sum_iL_i(x)t^i\sum_jL_j(x)s^jdx
 =\Gamma(\alpha+1)
 \frac{(1-t)^\alpha(1-s)^\alpha}{(1-ts)^{\alpha+1}}.
\]
Differentiation at $\alpha=0$ gives diagonal log moments $H_i-\gamma$
and off-diagonal moments $-1/|i-j|$.  Combining the Laguerre recurrence
for multiplication by $x$ gives $T_N$.  The same density proof for weight
$(1+f^2)e^{-x}$ shows the polynomial domain is a core.  Each interval
$(0,1)$ and $(1,\infty)$ maps bijectively onto $(-\infty,\gamma-1)$
under $f$, giving two positive absolutely continuous spectral measures.

The classical inequality $|e^{-x/2}L_j(x)|\le1$
\cite{ar:DLMFLaguerre} follows also from the normalized Hermite-function
identity $\langle h_j(\cdot+a),h_j\rangle=e^{-a^2/4}L_j(a^2/2)$ and
Cauchy--Schwarz.  Hence every normalized polynomial of degree below $N$
satisfies $e^{-x}|p(x)|^2\le N$.  At most half its mass lies in
$[1-\epsilon,1+\epsilon]$, $\epsilon=1/(4N)$.  Elsewhere
\[
 x-1-\log x\ge\epsilon-\log(1+\epsilon)
 \ge\epsilon^2/[2(1+\epsilon)].
\]
Integrating over at least half the mass proves \eqref{ar:spectralgap}.
\end{proof}
For $\ell_N=\operatorname{lcm}(1,\ldots,N-1)$, $\ell_1=1$,
$\det\{\ell_N(X-1)I_N-\ell_NT_N\}$ is an integer defining polynomial
of $\gamma_N$, with height at most $[\ell_N(1+5N)]^N$.
This is an upper bound for one defining polynomial, not a sharp minimal
height.  The first values are $\gamma_1=\gamma_2=0$, and $\gamma_3$ is
the largest root of $4X^3+14X^2+5X-1$.
The actual error bound \eqref{ar:spectralgap} excludes exponential
accuracy of raw finite sections in $N$, without excluding acceleration
or other operators.

\subsection{R06: reciprocal levels and one-zero excision}
The reciprocal construction in companion [001] supplies a real entire
function $H$ with
\[
 H(0)=2\operatorname{Cin}(1),\qquad
 H'(w)=\int_0^1e^{-sw}k(s)ds,\quad
 k(s)=J_0(2\sqrt{s(1-s)}),
\]
where $\operatorname{Cin}(1)=\int_0^1(1-\cos t)t^{-1}dt$.
Its inherited derivative-zero set is
$\{\pm2i\sqrt{\pi^2n^2-1}:n\ge1\}$, and its real restriction is an
increasing bijection.  These are inputs from the reciprocal construction.
Let $w_c$ denote the unique real solution $H(w_c)=2c$.

\begin{proposition}\label{ar:reallevels}
For every real $c$, all zeros of $H-2c$ are simple and every nonreal zero
lies in $\Re w<0$.  A right-half-plane zero occurs exactly when
$c>\operatorname{Cin}(1)$, and is then $w_c>0$.  A zero on the imaginary
axis occurs only for $c=\operatorname{Cin}(1)$, at zero.
\end{proposition}
\begin{proof}
For $u=2\sqrt{s(1-s)}\le1$, the Bessel series give
$J_0(u)\ge1-s+s^2$ and $0<J_1(u)\le u/2$ for $u>0$.
For $s\le1/2$ one has $k'(s)\le0$; for $s>1/2$,
$k'(s)\le2s-1$, so $sk'(s)-k(s)\le s^2-1<0$.
Thus $k(s)/s$ is strictly decreasing.  For $\sigma\ge0$, $y>0$,
\[
 \Im H(\sigma+iy)=\int_0^1e^{-\sigma s}\frac{k(s)}s\sin(ys)ds>0.
\]
Extend the decreasing weight by zero after $1$ and pair positive and
negative sine half-periods; all pairs are nonnegative and the first is
strictly positive.  The $1/s$ singularity is canceled by the sine.
Conjugation and the real bijection prove zero location.  The inherited
critical lattice is on the imaginary axis away from zero, where this same
inequality makes the $H$-value nonreal; real levels therefore have no
multiple zeros.
\end{proof}
For $c>\operatorname{Cin}(1)$ and $a=w_c$, removing the single upper
zero of the rotated function gives
\[
 \widehat E_c(z)=\frac{e^{-iz/2}(H(-iz)-2c)}{z-ia},\qquad
 \widehat\Theta_c(z)=\frac{z-ia}{z+ia}\,
 e^{iz}\frac{H(iz)-2c}{H(-iz)-2c}.
\]
The inherited Cartwright bounds are
$h_{\widehat E_c}(\theta)=\tfrac12|\sin\theta|$ and
$|\widehat E_c(x)|=O(\log(2+|x|)/(1+|x|))$.
All remaining zeros lie below the real axis.  Cartwright factorization,
together with $H(y)=\log y+O(1)$, $H(-y)\sim-e^y/y$, gives
$\widehat\Theta_c(iy)\sim-1/(y\log y)$ and zero exponential mean type.
Thus its reflection quotient is nonconstant inner and $\widehat E_c$ is
Hermite--Biehler of type $1/2$.
Logarithmic differentiation at zero gives the convergent positive trace
\[
 \sum_{\substack{H(\rho)=2c\\\rho\ne w_c}}
 \frac{-2\Re\rho}{|\rho|^2}
 =1+\frac{\sin1}{\operatorname{Cin}(1)-c}+\frac2{w_c}.
\]
The last term records the excised zero.  At the Euler level,
$\operatorname{Cin}(1)=\gamma-\operatorname{Ci}(1)$ gives
$1-\sin1/\operatorname{Ci}(1)+2/w_*$.  The same structure holds at
algebraic levels such as $c=1$; recovering $w_*$ spectrally does not
establish its algebraicity.

\subsection{R07: exact evaluation images of integer divisor filters}
For $k\ge1$, $d\ge2$, put $M=4^{k+d}$,
$a_j=4^{d-j}D(4^{k+j})/2\in\Z$.  The constraints
$\sum_{j=0}^dc_j=q$, $\sum jc_j=-kq$ give a rational value $N/(Mq)$,
where $N=Mq/2+\sum a_jc_j$.  Define
\[
 s_j=a_j-2a_{j+1}+a_{j+2},\quad
 g=\gcd(|s_0|,\ldots,|s_{d-2}|),\quad
 P=M/2+(k+1)a_0-ka_1.
\]
Then the image for fixed integer $q$ is exactly
\begin{equation}\label{ar:evaluationimage}
 N\in Pq+g\Z,
 \qquad \operatorname{den}(N/(Mq))=M|q|/\gcd(|N|,M|q|).
\end{equation}
For $g=0$ the image is the singleton $\{Pq\}$.
Indeed $q(k+1,-k,0,\ldots)$ is a particular solution and
$e_j-2e_{j+1}+e_{j+2}$ form an integer basis of the homogeneous kernel,
by successive coordinate elimination.  Bezout gives \eqref{ar:evaluationimage};
for $g>0$ the three-row map has Smith invariants $1,1,g$.
For $(k,d)=(1,3)$ the actual values are
$(D(4),D(16),D(64),D(256))=(8,50,280,1466)$, $(s_0,s_1)=(16,13)$
and $P=240$.  Every $p/256$ is realized by
\[
 c=(2,-1,0,0)+(p-240)(-4,13,-14,5).
\]
Such surjectivity supplies no privileged approximation of $\gamma$;
smallness, primitive height, and nonvanishing remain separate.

\subsection{R08: exact Sondow prefix rank and a larger positive cone}
Retain $I_n=\binom{2n}{n}\gamma-A_n+L_n$ and put
\[
 w_{n,t}=2\binom nt^2(H_t-H_{n-t}),\qquad
 \lambda_{p,n}=\sum_{t=0}^nw_{n,t}v_p((n+t)!).
\]
An integer vector $u$ cancels every logarithm exactly when
$\sum_nu_n\lambda_{p,n}=0$ for all relevant primes, by unique
factorization after clearing rational exponents.

\begin{proposition}\label{ar:sondowrank}
For every $M\ge1$, the logarithm matrix on columns $1,\ldots,M$ has
rank $\pi(2M)$ and nullity $M-\pi(2M)$.  On a general window $N,\ldots,M$
its rank is at least the number of odd primes in $[2N-1,2M-1]$.
If a nonzero kernel vector has last supported index $m\ge2$, then
$2m-1$ is composite.
\end{proposition}
\begin{proof}
The weights are antisymmetric under $t\mapsto n-t$, negative before the
midpoint and positive afterwards.  Every proper tail
$\sum_{t=r-n}^nw_{n,t}$, $n+1\le r\le2n$, is strictly positive.
For $n<p\le2n$, this tail at $r=p$ equals $\lambda_{p,n}$; for $p>2n$
the entry is zero.  For each odd prime choose column $(p+1)/2$.  In
increasing order this is a triangular minor with positive diagonal.
For a prefix add row $2$ and column $1$, whose entry is $2$; all prime
rows are covered.  The same observation at $p=2m-1$ gives the final claim.
\end{proof}
Let $U(z)=\sum_{j=0}^{M-N}u_{N+j}z^j$.  Since
$r=x(1-x)y(1-y)\in[0,1/16]$,
\[
 \sum_{n=N}^Mu_nI_n=
 \iint_{(0,1)^2}\frac{r^NU(r)}{(1-xy)(-\log(xy))}\,dx\,dy.
\]
Positivity of $U$ on that interval proves nonvanishing.  If $d=M-N$,
the Bernstein coefficients of $U(t/16)$ are exactly
\[
 b_k=\sum_{j=0}^k\frac{u_{N+j}}{16^j}\frac{\binom kj}{\binom dj}.
\]
Their strict positivity is less restrictive than first-coefficient
dominance.  The retained certificate on $1,\ldots,8$ has
$\min b_k=17363115$, fails first-coefficient dominance, and gives
\[
 57429190525504380\gamma-
 \frac{426310477993846627277}{12870}>50000.
\]
This enlarges the certified positive cone, not the approximation quality.
The lower bound follows from $I_1>3/1024$, obtained by restricting its
integral to $[1/4,3/4]^2$.  That prefix's two-dimensional logarithm kernel
also maps with rank two to the $(Q,A)$-plane; rank is different from
shortness of vectors in that quotient.

An unconstructed sufficient family for irrationality would combine exact
logarithm cancellation, strict positivity, and
\[
 d_{2M}\sum_{n=N}^M|u_n|16^{-(n-1)}\longrightarrow0,
 \qquad d_{2M}=\operatorname{lcm}(1,\ldots,2M).
\]
Indeed $I_n\le16^{-(n-1)}I_1$ and denominators of $A_n$ divide $d_{2M}$.
This condition has not been proved for an infinite family and addresses
irrationality rather than all algebraic degrees.

\subsection{R09: exact primitive cost of two-trace recovery}
For $q\ge1$, $x=c^{-1}$, let $U=T_2^{(q)}(c)$ and $V=T_4^{(q)}(c)$ be
the universal polyexponential inverse traces.  Their explicit expressions
are
\[
 U=x^2-2^{-q}x,
\quad V=x^4-2^{1-q}x^3+
 \left(\frac2{3^{q+1}}+\frac1{2^{2q+1}}\right)x^2-\frac{x}{6\cdot4^q}.
\]
Set $D_q=6\cdot24^q$, $N_q=4^{q+1}-6^q-3^{q+1}$ and
$B_q=4\cdot8^q-3\cdot6^q$.  Then
\begin{equation}\label{ar:tracerecovery}
 N_qx=D_qV-D_qU^2-B_qU
\end{equation}
is primitive of coefficient height exactly $D_q$ and, after division by
$N_q$, the unique total-degree-at-most-two expression in $U,V$ identically
equal to $x$.  No affine expression works.
Indeed $V-U^2=A_qU+C_qx$, with
\[
 A_q=\frac2{3^{q+1}}-\frac1{2^{2q+1}},\qquad
 C_q=-\frac1{6\cdot4^q}+\frac2{3\cdot6^q}-\frac1{2^{3q+1}}.
\]
Clearing gives $D_qA_q=B_q$, $D_qC_q=N_q$.  The number $N_q$ is odd
and congruent to $1$ modulo $3$, hence coprime to $D_q$.
The strict inequalities $0<A_q<1$, $|C_q|<1$ give the height.
Comparison of degrees eight, six, four, two and one in $x$ successively
fixes all coefficients of any quadratic recovery expression.
The unnormalized residual is asymptotic to $-4^{-q}/(6c)$, but after
integerization it is $N_q/c$, of size asymptotic to $6^q/|c|$.
This is optimality only in the stated low-degree coordinate problem.
At the algebraic level $c=2$, for example,
$T_2^{(1)}(2)=0$, $T_4^{(1)}(2)=1/288$, and all universal traces are
rational.  Inverse-branch complexity therefore does not imply arithmetic
complexity of these values.

\subsection{R10: raw-moment recovery and determinant information loss}
For independent formal variables $g,d$ over a characteristic-zero field
$K$, let
\[
 a=(g+1)d,\quad b=(g+1)(1-d),\quad c=(g+\tfrac32)d.
\]
Then $g=a+b-1$, $d=2(c-a)$ and $K[a,b,c]=K[g,d]$.
The complete relation ideal is generated by $2(A+B)(C-A)-A$, as the
same formulas define an inverse in the quotient.  At the actual moments
this recovers $\gamma$ and $\delta$ with coefficients of height at most
two, but the moments themselves are not rational approximants.

The first two determinant coordinates are instead
\[
 U=(g+1)d,\qquad
 V=(2g+3)^2(2d-1)-\frac{2g^2+5g+5}{2}d^2.
\]
They satisfy
\begin{equation}\label{ar:compression}
 [K(g,d):K(U,V)]=4,\qquad K[U,V]\cap K[g]=K.
\end{equation}
For the first assertion substitute $g=U/d-1$ to obtain
\[
 2d^4+(U-4)d^3+(2U^2-16U+2V+2)d^2
       +(-16U^2+8U)d+8U^2=0.
\]
Over $K(U)$ this is $A(d)+2Vd^2$, with coprime coefficients since
$A(0)=8U^2\ne0$.  It is irreducible as a primitive linear polynomial in
$V$; Gauss's lemma, using its constant leading coefficient as a polynomial
in $d$, proves irreducibility over $K(U,V)$.  Its degree is four.
For the second assertion, give $U,V$ weights $1,2$.  If a nonconstant
$Q(U,V)$ has highest weighted degree $M$, its coefficient of $d^M$ after
substitution is
\[
 (g+1)^M\sum_{i+2j=M}q_{ij}
 \left(-\frac{2g^2+5g+5}{2(g+1)^2}\right)^j\ne0.
\]
The rational function in parentheses is nonconstant and hence
transcendental over $K$.  Thus no nonconstant such $Q$ belongs to $K[g]$.
These are universal identities, not restrictions on accidental specialized
relations or on larger coordinate sets.

Finally all mixed positivity tests admit an algebraic-coordinate model:
for $\alpha\ge\gamma-1$,
\[
 k!(\alpha+k+1-H_k)=\int_0^\infty x^ke^{-x}
                  (x-\log x+\alpha-\gamma)dx
\]
has a positive almost-everywhere density.  Keeping the actual positive
Gompertz measure preserves the mixed Hankel positivity at $\alpha=0$.
The measure has changed; it does not preserve the defining unshifted Euler
integral.  This establishes insufficiency of those positivity tests alone,
not a contradiction in either arithmetic world.

% END INTEGRATED PART IV

\clearpage
\part{Infinite polynomial families and exact targets}\label{master:infinite}
% BEGIN INTEGRATED PART V
\section{Infinite polynomial families and exact arithmetic targets}
\label{inf:sec:families}

This section integrates the polynomial-family records of 9--10 September 2026. Its principal distinction is between a proved criterion, a terminating finite certificate, and a still unproved assertion that certificates succeed at unbounded complexity. Throughout, $H(P)$ is the maximum absolute coefficient of an integer polynomial. No criterion below is asserted to be satisfied by $\gamma$ at unbounded complexity.

\subsection{The degree-sensitive target is an equivalence}
\label{inf:sub:criterion}
For a nonzero polynomial $P$ with $P(x)\ne0$, put
\[
 B(P)=\log\max(2,H(P))+\deg P,
 \qquad \mathcal R_x(P)=\frac{-\log|P(x)|}{B(P)}.
\]

\begin{lemma}[Resultant lower bound]\label{inf:lem:resultant}
Let $\alpha$ be algebraic of degree $d$, with primitive minimal polynomial $f$, leading coefficient $a\ne0$, and conjugates $\alpha=\alpha_1,\ldots,\alpha_d$. Set
\[
 C_\alpha=|a|\prod_{i=2}^d\max(1,|\alpha_i|).
\]
For every $P\in\Z[X]$ of degree $r$ with $P(\alpha)\ne0$,
\begin{equation}\label{inf:eq:resultant}
 |P(\alpha)|\ge C_\alpha^{-r}\bigl((r+1)H(P)\bigr)^{-(d-1)}.
\end{equation}
If $0<\alpha<1$, then $C_\alpha=M(f)$, where $M$ denotes Mahler measure.
\end{lemma}
\begin{proof}
The integer $\operatorname{Res}(f,P)=a^r\prod_iP(\alpha_i)$ is nonzero. For $i\ge2$,
$|P(\alpha_i)|\le(r+1)H(P)\max(1,|\alpha_i|)^r$. Substitution into $|\operatorname{Res}(f,P)|\ge1$ proves the assertion.
\end{proof}

\begin{theorem}[Equivalent infinite-family target]\label{inf:thm:equivalence}
A real number $x$ is transcendental if and only if there are integer polynomials $P_n$, with $P_n(x)\ne0$, for which
\[
 \limsup_n\mathcal R_x(P_n)=\infty.
\]
For $x\in[0,1]$ this is also equivalent to existence, at unbounded integers $D\ge2$, of polynomials satisfying
\begin{equation}\label{inf:eq:box-target}
 \deg P_D,H(P_D)\le D,\qquad
 0<|P_D(x)|<2U_D,
 \quad U_D=\frac{D(D+1)}{(D+1)^{D+1}-1}.
\end{equation}
\end{theorem}
\begin{proof}
For algebraic $x$, Lemma~\ref{inf:lem:resultant} and $\log(r+1)\le r$ bound $\mathcal R_x(P)$ by a constant depending only on $x$.

Conversely, evaluate the $(D+1)^{D+1}$ polynomials with coefficient vectors in $\{0,\ldots,D\}^{D+1}$ at a transcendental $x\in[0,1]$. All values are distinct and lie in an interval of length at most $D(D+1)$. Two adjacent values therefore differ by a nonzero polynomial value of absolute value at most $U_D$. Since
\[
 \frac{-\log(2U_D)}{D+\log D}\sim\log D,
\]
these differences give the required family. For an arbitrary real $x$, the interval length is at most $D(D+1)\max(1,|x|)^D$; its additional logarithmic cost is only $O_x(D)$ and does not change the conclusion.
\end{proof}

\begin{remark}[Correction of an earlier formulation]
The earlier statement that this varying-degree criterion is stronger than necessary is replaced by Theorem~\ref{inf:thm:equivalence}. Fixed-degree or fixed-exponent criteria can be stronger than necessary, but the precise varying-degree condition above is an equivalence. Its converse uses transcendence to exclude collisions; it cannot be used to assume those collisions absent for $\gamma$.
\end{remark}

For linear forms, $0<|q_nx-p_n|\to0$, with integer $p_n,q_n$, implies irrationality. Roth's theorem gives transcendence if, for a fixed $\epsilon>0$, infinitely many distinct reduced fractions satisfy $0<|x-p_n/q_n|<q_n^{-2-\epsilon}$ \cite{pm:Roth1955}. A finite exponent improvement supplies neither infinitely many fractions nor the required uniform margin. In contrast, the varying-orbit criterion in Section~\ref{inf:sec:orbit} permits a specified vanishing margin; its normalization and hypotheses are different.

For completeness, a fixed-degree alternative uses the classical Nesterenko criterion. If, with $D$ fixed first,
\[
 \deg P_{D,n}\le D,\quad H(P_{D,n})\le e^{b_Dn+o(n)},\quad
 |P_{D,n}(x)|=e^{-a_Dn+o(n)},\qquad a_D,b_D>0,
\]
then
\[
 \dim_\Q\operatorname{span}_\Q(1,x,\ldots,x^D)\ge1+a_D/b_D
\]
by the nonoscillating case of \cite[Theorem 1]{inf:FischlerOsc}. The value asymptotic includes a lower estimate: positivity alone is insufficient. An unbounded set of ratios $a_D/b_D$ would prove transcendence, but no such family is obtained here.

\subsection{A finite certificate and a total output rule}
\label{inf:sub:canonical}
Suppose $[a,b]\subset[0,1]$ is a rational interval containing $x$, with midpoint $c$ and width $w$. For polynomials of degree at most $D$ and height at most $Q$, define
\[
 M_{D,Q}=\frac{QD(D+1)}2,
 \qquad U_{D,Q}=\frac{Q(D+1)}{(Q+1)^{D+1}-1}.
\]
Every such polynomial satisfies
\begin{equation}\label{inf:eq:derivative-budget}
 |P(x)-P(c)|\le M_{D,Q}w/2.
\end{equation}
Evaluate the entire coefficient grid $\{0,\ldots,Q\}^{D+1}$ at $c$, retaining repeated values. Let $g$ be the least adjacent gap.

\begin{proposition}[Finite separation and exact termination]\label{inf:prop:finite-test}
If $g>M_{D,Q}w/2$, every nonzero polynomial in the stated coefficient box is nonzero at $x$. If also $M_{D,Q}w/2<U_{D,Q}$, the minimum-gap difference satisfies
\[
 0<|P(x)|<2U_{D,Q}.
\]
Given rational enclosures of widths tending to zero, refinement until both tests hold terminates exactly when the box contains no nonzero polynomial vanishing at $x$.
\end{proposition}
\begin{proof}
Each polynomial in the symmetric coefficient box is a difference of two grid vectors, so its midpoint value has absolute value at least $g$. Equation~\eqref{inf:eq:derivative-budget} proves the first assertion. The minimum gap is at most $U_{D,Q}$, giving the second. If a relation exists, the corresponding midpoint difference is at most $M_{D,Q}w/2$, and the test cannot pass. Otherwise the finite collection of distinct evaluations at $x$ has a positive minimum gap $\delta$. The midpoint minimum gap is at least $\delta-M_{D,Q}w/2$; a sufficiently narrow interval passes both tests.
\end{proof}

Here is a completely specified total rule for $\gamma$. Let
\[
 G_D=(D+1)^{D+1},\quad M_D=\frac{D^2(D+1)}2,
 \quad 2^{k_D}>\frac{11M_D}{8U_D},
\]
where $k_D$ is the least positive integer satisfying the inequality. Put $N=2^{k_D}$ and, with $k=k_D$,
\[
 s_k=2\sum_{j=0}^{k-1}\frac1{(2j+1)3^{2j+1}},
 \qquad r_k=\frac3{4(2k+1)9^k},
\]
\begin{equation}\label{inf:eq:midpoint}
 a_D=H_N-k(s_k+r_k)-N^{-1},\quad
 b_D=H_N-ks_k,\quad c_D=(a_D+b_D)/2.
\end{equation}
The atanh series gives $s_k<\log2<s_k+r_k$, and harmonic integral comparison gives $0<H_N-\log N-\gamma<N^{-1}$. Consequently
\begin{equation}\label{inf:eq:midpoint-error}
 0<a_D<\gamma<b_D<1,\qquad
 b_D-a_D<\frac{U_D}{M_D},\qquad
 |c_D-\gamma|<\delta_D:=\frac1{D(G_D-1)}.
\end{equation}
Indeed $b_D-a_D=2^{-k}+kr_k<(11/8)2^{-k}$; the displayed choice of $k$ gives the width estimate. The elementary bounds $1/2<\gamma<2/3$, together with $k\ge6$, place these intervals inside $(0,1)$.

Order the grid polynomials by their rational values at $c_D$, breaking ties lexicographically by their coefficient vectors, and choose the lexicographically first pair realizing the least adjacent gap. Let $P_D$ be their difference, oriented to have nonnegative value at $c_D$. Then
\begin{equation}\label{inf:eq:minimum}
 |P_D(c_D)|=\min_{\substack{0\ne P\in\Z[X]\\\deg P,H(P)\le D}}|P(c_D)|,
 \qquad 0\le|P_D(\gamma)|<\tfrac32U_D.
\end{equation}
This rule always terminates, since it consists of finitely many rational operations. The assertion $P_D\ne0$ concerns its coefficient vector, not its value at $\gamma$.

\begin{proposition}[Full degree and primitivity]\label{inf:prop:full-degree}
For every $D\ge2$, the canonical output $P_D$ has degree $D$ and integer content $1$.
\end{proposition}
\begin{proof}
Bertrand's theorem supplies a prime $p$ with $N/2<p<N$ \cite{inf:Erdos}. The choice of $k$ implies $N/2>D$ and $N/2>2k+1$. In $H_N$, only the term $1/p$ has denominator divisible by $p$. The corrections in
\[
 c_D=H_N-ks_k-(2N)^{-1}-kr_k/2
\]
are $p$-integral, so $v_p(c_D)=-1$. Its reduced denominator exceeds $D$. The rational-root theorem therefore excludes a zero at $c_D$ of every nonzero polynomial of height at most $D$. Dividing a minimizer by content greater than $1$ would strictly decrease its nonzero value. If its degree were below $D$, multiplication by $X$ would also strictly decrease its value because $0<c_D<1$. Both contradict minimality.
\end{proof}

\subsection{Algebraic collapse, its effective onset, and the relation lattice}
\label{inf:sub:collapse}

\begin{theorem}[Exact dichotomy for the total family]\label{inf:thm:dichotomy}
If $\gamma$ is transcendental, every $P_D(\gamma)$ is nonzero. If $\gamma$ is algebraic, every sufficiently late $P_D(\gamma)$ is exactly zero. More generally, the latter conclusion holds at any algebraic $\alpha\in(0,1)$ for any family with degree and height at most $D$ and $|P_D(\alpha)|<2U_D$.
\end{theorem}
\begin{proof}
Let $d$ and $M=M(f)$ be the degree and Mahler measure of a hypothetical minimal polynomial. A nonzero value satisfies
\begin{equation}\label{inf:eq:collapse-bound}
 |P_D(\alpha)|\ge M^{-D}[D(D+1)]^{-(d-1)}.
\end{equation}
This is incompatible with $2U_D=\exp(-D\log D+O(\log D))$ for sufficiently large $D$. The transcendental case is the definition of transcendence.
\end{proof}

An exact sufficient condition for collapse is
\[
 (D+1)^{D+1}-1>2[D(D+1)]^dM^D.
\]
Two explicit, nonoptimal sufficient onsets are
\begin{align}
 D&\ge\max\{4d,\lceil4M^2\rceil\},\label{inf:eq:onset}\\
 D&\ge\left\lceil\max\{4M,16d\log(16dM)\}\right\rceil.\label{inf:eq:onset-linear}
\end{align}
For the first, use $2U_D<4(D+1)^{1-D}$; its ratio to the lower bound in \eqref{inf:eq:collapse-bound} is at most $4M^D/(D+1)^{D-2d+1}<4/2^D<1$. For the second, put $t=16dM$. At the prescribed onset $D_1$, one has $D_1+1\le t^2$ and $D_1\ge8d\log(D_1+1)$; monotonicity of $D/\log(D+1)$ extends this to larger $D$. The logarithm of the same comparison ratio is less than
\[
 \log4-D\log4+D/4<0.
\]
Since $M(f)\le\sqrt{d+1}H(f)$ by Jensen and Parseval, \eqref{inf:eq:onset} also yields the coefficient-only onset $D\ge\max\{4d,4(d+1)H(f)^2\}$. A finite nonvanishing certificate excludes a bounded region of degree \emph{and} height, not every algebraic number of a given degree.

\begin{proposition}[Saturated relation lattice and abundant collisions]\label{inf:prop:lattice}
For $D\ge d$, evaluation at algebraic $\alpha$ has kernel
\[
 \{P\in\Z[X]_{\le D}:P(\alpha)=0\}
   =f\Z[X]_{\le D-d}.
\]
This kernel is saturated of rank $D+1-d$. All nonzero Smith invariant factors of the matrix with columns $f,Xf,\ldots,X^{D-d}f$ equal $1$. Moreover, at most
\[
 V_D\le1+[D(D+1)]^dM^D
\]
distinct values occur on the coefficient grid $\{0,\ldots,D\}^{D+1}$.
\end{proposition}
\begin{proof}
Gauss's lemma proves the kernel formula. Saturation follows because $mP(\alpha)=0$ implies $P(\alpha)=0$. Alternatively, modulo every prime $p$, multiplication by the nonzero polynomial $f\bmod p$ is injective in $\mathbb F_p[X]$. Its columns remain independent, so no prime divides all maximal minors, proving the Smith assertion. The image has rank $d$, though for nonintegral $\alpha$ it need not have basis $1,\alpha,\ldots,\alpha^{d-1}$. For example, the image at $\alpha=1/2$ is $2^{-D}\Z$.

Distinct grid values are separated by the lower bound \eqref{inf:eq:collapse-bound} and lie in an interval of length $D(D+1)$. Counting separated points gives the bound for $V_D$.
\end{proof}
At least one fiber contains $\lceil G_D/V_D\rceil$ vectors. Subtracting one member produces at least $\lceil G_D/V_D\rceil-1$ distinct exact relations. Thus the algebraic branch permits many short relations; an assertion that an exact zero is unlikely does not contradict it.

\subsection{Canonical selection forces maximal multiplicity}
\label{inf:sub:multiplicity}
For a primitive irreducible polynomial $f$, define
\[
 \nu_D(f)=\max_{\substack{0\ne A\in\Z[X]\\\deg A,H(A)\le D}}v_f(A).
\]
Here $v_f(A)$ is the multiplicity of $f$ as a factor of $A$.

\begin{theorem}[Maximum admissible multiplicity]\label{inf:thm:multiplicity}
If $\gamma$ is algebraic with primitive minimal polynomial $f$ of degree $d$, then the canonical outputs satisfy
\[
 v_f(P_D)=\nu_D(f)
\]
for every sufficiently large $D$. In particular, with $L=\|f\|_1$,
\[
 f^{m_D}\mid P_D,\qquad
 m_D=\min\left\{\left\lfloor\frac{\log D}{\log L}\right\rfloor,
                  \left\lfloor\frac Dd\right\rfloor\right\}
     =\frac{\log D}{\log L}+O_f(1).
\]
\end{theorem}
\begin{proof}
Write $P_D=f^rB$ with $B(\gamma)\ne0$. Integer quotients follow from Gauss's lemma. The elementary Mahler estimates
\[
 H(B)\le2^{\deg B}M(B)\le2^DM(P_D)\le T_D:=2^DD(D+1)
\]
follow from expanding the product over roots and Jensen's unit-circle formula. They are uniform in $r$. Lemma~\ref{inf:lem:resultant} gives
\[
 |B(\gamma)|\ge\Lambda_D:=M(f)^{-D}((D+1)T_D)^{-(d-1)}
                         =\exp(-O_f(D+\log D)).
\]
By \eqref{inf:eq:midpoint-error},
\[
 |B(c_D)-B(\gamma)|\le\tfrac12D(D+1)T_D\delta_D
                  =\exp(-D\log D+O(D+\log D)),
\]
so $|B(c_D)|\ge\Lambda_D/2$ eventually.

Choose an admissible competitor $A=f^sC$ with $s=\nu_D(f)$. Again $H(C)\le T_D$. If $r<s$, exact minimality and $f(c_D)\ne0$ give
\[
 |B(c_D)|\le |f(c_D)|^{s-r}|C(c_D)|
 \le S(f)\delta_D(D+1)T_D,
 \qquad S(f)=\sum_{j\ge1}j|f_j|.
\]
For large $D$ the last expression is less than $\Lambda_D/2$, a contradiction. Thus $r=s$. Finally $f^{m_D}X^{D-dm_D}$ is admissible, since its height is at most $L^{m_D}\le D$; hence $\nu_D(f)\ge m_D$.
\end{proof}

Let $J_D=\gcd_{\Q[X]}(P_D,P_D')$, normalized to be monic. Under algebraicity, $f^{m_D-1}\mid J_D$ eventually. It follows that any one of the following assertions, if established on an unbounded set of canonical indices, proves transcendence:
\begin{itemize}
 \item $P_D$ is squarefree;
 \item a proved rational interval containing $\gamma$ has no real root of $J_D$;
 \item $\deg J_D=o(\log D)$;
 \item the number of roots of $J_D$ in such an interval, counted with multiplicity, is $o(\log D)$.
\end{itemize}
Sturm counts and squarefree factorization certify individual cases. Their unbounded success is unproved. If intervals may be freely refined, the second local criterion is equivalent in strength to transcendence: for a transcendental $\gamma$ each fixed $J_D$ has no zero at $\gamma$, whereas the algebraic branch forces the indicated repeated root.

The minimization rule is essential. At $x=1/2$,
\[
 (2X-1)(X^{D-1}-1),\qquad D\ge3,
\]
is primitive, squarefree, of degree $D$ and height $2$, yet vanishes at $x$. Squarefreeness plus the size bounds alone therefore does not suffice.

\subsection{Localized common factors, modular certificates, and infeasible jets}
\label{inf:sub:escapes}
For any target-small family, coprime finite blocks whose least index tends to infinity exclude eventual divisibility by one minimal polynomial. The more precise certificate is absence of a common real root in a proved interval containing $\gamma$. A persistent factor away from this interval is harmless. For example, for a transcendental $x\in[1/2,3/5]$, pigeonholing coefficient grids of degree $D-1$ gives $R_D$ with
\[
 0<|R_D(x)|<\frac{(5/2)D}{(D+1)^D-1}.
\]
Then $P_D=XR_D$ satisfies \eqref{inf:eq:box-target}, yet every member contains $X$. Thus global coprimality is not a necessary consequence of transcendence for arbitrary bounded families.

For fixed $N$, the monic gcd of $P_N,\ldots,P_M$ eventually stabilizes as $M$ increases, because it is a descending sequence of divisors of $P_N$. Equality at several consecutive stages cannot certify permanent stabilization: a later output can remove a factor. Under a hypothetical fixed ideal $(f)$, sums, products, Euclidean remainders and Bezout combinations remain in $(f)$. A Wronskian obeys
\[
 (fB)'(fC)-(fB)(fC)'=f^2(B'C-BC'),
\]
so it does not remove the relation either.

\begin{proposition}[Degree-preserving modular escape]\label{inf:prop:modular}
If reductions of integer polynomials $F_1,\ldots,F_s$ have gcd $1$ over $\mathbb F_p$, and at least one $F_i$ retains its degree, then their gcd over $\Q$ is $1$.
\end{proposition}
\begin{proof}
A primitive common factor $A$ divides every $F_i$ integrally. Degree preservation in one product forces the leading coefficient of $A$ to remain nonzero modulo $p$, so its reduction is a nonconstant common factor, a contradiction.
\end{proof}
The degree check cannot be omitted: $(pX+1)X$ and $(pX+1)(X+1)$ share $pX+1$ over $\Q$, but reduce to the coprime pair $X,X+1$. Alternatively, a nonzero reduction of the Sylvester determinant formed with the original degrees certifies coprimality.

Homogeneous congruences can be retained by pigeonholing. If a linear map $L:\mathbb F_p^{D+1}\to\mathbb F_p^r$ is prescribed, one residue fiber of $\{0,\ldots,Q\}^{D+1}$ has at least $\lceil(Q+1)^{D+1}/p^r\rceil$ members. Their nearest difference satisfies $L(P\bmod p)=0$ and
\[
 |P(x)|\le\frac{Q(D+1)}{\lceil(Q+1)^{D+1}/p^r\rceil-1}.
\]
This gives no prescribed nonzero residue. Selecting differences between distinct fibers requires a separate close-pair theorem. Moreover, a subset of $\mathbb F_p^{D+1}$ whose every nonzero difference has degree $D$ has at most $p$ members, since the leading-coefficient projection must be injective. Requiring every difference to have a full-degree irreducible reduction therefore destroys the large grid.

At the comparator $x=1/2$, every nonzero degree-$D$ integer polynomial value has size at least $2^{-D}$, while $2U_D<2^{-D}$ for $D\ge3$. Consequently every target-small output must be divisible by $2X-1$. Prescribing an odd leading coefficient, or irreducibility of degree at least $2$, is incompatible with the universal target there. These counterexamples concern constructions uniform in $x$; they do not exclude a future use of special identities of $\gamma$.

Differentiation escapes an ideal but does not preserve a small-value estimate. The stronger difficulty is the following feasibility obstruction.
\begin{proposition}[Identically stringent jets are impossible]\label{inf:prop:jet-obstruction}
Set $\epsilon_D=2U_D$ and let $r_D$ be the least positive integer such that $(5/3)^{r_D}\ge8D(D+1)$. For all sufficiently large $D$ and every $x\in[1/2,3/5]$, no nonzero polynomial of degree and height at most $D$ satisfies
\[
 |P^{(j)}(x)|<\epsilon_D\qquad(0\le j<r_D).
\]
\end{proposition}
\begin{proof}
Write $P=X^sR$ with $R(0)$ a nonzero integer, and put $K=r_D-1=O(\log D)$. Expanding $X^{-s}$ at $x$ gives, for $j\le K$,
\[
 \left|\frac{R^{(j)}(x)}{j!}\right|
 \le\eta:=\epsilon_D(K+1)(D+K)^K2^{D+K}.
\]
For $s=0$ the same upper bound is immediate. Remove the Taylor polynomial $T=\sum_{j=0}^K R^{(j)}(x)(X-x)^j/j!$ and put $Q=R-T$. Its coefficient norm satisfies
\[
 \|T\|_1\le\theta:=\eta(K+1)(8/5)^K
 =\exp(-D\log D+O(D+(\log D)^2))<1/2.
\]
Thus $Q$ has a root of multiplicity at least $r_D$ at $x$, $|Q(0)|\ge1/2$, and $\|Q\|_1\le D(D+1)+1/2$. The root product and Jensen's formula imply
\[
 \|Q\|_1\ge M(Q)\ge |Q(0)|x^{-r_D}
 \ge\tfrac12(5/3)^{r_D}\ge4D(D+1),
\]
a contradiction. Although $Q$ has real coefficients, these Mahler inequalities remain valid; integrality was needed only for $|R(0)|\ge1$.
\end{proof}

A formally sufficient degree-$d$ criterion demanding the same accuracy through derivative order $\lfloor D/d\rfloor$ is therefore empty on this interval for large $D$. Resultants would indeed force all these derivatives to vanish and hence a factor of degree exceeding $D$, but Proposition~\ref{inf:prop:jet-obstruction} shows that such a criterion cannot be used as an achievable universal construction. Unequal derivative budgets are outside this obstruction.

\subsection{Multiplicative recycling and common-denominator alternatives}
\label{inf:sub:products}
Define the additive cost $C(P)=\log M(P)+\deg P$ for nonconstant integer polynomials. The inequalities
\[
 H(P)\le2^{\deg P}M(P),\qquad M(P)\le(\deg P+1)H(P)
\]
show $C(P)\le2B(P)\le6C(P)$, while $C(PQ)=C(P)+C(Q)$.

\begin{proposition}[Primitive-factor reduction]\label{inf:prop:factor}
If $0<|P(x)|<1$, some primitive irreducible positive-degree factor $F$ of $P$ satisfies
\[
 \frac{-\log|F(x)|}{C(F)}\ge\frac{-\log|P(x)|}{C(P)}.
\]
Consequently the unbounded-ratio target is unchanged if all outputs are required to be primitive and irreducible.
\end{proposition}
\begin{proof}
Write $P=c\prod_iF_i^{m_i}$. The numerator equals $\sum_im_i[-\log|F_i(x)|]-\log|c|$ and the denominator equals $\sum_im_iC(F_i)+\log|c|$. The ratio is no larger than the greatest positive factor ratio. The comparison between $B$ and $C$ completes the reduction.
\end{proof}
Products and powers of a bounded-quality collection cannot amplify its normalized quality. This statement allows different factors and includes coefficient cancellation after expansion; it does not exclude addition or determinants that create new factors.

A different sufficient condition approximates powers simultaneously. Fix $D$ and suppose integer $q_n>0,p_{1,n},\ldots,p_{D,n}$ have nonzero errors
\[
 e_{j,n}=q_nx^j-p_{j,n},\qquad e_{1,n}\to0,\qquad e_{j+1,n}=o(e_{j,n}).
\]
Then $1,x,\ldots,x^D$ are linearly independent over $\Q$. Indeed an integer relation, multiplied by $q_n$, becomes an integer tending to zero and hence eventually zero. Dividing by the earliest error with a nonzero relation coefficient leaves that coefficient as a nonzero limit. Families for every fixed $D$ would prove transcendence. This elementary criterion is related to the vector criteria of \cite{inf:FischlerVectors}.

The rank-one moment curve imposes exact constraints. If $e_j=qx^j-p_j$, then
\begin{equation}\label{inf:eq:minor}
 \Delta_{i,j}=qp_{i+j}-p_ip_j
 =qx^ie_j+qx^je_i-qe_{i+j}-e_ie_j\in\Z.
\end{equation}
Under the stated hierarchy and $x\ne0$, $\Delta_{i,j}/(qe_i)\to x^j$ for $i<j$, and to $2x^i$ for $i=j$. These integers are eventually nonzero, implying
\[
 \liminf_nq_n|e_{i,n}|\ge|x|^{-j}\quad(i<j),\qquad
 \liminf_nq_n|e_{i,n}|\ge(2|x|^i)^{-1}\quad(i=j),
\]
whenever the required coordinates exist. In particular the earliest errors cannot also be $o(1/q_n)$.

Taking powers of one approximation $p_n/q_n$ with common denominator $q_n^D$ gives errors
\[
 jx^{j-1}q_n^{D-1}(q_nx-p_n)(1+o(1)),
\]
all on the same scale, and makes every minor \eqref{inf:eq:minor} vanish. A determinant descent is nevertheless an exact identity: for
\[
 q'=qp_2-p_1^2,\qquad p'=2(qp_3-p_1p_2),
\]
one has
\[
 q'x-p'=-3qxe_2+2qe_3-xe_1^2+2e_1e_2.
\]
Even when $q'\sim2xqe_1$ is smaller than $q$, the term $qe_2$ remains uncontrolled. No improved approximation follows without additional cancellation or divisibility.

\section{Positive realization and the complete two-pole integer lattice}
\label{inf:sec:positive}
This section uses the two-pole identities established in \eqref{ar:twopole}, restated here to fix normalization:
\[
\begin{gathered}
 T(R)=\int_1^\infty\lfloor t\rfloor R(t)\,dt,\quad
 r_0(t)=\frac1{t^2(t+1)},\quad d_k(t)=t^{-k}-(t+1)^{-k},\\
 T(r_0)=\gamma,\qquad T(d_k)=\frac1{k-1}.
\end{gathered}
\]
The new issue is the height of the \emph{output polynomial}, which need not include the denominators of exact null-mean corrections.

\subsection{The full lattice and its explicit basis}
\label{inf:sub:full-lattice}
Put $\ell(b)=\sum_{j=1}^mb_j/j$ and $d_m=\operatorname{lcm}(1,\ldots,m)$. The full integer lattice is
\[
 \Lambda_m=\{b\in\Z^m:\ell(b)\in\Z\}
 =\ker\left(b\longmapsto\sum_{j=1}^m(d_m/j)b_j\pmod{d_m}\right).
\]
The map is surjective, because the integers $d_m/j$ have gcd $1$: for each prime factor of $d_m$, choose $j$ equal to its largest power at most $m$. Thus $[\Z^m:\Lambda_m]=d_m$, whereas the coordinatewise sufficient sublattice $\operatorname{diag}(1,\ldots,m)\Z^m$ has index $m!$.

An explicit triangular basis is $v_1=e_1$ and, for $k\ge2$,
\[
 v_k=\begin{cases}
 pe_k-e_{k/p},&k=p^a,\\
 e_k-ae_u-be_w,&k=uw,\ (u,w)=1,\ u,w>1,\ aw+bu=1.
 \end{cases}
\]
Each $v_k$, $k\ge2$, has $\ell(v_k)=0$. The determinant of the column matrix is $\prod_{p^a\le m}p=d_m$, so these vectors form a basis of the entire lattice. Consequently
\[
 b=pv_1+\sum_{k=2}^mx_kv_k\quad(p,x_k\in\Z)
 \quad\Longrightarrow\quad T\left(qr_0-\sum_jb_jd_{j+1}\right)=q\gamma-p.
\]
The basis comparison is factorial versus exponential: $d_{2n}/d_n$ divides $\binom{2n}{n}$ by prime-power valuations, and iteration gives the sufficient bound $d_m<16^m$. This is a covolume fact, not a short-vector theorem for the intersection with a positivity cone.

For each prime $p\le m$, write $j=p^{s_j}u_j$ with $p\nmid u_j$ and $a=\lfloor\log_pm\rfloor$. The local integrality conditions are exactly
\[
 \sum_jb_jp^{a-s_j}u_j^{-1}\equiv0\pmod{p^a}.
\]
If $m/2<p\le m$, this forces $p\mid b_p$. For arbitrary integer $b$, with $S=\sum_j(d_m/j)b_j$,
\[
 \operatorname{den}\ell(b)=\frac{d_m}{\gcd(d_m,S)}.
\]
For a rational shape $\beta=(\beta_j)$ normalized by $q=1$, let $D_\beta=\operatorname{lcm}_j\operatorname{den}\beta_j$ and $p_0=\sum_j\beta_j/j$. The least integer multiplier making all residues and the rational output term integral is
\[
 Q_{\mathrm{full}}=\operatorname{lcm}(D_\beta,\operatorname{den}p_0).
\]
The coordinatewise rule uses $Q_{\mathrm{coord}}=\operatorname{lcm}_j\operatorname{den}(\beta_j/j)$. Primitive reduction of the final pair can be smaller still, so $Q_{\mathrm{full}}$ is not a lower bound on the primitive denominator.

\subsection{Null-mean corrections preserve output height}
\label{inf:sub:positive-lift}
The measure $d\eta(t)=\lfloor t\rfloor d_2(t)\,dt$ is a probability measure. Define
\[
 X(t)=\frac{r_0(t)}{d_2(t)}=\frac{t+1}{2t+1},\qquad
 V_k(t)=\frac{(k-1)d_k(t)}{d_2(t)}.
\]
Then $\mathbb EX=\gamma$ and $\mathbb EV_k=1$. The quantitative Bernstein construction of Theorem~\ref{ar:twosequence} gives rational $B_{m,k}$, $2\le k\le m^2+2$, for which
\[
 R_m=r_0-\sum_kB_{m,k}d_k,\quad
 0\le R_m\le\epsilon_md_2,\qquad\epsilon_m=\frac3{8m}.
\]
With $a_m=\sum_kB_{m,k}/(k-1)$, put
\[
 F_m=X-\sum_k\frac{B_{m,k}}{k-1}(V_k-1)
     =a_m+R_m/d_2.
\]
It follows exactly that $\mathbb EF_m=\gamma$ and $F_m\in[a_m,a_m+\epsilon_m]$; for $m\ge2$ this interval lies in $[0,1]$.

For $P(Z)=\sum_{j=0}^Dp_jZ^j$, define its symmetric multiaffine polarization
\[
 \mathcal P_D(x_1,\ldots,x_D)=\sum_{j=0}^D\frac{p_j}{\binom Dj}
      \sum_{|S|=j}\prod_{i\in S}x_i,
\]
and the rational tensor kernel
\[
 K_{m,P}(t_1,\ldots,t_D)=\prod_id_2(t_i)\,
                  \mathcal P_D(F_m(t_1),\ldots,F_m(t_D)).
\]
Independence yields the exact identity
\begin{equation}\label{inf:eq:positive-output}
 T^{\otimes D}(K_{m,P})=P(\gamma).
\end{equation}
The final polynomial remains $P\in\Z[Z]$, with its original degree and height; clearing all internal kernel denominators would charge a cost that the exact output identity does not require.

If $S(P)=\sum_{j\ge1}j|p_j|$, telescoping products give
\[
 |\mathcal P_D(x_1,\ldots,x_D)-P(a_m)|\le\epsilon_mS(P)
\]
throughout the range cube. Thus the rational test $P(a_m)>\epsilon_mS(P)$ proves strict positivity of the kernel. It eventually passes whenever $P(\gamma)>0$, since $|P(a_m)-P(\gamma)|\le\epsilon_mS(P)$. A sufficient conditional cost is $m>3S(P)/(4P(\gamma))$.

Conversely, a continuous nonnegative nonzero kernel has positive integral because the integration measure has full support. Therefore every nonzero integer polynomial has a positive lift for one sign if and only if $\gamma$ is transcendental. The construction is complete as a strict-sign semidecision procedure; proving termination for every nonzero $P$ would already solve the arithmetic question.

The same correction mechanism shows that the image of the positive two-pole cone is exactly the set of strict positive rational forms $q\gamma-p$, $q,p\in\Q$. For $q>0$, choose a certified rational lower endpoint $p/q<a<\gamma$ and a positive kernel with integral $\gamma-a$; multiplying it by $q$ and adding $(qa-p)d_2$ realizes the prescribed form. For $q<0$, use an upper endpoint and its positive error kernel; $q=0$ is immediate. Unconstrained existence of a positive kernel therefore adds no approximation exponent.

\subsection{A completed nonlinear family and a fixed-degree sum-of-squares bound}
\label{inf:sub:sos}
Let $A(Z)=26Z-15$ and $B(Z)=7Z-4$. The prior positive certificate gives $A(\gamma)>0$. Their exact relation is $7A-26B=-1$. The quadratic
\[
 Q(Z)=A(Z)^2+B(Z)^2=725Z^2-836Z+241
\]
has discriminant $-4$, hence is irreducible. With $a=15/26$ and a two-pole random variable $F\ge a$ of mean $\gamma$, its polarization is strictly positive:
\[
 725(a+u)(a+v)-418(2a+u+v)+241
 =\frac1{676}+\frac7{26}(u+v)+725uv>0\quad(u,v\ge0).
\]
Using $\gamma<5773/10000$ gives
\[
 0<Q(\gamma)<\frac{7141}{4000000}<1.
\]
Thus $Q^n$ has degree $2n$, height at most $1802^n$, and positive values bounded by $(7141/4000000)^n$. The same family is small at the rational point $577/1000$, where $Q=61/40000$. It does not prove irrationality. Likewise $(26X-15)^n$ has coefficient norm $41^n$ and equals $500^{-n}$ at that rational point. Proposition~\ref{inf:prop:factor} explains why powering does not improve normalized quality.

\begin{proposition}[Coprime integral sums of squares]\label{inf:prop:sos}
If $A,B\in\Z[X]$ are coprime over $\Q[X]$, $m=\max(\deg A,\deg B)\ge1$, and $P=A^2+B^2$, then for $0\le x\le1$,
\[
 P(x)\ge c_mH(P)^{-(2m-1)},
 \quad c_m=\left(((2m)!)^2C_m^{4m-2}(2m+1)^{2m-1}\right)^{-1},
 \quad C_m=(m+1)(2m)^m.
\]
\end{proposition}
\begin{proof}
A nonzero constant component gives $P\ge1$. Otherwise put $r=\deg A$, $s=\deg B$, and $H_0=\max(H(A),H(B))$. The adjugate of the Sylvester matrix gives $UA+VB=\operatorname{Res}(A,B)$ with degrees below $s,r$ and coefficient bounds $(r+s-1)!H_0^{r+s-1}$. Evaluation on $[0,1]$ yields
\[
 1\le(r+s)!H_0^{r+s-1}\sqrt{P(x)},
 \qquad P(x)\ge((2m)!)^{-2}H_0^{-(4m-2)}.
\]
Interpolate $A,B$ at $j/m$, $0\le j\le m$. Every Lagrange basis polynomial has coefficient norm at most $(2m)^m$, giving
\[
 H_0\le C_m\max_{[0,1]}\sqrt{P(x)}
      \le C_m\sqrt{(2m+1)H(P)}.
\]
Substitution proves the claimed bound.
\end{proof}
This is a fixed-degree result for coprime integral components. Its degree-dependent constant cannot be ignored when $m$ grows, and a common factor can retain the original small-value problem. For the displayed quadratic the sharper identity is $Q(x)=((725x-418)^2+1)/725$.

\section{An all-order Bernoulli--Stieltjes family and its rate obstruction}
\label{inf:sec:pade}
\subsection{Positive measure, exact Pad\'e bounds, and endpoint arithmetic}
\label{inf:sub:pade}
The two-pole kernel of Theorem~\ref{ar:twosequence} admits the gamma-free expression
\[
 g(u)=2\sum_{j\ge0}(1+ju)^{-3}-u^{-1},\quad g(0)=1,
 \qquad\gamma=\int_0^1(1-u)g(u)\,du.
\]
Let $b(t)=(e^{2\pi t}-1)^{-1}$ and $\rho(t)=2t^3b''(t)>0$. Then
\begin{equation}\label{inf:eq:stieltjes}
 H(v):=\frac{g(\sqrt v)-1}{\sqrt v}
       =\int_0^\infty\frac{\rho(t)}{1+vt^2}\,dt.
\end{equation}
To derive it, differentiate Binet's digamma integral \cite[5.9.10]{inf:DLMFGamma} twice and use $\psi''(x)=-2\zeta(3,x)$. The resulting formula is
\[
 \frac{g(u)-1}{u}=4\int_0^\infty
        \frac{t(3-u^2t^2)b(t)}{(1+u^2t^2)^3}\,dt.
\]
The rational factor is the second derivative of $2t^3/(1+u^2t^2)$. Two integrations by parts give \eqref{inf:eq:stieltjes}; boundary terms vanish using $b=O(t^{-1})$, $b'=O(t^{-2})$ at zero and exponential decay at infinity.

Push $\rho(t)dt$ forward by $x=t^2$ to a measure $\nu$. Its moments are rational:
\[
 \mu_k=2\int_0^\infty t^{2k+3}b''(t)\,dt
 =\frac{2(2k+3)(2k+2)!\zeta(2k+2)}{(2\pi)^{2k+2}}
 =(2k+3)|B_{2k+2}|.
\]
Both ordinary and shifted Hankel matrices are positive definite because $\nu$ and $x\nu$ have positive density. The moment problem is determinate: for $0<a<2\pi$, the series $\sum a^{2k}\mu_k/(2k)!$ converges. Any measure with these moments has finite $\cosh(a\sqrt x)$ integral. Symmetrizing its square-root pushforward gives a bilateral Laplace transform analytic near zero and fixed by the moments, hence fixes the measure. The same proof applies to $x\nu$. The formal moment series has radius zero; no convergence of that power series is asserted.

Let $L_n=[n-1/n]_H$ and $U_n=[n/n]_H$ for this formal moment series. Their coefficients are rational and all orders exist. If $p_n$ is the monic degree-$n$ orthogonal polynomial for $\nu$, the Gaussian quadrature formula gives
\begin{equation}\label{inf:eq:pade-error}
 L_n(v)=\sum_{j=1}^n\frac{w_j}{1+vx_j},\qquad
 H(v)-L_n(v)=\int\frac{(p_n(x)/p_n(-1/v))^2}{1+vx}\,d\nu(x)>0,
\end{equation}
with simple positive nodes and positive weights. Applying Gaussian quadrature to $x\nu$, with lower approximant $\widetilde L_n$, gives
\[
 U_n(v)=\mu_0-v\widetilde L_n(v)>H(v).
\]
This is Gauss--Radau quadrature including the node zero; its zero-node weight is $\int\widetilde p_n(x)^2d\nu(x)/\widetilde p_n(0)^2>0$, and other weights are the Gaussian weights divided by their positive nodes.

The quadrature measures have fixed mass, are tight by their first moments, and eventually reproduce each fixed moment. A higher fixed moment gives uniform integrability of lower powers. Every weak limit therefore has all the prescribed moments and equals the determinate measure. This proves $L_n(v),U_n(v)\to H(v)$ for $v>0$. Their transforms are bounded by $\mu_0$, so dominated convergence yields
\begin{equation}\label{inf:eq:pade-bracket}
 \alpha_n^-:=\int_0^1(1-u)(1+uL_n(u^2))\,du
 <\gamma<
 \alpha_n^+:=\int_0^1(1-u)(1+uU_n(u^2))\,du,
 \qquad\alpha_n^\pm\to\gamma.
\end{equation}

At a positive algebraic node $s$ the primitive integral is
\[
 \int_0^1\frac{u(1-u)}{1+su^2}\,du
 =\frac{\log(1+s)}{2s}-\frac1s+
          \frac{\arctan\sqrt s}{s\sqrt s}.
\]
All nodes and weights are algebraic. The logarithmic part is strictly positive, because the displayed logarithm and arctangent terms have positive weights. Writing the arctangent using logarithms of algebraic numbers makes it a nonzero algebraic-coefficient logarithmic linear form. Baker's theorem implies that this part, and hence each positive-order endpoint $\alpha_n^\pm$, is transcendental \cite{inf:Baker}. Their limit need not be transcendental. The first bounds are
\[
 -1+\frac34\log\frac43+\frac{\sqrt3\pi}4
 <\gamma<\frac7{18}+\frac{\log2}{12}+\frac\pi{24}.
\]
No integer-polynomial height may be assigned to these endpoints merely from their Pad\'e order.

\subsection{A quantitative all-order lower bound on the error}
\label{inf:sub:pade-rate}
\begin{theorem}[Rate obstruction for the unchanged lower family]\label{inf:thm:pade-rate}
For every $n\ge1$, writing $M=2n+1$,
\[
 \gamma-\alpha_n^-\ge
 \frac{3e^{-4/3}}{6400\pi^2M^4(\log M)^2}.
\]
The same lower bound applies to the full bracket width, but is not claimed for the individual upper error.
\end{theorem}
\begin{proof}
The error in \eqref{inf:eq:pade-error} is also the minimum of $\int q(x)^2(1+vx)^{-1}d\nu(x)$ over degree-$n$ polynomials with $q(-1/v)=1$. Indeed write a competitor as $p_n/p_n(-1/v)+(x+1/v)r$; the cross term vanishes by orthogonality.

Set $v=u^2$ and $Q(t)=q(t^2)$. It has degree at most $2n$ and $Q(i/u)=1$. By even extension,
\[
 H(u^2)-L_n(u^2)=\int_\R Q(t)^2w_u(t)\,dt,
 \qquad w_u(t)=\frac{\rho(|t|)}{2(1+u^2t^2)}.
\]
For $0<u\le1$,
\begin{equation}\label{inf:eq:hahn-comparison}
 w_u(t)\ge\frac1{2\pi}\operatorname{sech}^2(\pi t).
\end{equation}
For $t>0$, the ratio to $\operatorname{sech}^2(\pi t)$ is $\pi^2t^3\coth^3(\pi t)/(1+u^2t^2)$. If $t\le1$, use $t\coth(\pi t)\ge1/\pi$ and denominator at most $2$; if $t\ge1$, use $\coth(\pi t)\ge1$ and denominator at most $2t^2$.

The continuous Hahn normalization \cite[18.19.2--18.19.4; 18.20.9]{inf:DLMFHahn} gives an orthonormal basis for $\operatorname{sech}^2(\pi t)dt$:
\[
 \phi_k(t)=\sqrt{\frac{\pi(2k+1)}2}\,i^kF_k(2it),\qquad
 F_k(z)={}_3F_2\left(\begin{matrix}-k,k+1,(1+z)/2\\1,1\end{matrix};1\right).
\]
For $1\le y\le4/3$ we claim
\begin{equation}\label{inf:eq:bateman-bound}
 |F_k(-2y)|\le5(k+1)^{2y-1}.
\end{equation}
Put $a=1/2-y\in[-5/6,-1/2]$. A finite beta-integral expansion of the Legendre polynomial $\mathsf P_k$ gives
\[
 F_k(-2y)=1+\frac1{\Gamma(a)\Gamma(1-a)}
   \int_0^1s^{a-1}(1-s)^{-a}[\mathsf P_k(1-2s)-1],ds.
\]
The bracket vanishes to first order at zero, so the integral converges. Expand it termwise to recover the terminating hypergeometric series. The elementary Legendre estimate is
$|\mathsf P_k(1-2s)-1|\le\min(2,k(k+1)s)$; its derivative version follows by expressing $\mathsf P_k'$ as the sum of $(2j+1)\mathsf P_j$ of opposite parity and using $|\mathsf P_j|\le1$. With $T=k(k+1)$, split at $1/T$ to bound the integral by
\[
 \left(\frac1{a+1}+\frac2{-a}\right)T^{-a}\le10T^{-a}.
\]
Reflection gives $|\Gamma(a)\Gamma(1-a)|^{-1}\le1/\pi$, proving \eqref{inf:eq:bateman-bound}, including the immediate case $k=0$.

It follows that $\sum_{k=0}^{2n}|\phi_k(iy)|^2\le25\pi M^{4y}$. Expanding $Q$ in this basis and applying Cauchy--Schwarz at $iy$, together with \eqref{inf:eq:hahn-comparison}, yields
\[
 H(u^2)-L_n(u^2)\ge\frac1{50\pi^2M^{4/u}},\qquad3/4\le u\le1.
\]
Let $h=1/(4\log M)$; then $1-h\ge3/4$ and $M^{-4/u}\ge e^{-4/3}M^{-4}$ on $[1-h,1]$. Integrating $u(1-u)$ there, whose integral is at least $3h^2/8$, proves the theorem.
\end{proof}
Thus $-\log(\gamma-\alpha_n^-)\le4\log(2n+1)+2\log\log(2n+1)+O(1)$. A fixed exponential upper error $e^{-cn}$ cannot hold even on an unbounded subsequence. This is a lower bound, not a matching asymptotic, and does not exclude new cancellation between levels or moving evaluation points.

Two further restrictions retain their exact scope. Each positive-node integrand $w_ju(1-u)/(1+s_ju^2)$ has nonzero residue $w_j(1-z)/(2s_j)$ at $z=\pm i/\sqrt{s_j}$. Consecutive orthogonal polynomials have disjoint zero sets by their three-term recurrence. Therefore a nontrivial constant combination of two consecutive lower integrands still has a nonzero residue and has no rational antiderivative. This does not exclude a rational \emph{definite} integral from global logarithmic identities. Also a quadrature scheme supported on one fixed finite set $S$ cannot reproduce arbitrarily many moments: $\prod_{s\in S}(x-s)^2$ vanishes on that support but has positive $\nu$-integral. A fixed number field is not a fixed finite support.

\section{Bounded rational orbits and a vanishing approximation margin}
\label{inf:sec:orbit}
The orbit $xj/(j+1)$ stays in a fixed compact interval. Its conditional arithmetic analysis differs from that of $x/j$: translating the latter by $x$ would lose integer coefficients unless additional information about $x$ were already available.

\subsection{Exact-zero height has the sharp leading exponent}
\label{inf:sub:orbit-height}
Let $\alpha\in(0,1)$ be algebraic, with primitive irreducible polynomial
\[
 f(X)=a_dX^d+\cdots+a_0,\qquad a_d>0,\ a_0\ne0.
\]
Set
\[
 \beta_j=\alpha\frac j{j+1},\quad
 F_j(X)=j^df\left(\frac{j+1}jX\right),\quad
 f_j=F_j/g_j,
\]
where $g_j$ is the positive content of $F_j$.

\begin{lemma}[Content and leading coefficients]\label{inf:lem:orbit-content}
The $f_j$ are pairwise distinct primitive irreducible polynomials of degree $d$, and
\[
 g_j\mid |a_0a_d|,\qquad
 \operatorname{lc}(f_j)=\frac{a_d(j+1)^d}{g_j}
                  \ge\frac{(j+1)^d}{|a_0|}.
\]
\end{lemma}
\begin{proof}
Rational linear substitution preserves irreducibility. The norms
$N(\beta_j)=(j/(j+1))^dN(\alpha)$ are nonzero and distinct, hence the minimal polynomials differ. If a prime $p$ divides $j$, the leading coefficient of $F_j$ has valuation $v_p(a_d)$; if $p\mid j+1$, its constant coefficient has valuation $v_p(a_0)$. If neither is divisible by $p$, primitivity of $f$ implies that one coefficient of $F_j$ is a $p$-unit. Taking minimum coefficient valuations proves the content divisibility and the leading-coefficient bound.
\end{proof}

\begin{theorem}[Exact minimum-height exponent]\label{inf:thm:orbit-height}
For $D\ge dK$, let $h_{D,K}$ be the smallest coefficient height of a nonzero integer polynomial of degree at most $D$ vanishing at all $\beta_1,\ldots,\beta_K$. Then
\[
 \frac{((K+1)!)^d}{|a_0|^K}
 \le h_{D,K}\le\|f\|_1^K((K+1)!)^d,
 \qquad \log h_{D,K}=dK\log K+O_f(K).
\]
\end{theorem}
\begin{proof}
The product $S_K=\prod_{j=1}^Kf_j$ divides every such polynomial over $\Q$, hence integrally by Gauss's lemma. Its leading coefficient must divide the leading coefficient of the polynomial, giving the lower bound. Conversely $S_K$ itself has degree $dK$ and
\[
 H(S_K)\le\prod_j\|f_j\|_1\le\|f\|_1^K((K+1)!)^d.
\]
Stirling's elementary logarithmic estimate proves the last assertion.
\end{proof}

For $K=\lfloor c\log D/\log\log D\rfloor$, the height is $D^{dc+o(1)}$. Thus exact-zero polynomials within height $D$ exist eventually if $c<1/d$ and cannot exist if $c>1/d$. At the endpoint $c=1/d$ for this particular schedule, the upper construction still fits: writing $\ell=\log D$ and $L=\log\log D$,
\[
 \log H(S_K)\le\ell-\frac\ell L\log L+O_f(\ell/L)<\ell.
\]
This leading-coefficient argument improves an earlier comparison using $M(R)\le D\sqrt{D+1}$, which lost an unnecessary exponent $1/2$.

\subsection{A fixed gain above one is sufficient}
\label{inf:sub:kappa}
If $\deg R,H(R)\le D$ and $R(\beta_j)\ne0$, the resultant and
$M(f_j)\le M(f)(j+1)^d$ imply
\begin{equation}\label{inf:eq:orbit-resultant}
 |R(\beta_j)|\ge[M(f)(j+1)^d]^{-D}[D(D+1)]^{-(d-1)}.
\end{equation}

\begin{theorem}[Conditional orbit criterion]\label{inf:thm:kappa}
Fix $\kappa>1$. For every positive integer $d$, put
\[
 c_d=\frac{1+\kappa}{2d},\qquad
 K_{D,d}=\left\lfloor c_d\frac{\log D}{\log\log D}\right\rfloor.
\]
Suppose that, for each $d$, at unbounded integers $D$ there is a nonzero $R_D^{(d)}\in\Z[X]$ with degree and height at most $D$ such that
\begin{equation}\label{inf:eq:orbit-hypothesis}
 \max_{1\le j\le K_{D,d}}
 \left|R_D^{(d)}\left(\gamma\frac j{j+1}\right)\right|
 \le\exp\left(-\kappa\frac{D\log D}{K_{D,d}}\right).
\end{equation}
Then $\gamma$ is transcendental.
\end{theorem}
\begin{proof}
Assume $\gamma$ algebraic of degree $d$ and use that degree's family. Uniformly for $j\le K_{D,d}$, the negative logarithm of the nonzero lower bound \eqref{inf:eq:orbit-resultant} is at most $(d+o(1))D\log\log D$. The requested exponent is $(\kappa/c_d+o(1))D\log\log D$, and $\kappa/c_d>d$. Thus every value is eventually exactly zero. Theorem~\ref{inf:thm:orbit-height} then forces
\[
 H(R_D^{(d)})\ge D^{dc_d+o(1)}=D^{(1+\kappa)/2+o(1)}>D,
\]
a contradiction.
\end{proof}
The theorem proves sufficiency, not existence or an equivalence. The $d=1$ instance would exclude rationality only. The unbounded sets of successful $D$ need not be common to the different fixed degrees. The numerical choice $\kappa=1.03$ is permitted here with $c_d=1.015/d$; it is not the exponent of any earlier isolated linear-form certificate.

\subsection{A margin tending to zero can also suffice}
\label{inf:sub:vanishing}
Let
\[
 \ell=\log D,\quad L=\log\log D,\quad T=\log\log\log D,
 \qquad w_D\to\infty,\quad w_D=o(L),
\]
and define, for sufficiently large $D$,
\begin{equation}\label{inf:eq:vanishing-schedule}
 \kappa_D=1+\frac{2w_D}{L},\qquad
 K_{D,d}=\left\lfloor\frac\ell{d(L-T-w_D)}\right\rfloor.
\end{equation}

\begin{theorem}[Vanishing-margin version]\label{inf:thm:vanishing}
The conclusion of Theorem~\ref{inf:thm:kappa} holds if its parameters are replaced by \eqref{inf:eq:vanishing-schedule}.
\end{theorem}
\begin{proof}
Under an algebraic hypothesis of degree $d$, put $B_D=L-T-w_D$ and $K=K_{D,d}$. Since $B_D\sim L$,
\[
 K=\frac\ell{dB_D}+O(1),\qquad\log K=L-T-\log d+o(1).
\]
The negative logarithm of the nonzero-value lower bound is at most
\[
 A_D=D\log M(f)+dD\log(K+1)+(d-1)\log[D(D+1)]
     =dD(L-T)+O_f(D).
\]
The demanded exponent, including the floor error, is
\begin{align*}
 E_D&=\kappa_DD\ell/K\\
 &=dD\left(L-T+w_D-\frac{2w_D(T+w_D)}L\right)
       +O_d(DL^2/\ell).
\end{align*}
Hence $E_D-A_D=(d+o_f(1))Dw_D>0$, and all orbit values vanish. On the other hand,
\begin{align*}
 d\log((K+1)!)-K\log|a_0|-\ell
 &=\frac\ell{B_D}\left(w_D-\log d-1-\frac{\log|a_0|}d+o(1)\right)
       +O_f(L)\\
 &=(1+o_f(1))\frac{\ell w_D}L>0.
\end{align*}
The leading-coefficient lower bound again contradicts height at most $D$.
\end{proof}
A prescribed margin $\epsilon_D=\kappa_D-1$ can therefore be used when
$\epsilon_D\to0$ and $\epsilon_D\log\log D\to\infty$, by choosing $w_D=\epsilon_DL/2$. An explicit alternative schedule is
\[
 \widetilde\kappa_D=1+4T/L,\qquad
 \widetilde K_{D,d}=\left\lfloor(1+2T/L)\frac\ell{dL}\right\rfloor.
\]
Its exponent exceeds the resultant threshold by $(3d+o_f(1))DT$, while the exact-zero logarithmic height exceeds $\ell$ by $(1+o_f(1))\ell T/L$. Both inequalities survive integer rounding. Neither schedule has been achieved by a gamma-specific infinite family.

\subsection{Universal counting and a finite success that does not extend}
\label{inf:sub:orbit-feasibility}
For any fixed $x\in(0,1)$, evaluate $\{0,\ldots,D\}^{D+1}$ at $K$ points $xj/(j+1)$. Its images lie in $[0,D/(1-x)]^K$. Partition into $t^K<G_D$ boxes, taking $t$ maximal. Two coefficient vectors lie in one box, giving a nonzero difference with
\[
 \max_j|R(xj/(j+1))|\le\frac{D}{(1-x)t}
 =\exp\left(-(1+o(1))\frac{D\log D}{K}\right)
\]
for the above scales of $K$. Zero values are allowed. The notation $o(1)$ supplies neither a positive sign nor the necessary rate for the margin in Theorem~\ref{inf:thm:vanishing}.

The archived finite search used the older schedule $\kappa=2$, $d=1$, $K=\lfloor(7/4)\log D/\log\log D\rfloor$. At the four selected degrees $K=4$. The candidate
\[
 R_8=X^6(2X-1)^2,\qquad H(R_8)=4,
\]
was certified by rational arithmetic to satisfy
\[
 \max_{j\le4}|R_8(\gamma j/(j+1))|<173/10^6<1/4096.
\]
The selected candidates $X^9(2X-1)^3$, $X^{13}(2X-1)^3$ and $X^{21}(2X-1)^3$ at $D=12,16,24$ failed that same target. Failure of selected candidates is not a nonexistence theorem.

There is, however, an all-order obstruction to extending a success just by powers of $X$. For fixed nonzero $Q\in\R[X]$ and $x\in(0,1)$, set $R_D=X^{D-\deg Q}Q$ and $K_D\sim c\log D/\log\log D$. Then
\[
 \frac{-K_D\log\max_{j\le K_D}|R_D(xj/(j+1))|}{D\log D}\longrightarrow0.
\]
Indeed, if $Q$ has a root of multiplicity $r\ge0$ at $x$, its value at $xK/(K+1)$ is comparable to $K^{-r}$. Boundedness on $[0,x]$ and this last point give
\[
 -\log\max_j|R_D(xj/(j+1))|
 =D\log(1/x)+O(D/K+\log K).
\]
Normalization proves the claim. A construction whose factor structure changes with $D$ is not covered.

\section{Bernoulli convolutions: an explicitly conditional arithmetic branch}
\label{inf:sec:fractal}
Let
\[
 \nu_x=\operatorname{Law}\left(\sum_{j\ge0}\varepsilon_jx^j\right),
 \qquad\Pr(\varepsilon_j=\pm1)=\tfrac12,
\]
and denote by $\mathcal F$ the additional hypothesis $\dim_H\nu_\gamma<1$. All entropies in this section use natural logarithms.

The imported facts are the following: transcendental parameters in $(1/2,1)$ have full dimension; dimension drop produces height-one polynomials of arbitrarily strong exponential smallness; and for algebraic $x$ the dimension is $\min(1,h_x/\log(1/x))$, where $h_x$ is entropy per digit. These are recorded with their sources in \cite[Theorems 3, 6, 8]{inf:Varju}, using Hochman's theorem \cite{inf:Hochman}. They are external theorems, not new assertions proved by the present program.

\begin{proposition}[Arithmetic consequences of dimension drop]\label{inf:prop:fractal-branch}
If $\mathcal F$ holds, then $\gamma$ is an algebraic unit, is a root of a nonzero polynomial with coefficients in $\{-1,0,1\}$, and satisfies
\[
\begin{gathered}
 \frac12<|\sigma(\gamma)|<2\quad\text{for every conjugate},\qquad |N(\gamma)|=1,\\
 [\Q(\gamma):\Q]\ge4,\qquad \gamma^{-1}\le M(\gamma)<16.
\end{gathered}
\]
In particular every non-Archimedean absolute value of $\gamma$ equals $1$.
\end{proposition}
\begin{proof}
Varj\'u's theorem first forces algebraicity. Let $d$ and $M$ be its hypothetical degree and Mahler measure. For every fixed $0<\theta<1/M$, dimension drop gives nonzero polynomials $Q_n$ of degree at most $n$ and coefficients in $\{-1,0,1\}$ with $|Q_n(\gamma)|<\theta^n$ for large $n$. Nonzero values would instead be at least $M^{-n}(n+1)^{1-d}$ by the resultant. Thus some $Q_n(\gamma)$ equals zero.

Remove any power of $X$ and change sign to make its leading coefficient $1$. Its constant coefficient is $\pm1$, hence $\gamma$ and $\gamma^{-1}$ are algebraic integers. If a root $z$ of this polynomial had $|z|\ge2$, its leading term would strictly exceed the sum of all lower terms in modulus. Applying the same argument to the reciprocal gives $|z|>1/2$. All conjugates obey these bounds. The unit norm identity yields
\[
 M(\gamma)=\prod_{|\sigma(\gamma)|<1}|\sigma(\gamma)|^{-1}\ge\gamma^{-1}.
\]

For the upper Mahler bound, the entropy estimate of Breuillard--Varj\'u \cite[Proposition 13]{inf:BV}, specialized to algebraic units and a fair sign $E$, gives $h_x\ge\Phi(M(x))$, where for an independent standard Gaussian $G$,
\[
 \Phi(a)=\sup_{t>0}\{h_{\rm diff}(taE+G)-h_{\rm diff}(tE+G)\}.
\]
Choose $t=a^{-1/2}$. The difference equals
$I(E;\sqrt aE+G)-I(E;a^{-1/2}E+G)$. A sign decision in the first channel has error probability at most $e^{-a/2}$. Its binary conditional entropy is at most $2e^{-a/4}$, since $h_b(p)\le2\sqrt p$. The second information is at most $1/(2a)$ by comparison of its two shifted Gaussian components with $N(0,1)$ and nonnegativity of mixture relative entropy. Hence
\[
 \Phi(a)\ge\log2-2e^{-a/4}-\frac1{2a}.
\]
For $a\ge16$, use $e^4>50$ and $\log(11/10)>1/11>57/800$ to obtain
\[
 \Phi(a)>\log2-57/800>\log(20/11).
\]
Since $\gamma>11/20$, $M(\gamma)\ge16$ would force $h_\gamma>\log(1/\gamma)$ and therefore full dimension, contradicting $\mathcal F$.

For the degree bound, the elementary harmonic/logarithm enclosure certifies $577/1000<\gamma<578/1000$. A monic unit polynomial of degree at most $3$ whose roots all have modulus below $2$ belongs to one of the following larger coefficient boxes:
\[
 X+c\ (c=\pm1),\quad X^2+aX+c\ (|a|\le3,c=\pm1),\quad
 X^3+aX^2+bX+c\ (|a|\le5,|b|\le7,c=\pm1).
\]
These contain $2$, $14$, and $330$ polynomials. At $m=231/400$ and $r=1/2000$, the rational lower bound
\[
 |P(x)|\ge|P(m)|-r\sum_jj|p_j|
\]
is positive for all of them throughout the gamma interval. Direct finite evaluation gives the respective minima $211/500$, $13999/160000$, and $677991/64000000$. This excludes all three degrees under the unit and conjugate restrictions.
\end{proof}
The final finite enumeration is reproducible in the archived \texttt{verify\_fractal\_subworlds.py}; it does not restrict the degree of $\gamma$ without the hypothesis $\mathcal F$.

There are actual algebraic dimension-drop models in the wider interval $(1/2,3/5)$. Let $\beta$ be the real root of $X^3-X^2-X-1$ and $\lambda=\beta^{-1}$. The discriminant is $-44$, and the other conjugates of $\beta$ have modulus $\beta^{-1/2}<1$, so $\beta$ is a Pisot unit. Moreover $\lambda\in(1/2,3/5)$, $M(\lambda)=\beta<2$, and there is no unit-circle conjugate. The strict entropy bound of \cite[Proposition 7]{inf:BV} gives $h_\lambda<\log M(\lambda)=\log(1/\lambda)$, hence dimension drop. This parameter lies outside the narrower gamma interval. The inference uses the entropy theorem, not the invalid general implication from singularity to dimension drop.

Conversely, full dimension does not force transcendence: $4/7$ is algebraic in the same broad interval and cannot be a root of any nonzero $\{-1,0,1\}$ polynomial by the rational-root theorem, so the imported dimension theorem gives full dimension. Neither $\mathcal F$ nor its negation has been established for $\gamma$ here. Proving $\mathcal F$ would support algebraicity, whereas excluding it alone would not prove transcendence.

\subsection{The exact logical closure and the finite record}
\label{inf:sub:status}
Let A mean algebraicity of $\gamma$, B its transcendence, and C the definitions and proved statements used above. The orbit result has the form
\[
 \mathrm C+\mathrm A+\text{orbit-family existence}\ \Longrightarrow\ \bot.
\]
The fractal result gives $\mathrm C+\mathrm B+\mathcal F\Longrightarrow\bot$, and excludes combinations of $\mathcal F$ with nonunit behavior, a forbidden conjugate, degree at most $3$, or Mahler measure at least $16$. None is a contradiction from A alone or B alone. A finite implication registry can trace minimal sets of extra assumptions; it is not a general theorem-discovery engine, a consistency proof, or a proof of independence.

The archived finite records preserve several different experiments without identifying their outputs:
\begin{itemize}
 \item Full-box separation at $D=2,3,4,5$ was certified using $27$, $256$, $3125$, and $46656$ grid values. The associated positive-gap polynomials are $2X^2-X$, $3X^3-3X^2-X+1$, $3X^4+3X^3-4X^2-X+1$, and $-3X^5+3X^4-5X^3+2X^2+2X-1$. These exclude only the respective simultaneous degree-and-height boxes.
 \item The stricter canonical midpoint rule was independently executed for $D=2,3,4$, with $k_D=6,10,14$. Its outputs are the first three polynomials just listed and are squarefree. Their midpoint denominators retain the primes $37$, $521$, and $8209$. The $D=5$ finite-box experiment is not asserted to be that canonical output.
 \item A separate primitive-irreducible search produced degrees $1,2,4,6,8,10,12$ and certified $\mathcal R_\gamma(P)>0,1,3,4,5,6,6$, respectively. Irreducibility was checked at primes $2,3,17,11,2,19,101$. These do not satisfy the same height-at-most-$D$ schedule. A certified upper value $U$ and the rational inequality $U(3^{\deg P}\max(2,H(P)))^s<1$ imply $\mathcal R_\gamma(P)>s$, since $e<3$.
 \item The full-lattice positive candidate at pole order $8$ had $(q,p)=(104,60)$ and coefficient vector $(52,17,3,-16,25,-21,7)$. It is exactly half the previous $(208,120)$ kernel and reduces to the same primitive form $26\gamma-15$, so it is not counted as a new approximation. Other archived positive candidates reduce to $7\gamma-4$ and $201\gamma-116$; solver timeouts certify no optimum.
 \item The first six lower and upper Pad\'e identities were checked exactly; the sixth bracket width is below $3.772\cdot10^{-6}$. The all-order existence and rate claims rest on the proofs above, not this sample.
\end{itemize}
The source packages retain exact coefficients, rational intervals and independent check scripts. Their reported finite verification counts are experimental provenance, not a substitute for any quantified theorem. The incomplete assertions remain unbounded nonvanishing of the canonical outputs, an infinite local multiplicity escape, a common-denominator separated-error family, and the orbit improvement required in \eqref{inf:eq:orbit-hypothesis} or \eqref{inf:eq:vanishing-schedule}.

% END INTEGRATED PART V

\clearpage
\part{The two alternatives and local arithmetic}\label{master:worlds}
% BEGIN INTEGRATED PART VI
\section{The two arithmetic cases and the obligation to preserve relations}
\label{wr:sec-worlds}

This part consolidates the two-world audit, the algebraic-system notes, the
one-prime investigation, and their cross-field extensions. Throughout, the
actual real number under investigation is fixed. We use the labels
\(\mathrm A_x:x\in\Qbar\) and \(\mathrm B_x:x\notin\Qbar\), and specialize
to \(x=\gamma\) only after proving the general statements. These labels do
not describe newly constructed models of the full definition of Euler's
constant. None of the conclusions below determines which case holds for
\(\gamma\).

The logical bookkeeping is simple but essential. If a body of established
facts \(\mathcal T\) proves both \(\mathrm A_\gamma\Rightarrow C\) and
\(\neg C\), it proves transcendence. If instead the derivation uses an
additional hypothesis \(H\), then a contradiction proves
\(\neg\mathrm A_\gamma\lor\neg H\). One cannot discard \(H\) without a
separate proof. Several unsuccessful transfers in this part fail precisely
at this point.

\subsection{The relation ideal, derivations, and formal motions}

For a real number \(x\), put
\[
 I_x=\{P\in\Q[X]:P(x)=0\},\qquad R_x=\Q[x],\qquad K_x=\Q(x).
\]
The evaluation presentation \(R_x\simeq\Q[X]/I_x\) always holds. In the
algebraic case, let \(f\in\Z[X]\) be the primitive irreducible minimal
polynomial, with positive leading coefficient, and write \(d=\deg f\).
Then \(I_x=(f)\), \(R_x=K_x\), and \([K_x:\Q]=d\). In the transcendental
case these objects are respectively \(0,\Q[X],\Q(X)\). In particular,
\begin{equation}\label{wr:eq-rank}
 \dim_\Q\langle1,x,\ldots,x^D\rangle
 =\begin{cases}\min(D+1,d),&\mathrm A_x,\\ D+1,&\mathrm B_x.\end{cases}
\end{equation}
The ground field in this formula cannot be replaced by \(\R\): real
coefficient evaluation has rank one in either case.

\begin{proposition}[Equivalent forms of algebraic freedom]\label{wr:prop-freedom}
For a real number \(x\), the following conditions are equivalent:
\begin{enumerate}
\item \(x\) is transcendental;
\item a \(\Q\)-derivation \(D:K_x\to K_x\) satisfies \(Dx=1\);
\item \(\dim_{K_x}\Omega_{K_x/\Q}=1\);
\item \(P(x)=0\) implies \(P'(x)=0\) for every \(P\in\Q[X]\);
\item a \(\Q\)-algebra homomorphism \(\Phi:K_x\to K_x[[t]]\) has constant
term the identity and satisfies \(\Phi(x)=x+t\);
\item a \(\Q\)-automorphism of \(K_x\) sends \(x\) to \(x+1\).
\end{enumerate}
In the algebraic case every \(\Q\)-derivation is zero, and every formal
self-deformation with identity constant term is constant.
\end{proposition}
\begin{proof}
In the transcendental case use ordinary differentiation and substitution
in \(\Q(X)\). For example,
\[
 D\!\left(\frac{P(x)}{Q(x)}\right)
 =\frac{P'(x)Q(x)-P(x)Q'(x)}{Q(x)^2},\qquad
 \Phi\!\left(\frac{P(x)}{Q(x)}\right)
 =\frac{P(x+t)}{Q(x+t)}.
\]
The latter denominator is invertible as a formal power series. In the
algebraic case, separability gives \(f'(x)\ne0\); differentiating the
relation gives \(f'(x)Dx=0\), and the same equation for K\"ahler
differentials gives \(\Omega_{K_x/\Q}=0\). If a nonzero ideal of relations
were stable under formal differentiation, differentiating its element of
least positive degree would contradict minimality. If
\(\Phi(x)-x=ct^m+O(t^{m+1})\), its first nonzero coefficient in
\(f(\Phi(x))\) would be \(cf'(x)\). Finally, an automorphism sending
\(x\) to \(x+1\) would make all the distinct numbers \(x+n\) conjugates
of \(x\), which is impossible in finite degree. These arguments also
prove the asserted rigidity; the standard module formulation is recorded
in \cite{wr:stacks-differentials,wr:stacks-etale}.
\end{proof}

A derivation on a real subfield which is continuous for the usual topology
must vanish: approximate every element by rational numbers. Likewise an
order-preserving \(\Q\)-automorphism of a real subfield is the identity,
because a rational number strictly between an element and its purported
image contradicts order preservation. Consequently neither continuity
nor order preservation may be silently added to the maps in
Proposition~\ref{wr:prop-freedom}.

There is a useful quantitative form of this warning. For
\begin{equation}\label{wr:eq-harmonic-approximant}
 s_n=2\sum_{j=0}^{n-1}\frac{1}{(2j+1)3^{2j+1}},\quad
 \rho_n=\frac{3}{4(2n+1)9^n},\quad
 u_n=H_{2^n}-ns_n,
\end{equation}
the positive tail of \(\log2=2\operatorname{arctanh}(1/3)\) and integral
comparison for harmonic numbers give
\begin{equation}\label{wr:eq-real-approx}
 0<u_n-\gamma<2^{-n}+n\rho_n.
\end{equation}
Under \(\mathrm B_\gamma\), the derivation with \(D\gamma=1\) has
\[
 \frac{|D(u_n-\gamma)|}{|u_n-\gamma|}
 >\frac1{2^{-n}+n\rho_n}=(1+o(1))2^n.
\]
This is the expected discontinuity, not a contradiction. More generally,
for every fixed \(m\),
\[
 \{(u,Du,\ldots,D^mu):u\in\Q[\gamma]\}
 \quad\hbox{is dense in }\R^{m+1}.
\]
Indeed the map from the coefficients of a polynomial of degree at most
\(m\) to its jet at \(\gamma\) is an invertible real linear map; apply it
to the dense set of rational coefficient vectors. No integer height
bound is present in this density statement.

One can also express the missing assertion as a finite-box estimate:
\(\mathrm B_\gamma\) is equivalent to the existence, for every \(D\), of
some finite \(C_D\) such that
\[
 |P'(\gamma)|\le C_D|P(\gamma)|
 \quad\bigl(P\in\Z[X],\ \deg P,H(P)\le D\bigr).
\]
In case B take the maximum over the finite box; in case A a sufficiently
large box contains \(f\), for which the inequality fails. Merely renaming
this existence assertion does not supply a bound or prove it.

\subsection{Counting distinct values and the relation lattice}

Define
\[
 W_D(x)=\#\left\{\sum_{j=0}^D a_jx^j:
                 (a_0,\ldots,a_D)\in\{0,\ldots,D\}^{D+1}\right\},
 \qquad 0<x<1.
\]
\begin{proposition}[Exact counting dichotomy]\label{wr:prop-counting}
If \(x\) is transcendental then \(W_D(x)=(D+1)^{D+1}\). If \(x\) is
algebraic with primitive minimal polynomial \(f\) of degree \(d\), then
\begin{equation}\label{wr:eq-counting-bound}
 W_D(x)\le1+[D(D+1)]^d M(f)^D.
\end{equation}
Consequently
\[
 x\notin\Qbar
 \quad\Longleftrightarrow\quad
 \limsup_{D\to\infty}\frac{\log W_D(x)}D=\infty.
\]
\end{proposition}
\begin{proof}
Equality of two grid values is a polynomial relation, proving the first
assertion. In the algebraic case a nonzero value of an integer polynomial
\(P\) of degree and height at most \(D\) obeys
\[
 |P(x)|\ge M(f)^{-D}[D(D+1)]^{-(d-1)}.
\]
To see this, use the nonzero integer \(\operatorname{Res}(f,P)\), and
bound each conjugate value other than \(P(x)\) by
\(D(D+1)\max(1,|x_i|)^D\). Since \(0<x<1\), multiplication gives the
stated bound. All grid values lie in an interval of length at most
\(D(D+1)\); divide that length by the separation bound. The last claim
follows because the algebraic bound is exponential times a polynomial,
whereas the transcendental count has logarithm asymptotic to \(D\log D\).
\end{proof}

The normalized quantity \(\log W_D(x)/((D+1)\log(D+1))\) therefore tends
to zero in A and is identically one in B. It counts exact values, not
coefficient vectors or cells at a prescribed resolution.

In case A and \(D\ge d\), Gauss's lemma gives the exact integer relation
lattice
\begin{equation}\label{wr:eq-relation-lattice}
 \{P\in\Z[X]_{\le D}:P(x)=0\}=f\Z[X]_{\le D-d}.
\end{equation}
Its basis \(f,Xf,\ldots,X^{D-d}f\) has covolume, in its own real span,
at most \(\|f\|_2^{D+1-d}\) by Hadamard's inequality. It is saturated:
\(mP(x)=0\) with \(m\ne0\) implies \(P(x)=0\). Equivalently its nonzero
Smith invariant factors are one. Reduction of the multiplication matrix
modulo any prime is injective, since \(f\bmod p\) is a nonzero
polynomial. Writing \(L=\|f\|_1\), the height-\(D\) box contains at least
\begin{equation}\label{wr:eq-abundant-relations}
 (2\lfloor D/L\rfloor+1)^{D-d+1}-1
\end{equation}
nonzero exact relations: multiply \(f\) by every coefficient vector in
\([ -\lfloor D/L\rfloor,\lfloor D/L\rfloor]^{D-d+1}\). Thus an abundance
of short exact relations is compatible with A. For \(x=1/2\), the basis
from \(f=2X-1\) has Gram matrix with diagonal 5 and adjacent entries
\(-2\); its covolume is \(\sqrt{(4^{D+1}-1)/3}\), while
\(W_D(1/2)\le2D\,2^D+1\).

Certified interval packing can give lower bounds for \(W_D\). If
\(|c-x|\le\varepsilon\), two grid polynomials change their difference
by at most \(B_D\varepsilon\), where \(B_D=D^2(D+1)/2\). In A, if
\(2B_D\varepsilon\) is smaller than the resultant separation above,
clustering sorted midpoint values across gaps at most
\(B_D\varepsilon\) recovers their exact value classes. In B, a sufficiently
fine precision separates any fixed finite grid, but no uniform bound on
that precision is obtained here. Such a uniform packing theorem at
\(\gamma\) would be additional arithmetic input.

\subsection{Fractal dimension is a different observable}

Section~\ref{inf:sec:fractal} records the imported Varj\'u and Hochman
theorems and the corresponding conditional arithmetic conclusions.
For comparison with \(W_D\), retain the rational example \(x=4/7\).
Differences of distinct binary words cannot vanish by the rational-root
theorem, and their length-\(n\) separation is at least \(7^{-(n-1)}\).
Thus Hochman's theorem \cite{inf:Hochman} gives full Bernoulli-convolution dimension,
while
\[
 W_D(4/7)\le1+\tfrac73D\,7^D.
\]
The latter follows by counting multiples of \(7^{-D}\) in an interval
of length less than \(7D/3\). Full dimension and vanishing normalized
exact counting growth therefore coexist at a rational parameter.

The scale mismatch can be seen without a dimension theorem. The law of
\(\sum_{j=0}^D A_jx^j\), with digits uniform in \(\{0,\ldots,D\}\), has
support in \([0,D/(1-x)]\). At interval resolution \(e^{-cD}\) its
Shannon entropy is at most
\[
 \log\left(2+\frac{De^{cD}}{1-x}\right)=cD+O_x(\log D)
\]
for every \(x\in(0,1)\). Fixed exponential resolution cannot read the
\(D\log D\) exact atomic count. A dimension-zero exceptional set also
need not exclude a specified computable parameter.

\section{One-prime character counts and optimal finite camouflage}
\label{wr:sec-local}

All homomorphisms in this section preserve the identity and fix the
integers. Set
\[
 n_p(x)=\#\operatorname{Hom}_\Z(\Z[x],\mathbb F_p),\qquad
 N_{p,k}(x)=\#\operatorname{Hom}_\Z(\Z[x],\Z/p^k\Z).
\]
For algebraic \(x\), primitivity and Gauss's lemma give
\(\Z[x]\simeq\Z[X]/(f)\), even when \(f\) is not monic. Thus these
counts are precisely the respective root counts of \(f\). For
transcendental \(x\) the source is the free polynomial ring, so
\begin{equation}\label{wr:eq-free-counts}
 n_p(x)=p,\qquad N_{p,k}(x)=p^k.
\end{equation}
In either case, a proposed image \(a\) in a target ring \(S\) defines an
actual homomorphism exactly when
\begin{equation}\label{wr:eq-relation-preservation}
 P(x)=0\quad\Longrightarrow\quad P(a)=0\text{ in }S
 \qquad\text{for every }P\in\Z[X].
\end{equation}
This implication is the unresolved input at \(x=\gamma\).

\subsection{All primes and exact eventual counts at one prime}

\begin{proposition}[Prime characters]\label{wr:prop-prime-characters}
For every real \(x\),
\[
 \bigl[\operatorname{Hom}_\Z(\Z[x],\mathbb F_p)\ne\varnothing
       \text{ for every prime }p\bigr]
 \Longleftrightarrow x\in\Z\text{ or }x\notin\Qbar.
\]
If the condition is imposed only at all sufficiently large primes, the
right side becomes \(x\in\Q\) or \(x\notin\Qbar\). Requiring only
infinitely many primes does not distinguish the cases.
For algebraic \(x\),
\[
 \sum_{p\le X}n_p(x)\sim\operatorname{Li}(X),
\]
whereas for transcendental \(x\) the sum is asymptotic to
\(X^2/(2\log X)\).
\end{proposition}
\begin{proof}
For \(x=a/b\) in lowest terms, \(b>0\), a prime dividing \(b>1\)
prevents any map: the relation \(bx=a\) would become \(0=\bar a\).
At all other primes the rational case has a unique image. In degree
\(d>1\), the Galois group of the splitting field acts transitively on
the roots of \(f\). Its average number of fixed points is one, while
the identity has \(d>1\) fixed points. It therefore contains a
fixed-point-free element. Chebotarev supplies a positive-density set of
primes with this Frobenius class, at which \(f\) has no root. Exclude
the finitely many leading-coefficient and discriminant primes in making
the Frobenius identification. The same theorem applied to the identity
class supplies infinitely many primes with roots. Averaging the fixed
point character gives the algebraic asymptotic; in the transcendental
case use \eqref{wr:eq-free-counts} and partial summation of the prime
number theorem. The needed Frobenius correspondence and density theorem
are recorded in \cite{wr:milne,wr:ekt}.
\end{proof}

More generally, if \(m_r\) is the number of Galois orbits on ordered
\(r\)-tuples of roots, then
\(\pi(X)^{-1}\sum_{p\le X}n_p(x)^r\to m_r\). In particular \(m_1=1\)
and the limiting variance is \(m_2-1\). These are conditional
fingerprints, not measured values of \(n_p(\gamma)\).

\begin{theorem}[Exact eventual one-prime count]\label{wr:thm-eventual}
Let \(x\) be algebraic with primitive minimal polynomial \(f\), and
choose \(A,B\in\Z[X]\), \(M\in\Z\setminus\{0\}\), such that
\(Af+Bf'=M\). For a fixed prime \(p\), put \(s=v_p(M)\). Then
\begin{equation}\label{wr:eq-eventual}
 k>2s\quad\Longrightarrow\quad
 N_{p,k}(x)=
 \sum_{\substack{\alpha\in\Z_p\\f(\alpha)=0}}
 p^{v_p(f'(\alpha))}\le d p^s.
\end{equation}
The sum is zero if \(f\) has no integral \(p\)-adic root. In particular
\(N_{p,k}(x)\) is eventually constant.
\end{theorem}
\begin{proof}
The B\'ezout identity and \(f(a)\equiv0\pmod{p^k}\) imply
\(t=v_p(f'(a))\le s\). Since \(k>2t\), the strong Hensel lemma
\cite[Theorem~4.1]{wr:hensel} gives a root \(\alpha\in\Z_p\) with
\(v_p(a-\alpha)\ge k-t>s\). As \(f'\) has integral coefficients,
\(v_p(f'(\alpha))=t\). Conversely, the residues in
\[
 a\equiv\alpha\pmod{p^{k-s_\alpha}},\qquad
 s_\alpha=v_p(f'(\alpha)),
\]
all solve the congruence, by Taylor expansion; each such ball contains
\(p^{s_\alpha}\) residues modulo \(p^k\). The balls are disjoint. If
two overlapped, distinct roots \(\alpha,\beta\) would satisfy
\(v_p(\alpha-\beta)>s\); division by both monic linear factors in
\(\Z_p[X]\) would then force \(v_p(f'(\alpha))>s\). Summing the
disjoint contributions proves the formula.
\end{proof}

\begin{corollary}[One-prime dichotomy]\label{wr:cor-one-prime}
For any fixed prime \(p\) and every real \(x\),
\[
 x\notin\Qbar\ \Longleftrightarrow\ \sup_kN_{p,k}(x)=\infty,
 \qquad
 \lim_{k\to\infty}\frac{\log_p(1+N_{p,k}(x))}{k}
 =\begin{cases}0,&x\in\Qbar,\\1,&x\notin\Qbar.\end{cases}
\]
Also
\[
 \operatorname{Hom}_\Z(\Z[x],\Z_p)
 \simeq\varprojlim_k\operatorname{Hom}_\Z(\Z[x],\Z/p^k\Z),
\]
and this set is infinite if and only if \(x\) is transcendental.
\end{corollary}
\begin{proof}
Combine \eqref{wr:eq-free-counts} and Theorem~\ref{wr:thm-eventual}.
The inverse-limit statement follows by taking componentwise ring
operations. Its left side has at most \(d\) elements in degree \(d\)
and is in bijection with \(\Z_p\) in the transcendental case.
\end{proof}

Finite-level cardinalities need not count infinite branches. For
\(x=\sqrt{17}\), \(N_{2,k}(x)=4\) for \(k\ge3\), although only two
maps to \(\Z_2\) exist. Indeed \(-2(X^2-17)+X(2X)=34\), and both
\(2\)-adic roots have derivative valuation one. The roots modulo 8
are \(1,3,5,7\), whereas the roots modulo 16 are \(1,7,9,15\);
only two of the former four branches lift at that step.

There is a second counting trap. If \(J\) consists of finitely many
verified relations and \(J\subset I=\ker(\Z[X]\to\Z[x])\), then
\begin{equation}\label{wr:eq-map-inclusion}
 \operatorname{Hom}(\Z[X]/I,S)
 \hookrightarrow\operatorname{Hom}(\Z[X]/J,S).
\end{equation}
Thus the number of not-yet-excluded images is an \emph{upper} bound on
actual maps. It cannot provide the lower bound required by
Corollary~\ref{wr:cor-one-prime}.

\subsection{A height-free bound valid at every depth}

\begin{theorem}[Primitive polynomial root bound]\label{wr:thm-height-free}
Let \(f\in\Z_p[X]\) have positive degree \(d\), with at least one unit
coefficient. For every \(k\ge1\),
\begin{equation}\label{wr:eq-height-free}
 \#\{a\bmod p^k:f(a)\equiv0\pmod{p^k}\}
 \le\min\{p^k,d p^{k-\lceil k/d\rceil}\}.
\end{equation}
It follows that \(N_{p,k}(x)\) satisfies this bound for every algebraic
\(x\) of degree \(d\), including nonintegral numbers and bad primes.
\end{theorem}
\begin{proof}
Factor \(f(X)=c\prod_{i=1}^d(X-\alpha_i)\) in \(\C_p\). Multiplicativity
of the Gauss norm gives
\(1=|c|_p\prod_i\max(1,|\alpha_i|_p)\). For \(a\in\Z_p\), the
factors from roots outside the unit disc cancel, leaving
\[
 |f(a)|_p=\prod_{|\alpha_i|_p\le1}|a-\alpha_i|_p.
\]
If this product is at most \(p^{-k}\), one of its at most \(d\) factors
has valuation at least \(k/d\). Any two integral points in the ball
around that root differ by an element of valuation at least \(k/d\),
hence are congruent modulo \(p^{\lceil k/d\rceil}\). Such a ball
contains at most \(p^{k-\lceil k/d\rceil}\) residues modulo \(p^k\).
Take the union over at most \(d\) roots. Primitivity supplies the unit
coefficient for a minimal polynomial at every prime.
\end{proof}

If \(M_0\) distinct genuine maps at depth \(k\) satisfy
\[
 M_0>D p^{k-\lceil k/D\rceil},
\]
then algebraic degree at most \(D\) is excluded. This is a criterion
without a height parameter, but no such lower bound for the actual
\(\gamma\)-ring is established here.

\subsection{The exact minimum degree of finite camouflage}

For \(M\ge2\), define \(\mu(M)=\min\{m\ge0:M\mid m!\}\). The
relevant classical mechanism is fixed divisors, or the elementary
\(\Z\)-case of generalized factorials \cite{wr:bhargava}.

\begin{lemma}[Fixed-divisor bound]\label{wr:lem-fixed-divisor}
For \(0\ne f\in\Z[X]\) of degree \(d\), let \(c(f)\) be its coefficient
content and let \(\operatorname{fd}(f)=\gcd\{f(a):a\in\Z\}\). Then
\(\operatorname{fd}(f)\mid d!c(f)\).
\end{lemma}
\begin{proof}
The Newton expansion
\[
 f(X)=\sum_{j=0}^d\Delta^jf(0)\binom Xj
     =\sum_{j=0}^d\frac{\Delta^jf(0)}{j!}(X)_j
\]
has every \(\Delta^jf(0)\) divisible by the fixed divisor. Multiplication
by \(d!\) shows that this divisor divides every coefficient of \(d!f\),
whose gcd is \(d!c(f)\).
\end{proof}

\begin{theorem}[Optimal finite camouflage]\label{wr:thm-camouflage}
Fix finitely many primes \(S\ne\varnothing\), depths \(L_p\ge1\), and
a nonempty open real interval \(J\). Put \(M=\prod_{p\in S}p^{L_p}\). The least
possible degree of an algebraic \(x\in J\) satisfying
\begin{equation}\label{wr:eq-camouflage}
 N_{p,k}(x)=p^k\qquad(p\in S,\ 1\le k\le L_p)
\end{equation}
is exactly \(\mu(M)\). This minimum is attained by an algebraic integer.
Allowing nonmonic minimal polynomials cannot reduce it.
\end{theorem}
\begin{proof}
If \(f\) is the primitive minimal polynomial of any such \(x\), then
\eqref{wr:eq-camouflage} says that all integer values of \(f\) are
multiples of \(M\). Lemma~\ref{wr:lem-fixed-divisor} gives \(M\mid d!\),
so \(d\ge\mu(M)\).

For existence let \(m=\mu(M)\ge2\), \(Q(X)=(X)_m\), and choose a prime
\(\ell\notin S\). Coefficientwise Chinese remaindering produces a
monic degree-\(m\) polynomial \(f_0\) congruent to \(Q\) modulo \(M\),
with every lower coefficient divisible by \(\ell\), and constant
coefficient congruent to \(\ell\) modulo \(\ell^2\). Choose a rational
subinterval \((u,v)\subset J\) and \(a/b\in(u,v)\), \(b>0\). Set
\[
 f_A(X)=f_0(X)+M\ell^2 A(bX-a),\qquad A\in\Z_{>0}.
\]
This preserves the degree, monicity, congruences, and Eisenstein
irreducibility at \(\ell\). For large \(A\), its signs at \(u,v\) are
opposite, so it has a real root in \(J\). Since
\((a)_m=m!\binom am\) is divisible by \(M\) for every integer \(a\),
all prescribed residue classes are roots. The presentation of
\(\Z[x]\) proves \eqref{wr:eq-camouflage}.
\end{proof}

The construction also reveals a separating measurement outside the
prescribed list: \(N_{\ell,1}(x)=1\) and \(N_{\ell,2}(x)=0\).
Modulo \(\ell\) its polynomial is \(X^m\); a root modulo \(\ell^2\)
would be divisible by \(\ell\), which leaves the nonzero constant term
modulo \(\ell^2\).

Legendre's formula gives
\[
 \mu(p^L)=(p-1)L+O_p(\log(L+1)),\qquad
 \mu\!\left(\prod p^{L_p}\right)=\max_p\mu(p^{L_p}).
\]
For example, the optimal degrees at \(p=2\) and depths \(3,7,10,100\)
are \(4,8,12,104\), respectively. The earlier construction of degree
\(p^L\) was valid but is superseded as an optimization claim.
An exact degree-four example in
\((577215/10^6,577216/10^6)\) is the root of
\[
 X^4+18X^3+27X^2+144000018X-83119056.
\]
It is Eisenstein at 3 and is a null polynomial modulo 8. The earlier
fourth-degree example
\(X^4-6X^3+3X^2+28794X-16620\), with a root in
\((577/1000,578/1000)\), instead has counts \(2,4,0,0,\ldots\).
Both are finite-data comparisons, not alternate realizations of
\(\gamma\)'s full definition.

\begin{corollary}[Exact full-residue depth]\label{wr:cor-depth}
The first depth at which full counts exclude every algebraic degree
\(\le D\) is
\[
 k_D=v_p(D!)+1.
\]
For \(p=2\), this is \(D-s_2(D)+1\), where \(s_2\) is binary digit
sum. In particular, \(N_{2,9}(\gamma)=512\), if proved by genuine
relation-preserving maps, would exclude degree at most 10.
\end{corollary}
\begin{proof}
Full counts in degree \(d\le D\) force \(p^k\mid d!\mid D!\).
Conversely, for every \(k\le v_p(D!)\),
Theorem~\ref{wr:thm-camouflage} supplies a witness of degree at most
\(D\) in any desired interval.
\end{proof}
No assertion about \(N_{2,9}(\gamma)\) is proved by this corollary.

\section{Real limits, arithmetic places, and invariant analytic data}
\label{wr:sec-invariants}

\subsection{The original approximants and denominator support}

The real convergence in \eqref{wr:eq-real-approx} has a different
\(2\)-adic behavior:
\begin{equation}\label{wr:eq-u-valuation}
 v_2(u_n)=-n.
\end{equation}
Indeed \(1/2^n\) is the unique term of \(H_{2^n}\) with valuation
\(-n\); every denominator in \(s_n\) is odd. Thus the original upper
approximants cannot be reduced directly to integral residue rings at
2. The later simultaneous real/\(2\)-adic construction repairs this
particular defect but still requires relation preservation. Equation
\eqref{wr:eq-u-valuation} is retained as the exact scope of the older
failed construction, not as an obstruction to all approximations.

For a nonzero reduced rational \(r=a/b\), \(|r|\le B\), and prescribed
\(v_p(r)\ge k_p\ge0\) at finitely many primes, one necessarily has
\begin{equation}\label{wr:eq-denominator-budget}
 b\ge B^{-1}\prod_p p^{k_p}.
\end{equation}
The numerator is divisible by that product, and \(|a|\le Bb\).
Moreover a sequence with denominator primes in a fixed finite set
\(S\), which converges in \(\R\) and at every place of \(S\), is
eventually constant. Its differences are small at all these places,
whereas a nonzero element of \(\Z[S^{-1}]\) satisfies
\(|r|_\infty\prod_{p\in S}|r|_p\ge1\), by the product formula and
integrality outside \(S\).

Uniform support is indispensable. The sequence
\begin{equation}\label{wr:eq-all-places-counterexample}
 r_n=\frac{(n!)^2}{2(n!)^2+1}
\end{equation}
converges to \(1/2\) in \(\R\) and to 0 in every fixed \(\Q_p\).
For large \(n\), the denominator is a \(p\)-adic unit and the numerator
valuation is \(2v_p(n!)\to\infty\). Nevertheless no unital map
\(\Z[1/2]\to\mathbb F_2\) exists. Even convergence at every finite
place does not imply descent of the real limit's relations.
The assertion concerns separate convergence at each place, not
convergence in the adelic restricted-product topology. More generally,
finite Chinese remaindering with sufficiently large denominators,
followed by a diagonal construction, realizes arbitrary prescribed
\(p\)-adic integer limits together with an arbitrary real limit.

\subsection{Exact valuations of the canonical rational midpoint}

Use the midpoint \(c_D\) of \eqref{inf:eq:midpoint}, put
\(N=2^k\), and write \(c_D=A_D/B_D\) in lowest terms, \(B_D>0\).
To avoid confusing the small logarithmic tail with a rational
approximation, denote that tail bound here by \(\rho_k\), as in
\eqref{wr:eq-harmonic-approximant}. Then
\[
 c_D=H_N-ks_k-\frac1{2N}-\frac{k\rho_k}{2}.
\]
The canonical schedule has \(k\ge6\) and \(N/2>\max(D,2k+1)\).
Let \(\Pi_N=\prod_{N/2<p<N}p\), over primes in the indicated interval.
Only the harmonic term \(1/p\) is nonintegral at such a prime, whereas
\(-1/(2N)\) is the unique term with smallest 2-adic valuation.
Consequently
\begin{equation}\label{wr:eq-midpoint-valuations}
 v_2(c_D)=-(k+1),\qquad v_p(c_D)=-1\quad(N/2<p<N).
\end{equation}
If \(\ell_N=\operatorname{lcm}(1,\ldots,N)\), inspection of all
explicit denominators gives
\begin{equation}\label{wr:eq-midpoint-denominator}
 2N\Pi_N\mid B_D,\qquad B_D\mid8\cdot9^k\ell_N,\qquad
 N^2<B_D\le8\cdot9^k4^{N-1}.
\end{equation}
The strict lower bound uses Bertrand's postulate; the last elementary
upper bound follows by iterating
\(\ell_{2n}/\ell_n\mid\binom{2n}{n}\le4^n\).

For every polynomial \(P(X)=\sum_{j=0}^D h_jX^j\) of degree exactly
\(D\), with \(H(P)\le D\), the leading term is uniquely smallest at
these places. Hence
\begin{equation}\label{wr:eq-midpoint-polynomial}
 v_p(P(c_D))=-D\ (N/2<p<N),\qquad
 v_2(P(c_D))=v_2(h_D)-D(k+1).
\end{equation}
Equivalently the homogenized integer numerator
\[
 T_P(A_D,B_D)=\sum_{j=0}^D h_jA_D^jB_D^{D-j}
\]
is coprime to \(\Pi_N\) and has 2-adic valuation \(v_2(h_D)\).
In particular
\[
 |P(c_D)|\ge B_D^{-D}\ge(8\cdot9^k4^{N-1})^{-D}.
\]
This lower bound is compatible with the canonical small real value.
Its guaranteed logarithmic scale is \(-O(DN)\), whereas the target
has scale \(-D\log D+O(\log D)\). Large denominator valuations permit
small real values in the product formula; they are not numerator
divisibilities.

Under a hypothetical factorization \(P_D=f^rC\), the large-prime
valuations satisfy
\[
 -r\deg f-(D-r\deg f)=-D.
\]
They impose no inconsistent constraint on \(r\). A concrete comparison makes this compatibility explicit:
for \(D\ge3\), put
\[
 c^*=\frac12+\frac1{2N\Pi_N},\quad
 r=\left\lfloor\frac{\log D}{\log3}\right\rfloor,\quad
 Q_D(X)=(2X-1)^rX^{D-r}.
\]
Then \(\deg Q_D=D\), \(H(Q_D)\le3^r\le D\),
\(Q_D(1/2)=0\), while
\[
 0<Q_D(c^*)=(N\Pi_N)^{-r}(c^*)^{D-r}<2U_D.
\]
Its valuations coincide with \eqref{wr:eq-midpoint-valuations} and
\eqref{wr:eq-midpoint-polynomial}, with leading coefficient \(2^r\).
The comparator is not the actual harmonic midpoint or the canonical
minimizer. It shows that these particular real and local properties
can coexist with a repeated algebraic factor; an additional estimate
using the actual harmonic numerator would be required to separate them.

\subsection{All Hankel tests and a freely movable mean}

Let \(E_0\) have the standard exponential distribution and define
\[
 G_a=-\log E_0+a-\gamma\qquad(a\in\R).
\]
The Gamma integral and its Taylor expansion \cite{wr:dlmf} give
\begin{equation}\label{wr:eq-gumbel-mgf}
 M_a(z)=\mathbb E(e^{zG_a})
 =e^{(a-\gamma)z}\Gamma(1-z)
 =\exp\left(az+\sum_{n\ge2}\frac{\zeta(n)}n z^n\right),
 \qquad |z|<1.
\end{equation}
Thus \(\mathbb E G_a=a\). If \(m_n(a)=\mathbb E(G_a^n)\),
\begin{equation}\label{wr:eq-moment-recurrence}
 m_0(a)=1,\qquad
 m_{n+1}(a)=a m_n(a)+
 \sum_{j=1}^n\binom nj j!\zeta(j+1)m_{n-j}(a).
\end{equation}

\begin{proposition}[Shift invariance of all principal Hankel determinants]
\label{wr:prop-hankel-shift}
For every \(N\ge1\) and every real \(a\),
\[
 \Delta_N(a)=\det(m_{i+j}(a))_{0\le i,j<N}=Delta_N(0)>0.
\]
In particular,
\[
 \Delta_2=\zeta(2),\qquad
 \Delta_3=2\zeta(2)^3+6\zeta(2)\zeta(4)-4\zeta(3)^2.
\]
\end{proposition}
\begin{proof}
Translation of the polynomial basis \(1,X,\ldots,X^{N-1}\) is a
triangular change of basis with diagonal entries one. The corresponding
moment matrices are congruent by this matrix, preserving determinant.
The density of \(G_a\) is positive on all of \(\R\), so the integral
of the square of a nonzero polynomial is strictly positive. Its moment
Gram matrix is therefore positive definite.
\end{proof}

The entire infinite family is compatible with every mean. For example,
when \(a=1/2\), \(P(X)=2X-1\) satisfies \(P(a)=0\), while
\(\mathbb E(P(G_a)^2)=4\zeta(2)>0\). Positivity of a random polynomial's
square is not nonvanishing at its mean. The later tensor-statistic
construction refines this distinction by computing the exact variance
terms and their signed cancellation; it does not turn this older
Hankel family into an arithmetic detector.

\subsection{Renormalization, multiplication, and the full L\'evy measure}

\begin{proposition}[Complete maximum fixed-point family]\label{wr:prop-max-fixed}
Let \(F\) be a continuous distribution function with
\(0<F(t)<1\) for every real \(t\). Then
\[
 F(t+\log m)^m=F(t)\quad(m=2,3,\ldots)
\]
holds if and only if \(F(t)=\exp(-C e^{-t})\) for some \(C>0\).
\end{proposition}
\begin{proof}
Put \(h=-\log F>0\). The equations give
\(h(t+\log(m/n))=(n/m)h(t)\) for positive integers \(m,n\).
Density of the positive rationals and continuity extend this to
\(h(t+s)=e^{-s}h(t)\) for all real \(s\). Take \(t=0\).
Direct substitution proves the converse.
\end{proof}
The mean in this family is \(\gamma+\log C\), so every real mean
occurs. Equivalently, for independent copies,
\(\max(G_{a,1},\ldots,G_{a,m})-\log m\) has the law of \(G_a\) for
all \(a\) and all \(m\).

Gauss multiplication applied to \eqref{wr:eq-gumbel-mgf} gives
\begin{align}
 \prod_{r=0}^{m-1}M_a\!\left(\frac{z+r}{m}\right)
 &= (2\pi)^{(m-1)/2}m^{z-1/2}
 e^{(a-\gamma)(m-1)/2}M_a(z),\label{wr:eq-mult-un}\\
 \prod_{r=0}^{m-1}
 \frac{M_a((z+r)/m)}{M_a(r/m)}
 &=m^zM_a(z).\label{wr:eq-mult-normalized}
\end{align}
These identities hold initially where the factors are finite and then
meromorphically. Normalization at \(z=0\) removes the entire drift
factor; all the normalized identities therefore hold at every real
mean.

There is also a construction which never uses a decimal value of
\(\gamma\). For independent standard exponentials \(E_n\), the sums
\[
 S_{N,a}=a+\sum_{n=1}^N\frac{E_n-1}{n}
\]
are Cauchy in \(L^2\), and their characteristic functions converge to
\(M_a(it)\) by the Gamma product. Equivalently,
\(\max(E'_1,\ldots,E'_N)-H_N+a\) has the same finite law: both
uncentered moment generating functions are
\(\prod_{n=1}^N n/(n-z)\). This construction works at rational \(a\)
just as well as at \(a=\gamma\).

\begin{proposition}[Infinite divisibility with fixed L\'evy measure]
\label{wr:prop-levy}
Every \(G_a\) is infinitely divisible, and for \(\Re z<1\),
\begin{equation}\label{wr:eq-levy}
 \log M_a(z)=az+
 \int_0^\infty\frac{e^{zx}-1-zx}{x(e^x-1)}\,dx.
\end{equation}
Only the drift \(a\) varies; the positive L\'evy measure is
\(\mathbf1_{x>0}\,dx/[x(e^x-1)]\).
\end{proposition}
\begin{proof}
Near zero the compensated integrand is bounded; at infinity it is
integrable in the stated half-plane. Expanding for \(|z|<1\) and using
\[
 \int_0^\infty\frac{x^{n-1}}{e^x-1}\,dx=(n-1)!\zeta(n)
 \qquad(n\ge2)
\]
proves the identity, followed by analytic continuation. To construct
convolution roots directly, take independent Gamma variables
\(V_{n,q}\) of shape \(1/q\) and scale \(1/n\). The centered sum
\[
 a/q+\sum_{n\ge1}\left(V_{n,q}-\frac1{qn}\right)
\]
converges in \(L^2\), with variance \(\zeta(2)/q\). The sum of \(q\)
independent copies has the law of the limit of \(S_{N,a}\), because
Gamma shapes add.
\end{proof}

The exponentiated variable is \(e^{-G_a}=\lambda_a E_0\), where
\(\lambda_a=e^{\gamma-a}\). Its Laplace transform and moments satisfy
\[
 L_a(s)=\frac1{1+\lambda_as},\qquad
 (-1)^jL_a^{(j)}(s)=\frac{j!\lambda_a^j}{(1+\lambda_as)^{j+1}}>0,
 \qquad \mathbb E(e^{-nG_a})=n!\lambda_a^n.
\]
Thus complete monotonicity and the positive moment recursion also
allow every mean. Gamma reflection gives two further shift-invariant
identities:
\[
 M_a(z)M_a(-z)=\frac{\pi z}{\sin\pi z},\qquad
 |M_a(it)|^2=\frac{\pi t}{\sinh\pi t}.
\]
Their modulus discards the phase which carries the location parameter.

Within this family, the normalization selecting the actual constant is
exactly
\begin{equation}\label{wr:eq-normalization}
 a=\gamma\ \Longleftrightarrow\ M_a(-1)=1
 \Longleftrightarrow\lambda_a=1
 \Longleftrightarrow F_a(0)=e^{-1}
 \Longleftrightarrow e^{(a-\gamma)(m-1)/2}=1
\end{equation}
for one fixed integer \(m\ge2\). These are several forms of one
normalization, not independent arithmetic constraints. Subtracting
\(\log N\) from the finite maximum selects \(\gamma\); subtracting
\(H_N-a\) selects \(a\). Prescribing the divergent centering term alone
leaves a finite location freedom.

Hermite--Lindemann implies that among algebraic real \(a\), at most one
can have \(\lambda_a\) algebraic: two such values would give the
algebraic quotient \(e^{b-a}\) for nonzero algebraic \(b-a\). Under A,
that unique value is precisely \(a=\gamma\). This is compatible with
A. Proving \(e^\gamma\) algebraic would instead prove \(\gamma\)
transcendental, since \(\gamma\ne0\), but none of the preceding
identities proves that assertion.

Finally, none of the functions \(M_a\) is itself an arithmetic
E-function or G-function. Its first two Taylor coefficients are
\[
 [z]M_a=a,\qquad [z^2]M_a=\frac{a^2+\zeta(2)}2.
\]
If both were algebraic, \(\zeta(2)=\pi^2/6\) would be algebraic. The
same argument survives any nonzero algebraic rescaling of \(z\).
This excludes a direct invocation of arithmetic E/G-function theorems
on this exact family, not every possible transformation of it.

\subsection{The earlier curvature and capacity entry tests}

The subsequent tensor-category classification supersedes the narrow
rank-one twist test, but the original all-prime formula remains useful.
With \(\Phi=\Ein\) and
\[
 Y_a=(1,e^{-z},\Phi-a)^T,\qquad
 A(z)=\begin{pmatrix}0&0&0\\0&-1&0\\z^{-1}&-z^{-1}&0\end{pmatrix},
 \qquad\nabla=\partial-A,
\]
the \(p\)-curvature in characteristic \(p\), including \(p=2\), is
\begin{equation}\label{wr:eq-curvature}
 \psi_p=
 \begin{pmatrix}
 0&0&0\\0&1&0\\
 z^{-p}&-z^{-p}\displaystyle\sum_{j=0}^{p-1}z^j/j!&0
 \end{pmatrix}.
\end{equation}
For completeness, \(\nabla e_1=-z^{-1}e_3\),
\(\nabla e_2=e_2+z^{-1}e_3\), and \(\nabla e_3=0\). Wilson's
congruence gives the first column; iteration gives the second lower
entry as \(\sum_{j=0}^{p-1}(-1)^j j!z^{-j-1}\), which is the displayed
truncated exponential. The matrix has rank two, kernel generated by
\(e_3\), and both characteristic and minimal polynomials
\(T^2(T-1)\).

Tensoring with any rank-one connection shifts the eigenvalues to
\(c_p,c_p,c_p+1\). A fixed nonconstant rational pullback \(z=h(t)\)
multiplies their difference by \(h'(t)^p\), nonzero at all good primes;
a rational gauge conjugates curvature. Therefore these operations
cannot make the entire module globally nilpotent. The connections and
transformations must be defined over one fixed number field, with
only finitely many primes excluded. These are the standard tensor and
pullback rules \cite{at:Katz1982}; no general \(p\)-curvature conjecture is
assumed. The pure logarithmic subconnection spanned by \(e_1,e_3\)
does have nonzero nilpotent curvature, and its arbitrary integration
constant does not acquire an arithmetic classification from that fact.
The scalar \(a\) does not occur in \(A(z)\); changing it changes a
framed solution, not the differential module.

For the capacity route, the direct fixed-scaling candidate also fails
an explicit input test. If \(B,C\) are nonzero integers and \(p\nmid BC\),
then the \(z^p\) coefficient of \(B\Ein(Cz)\) has valuation \(-2\).
Thus no fixed integral scaling makes its coefficients all integral.
Multiplying coefficient \(n\) by \(n!\), the inverse Borel convention
used here, gives \(\log(1+z)\); the same coefficient test now has
valuation \(-1\). Differentiation gives a rational function but removes
the constant \(-a\) under investigation. These facts test only this
specific sequence of transformations. They do not prohibit an adelic,
nonlinear, multivariable, or moving-denominator capacity construction.

The analytic relationship to the A/B split remains conditional. With
\(\eta=\Ein(1)\) and \(\delta=\int_0^\infty e^{-t}/(1+t)\,dt\),
\(\eta=\gamma+\delta/e\). The known algebraic independence of
\(e,\eta\) \cite{at:Rivoal2012} implies, under A, the algebraic independence
of \(e,\delta\), because
\(\Qbar(e,\eta)=\Qbar(e,\delta)\). It also forces
\[
 \delta\in(\mathbf E\cap\mathbf D)\setminus\Qbar,
\]
where the two rings have the specified meanings of
\cite{pm:FischlerRivoal2024}. This contradicts the \emph{conjectural}
intersection equality, not a proved equality. Here it is \(\delta\in
\mathbf D\) that is used; division by \(e\) is not automatically
allowed merely because \(\mathbf D\) is a ring. The earlier displayed
Stokes invariance and the later curvature computation do not add the
missing arithmetic lifting theorem.

\section{A feasibility obstruction for replication near zero}
\label{wr:sec-orbit-obstruction}

This section concerns new auxiliary polynomials, not the canonical
minimum-gap outputs. Ordinary simultaneous pigeonholing on the grid
\(\{0,\ldots,D\}^{D+1}\), at \(K\) points in \((0,1)\), supplies a
nonzero polynomial of degree and height at most \(D\) with error
\[
 \exp\bigl(-(1/c+o(1))D\log\log D\bigr),
 \qquad K=\left\lfloor c\frac{\log D}{\log\log D}\right\rfloor.
\]
For the rational dilation orbit of a hypothetical degree-\(d\) algebraic
number, the resultant comparison forces exact zeros when \(c<1/d\),
whereas the product of the different minimal-polynomial copies exceeds
the Mahler budget only when \(c>3/(2d)\). These two windows do not
intersect. Trying to bridge the gap by an arbitrary stronger error
requirement can make the proposed construction impossible even without
assuming algebraicity.

\begin{theorem}[Universal obstruction for the orbit tending to zero]
\label{wr:thm-orbit-no-go}
Fix \(x\in(0,1)\), \(c>3/2\), and \(\kappa>c\), and put
\(K=\lfloor c\log D/\log\log D\rfloor\). For all sufficiently large
\(D\), no nonzero \(P\in\Z[X]\) with \(\deg P,H(P)\le D\) satisfies
\begin{equation}\label{wr:eq-orbit-impossible}
 \max_{1\le j\le K}|P(x/j)|
 \le\exp\left(-\kappa\frac{D\log D}{K}\right).
\end{equation}
\end{theorem}
\begin{proof}
Write \(P=X^sR\), where \(R(0)\) is a nonzero integer. If the right side
of \eqref{wr:eq-orbit-impossible} is \(\varepsilon\), then
\(|R(x/j)|\le\varepsilon(K/x)^D=:\eta\). Interpolate these values by
\(T\) of degree at most \(K-1\). Node separation is at least \(x/K^2\),
so the Lagrange basis yields
\[
 \|T\|_1\le K\eta(2K^2/x)^{K-1}=:\theta_D.
\]
Consequently
\[
 \log\theta_D
 \le-\kappa\frac{D\log D}{K}+D\log(K/x)+O_x(K\log K)
 =-(\kappa/c-1+o(1))D\log\log D\to-\infty.
\]
The real polynomial \(Q=R-T\) has all \(K\) roots \(x/j\), with
\(|Q(0)|\ge1/2\) for large \(D\), while
\[
 M(Q)\le\|Q\|_2\le D\sqrt{D+1}+1/2.
\]
The constant-term form of the Mahler product gives independently
\[
 M(Q)\ge |Q(0)|\prod_{j=1}^K\frac jx
 \ge\tfrac12K!x^{-K}=D^{c+o(1)}.
\]
This contradicts the upper bound \(D^{3/2+o(1)}\). Integrality is used
before interpolation, to bound the nonzero constant coefficient below;
the Mahler inequalities themselves also hold for real or complex
coefficients.
\end{proof}

The contradiction eliminates the specified demand, not either arithmetic
case. Replacing the nodes by \(xj/(j+1)\) removes this particular
geometric obstruction: their inverse product is only
\((K+1)x^{-K}\). A real-coefficient polynomial with these roots and
constant coefficient one has coefficient norm at most
\((1+2/x)^K=D^{o(1)}\). Its coefficients need not be integers. Thus this
observation neither constructs the desired integer family nor proves
its feasibility. The bounded-orbit criterion of Section~\ref{inf:sec:orbit}
retains the arithmetic factor cost and the missing simultaneous
approximation assumption.

There is also no free recovery of the obstruction by the translation
\(Y=x-X\). For transcendental \(x\), translating an integer polynomial
introduces coefficients involving \(x\); the new constant term is
uncontrolled. For instance, if
\(B_K(X)=\prod_{j=1}^K(X-xj/(j+1))\), then
\(B_K(x)=x^K/(K+1)!\), exactly compensating the apparent factorial from
the translated roots.

A related audit of the canonical multiplicity theorem is worth retaining.
If \(f^{r_D}\mid P_D\), \(H(P_D)\le D\), and \(\deg P_D\le D\), then
\[
 r_D\le\frac{\log D+\frac12\log(D+1)}{\log M(f)}.
\]
For \(0<x<1\) algebraic, \(M(f)>1\); the usual available competitors
\(f^mX^{D-md}\), with \(\|f\|_1^m\le D\), already allow
\(m\) of order \(\log D\). Hence the lower logarithmic multiplicity
and this upper logarithmic budget are compatible. Similarly, after
removing roots at zero from an integer polynomial \(R\), the number
of its roots in \(|z|\le b<1\), counted with multiplicity, is at most
\[
 \frac{\log D+\frac12\log(D+1)-\log|R(0)|}{\log(1/b)}.
\]
This is \(O(\log D)\), not the smaller bound that would contradict the
canonical algebraic branch.

\section{Effective nonvanishing tests and the limits of generic world search}
\label{wr:sec-effective-logic}

\subsection{What can be certified by exact rational computation}

A computable real \(x\) is supplied with rational approximants \(q_n\)
and a guaranteed bound \(|q_n-x|\le2^{-n}\). Choose a rational \(R\)
containing these values in \([-R,R]\). For
\(P(X)=\sum_{j=0}^d a_jX^j\), put
\(L_P=\sum_{j=1}^d j|a_j|R^{j-1}\). Then
\begin{equation}\label{wr:eq-semidecision}
 |P(q_n)|>L_P2^{-n}\quad\Longrightarrow\quad P(x)\ne0.
\end{equation}
The test is exact and terminates precisely when \(P(x)\ne0\): once
\(L_P2^{-n}<|P(x)|/3\), the inequality holds. A zero value can never
pass. If \([a_n,b_n]\subset[0,1]\) encloses \(\gamma\), its midpoint
\(r_n\) gives the common finite-box certificate
\begin{equation}\label{wr:eq-box-certificate}
 |P(r_n)|>\frac{HD(D+1)}4(b_n-a_n)
 \quad\Longrightarrow\quad P(\gamma)\ne0
 \quad(\deg P\le D,\ H(P)\le H).
\end{equation}

\begin{proposition}[Finite-box termination and least survivor]
\label{wr:prop-survivor}
The parallel test of all nonzero integer polynomials of degree at most
\(D\) and height at most \(H\) terminates if and only if none vanishes
at \(x\). Consequently termination for every pair \((D,H)\) is
equivalent to transcendence of \(x\).

Effectively enumerate the primitive irreducible integer polynomials,
with positive leading coefficient, as \(f_1,f_2,\ldots\), and delete
candidates whenever \eqref{wr:eq-semidecision} certifies nonvanishing.
Let \(m_s(x)\) be the least index not deleted after stage \(s\), with
all candidates eventually tested. Then \(m_s\) is computable and
nondecreasing, and
\[
 \begin{array}{ll}
 x\in\Qbar:&m_s(x)\text{ eventually equals the minimal-polynomial index},\\
 x\notin\Qbar:&m_s(x)\longrightarrow\infty.
 \end{array}
\]
\end{proposition}
\begin{proof}
There are finitely many tests in a box, so if all its values are
nonzero their finite termination times have a finite maximum.
Conversely a zero value cannot be certified. For the survivor claim,
all candidates preceding a true minimal polynomial are eventually
deleted and that polynomial never is. In the transcendental case each
fixed finite initial list is eventually deleted. Permanent deletion
makes the least surviving index nondecreasing.
\end{proof}

Long observed stability does not certify eventual stability. If
algebraicity and explicit degree/height bounds are supplied independently,
exact isolation of the finitely many possible roots does identify the
number from sufficiently narrow intervals. This is a result under a
finite-candidate promise, not a proof of that promise for \(\gamma\).
The archived degree-3, height-4 calculation checked 2,927 primitive
polynomials and certified \(|P(\gamma)|>10^{-4}\) throughout that box.
It is an independent reproduction of a subset of the earlier degree-4,
height-4 verification, not a new exclusion range. The exact scripts and
certificates are retained in the accompanying archive.

\subsection{Two explicit reductions from the halting problem}

\begin{theorem}[No uniform arithmetic semidecision for computable reals]
\label{wr:thm-no-generic-decider}
On inputs consisting of a program for a fast rational Cauchy name of a
real number, neither algebraicity nor transcendence has a semidecision
procedure valid for every such input. This remains true under the
promise that the input is either rational or transcendental.
\end{theorem}
\begin{proof}
Fix a computable transcendental \(\tau\in(0,1)\), and let \(H_e=t\)
mean that a machine \(e\) first halts at step \(t\ge1\). Define
\[
 x_e=\begin{cases}0,&e\text{ never halts},\\2^{-t}\tau,&H_e=t.
 \end{cases}
\]
These reals have fast names uniformly in \(e\): to precision \(2^{-n}\)
simulate \(n\) steps; if the machine has not halted, any future value
is smaller than \(2^{-n}\), and otherwise compute the indicated multiple
of \(\tau\). Here \(x_e\) is algebraic exactly when \(e\) never halts.
A semidecision procedure for algebraicity would enumerate nonhalting
machines, contradicting the undecidability of halting.

For the other direction choose rational \(q_t\) with
\(|q_t-\tau|\le2^{-t-2}\), and put
\[
 y_e=\begin{cases}\tau,&e\text{ never halts},\\q_t,&H_e=t.
 \end{cases}
\]
To precision \(2^{-n}\), simulate \(n+2\) steps. If no halt has occurred,
output \(q_{n+2}\); even a later halt changes the value by less than
the error budget. If a halt has occurred, output the corresponding
rational exactly. Thus \(y_e\) is uniformly computable and is
transcendental exactly when \(e\) never halts. This excludes a uniform
semidecision for transcendence. Both constructions have the stated
rational/transcendental promise. Nonzero rational affine changes put
them in any prescribed rational open interval.
\end{proof}

One may take \(\tau=\sum_{m\ge1}10^{-m!}\). Its partial sums
\(p_n/q_n\), with \(q_n=10^{n!}\), satisfy
\(0<\tau-p_n/q_n<2q_n^{-n-1}\). If a fixed nonzero integer polynomial
\(P\) of degree \(d\) vanished at \(\tau\), sufficiently late partial
sums would not be its roots, so \(|P(p_n/q_n)|\ge q_n^{-d}\),
contradicting the derivative bound times the displayed tail. This also
makes the auxiliary transcendental used in the reduction explicit.
The general undecidability background is \cite{wr:turing}; the fast
names and reductions above are supplied directly.

Theorem~\ref{wr:thm-no-generic-decider} does not exclude a proof for the
specific constant \(\gamma\), nor establish independence of its
arithmetic nature from any axiom system.

\subsection{Finite observations, infinitesimals, and proof search}

Any finite set of strict polynomial inequalities true at a real point
persists in a neighborhood. That neighborhood contains rationals,
algebraic numbers of arbitrarily high degree, and transcendental
numbers. Hence these finite tests alone cannot classify the point.
This does not apply to an exact algebraic equality or assert that the
comparison points satisfy \(\gamma\)'s full analytic definition.

An explicit non-Archimedean example locates a different possible
confusion. In the ordered Laurent series field \(\R((t))\), with \(t\)
a positive infinitesimal, \(c=\gamma+t\) belongs to every strict
rational interval containing \(\gamma\), but \(c\ne\gamma\) and is
transcendental over \(\R\). In the actual complete real field the same
shrinking intervals determine one element. Passing to an infinitesimal
field changes that uniqueness property; it does not construct another
Euler constant satisfying all the original definitions. Likewise the
formal map of Proposition~\ref{wr:prop-freedom} may preserve the
inherited order into \(K_x((t))\), although a nontrivial real-field
self-automorphism cannot preserve order.

Finally, fix a consistent effectively axiomatized classical theory
\(T\) in which the algebraicity assertion \(A\) is expressible.
Enumerating contradiction proofs in \(T+A\) and \(T+\neg A\) in
parallel terminates exactly when \(T\) decides \(A\), because the
deduction theorem gives
\[
 T+A\vdash\bot\ \Longleftrightarrow\ T\vdash\neg A,
 \qquad
 T+\neg A\vdash\bot\ \Longleftrightarrow\ T\vdash A.
\]
The excluded-middle theorem alone does not guarantee termination. This
is conditional proof-search bookkeeping, not a claim that either
extension is inconsistent or that Euler's constant is independent of
ZFC.

\begin{remark}[What this part supplies]\label{wr:rem-status}
The completed general results are the exact local-count dichotomy,
the height-free bound, the optimal finite-camouflage degree, the
formal-motion classification, the analytic invariant families, the
specific replication obstruction, and the effective-test theorems.
Their applications to \(\gamma\) are equivalences, sufficient
conditions, or explicitly limited failures. The missing positive input
for the local route remains an unbounded family of \emph{actual}
relation-preserving maps. The missing input for value counting is a
superexponential lower bound on \emph{distinct evaluated values}.
Neither follows from a collection of surviving candidates, from the
number of coefficient vectors, or from shift-invariant positivity.
\end{remark}

% END INTEGRATED PART VI

\clearpage
\part{Arithmetic transfer and its obstructions}\label{master:transfer}
% BEGIN INTEGRATED PART VII
\section{Arithmetic transfer: notation and scope}
\label{at:sec:introduction}

This part preserves the arithmetic-transfer revision as a coherent
collection of four mechanisms with explicit hypotheses and proofs. The
arithmetic question is whether a nonzero polynomial $P\in\Q[X]$ can satisfy
$P(\gamma)=0$. Calling the algebraic and transcendental alternatives two
``worlds'' introduces no additional axiom: to exclude one alternative one
must derive a contradiction from it and independently established facts.
None of the constructions below produces that contradiction.

Throughout,
\begin{align}
 \gamma&=\lim_{n\to\infty}(H_n-\log n),& H_n&=\sum_{j=1}^n\frac1j,\notag\\
 \Ein(z)&=\sum_{n\ge1}\frac{(-1)^{n+1}z^n}{n\,n!},&
 E_1(x)&=\int_x^\infty\frac{e^{-t}}t\,dt\quad(x>0),\label{at:eq:notation}\\
 \delta&=\int_0^\infty\frac{e^{-t}}{1+t}\,dt=eE_1(1),&
 \eta&=\Ein(1)=\gamma+\delta/e.\notag
\end{align}
On the positive real branch, and by continuation on compatible branches,
$\Ein(z)=\gamma+\log z+E_1(z)$; see \cite{at:DLMF}.
We use $|p|_p=p^{-1}$. Congruences for rational numbers with denominator
prime to $p$ are interpreted in $\Z_p$. A common denominator for algebraic
numbers is the least positive rational integer making them all algebraic
integers. Differential modules have derivation $\partial=d/dz$ over
$K=\Qbar(z)$.

There are four separate assertions in this part. The approximation and
grid results concern the preservation of polynomial relations under
different completions. The differential classification concerns the finite
tensor category of one specified module, not arbitrary transformations of
functions. The statistical result concerns kernels valid for an entire
location family, not every possible identity special to $\gamma$. The
denominator result concerns an explicitly defined formal quotient; no
closure under division at a zero of a Borel sum is assumed.

\paragraph{Attribution and scope.}
The approximation argument is an elementary instance of simultaneous
approximation in different completions. The statistical constructions use
classical U-statistics and finite differences \cite{at:Hoeffding1948}; their
identities are proved directly. The differential argument uses the
primitive-dependence theorem and differential Galois theory
\cite{at:Srinivasan2008,at:Katz1982}. The factorial-denominator argument uses the
regular-singularity theorem for $G$-operators \cite{pm:FischlerRivoal2016}.
Its prime-discrepancy statement is also a consequence of
\cite[Theorem 1]{at:Ernvall2018}, as proved below. We make no claim of research
priority for these auxiliary results. The mixed-series comparison uses
the arithmetic framework of \cite{pm:FischlerRivoal2024}; it does not prove
the missing lifting assertion. This part is a research note documenting
constructions and obstructions, not a claimed solution of the constant
problems.

\section{Explicit real and \texorpdfstring{$2$}{2}-adic approximation}
\label{at:sec:approximation}

\begin{theorem}[Prescribed simultaneous limits]
\label{at:app:construction}
Let $x\in\R$, $a\in\Z_2$, and let $[l_n,u_n]$ be rational intervals
containing $x$ with width at most $(2/9)^n$. Write $c_n=(l_n+u_n)/2$,
$t_n=2^n$, $q_n=9^n$, and let $a_n$ be the residue of $a$ modulo $t_n$.
Choose $0\le b_n<t_n$ with $b_n\equiv q_na_n\pmod{t_n}$ and set
\begin{equation}
 r_n=\frac{b_n+t_n\operatorname{round}((q_nc_n-b_n)/t_n)}{q_n}.
 \label{at:eq:rounding}
\end{equation}
Then $r_n\in\Z[1/3]$, and
\begin{equation}
 |r_n-x|\le(2/9)^n,\qquad r_n\equiv a\pmod{2^n}.
 \label{at:eq:rounding-bounds}
\end{equation}
In particular, $r_n\to x$ in $\R$ and $r_n\to a$ in $\Q_2$.
\end{theorem}
\begin{proof}
Rounding gives $|r_n-c_n|\le t_n/(2q_n)$, while the interval gives
$|c_n-x|\le t_n/(2q_n)$. Their sum is the real estimate. Since $q_n$
is odd, it is a unit in $\Z_2$, and the numerator congruence gives the
$2$-adic estimate. The residues are compatible because they are residues
of one fixed $a$. Either tie convention for rounding is valid.
\end{proof}

For computable $x$ and $a$ the construction is an algorithm. A
noncomputable $a$ requires its residues as oracle input. Thus no
computability assertion about arbitrary $2$-adic integers is implicit.

\begin{lemma}[Rational intervals for $\gamma$]
\label{at:app:intervals}
For $m\ge1$, define
\begin{equation}
 s_m=2\sum_{j=0}^{m-1}\frac1{(2j+1)3^{2j+1}},\quad
 R_m=\frac3{4(2m+1)9^m},\quad
 U_m=H_{2^m}-ms_m,\quad \beta_m=2^{-m}+mR_m.
\end{equation}
Then $U_m-\beta_m<\gamma<U_m$, and $\beta_{3n}<(2/9)^n$ for $n\ge1$.
\end{lemma}
\begin{proof}
The expansion $\log2=2\sum_{j\ge0}((2j+1)3^{2j+1})^{-1}$ has a
positive tail less than $R_m$, by a geometric-series bound. The integral
comparison for the harmonic sum gives
$0<H_N-\log N-\gamma<1/N$. Combining the two bounds proves the interval.
Finally
\[
 \frac{\beta_{3n}}{(2/9)^n}
 =\left(\frac9{16}\right)^n+
 \frac{9n}{4(6n+1)162^n}
 \le\frac9{16}+\frac3{8\cdot162}<1.
\]
\end{proof}
Consequently, Theorem~\ref{at:app:construction} gives explicit rational
approximants to $\gamma$ without reading its decimal expansion or assuming
its arithmetic nature. The harmonic sum at $m=3n$ has $2^{3n}$ terms;
the construction is not asserted to be computationally efficient.

\begin{proposition}[Failure of relation preservation]
\label{at:app:counterexample}
The simultaneous limits in Theorem~\ref{at:app:construction}, including
compatibility at every level, do not imply
$P(x)=0\Rightarrow P(a)=0$ for $P\in\Z[X]$.
\end{proposition}
\begin{proof}
Take $x=c_n=1/3$ and $a=0$. Then
\[
 r_n=\frac{2^n}{9^n}\operatorname{round}\left(\frac{9^n}{3\cdot2^n}\right)
 \longrightarrow\frac13\text{ in }\R,
 \qquad r_n\longrightarrow0\text{ in }\Q_2.
\]
For $P(X)=3X-1$ we have $P(x)=0$, but
$P(r_n)\equiv-1\pmod{2^n}$ and $P(a)=-1$. Evaluation therefore
defines a homomorphism on $\Z[X]$, but it does not descend to
$\Z[x]=\Z[X]/\ker(P\mapsto P(x))$ with $x\mapsto a$.
The choice $x=1/2$, $a=0$, $P=2X-1$ gives another example.
\end{proof}

\begin{proposition}[Denominator escape]
\label{at:app:escape}
Suppose $r_n\in\Z[1/3]$ converges to $x$ in $\R$ and to $a$ in
$\Q_2$. Unless $x\in\Z[1/3]$ and $a$ is its natural $2$-adic image,
the exponents of $3$ in the reduced denominators of $r_n$ tend to infinity.
For the construction at $x=\gamma$, $a=0$, one has
$|r_n|_3>2^n$ for all sufficiently large $n$.
\end{proposition}
\begin{proof}
If infinitely many reduced denominators have exponent at most $e$, that
subsequence lies in the discrete real set $3^{-e}\Z$. Real convergence
makes it bounded, so a further subsequence is constant. Uniqueness of
both limits yields precisely the stated exception. At $x=\gamma$ and
$a=0$, eventually $0<r_n<1$. Write $r_n=A'_n/3^{e_n}$ in lowest terms.
The numerator is a positive multiple of $2^n$. Hence
$|r_n|_3=3^{e_n}>A'_n\ge2^n$.
\end{proof}

\section{What is sufficient to transfer a relation?}
\label{at:sec:grid}

\begin{proposition}[An integer-zero sufficient condition]
\label{at:grid:transfer}
Let $r_n=A_n/q_n$ be reduced rationals with $q_n>0$ and $p\nmid q_n$,
$r_n\to x$ in $\R$, and $r_n\to a$ in $\Q_p$. If, for an integer
$d\ge1$, $q_n^d|r_n-x|\to0$, then every $P\in\Z[X]$ of degree at most
$d$ satisfies
\[
 P(x)=0\quad\Longrightarrow\quad P(r_n)=0\text{ eventually}
 \quad\Longrightarrow\quad P(a)=0.
\]
\end{proposition}
\begin{proof}
For $e=\deg P\le d$, the number $q_n^eP(r_n)$ is an integer. On a
fixed compact real interval containing $x$ and the tail of the sequence,
bound $|P'|$ by $C_P$. Under $P(x)=0$ its absolute value is at most
$C_Pq_n^d|r_n-x|\to0$, so it eventually vanishes. Polynomial continuity
in $\Q_p$ proves the last implication. Constant polynomials cause no
difficulty.
\end{proof}

This is a sufficient condition for this direct procedure, not a necessary
condition for arbitrary proofs. If infinitely many $r_n$ are distinct and
the condition holds for every fixed $d$, then $x$ is already
transcendental: a nonzero polynomial has finitely many roots.

\begin{theorem}[Common-denominator grid bound]
\label{at:grid:bound}
For $M\ge2$ distinct rationals $r_j=A_j/q$ with integer numerators and
the same positive denominator, every $x\in\R$ satisfies
\begin{equation}
 \max_{1\le j\le M}|r_j-x|\ge\frac{M-1}{2q}.
 \label{at:eq:grid-bound}
\end{equation}
For the $2^n$ distinct residues in construction~\eqref{at:eq:rounding},
at each $n$ this gives
\begin{equation}
 9^{nd}\max_{a_n\bmod2^n}|r_{n,a_n}-\gamma|
 \ge\frac{2^n-1}{2}\,9^{n(d-1)}\qquad(d\ge1).
\end{equation}
\end{theorem}
\begin{proof}
The largest and smallest of $M$ distinct integer numerators differ by
at least $M-1$. Their rational values have that difference divided by
$q$, and at least one lies at distance half the difference or more from
$x$. Distinct residues in the construction force distinct numerators.
\end{proof}

The theorem excludes a uniformly vanishing integer-zero budget for the
entire common-denominator grid. It does not exclude a favorable single
sequence, estimates using individual reduced denominators, or a different
relation-preserving construction. For one prescribed residue $b$ modulo
$t$ and a fixed denominator $q$, the exact minimax identity is
\[
 \sup_{x\in\R}\inf_{m\in\Z}
 \left|x-\frac{b+tm}{q}\right|=\frac{t}{2q}.
\]
Indeed the admissible points are an equally spaced lattice, and its
midpoints attain the worst distance.

\begin{proposition}[Finite denominator support]
\label{at:grid:finite-support}
Suppose a rational sequence converges in $\R$ and in every $\Q_p$,
and the prime factors of all sufficiently late reduced denominators
belong to a fixed finite set $S$. Then the sequence is eventually constant.
\end{proposition}
\begin{proof}
Convergence in each $\Q_p$, $p\in S$, bounds the denominator exponent
at that prime. A fixed positive integer $B$ therefore clears all late
denominators. The sequence $Br_n$ is a convergent real sequence of
integers and is eventually constant.
\end{proof}
The support condition cannot be omitted. The rationals
$r_n=(n!)^2/(2(n!)^2+1)$ converge to $1/2$ in $\R$ and to $0$ in
every fixed $\Q_p$. Controlling each fixed place separately does not
control primes that enter the denominators with $n$.

\section{Globally nilpotent objects in the tensor category of the Ein connection}
\label{at:ten:section}

Put $k=\Qbar$, $K=k(z)$, and $\partial=d/dz$.  Let $M$ be the
differential module represented by
\begin{equation}\label{at:ten:connection}
 Y'=A(z)Y,\qquad
 A(z)=\begin{pmatrix}
 0&0&0\\
 0&-1&0\\
 z^{-1}&-z^{-1}&0
 \end{pmatrix}.
\end{equation}
We write $\langle M\rangle_\otimes$ for the category obtained from $M$
by finitely many tensor products, duals, finite direct sums, subobjects,
and quotients over $K$.  Every object and every finite collection of
morphisms in this category descend to a number field.  An object is
\emph{globally nilpotent} if, for a number-field model, its
$p$-curvature is nilpotent at all but finitely many finite places of
good reduction.  This property is unchanged by finite extension of
the number field or by a rational change of basis.

We give a concrete application of standard Picard--Vessiot and
tensor-category theory.  No priority claim is made for the resulting
classification.  The classification concerns the differential module;
the additional analytic normalization containing $\gamma$ is discussed
separately at the end of the section.

\subsection{Primitive independence and the Picard--Vessiot field}

On a simply connected domain in $\C\setminus\{0\}$ containing a
positive real interval, choose the branches
\[
 E=e^{-z},\qquad L=\log z,\qquad J=E_1(z).
\]
Here $E_1$ is continued from its usual positive-real definition.  They
satisfy
\begin{equation}\label{at:ten:derivatives}
 E'=-E,\qquad L'=z^{-1},\qquad J'=-E/z,
 \qquad \Ein(z)=\gamma+L+J;
\end{equation}
the last normalization agrees with \cite[Eq.~6.2.4]{at:DLMF}.

\begin{lemma}\label{at:ten:no-primitive}
If $b,c\in\C$ and $R\in\C(z,E)$ satisfy
\[
 R'=b/z-cE/z,
\]
then $b=c=0$.
\end{lemma}

\begin{proof}
There is no nonzero algebraic function $r$ over $\C(z)$ satisfying
$r'=\lambda r$ for a nonzero constant $\lambda$.  Indeed, in a finite
extension containing $r$, the norm $q$ of $r$ would satisfy
$q'/q=n\lambda$, where $n$ is the extension degree.  The logarithmic
derivative of a nonzero rational function tends to zero at infinity,
which is impossible.  In particular, $E$ is transcendental over
$\C(z)$.

View $R$ as a rational function of the independent variable $E$, and
factor its denominator over an algebraic closure of $\C(z)$.  At a
nonzero finite pole $E=r(z)$, let
$a(z)(E-r)^{-s}$ be its leading principal part.  Its derivative has
coefficient $s a(r+r')$ at pole order $s+1$.  The prescribed derivative
has no such pole, so $r'=-r$, contrary to the preceding paragraph.
Consequently the only possible finite $E$-pole is zero, and
\[
 R=\sum_{n=u}^{v} r_n(z)E^n,\qquad r_n(z)\in\C(z).
\]
For $n\ne0,1$, comparison of coefficients gives
$r_n'-n r_n=0$, hence $r_n=0$.  At $n=0$ it gives $r_0'=b/z$;
the residue of a rational derivative is zero, so $b=0$.
At $n=1$ it gives $r_1'-r_1=-c/z$.  A finite pole of $r_1$ would
produce, upon differentiation, a pole of strictly higher order than
the right-hand side permits, including at zero.  Thus $r_1$ is a
polynomial, which forces $c=0$.
\end{proof}

\begin{lemma}\label{at:ten:independence}
The functions $E,L,J$ are algebraically independent over $\C(z)$.
Moreover, the constant field of
\[
 F=K(E,L,J)
\]
is exactly $k$.
\end{lemma}

\begin{proof}
The transcendence of $E$ was proved above.  Both $L$ and $J$ are
antiderivatives over $\C(z,E)$.  Their meromorphic-function extension
has constant field $\C$, so the no-new-constants hypothesis in the
Kolchin--Ostrowski theorem is satisfied.  If $L,J$ were algebraically
dependent over $\C(z,E)$, that theorem would give
$bL+cJ\in\C(z,E)$ for some $(b,c)\ne(0,0)$; see
\cite[Theorem~2.5]{at:Srinivasan2008}.  Differentiating contradicts
Lemma~\ref{at:ten:no-primitive}.

Now an element of $F$ with zero derivative is a complex constant $d$.
Write it as a quotient of polynomials in $E,L,J$ with coefficients in
$k(z)$.  Clearing the denominator and using the independence just
proved, comparison of any nonzero denominator coefficient shows that
$d\in k(z)$.  Since $d$ is constant, $d\in k$.
\end{proof}

\begin{proposition}\label{at:ten:group}
The field $F$ is a Picard--Vessiot field for $M$ over $K$, with
fundamental matrix
\begin{equation}\label{at:ten:fundamental}
 \Phi=\begin{pmatrix}1&0&0\\0&E&0\\L&J&1\end{pmatrix}.
\end{equation}
Its differential Galois group consists of the transformations
\begin{equation}\label{at:ten:automorphisms}
 E\longmapsto tE,\qquad L\longmapsto L+b,\qquad
 J\longmapsto tJ+c,
 \qquad t\in k^\times,\quad b,c\in k.
\end{equation}
It is isomorphic to $\Ga\times(\Ga\rtimes\Gm)$, with the action
specified below.
\end{proposition}

\begin{proof}
Direct differentiation gives $\Phi'=A\Phi$ and $\det\Phi=E$.
Its entries and inverse determinant generate $F$, which has no new
constants by Lemma~\ref{at:ten:independence}.  Thus it is a
Picard--Vessiot field.  Every differential automorphism must have
the form \eqref{at:ten:automorphisms}, by successive use of
\eqref{at:ten:derivatives}.  Conversely, algebraic independence shows
that every transformation in \eqref{at:ten:automorphisms} is a
well-defined differential automorphism.

In the convention $\sigma(\Phi)=\Phi g_\sigma$, its matrix is
\[
 g(t,b,c)=\begin{pmatrix}1&0&0\\0&t&0\\b&c&1\end{pmatrix}.
\]
Set $T(t)=g(t,0,0)$, $U_b(b)=g(1,b,0)$, and
$U_c(c)=g(1,0,c)$.  Then
\begin{equation}\label{at:ten:conjugation}
 T(t)U_c(c)T(t)^{-1}=U_c(c/t),\qquad
 [T,U_b]=[U_b,U_c]=1.
\end{equation}
These identities specify the asserted semidirect product.
\end{proof}

\subsection{The classification}

The solution functor in $F$ identifies $\langle M\rangle_\otimes$
with the finite-dimensional algebraic representations of this
Picard--Vessiot group.  In particular, invariant subspaces correspond
to differential submodules, and quotients correspond to quotients.
This is the standard tensor interpretation of differential Galois
theory; see \cite[Sections~I and IV]{at:Katz1982}.  Let $F_W\subset F$
denote the Picard--Vessiot subfield generated by a fundamental matrix
for an object $W$ under this identification.

\begin{theorem}\label{at:ten:classification}
For an object $W\in\langle M\rangle_\otimes$, the following are
equivalent:
\begin{enumerate}
 \item $W$ is globally nilpotent.
 \item The torus $T$ acts trivially on its solution representation.
 \item Over $K$, the module $W$ is isomorphic to
       $Y'=(N/z)Y$ for a constant nilpotent matrix $N$ over $k$.
 \item $F_W\subset K(L)$.
\end{enumerate}
In case (3), a suitable $K$-basis gives the fundamental matrix
$\exp(NL)$, a finite polynomial in $L$.  In an arbitrary initial
basis, a rational gauge factor can also occur.
\end{theorem}

\begin{proof}
First consider an arbitrary finite-dimensional algebraic
representation of the group in Proposition~\ref{at:ten:group}.
The two additive subgroups act as $\exp(bB)$ and $\exp(cC)$ for
commuting nilpotent matrices $B,C$.  Their common kernel is nonzero:
$\ker B$ is $C$-stable and the restriction of $C$ to it is nilpotent.
By \eqref{at:ten:conjugation} this common kernel is $T$-stable.
Since a torus representation is diagonalizable, it contains a
$T$-eigenline.  That line is invariant under the whole group.
Repeating in the quotient produces a full invariant flag.
Each one-dimensional factor has trivial additive action and torus
character $t\mapsto t^m$ for some integer $m$.  Under the differential
Galois correspondence it is the rank-one module with solution
$E^m=e^{-mz}$.

Suppose (1) holds.  The flag just constructed, its inclusions and
quotients, and the rational changes of basis identifying its
one-dimensional factors with the modules $E^m$, involve only
finitely many algebraic coefficients.  Choose one number field
containing them and a model for $W$.  After removing finitely many
places, these matrices, their required nonzero minors, and their
connection identities reduce correctly.  Consequently the flag
survives at the generic point in good characteristic, and
$p$-curvature on a factor is induced from the $p$-curvature on $W$.
Nilpotence therefore passes to every factor.

For the rank-one connection $\nabla_\partial=\partial+m$, however,
its curvature at a place of residue characteristic $p$ is
\[
 \psi_p=\nabla_\partial^p-\nabla_{\partial^{[p]}}=m^p.
\]
Here $\partial^{[p]}=0$ on the rational function field in
characteristic $p$.  If $m\ne0$, this scalar is nonzero at all places
except those over the finitely many rational primes dividing $m$.
Thus every factor has $m=0$.  The torus is diagonalizable, and all
of its weights on a representation are its weights on these flag
factors.  Hence $T$ acts trivially, proving (1)$\Rightarrow$(2).

If (2) holds, \eqref{at:ten:conjugation} implies
$\rho(U_c(c))=\rho(U_c(c/t))$ for every $t\ne0$.
The algebraic morphism $c\mapsto\rho(U_c(c))$ is then constant on
$k^\times$, hence on $k$, and its value at zero is the identity.
Therefore $U_c$ also acts trivially.  The representation factors
through the remaining additive group $U_b$.  It has the form
$b\mapsto\exp(bN)$ for a nilpotent constant matrix $N$.
The module $Y'=(N/z)Y$ has fundamental matrix
\[
 \exp(NL)=\sum_{j=0}^{r-1}\frac{N^jL^j}{j!},\qquad N^r=0,
\]
and precisely that representation.  The tensor correspondence now
gives (3), and also gives (4).  Conversely, (4) implies that the
subgroup fixing $L$, in particular $T$, acts trivially, so it
implies (2).

Finally, reduce a number-field model of $N$ at any good place.
For $\nabla_\partial=\partial-N/z$, direct iteration on constant
vectors gives
\[
 \psi_p=z^{-p}(N-N^p).
\]
Indeed, in characteristic $p$ the relevant polynomial identity is
$\prod_{j=0}^{p-1}(X+j)=X^p-X$; the general formula follows from
the linearity of $p$-curvature over the rational function field.
The matrix $N-N^p=N(I-N^{p-1})$ is nilpotent.  Thus (3) implies (1),
completing the equivalences.
\end{proof}

\subsection{Invariants and the remaining comparison constant}

\begin{proposition}\label{at:ten:invariants}
For $\mathcal R=K[E,E^{-1},L,J]$, one has
\[
 \mathcal R^T=K[L,J/E],\qquad
 \mathcal R^{\langle T,U_c\rangle}=K[L],\qquad
 F^{\langle T,U_c\rangle}=K(L).
\]
\end{proposition}

\begin{proof}
A monomial $E^aJ^bL^c$, with $a\in\mathbb Z$ and $b,c\ge0$,
has torus weight $a+b$.  The weight-zero monomials give the first
equality.  Invariance under every translation $J\mapsto J+c$
first removes dependence on $J$; invariance under every dilation of
$E$ then removes dependence on $E$.  This proves the second equality.
The same translation and dilation arguments for rational functions
prove the assertion about $F$.
\end{proof}

In particular, the single function $J/E$ has torus weight zero, but
this does not make its generated representation a trivial torus
representation.  In fact,
\[
 (J/E)'=J/E-1/z,\qquad
 \sigma(J/E)=J/E+c/(tE).
\]
Closure under the group action reintroduces the nonzero-weight
function $E^{-1}$.  This distinction prevents a weight-zero
normalization of one function from evading
Theorem~\ref{at:ten:classification}.

\begin{remark}[Scope and framing]\label{at:ten:framing}
For every complex constant $a$, the distinguished analytic solution
satisfies
\begin{equation}\label{at:ten:comparison}
 \begin{pmatrix}1\\E\\\Ein(z)-a\end{pmatrix}
 =\Phi(z)\begin{pmatrix}1\\1\\\gamma-a\end{pmatrix}.
\end{equation}
Thus $\gamma-a$ is a coordinate of an analytic comparison vector,
not a coefficient of \eqref{at:ten:connection}.  The classification of
globally nilpotent objects does not determine this coordinate.

For example, both $K(E,L,J)$ and $K(E,L,J+\gamma)$ are
Picard--Vessiot fields with constant field $k$, by the same
independence argument.  In their common analytic overfield their
compositum contains the constant $(J+\gamma)-J=\gamma$.
The assertion that each individual model has no new constants does
not assert that their analytically specified compositum has none.
Replacing this analytic comparison by an abstract
Picard--Vessiot isomorphism would discard the datum under study.

The theorem classifies all the stated finite tensor constructions
within this fixed module.  It does not cover adjoining independent
functions, differentiating an auxiliary parameter, applying a
Fourier--Laplace transform, taking an infinite limit, or imposing
new arithmetic constraints on a comparison vector.  In particular,
it proves neither the algebraicity nor the transcendence of
$\gamma$.
\end{remark}

\section{Independent samples, exact cancellation, and positivity}
\label{at:sec:statistics}

Let $T$ be a standard exponential random variable and let
$G_a=-\log T+a-\gamma$. Its density is strictly positive on $\R$.
The Gamma integral gives $\E\exp(t(-\log T))=\Gamma(1-t)$ for $t<1$;
differentiating the logarithmic expansion at $t=0$ gives
\begin{equation}
 \E G_a=a,\qquad\Var(G_a)=v:=\zeta(2).
\end{equation}
These identities follow from \cite[Sections 5.2 and 5.7]{at:DLMF}. At
$a=\gamma$ the sample is just $-\log T$, with no numerical input of
$\gamma$. The translated family is an analytical comparison family,
not an assertion that all its means equal Euler's constant.

\begin{theorem}[Unbiased polynomial and its second moment]
\label{at:sta:moments}
Let $P(X)=\sum_{k=0}^d p_kX^k\in\Q[X]$ have degree $d\ge1$. For $N\ge d$,
define the symmetric multi-affine polynomial
\begin{equation}
 U_{P,N}(x_1,\ldots,x_N)=
 \sum_{k=0}^d\frac{p_k}{\binom Nk}
 \sum_{\substack{I\subseteq\{1,\ldots,N\}\\ |I|=k}}\prod_{i\in I}x_i.
 \label{at:eq:unbiased}
\end{equation}
If its arguments are independent copies of $G_a$, then
\begin{align}
 \E_a U_{P,N}&=P(a),\label{at:eq:mean}\\
 \E_a U_{P,N}^2&=P(a)^2+
 \sum_{r=1}^d\left(\frac{P^{(r)}(a)}{r!}\right)^2
 \frac{v^r}{\binom Nr}.
 \label{at:eq:second-moment}
\end{align}
The same formula holds for any independent identically distributed
sample with mean $a$ and finite variance $v$.
\end{theorem}
\begin{proof}
Independence makes the expectation of each $k$-fold product equal to
$a^k$, proving~\eqref{at:eq:mean}. Put $Z_i=G_{a,i}-a$ and
\[
 V_{r,N}=\binom Nr^{-1}\sum_{|I|=r}\prod_{i\in I}Z_i,
 \qquad V_{0,N}=1.
\]
For a fixed $r$-element subset, counting the $k$-element subsets
containing it gives
$\binom{N-r}{k-r}/\binom Nk=\binom kr/\binom Nr$. Expanding each
factor as $a+Z_i$ therefore yields
\begin{equation}
 U_{P,N}=\sum_{r=0}^d\frac{P^{(r)}(a)}{r!}V_{r,N}.
 \label{at:eq:hoeffding}
\end{equation}
Products belonging to different subsets have a centered variable to
its first power, and hence expectation zero. Equal subsets contribute
$v^r$. Thus $\E V_{r,N}V_{s,N}=0$ for $r\ne s$ and
$\E V_{r,N}^2=v^r/\binom Nr$. Squaring~\eqref{at:eq:hoeffding} proves
the result. Only second moments are used, since each variable occurs
with exponent at most two after squaring.
\end{proof}

This is a direct U-statistic calculation in the sense of
\cite{at:Hoeffding1948}. At a real root $a$ of multiplicity $m$ of $P$,
the first surviving term in~\eqref{at:eq:second-moment} gives
\begin{equation}
 \E_a U_{P,N}^2\sim
 \frac{P^{(m)}(a)^2v^m}{m!}\,N^{-m}.
 \label{at:eq:root-asymptotic}
\end{equation}
It is strictly positive for every $N$, even if $P(a)=0$. For example,
$P(X)=qX-p$ and $a=p/q$ give $\E_a U_{P,N}^2=q^2v/N$.
Positive expectations tending to zero do not contradict rationality
of their mean.

\begin{theorem}[Exact integer jackknife]
\label{at:sta:jackknife}
For any integer $N_0\ge d$,
\begin{equation}
 P(a)^2=\sum_{j=0}^d(-1)^{d-j}\binom dj\binom{N_0+j}d
 \E_a U_{P,N_0+j}^2.
 \label{at:eq:jackknife}
\end{equation}
The weights are integers and do not depend on $a$ or on the sample
variance. For $N_0=d$, their absolute sum is at most
$(3+2\sqrt2)^d$.
\end{theorem}
\begin{proof}
Let $F(N)=\E_aU_{P,N}^2$ and let $(N)_d=N(N-1)\cdots(N-d+1)$.
Equation~\eqref{at:eq:second-moment} implies
\[
 (N)_dF(N)=P(a)^2(N)_d+
 \sum_{r=1}^d\left(\frac{P^{(r)}(a)}{r!}\right)^2
 v^r r!(N-r)_{d-r}.
\]
The right side is a polynomial of degree at most $d$, whose coefficient
of $N^d$ is $P(a)^2$. Its $d$th forward difference is $d!P(a)^2$.
Dividing that difference by $d!$ proves~\eqref{at:eq:jackknife}.
For $N_0=d$ the absolute sum satisfies
\[
 \sum_{j=0}^d\binom dj\binom{d+j}d
 =[t^d](1+t)^d(2+t)^d
 \le\left(\frac{(1+\sqrt2)(2+\sqrt2)}{\sqrt2}\right)^d
 =(3+2\sqrt2)^d.
\]
The coefficient identity follows by summing
$\binom dj(1+t)^{d+j}$ before taking the coefficient. The inequality
uses nonnegative coefficients evaluated at $t=\sqrt2$.
\end{proof}

For $d=1,2,3,4$, the weights at sample sizes $d,\ldots,2d$ are
respectively
\[
 (-1,2),\quad(1,-6,6),\quad(-1,12,-30,20),\quad
 (1,-20,90,-140,70).
\]
If $P\in\Z[X]$, multiplication by $((2d)!)^2$ clears all polynomial
kernel denominators in~\eqref{at:eq:jackknife} at $N_0=d$:
each $\binom Nk$ divides $N!$, and $N!$ divides $(2d)!$.
This is a nonoptimal denominator bound, not an arithmetic lower bound
for the expectation.

\begin{proposition}[Signs required for universal elimination]
\label{at:sta:signs}
A linear combination of second moments that preserves their constant
term and cancels their $1/N$ term identically must use a negative weight.
\end{proposition}
\begin{proof}
For sample sizes $N_j>0$, the two coefficient conditions are
$\sum_jw_j=1$ and $\sum_jw_j/N_j=0$. They cannot hold with every
$w_j\ge0$. These are conditions for universal coefficient cancellation;
no independence assumption about special values of $P^{(r)}(a)$ is needed.
\end{proof}

Independent blocks offer a different exact identity:
\begin{equation}
 \E_a[U_{P,N}(X)U_{P,N}(Y)]=P(a)^2.
 \label{at:eq:independent-product}
\end{equation}
Its integrand need not be nonnegative. In fact $U_{P,N}$ is
nonconstant, since $U_{P,N}(x,\ldots,x)=P(x)$. A nonconstant
multi-affine polynomial takes both signs: fix all but one variable
where the remaining affine slope is nonzero. The positive density of
the sample law gives each sign positive probability, so independent
blocks make the product negative with positive probability.

\begin{theorem}[Finite-sample nonnegative-kernel obstruction]
\label{at:sta:positive}
Fix a finite sample size $N$ and a measurable function
$K:\R^N\to[0,\infty]$ independent of $a$. Either $\E_aK=0$ for every
$a\in\R$, or $\E_aK>0$ for every $a\in\R$. Consequently, if a
nonzero real polynomial $P$ has a real root, no such kernel satisfies
$\E_aK=P(a)^2$ for every real $a$.
\end{theorem}
\begin{proof}
If $\E_rK=0$, then $K=0$ almost surely under the law of the sample
with mean $r$. Its density is strictly positive everywhere, so $K=0$
Lebesgue almost everywhere. Every other sample law is equivalent to
Lebesgue measure, proving $\E_aK=0$ for all $a$. Apply this at a root
of $P$ and then at a point where $P$ is nonzero.
\end{proof}
The proof is an absolute-continuity argument; compare the broader
constrained-estimation discussion in \cite{at:Somekh2016}. It excludes a
single kernel valid across the whole location family. It does not
exclude a new arithmetic identity proved only at $a=\gamma$.

\begin{remark}[The infinite-sample boundary]
For infinitely many independent samples the strong law gives
$\overline G_N\to a$ almost surely. The nonnegative functional
$P(\lim_N\overline G_N)^2$, defined as zero if the limit does not exist,
has expectation $P(a)^2$ for every $a$. There is no contradiction with
Theorem~\ref{at:sta:positive}: different means give mutually singular
infinite product laws. If $P(a)=0$, this functional is zero almost
surely and provides no nonvanishing certificate.
\end{remark}

\begin{proposition}[Quadratic algebraic-coefficient kernels]
\label{at:sta:quadratic}
Let $K(x)=x^{\mathsf T}Ax+b^{\mathsf T}x+c$ have real algebraic
coefficients, with $A$ symmetric, and suppose $K\ge0$ on $\R^N$.
If the polynomial $a\mapsto\E_aK$ has algebraic coefficients, then
$K$ is a nonnegative constant.
\end{proposition}
\begin{proof}
With $\mathbf1=(1,\ldots,1)^{\mathsf T}$,
\[
 \E_aK=a^2\mathbf1^{\mathsf T}A\mathbf1
 +a b^{\mathsf T}\mathbf1+c+\zeta(2)\tr A.
\]
Since $\zeta(2)=\pi^2/6$ is transcendental, the constant coefficient
can be algebraic only when $\tr A=0$. Nonnegativity of $K$ forces
$A$ to be positive semidefinite, so $A=0$. A globally nonnegative
affine function has $b=0$, and then $c\ge0$.
\end{proof}

\section{Exact denominator growth after division of Euler's series}
\label{at:den:section}

The distinction between an analytic zero and arithmetic divisibility can be
made explicit for Euler's factorial series.  The construction in this section
has two purposes: to compute the denominator loss exactly, and to identify
the additional assertion that would be needed to deduce transcendence from
a Borel-summed zero.  The weighted prime nonagreement result below is also a
consequence of an existing global relation theorem; no priority claim is made
for that consequence.

\subsection{Normalization and exact common denominators}

Put
\begin{equation}\label{at:den:def-series}
 f(z)=\sum_{n\geq0}(-1)^n n!z^n,
 \qquad S_n=\sum_{j=0}^n(-1)^j j!,
 \qquad
 g_a(z)=\frac{f(z)-a}{z-1}
       =\sum_{n\geq0}(a-S_n)z^n,
\end{equation}
where $a\in\Qbar$ and the quotient is taken in $\Qbar[[z]]$.  Its formal
Borel transform is
\begin{equation}\label{at:den:borel}
 G_a(z)=\sum_{n\geq0}b_n(a)z^n,
 \qquad b_n(a)=\frac{a-S_n}{n!}.
\end{equation}
For an algebraic number $b$, let $\operatorname{den}(b)$ be the least positive
integer $D$ for which $Db$ is an algebraic integer.  Write $D_N(a)$ for the
least positive integer that makes all of $b_0(a),\ldots,b_N(a)$ algebraic
integers.

We use the usual arithmetic Gevrey convention: a formal series
$\sum c_nz^n$ is an arithmetic Gevrey series of order $1$ when
$\sum(c_n/n!)z^n$ is a $G$-function; see
\cite{pm:FischlerRivoal2024} for this class and its summation conventions.
In particular, a bound of the form $|c_n|\leq C^{n+1}n!$ alone does
not establish this arithmetic condition.  A $G$-function must also have
algebraic coefficients, exponentially bounded conjugates and common
denominators, and satisfy a linear differential equation over $\Qbar(z)$.

\begin{proposition}[Exact denominator formula]\label{at:den:exact}
For $n\geq1$,
\begin{equation}\label{at:den:recurrence}
 nb_n(a)=b_{n-1}(a)-n(-1)^n,
 \qquad b_0(a)=a-1.
\end{equation}
Consequently,
\begin{equation}\label{at:den:nested}
 D_N(a)=\operatorname{den}(b_N(a)),
 \qquad D_{N-1}(a)\mid D_N(a).
\end{equation}
If $a=A/B\in\Q$, where $A,B\in\Z$, $B>0$, and $\gcd(A,B)=1$, then
\begin{equation}\label{at:den:rational-exact}
 D_N(A/B)=\frac{BN!}{\gcd(BN!,A-BS_N)}.
\end{equation}
The convention $\gcd(m,0)=m$ makes the formula valid when $A-BS_N=0$.
\end{proposition}

\begin{proof}
The identity $S_n-S_{n-1}=(-1)^n n!$ gives
\eqref{at:den:recurrence}.  If $Db_n(a)$ is an algebraic integer, then
\[
 Db_{n-1}(a)=nDb_n(a)+Dn(-1)^n
\]
is an algebraic integer as well.  Applying this implication successively
shows that the denominator of $b_N(a)$ clears all earlier coefficients.
The converse is part of the definition of a common denominator.  The
divisibility assertion follows, and reduction of
$(A-BS_N)/(BN!)$ gives \eqref{at:den:rational-exact}.
\end{proof}

\begin{theorem}[Factorial denominators for nonrational algebraic parameters]
\label{at:den:algebraic-factorial}
For every $a\in\Qbar\setminus\Q$,
\begin{equation}\label{at:den:factorial-growth}
 \log D_N(a)=\log N!+O_a(1).
\end{equation}
More explicitly, let $K=\Q(a)$, choose an integral basis
$1,\omega_1,\ldots,\omega_{d-1}$ of $\mathcal O_K$, and write
\[
 a=\frac{c_0+c_1\omega_1+\cdots+c_{d-1}\omega_{d-1}}{B},
 \qquad B>0,\quad c_i\in\Z.
\]
Then
\begin{equation}\label{at:den:integral-basis-formula}
 D_N(a)=
 \frac{BN!}{\gcd(BN!,c_0-BS_N,c_1,\ldots,c_{d-1})}.
\end{equation}
In particular, for $h=\gcd(c_1,\ldots,c_{d-1})>0$,
\begin{equation}\label{at:den:factorial-two-sided}
 \frac{BN!}{h}\leq D_N(a)\leq BN!.
\end{equation}
\end{theorem}

\begin{proof}
First, an integral basis containing $1$ exists.  Indeed,
$\mathcal O_K/\Z$ is torsion-free: if $m\beta\in\Z$ with
$\beta\in\mathcal O_K$ and $m\in\Z\setminus\{0\}$, then $\beta$ is a
rational algebraic integer and therefore belongs to $\Z$.  Thus $1$ is a
primitive vector in the free abelian group $\mathcal O_K$ and extends to a
basis.

In that basis the coordinates of $b_N(a)$ are
\[
 \frac{c_0-BS_N}{BN!},
 \frac{c_1}{BN!},\ldots,\frac{c_{d-1}}{BN!}.
\]
Integrality is coordinatewise, so Proposition~\ref{at:den:exact} gives
\eqref{at:den:integral-basis-formula}.  Since $a$ is not rational, at least
one $c_i$ with $i\geq1$ is nonzero.  The gcd in the denominator of
\eqref{at:den:integral-basis-formula} is therefore bounded above by the fixed
positive integer $h$.  This proves \eqref{at:den:factorial-two-sided} and
\eqref{at:den:factorial-growth}.
\end{proof}

\subsection{The Borel differential equation and its monodromy}

\begin{proposition}\label{at:den:borel-ode}
The germ $G_a$ is holomorphic at $0$ and satisfies
\begin{equation}\label{at:den:inhomogeneous}
 G_a'(z)=G_a(z)+\frac{1}{(1+z)^2},
 \qquad G_a(0)=a-1.
\end{equation}
Analytic continuation once counterclockwise around $-1$ changes it by
\begin{equation}\label{at:den:monodromy}
 \Delta G_a(z)=-2\pi i e^{z+1}.
\end{equation}
For every $a\in\Qbar$, the function $G_a$ is not a $G$-function, and
$g_a$ is not an arithmetic Gevrey series of order $1$.
\end{proposition}

\begin{proof}
The estimate $\sum_{j=0}^n j!\leq3n!$ shows that, for every embedding
$\sigma:\Q(a)\hookrightarrow\C$,
\begin{equation}\label{at:den:conjugate-bound}
 |\sigma(b_n(a))|\leq |\sigma(a)|+3.
\end{equation}
Thus the germ converges for $|z|<1$.  Summing
\eqref{at:den:recurrence} gives \eqref{at:den:inhomogeneous}, hence
\begin{equation}\label{at:den:integral-solution}
 G_a(z)=e^z\left(a-1+\int_0^z\frac{e^{-t}}{(1+t)^2}\,dt\right)
\end{equation}
along any path avoiding $-1$.  The integrand has residue $-e$ there;
the residue theorem gives \eqref{at:den:monodromy}.

Suppose that $G_a$ were a $G$-function.  The minimal differential
operator annihilating a $G$-function is regular singular at every point,
including infinity; see \cite[Section~1]{pm:FischlerRivoal2016} for this
standard consequence of the theory of $G$-operators.  Analytic
continuations of $G_a$, and therefore their differences, solve the same
homogeneous equation.  Equation~\eqref{at:den:monodromy} would make $e^z$
a solution.  This is incompatible with regular singularity at infinity:
after finite ramification, such solutions have power--logarithm local
expansions, whereas $e^z$ has exponential growth on a ray to infinity.
The contradiction proves the assertion about $G_a$; the statement about
$g_a$ follows from the Borel normalization.
\end{proof}

\begin{remark}[An explicit minimal operator]\label{at:den:minimal-operator}
The homogeneous operator
\begin{equation}\label{at:den:operator}
 \left(\partial+\frac{2}{1+z}\right)(\partial-1),
 \qquad \partial=\frac{d}{dz},
\end{equation}
annihilates $G_a$ and has minimal order $2$, even over $\C(z)$.
For if $G_a'=rG_a$ with $r\in\C(z)$, then
$(r-1)G_a=(1+z)^{-2}$.  The case $r=1$ is impossible, and otherwise
$G_a$ would be rational.  A rational solution $R$ of
$R'-R=(1+z)^{-2}$ can have no finite pole except $-1$.  At $-1$ its
principal part would have to be $-1/(1+z)$, leaving a nonzero simple
pole in $R'-R$, a contradiction.  This also gives a direct check of
minimality.  The coefficient of $\partial$ in \eqref{at:den:operator} is
$-1+2/(1+z)$, which fails the Fuchs condition at infinity.
\end{remark}

\begin{remark}[Cancellation on one ray does not remove monodromy]
\label{at:den:gompertz-branch}
For $x>0$ put $E_1(x)=\int_x^\infty e^{-t}\,dt/t$ and
$\delta=eE_1(1)$.  Formula~\eqref{at:den:integral-solution} gives, on the
positive real axis,
\begin{equation}\label{at:den:closed-form}
 G_a(z)=e^z(a-\delta)-\frac{1}{1+z}+e^{z+1}E_1(z+1).
\end{equation}
The term $e^z(a-\delta)$ vanishes at $a=\delta$.  The monodromy
\eqref{at:den:monodromy} does not.  Thus removal of an exponentially
growing term on one branch is insufficient to establish the arithmetic
condition governing all branches and normalized denominators.
\end{remark}

The coefficients of $G_a$ belong to a fixed number field,
\eqref{at:den:conjugate-bound} supplies the conjugate bound, and
\eqref{at:den:operator} supplies holonomicity.  Hence its failure to be a
$G$-function is exactly the failure of an exponential bound on $D_N(a)$.
This already implies
$\limsup_{N\to\infty}\log D_N(a)/N=+\infty$.
Obtaining a limit rather than a limsup for rational $a$ requires the
following argument.

\subsection{Rational parameters and weighted prime nonagreement}

For each prime $p$, the alternating series $f$ in
\eqref{at:den:def-series} converges in $\Q_p$ at $z=1$.
Set
\begin{equation}\label{at:den:epadics}
 E_p=\sum_{n\geq0}(-1)^n n!\in\Z_p,
 \qquad
 \mathcal E(a)=\{p:E_p\ne a\text{ in }\Q_p\}
 \quad(a\in\Q).
\end{equation}

\begin{theorem}[Weighted nonagreement and denominator growth]
\label{at:den:rational-growth}
For every $a\in\Q$,
\begin{equation}\label{at:den:weighted-divergence}
 \sum_{p\in\mathcal E(a)}\frac{\log p}{p-1}=+\infty
\end{equation}
and
\begin{equation}\label{at:den:superexponential}
 \lim_{N\to\infty}\frac{\log D_N(a)}{N}=+\infty.
\end{equation}
The first assertion is a consequence of the global relation theorem
of Ernvall-Hyt\"onen, Matala-aho, and Sepp\"al\"a
\cite[Theorem~1]{at:Ernvall2018}.  A separate proof through
Proposition~\ref{at:den:borel-ode} is included to connect that assertion
directly to the exact denominators.
\end{theorem}

\begin{proof}
Write $a=A/B$ in lowest terms with $B>0$.  If
$p\notin\mathcal E(a)$, then $a=E_p\in\Z_p$, so $p\nmid B$ and
\begin{equation}\label{at:den:tail-valuation}
 v_p(a-S_N)=v_p(E_p-S_N)\geq v_p((N+1)!).
\end{equation}
Indeed, every term of the $p$-adically convergent tail is divisible by
$(N+1)!$ in $\Z_p$.  Such a prime divides none of the $D_N(a)$, and
\eqref{at:den:rational-exact} yields
\begin{equation}\label{at:den:support-divisibility}
 D_N(a)\mid B\prod_{p\in\mathcal E(a)}p^{v_p(N!)}.
\end{equation}
The product is finite for each $N$.  If the sum in
\eqref{at:den:weighted-divergence} had finite value $C$, Legendre's
bound $v_p(N!)\leq N/(p-1)$ would give
\[
 \log D_N(a)\leq\log B+CN.
\]
Together with the conjugate bound and the differential equation, this
would make $G_a$ a $G$-function, contradicting
Proposition~\ref{at:den:borel-ode}.  This proves
\eqref{at:den:weighted-divergence}.

Fix now $p\in\mathcal E(a)$ and put
$t_p=v_p(A-BE_p)<\infty$.  Convergence in $\Q_p$ gives
$v_p(A-BS_N)=t_p$ for all sufficiently large $N$.  Increasing $N$
further if necessary, \eqref{at:den:rational-exact} therefore gives
\begin{equation}\label{at:den:stable-valuation}
 v_p(D_N(a))=v_p(B)+v_p(N!)-t_p.
\end{equation}
For every finite $T\subset\mathcal E(a)$ it follows that
\begin{equation}\label{at:den:finite-liminf}
 \liminf_{N\to\infty}\frac{\log D_N(a)}{N}
 \geq\sum_{p\in T}\frac{\log p}{p-1},
\end{equation}
since $v_p(N!)/N\to1/(p-1)$ at each fixed prime.  The right-hand side
can be made arbitrarily large by
\eqref{at:den:weighted-divergence}.  This proves
\eqref{at:den:superexponential}.
\end{proof}

\begin{remark}[Derivation from the existing global relation theorem]
\label{at:den:prior-global-relation}
For completeness, the precise implication from
\cite[Theorem~1]{at:Ernvall2018} can be checked without any prime-density
claim.  Specializing that theorem to the argument $-1$, a set $T$ of
primes has infinitely many members with $A-BE_p\ne0$ whenever, for
every finite $S\subset T$,
\begin{equation}\label{at:den:prior-condition}
 \limsup_{n\to\infty}
 4^n n!\prod_{p\in T\setminus S}|n!|_p^2=0.
\end{equation}
If \eqref{at:den:weighted-divergence} were false, with finite sum $C$,
take $T$ to be the complement of $\mathcal E(a)$.  The product formula
and Legendre's bound imply
\begin{align*}
 4^n n!\prod_{p\in T\setminus S}|n!|_p^2
 &=\frac{4^n}{n!}
   \prod_{p\in\mathcal E(a)\cup S}|n!|_p^{-2}\\
 &\leq\frac{\exp\bigl((\log4+2C+2C_S)n\bigr)}{n!}
 \longrightarrow0,
 \qquad C_S=\sum_{p\in S}\frac{\log p}{p-1}.
\end{align*}
The cited theorem would then give a prime of nonagreement inside the
set $T$ of agreement primes, a contradiction.  Accordingly,
\eqref{at:den:weighted-divergence} is recorded here as a known consequence
and as an input to the explicit denominator conclusion, not as a new
prime-distribution theorem.
\end{remark}

\begin{corollary}\label{at:den:all-algebraic}
For every $a\in\Qbar$ and every real $C>1$, there is $N_0(a,C)$
such that $D_N(a)>C^N$ for every $N\geq N_0(a,C)$.  Equivalently,
\[
 \lim_{N\to\infty}\frac{\log D_N(a)}N=+\infty.
\]
\end{corollary}

\begin{proof}
For nonrational algebraic $a$, use
Theorem~\ref{at:den:algebraic-factorial}.  For rational $a$, use
Theorem~\ref{at:den:rational-growth}.
\end{proof}

\begin{remark}
For rational $a$, these arguments do not establish
$\log D_N(a)\sim\log N!$.  The all-$N$ superexponential conclusion
must be distinguished from that stronger factorial asymptotic.
Nor does the prime $p$ in \eqref{at:den:epadics} remain fixed when a
rational candidate $a$ is changed.
\end{remark}

\subsection{Finite certificates and the exact scope of the infinite family}

\begin{proposition}[Finite certificates of nonagreement]
\label{at:den:certificates}
Let $a=A/B\in\Q$ be in lowest terms, with $B>0$.  A triple
$(p,N,r)$, where $p$ is prime, $N\geq0$, and $r\geq1$, certifies
$E_p\ne a$ if
\begin{equation}\label{at:den:certificate-test}
 v_p(B)+v_p((N+1)!)\geq r,
 \qquad A-BS_N\not\equiv0\pmod{p^r}.
\end{equation}
Conversely, every prime of $\mathcal E(a)$ admits such a certificate.
For each positive rational $R$, there is a terminating procedure
that produces a finite certified set $T_R\subset\mathcal E(a)$ with
\begin{equation}\label{at:den:certificate-weight}
 \sum_{p\in T_R}\frac{\log p}{p-1}>R.
\end{equation}
\end{proposition}

\begin{proof}
The tail bound gives
$B(E_p-S_N)\equiv0\pmod{p^r}$ under the first condition in
\eqref{at:den:certificate-test}.  The second condition therefore implies
$A-BE_p\ne0$.  Conversely, if
$t_p=v_p(A-BE_p)<\infty$, take $r=t_p+1$ and then take $N$ large
enough that the tail has valuation at least $r$.

Enumerate primes $p$, integers $N,r$, and check
\eqref{at:den:certificate-test} using integer arithmetic.  Remove repeated
primes.  Simultaneously refine rational lower bounds for the finitely
many logarithms already encountered.  By
\eqref{at:den:weighted-divergence}, some finite certified set has total
weight strictly greater than $R$; its sufficiently accurate rational
lower bound will eventually certify the same strict inequality.
Thus the procedure terminates.
\end{proof}

This is an effective infinite family of finite certificates concerning
different $p$-adic values.  It is not a family of certificates proving
the irrationality of one fixed $E_p$, of the real constant $\delta$,
or of $\gamma$.

\subsection{What zero lifting would have to add}

The positive-direction Borel sum of $f$ at $1$ is
\[
 f_0(1)=\int_0^\infty\frac{e^{-t}}{1+t}\,dt=\delta.
\]
In particular, the restriction of the zero-lifting conjecture
\cite[Conjecture~2]{pm:FischlerRivoal2024} to this single family is the
assertion
\begin{equation}\label{at:den:delta-lifting}
 \text{for every }a\in\Qbar,\qquad
 \delta=a\ \Longrightarrow\
 g_a\text{ is arithmetic Gevrey of order }1.
\end{equation}

\begin{proposition}\label{at:den:delta-equivalence}
Assertion~\eqref{at:den:delta-lifting} is equivalent to the
transcendence of $\delta$.
\end{proposition}

\begin{proof}
If $\delta$ is transcendental, every antecedent in
\eqref{at:den:delta-lifting} is false.  If $\delta$ is algebraic,
take $a=\delta$.  The conclusion of
\eqref{at:den:delta-lifting} contradicts
Proposition~\ref{at:den:borel-ode}.
\end{proof}

Thus the exact denominator calculation establishes the obstruction
on the conclusion side of \eqref{at:den:delta-lifting}; it does not
establish the implication.  Ordinary Gevrey growth and formal
division do not supply that missing arithmetic assertion.

For comparison with the earlier mixed-function formulation for
$\gamma$, define
\begin{equation}\label{at:den:gamma-mixed}
 \widehat D(z)=\sum_{n\geq0}(-1)^n n!z^{-n-1},
 \qquad F_a(z)=e^z(\Ein(z)-a),
 \qquad\Psi_a(z)=F_a(z)-\widehat D(z),
\end{equation}
where
\[
 \Ein(z)=\sum_{n\geq1}\frac{(-1)^{n-1}z^n}{n\,n!}.
\]
Here a mixed arithmetic Gevrey function means a bilateral formal sum
$H(z)+h(1/z)$ with $H$ an $E$-function and $h$ arithmetic Gevrey of
order $1$, as in \cite[Definition~1]{pm:FischlerRivoal2024}.  Multiplication
by a polynomial is well-defined coefficientwise.  The expression
``division by $z-1$ within the mixed class'' means the existence of a
mixed function $\Theta$ satisfying $(z-1)\Theta=\Psi_a$; it does not
presuppose a field structure on all bilateral sums.

On the positive real axis,
$\Ein(z)=\gamma+\log z+E_1(z)$ and
$\widehat D_0(z)=e^zE_1(z)$.  Consequently,
\begin{equation}\label{at:den:gamma-sum}
 \Psi_{a,0}(z)=e^z(\gamma+\log z-a),
 \qquad\Psi_{a,0}(1)=e(\gamma-a).
\end{equation}

\begin{proposition}[The earlier mixed zero-lifting equivalence]
\label{at:den:gamma-equivalence}
For every $a\in\Qbar$, there is no mixed arithmetic Gevrey function
$\Theta$ such that $(z-1)\Theta=\Psi_a$.  Consequently the assertion
\begin{equation}\label{at:den:gamma-lifting}
 \text{for every }a\in\Qbar,\qquad
 \Psi_{a,0}(1)=0\ \Longrightarrow\
 \Psi_a\text{ is divisible by }z-1\text{ in the mixed class}
\end{equation}
is equivalent to the transcendence of $\gamma$.
\end{proposition}

\begin{proof}
Suppose $\Theta=H(z)+h(1/z)$ exists.  Absorb the constant term of
$h$ into $H$ and write $h(1/z)=\sum_{j\geq1}h_jz^{-j}$, where
$h_j\in\Qbar$.  Comparing the nonnegative powers of $z$ in
$(z-1)\Theta=F_a-\widehat D$ gives the entire-function identity
\[
 F_a(z)=(z-1)H(z)+h_1.
\]
Hence $F_a(1)=h_1$ is algebraic.  However, $e$ and $\Ein(1)$ are
algebraically independent, as follows from
\cite[Section~5(ii)]{at:Rivoal2012}.  Therefore
$F_a(1)=e(\Ein(1)-a)$ cannot be algebraic: an algebraic value $b$
would give the nonzero relation $X(Y-a)-b=0$ at
$(X,Y)=(e,\Ein(1))$.  This proves nondivisibility.

By \eqref{at:den:gamma-sum}, the antecedent of
\eqref{at:den:gamma-lifting} is exactly $\gamma=a$.  If $\gamma$ is
transcendental all such antecedents are false; if it is algebraic,
choosing $a=\gamma$ contradicts the nondivisibility just proved.
\end{proof}

Proposition~\ref{at:den:gamma-equivalence} restates the bridge already
identified in the preceding manuscript and specializes the mixed
zero-lifting framework of \cite{pm:FischlerRivoal2024}.  It is not a new
transcendence result.  The additional content of the present section
is the explicit denominator analysis of the Euler component and its
connection with prime nonagreement.  Neither equivalence proves its
own zero-lifting antecedent-to-conclusion assertion.

\section{Remaining arithmetic assertions}
\label{at:sec:remaining}

The preceding sections complete several auxiliary constructions. Their
endpoints do not identify which polynomials vanish at $\gamma$.
The following distinctions fix the logical status of the program.

\begin{enumerate}
\item \textbf{Relations and completions.}
To obtain a homomorphism $\Z[\gamma]\to\Z_p$ sending $\gamma$ to $a$,
one must prove $P(\gamma)=0\Rightarrow P(a)=0$ for every
$P\in\Z[X]$. The explicitly constructed simultaneous limits do not
give this implication. The counterexamples in
Proposition~\ref{at:app:counterexample} establish this failure without
any assumption about $\gamma$.
\item \textbf{Modules and comparison vectors.}
Theorem~\ref{at:ten:classification} classifies the stated differential
modules. It does not determine the analytic comparison coordinate
$\gamma-a$ in~\eqref{at:ten:comparison}. An arithmetic theorem about that
coordinate would be an additional result.
\item \textbf{Exact cancellation and nonvanishing.}
Theorem~\ref{at:sta:jackknife} cancels the second-moment correction terms
exactly. Its signed combination does not retain positivity. Theorem
\ref{at:sta:positive} excludes a nonnegative replacement valid across the
whole finite-sample location family; an identity special to the
actual Gamma integral would require additional information.
\item \textbf{Analytic zeros and arithmetic division.}
Corollary~\ref{at:den:all-algebraic} describes the normalized denominator
loss of a formal quotient. The implication from a zero of a Borel sum
to division in the arithmetic class is unproved. Propositions
\ref{at:den:delta-equivalence} and \ref{at:den:gamma-equivalence} explain
why, for these specific families, proving that implication already
settles the corresponding transcendence question.
\end{enumerate}

These are limitations of the specified constructions, not impossibility
theorems for every method of proving transcendence. The paper makes no
claim that the missing implications must follow after further numerical
calculation or after enlarging the number of fields considered. Neither
arithmetic alternative for $\gamma$ has been excluded here.


\section{Transfer certificates and reproducibility}
\label{at:app:verification}

The accompanying scripts use only the Python standard library. Their
finite checks are separated from the proofs of the infinite assertions.
They should be run with assertions enabled:
\begin{verbatim}
python3 verify_adelic_candidates.py
python3 verify_positive_tensor.py
python3 verify_zero_lifting.py
\end{verbatim}

\paragraph{Simultaneous approximation.}
The first script checks 25 rational candidates for $\gamma$ at five
levels and five prescribed $2$-adic targets. It checks the certified
real bounds, congruences, compatibility, and the denominator cost for
target zero. It also checks 40 rational counterexamples, 20 each for
$x=1/3$ and $x=1/2$. The proof uses $m=3n$; the implementation chooses
the first $m\le3n$ giving sufficient interval width.
\paragraph{Statistical identities.}
The second script checks 90 exact cancellation identities for degrees
up to 12 and 14 exhaustive finite-distribution cases for an asymmetric
three-point law. The direct algebraic proofs in
Section~\ref{at:sec:statistics} apply to all degrees and sample sizes.
\paragraph{Factorial denominators.}
The third script checks 2,173 coefficient and common-denominator cases
for 53 rational parameters at $N=0,\ldots,40$, 41 denominator cases
for $\sqrt2$, and 18 explicit finite $p$-adic nonagreement certificates.
The limits in Section~\ref{at:den:section} are proved there; they are not
inferred from these finite computations.

The approximation and zero-lifting scripts write JSON certificates.
The statistical script prints its check record. A SHA-256 manifest
identifies the delivered versions. Exact arithmetic prevents rounding
errors in these checks but does not turn them into formal proofs of
the general theorems. Independent mathematical review was performed
for the main arguments; no proof-assistant certification or external
peer review is claimed.

\section{Transfer revision record}
\label{at:app:revision}

Revision 2 integrates the five proof notes into a standalone paper. It
makes the computability input, reduced versus common denominator
distinction, finite tensor-category restriction, and finite-sample
location-family restriction explicit. It separates a denominator
limsup from the all-index limit, treats nonrational algebraic parameters
using an integral basis, and defines mixed division by a polynomial
without assuming a field of unrestricted bilateral series. The
prime-discrepancy result is attributed to its existing global-relation
source, and the Katz reference is identified as the 1982 paper.

The previously reported finite checks are retained, with the same
counts. This revision adds a complete exposition and corrects scope
and bibliographic details; it adds no claim of irrationality or
transcendence for $\gamma$ or $\delta$.


% END INTEGRATED PART VII
\clearpage
\part{Unified targets, corrections, and the source record}
\label{master:closing}
\section{What would complete an arithmetic argument?}
\label{master:remaining}

The constructions in the preceding parts have several outputs that must
not be identified with one another. A nonzero vector in an integer kernel
can have zero image in the target quotient. A nonzero target polynomial
can vanish at an algebraic argument. A polynomial that does not vanish can
still have an insufficiently small value relative to its primitive height.
These are separate assertions, and the surviving targets below keep them
separate.

\paragraph{Primitive rational forms.}
For the positive-moment route, a sufficient completed output would be
coprime integers $p_n,q_n$, with $q_n\to\infty$, and one fixed
$\varepsilon>0$ such that
\begin{equation}\label{master:roth-output}
 0<|q_n\gamma-p_n|<q_n^{-1-\varepsilon}
 \quad\text{for infinitely many }n.
\end{equation}
The strict inequality excludes rational $\gamma$ once the denominators are
large, and Roth's theorem excludes every irrational algebraic $\gamma$.
Thus this is a transcendence criterion. The one certified exponent
$103/100$ in the positive-moment part is an instance at a finite target;
it is not an infinite family satisfying \eqref{master:roth-output}.
The denominator here is the reduced target denominator, after cancellation
of every common factor. The size of an auxiliary residue vector is not a
substitute for it.

\paragraph{Small polynomial values.}
Part~\ref{master:infinite} states the degree-sensitive polynomial criteria
and proves what algebraicity would force for the corresponding total
selection rule. Producing nonzero \emph{polynomials} under that rule is not
the unresolved step: their \emph{values} may eventually vanish if
$\gamma$ is algebraic. A successful completion must force escape from the
relation ideal with the stated degree and height budgets, or independently
certify sufficiently many nonzero values. Neither finite modular tests nor
selection from a shrinking real interval automatically supplies this.

\paragraph{Local arithmetic.}
Part~\ref{master:worlds} turns the two alternatives into exact statements
about maps from $\Z[\gamma]$ to residue rings. Those statements are useful
because they specify an observable which would distinguish the cases.
Part~\ref{master:transfer} constructs prescribed real and $2$-adic limits,
but also proves that these two limits need not preserve polynomial
relations. A completion by this route must establish the missing
well-defined map from the actual ring $\Z[\gamma]$, with enough depth or
unbounded counts to apply the local criterion. Choosing residues freely in
a rational approximation does not establish that map.

\paragraph{Differential and mixed-series arithmetic.}
The differential reductions classify objects in a specified tensor category;
the factorial-quotient theorems classify denominators of a specified
division operation. They are completed results about those objects.
Their relation to \eqref{master:target} still requires an arithmetic
specialization or zero-lifting assertion. In particular, a common analytic
value of an $E$-series and a summed divergent series cannot simply be divided
out in the arithmetic category. Assuming that operation would assume the
very bridge the construction is intended to prove.

\paragraph{The two-case strategy.}
Let A denote ``$\gamma$ is algebraic'' and B denote ``$\gamma$ is
transcendental''. Deductions made under A can exclude A only when every
additional premise in the contradiction has independent justification.
An attractive property obtained in B cannot be imported as an established
fact into A. Conversely, the fact that many finite observations can be
modeled in both cases does not establish that either case is consistent
with every theorem. The local camouflage results quantify a particular
finite limitation; they do not settle the original alternative.

\section{Corrections carried into the cumulative edition}
\label{master:corrections}

\begin{enumerate}[leftmargin=*,itemsep=6pt]
\item \textbf{Construction tests have a specified scope.}
Independence from unknown digits, a fixed target space, favorable error
versus height, and value nonvanishing are useful tests for an explicit
integer-form strategy. They are not necessary conditions for every
possible proof of transcendence. A construction from certified rational
intervals is not circular merely because the intervals locate $\gamma$.
It still needs its own arithmetic nonvanishing argument.

\item \textbf{Dirichlet gives existence, not an upper ceiling.}
A Dirichlet exponent cannot be used as an unconditional upper bound on
the approximation exponent of an unknown number. In particular, an
asserted universal bound such as $\tau_R\le2$ does not follow from the
pigeonhole principle. Any elimination estimate that uses such a bound
must state it as an additional hypothesis, not as a completed obstruction.

\item \textbf{Three nonvanishing statements differ.}
An integer coefficient vector, its primitive target, and the evaluated
linear form can vanish independently at successive stages. Positivity
is a value certificate only when it applies to the actual final form.
Siegel or Minkowski bounds on a kernel vector do not alone locate that
vector in a positive cone or outside the target kernel.

\item \textbf{Finite improvements have not been extrapolated.}
The finite $103/100$ exponent, interval contractions, numerical
asymptotic experiments, and exploratory large-$N$ thresholds are retained
with their actual status. A fixed-margin infinite family with unbounded
primitive denominator is not inferred from a finite run. Explicit horizon
claims require effective inequalities and constants, not only an
asymptotic comparison.

\item \textbf{Target height is used after reduction.}
A large common denominator or a large representation coefficient is not
by itself a lower bound for reduced rational height or minimal-polynomial
height. Upper bounds, lower bounds, and exact asymptotics are distinguished
throughout the moment, Todd, and moving-parameter constructions.

\item \textbf{A strict Stokes transition is distinguished from a comparison
entry.}
The elementary lateral jump of the exponential integral is determined by
$2\pi i$. The appearance of $\gamma$ in a regular-to-asymptotic comparison
requires the relevant normalizations or framing. The phrase ``$\gamma$ is
a Stokes constant'' is not used to infer an arithmetic theorem for the
unframed jump matrix.

\item \textbf{Finite camouflage uses the sharp degree.}
The earlier degree $p^L$ construction is replaced in the operative
statement by the minimum $\mu(p^L)=\min\{m:p^L\mid m!\}$.
The older construction remains a valid, nonoptimal example in its source
record. The exact root-count theorem concerns the relation ideal, not a
selected polynomial with an unrelated algebraic root.

\item \textbf{The four transfer mechanisms retain their corrected
hypotheses.}
Approximation denominators are positive; the statistical polynomial has
exact degree $d\ge1$; the cleared second-moment polynomial has degree
\emph{at most} $d$; factorial-series sign and argument conventions are
fixed explicitly. Mixed divisibility means existence of a series in the
stated class satisfying the product identity. No closure under analytic
division is silently added.
\end{enumerate}

These corrections do not turn an unfinished transcendence proof into a
finished one. They ensure that the positive results and the limits of each
construction can be assessed without inheriting an earlier overstatement.

\section{Coverage and reproducibility}
\label{master:coverage}

The companion Korean guide contains a file-by-file coverage table. Each
entry identifies its destination in this manuscript or the preserved
source calculation and records whether the content is proved, conditional,
finite, experimental, or superseded. The accompanying archive contains:
\begin{enumerate}[leftmargin=*]
\item this single-file LaTeX source and its compiled PDF;
\item the Korean reading guide, source-claim ledgers, and a machine-readable
manifest with SHA-256 digests;
\item all 59 identified source files unchanged, including the 27 original
research ZIP packages, their proof notes, programs, and certificates;
\item the editorial section fragments and assembly script used for this
edition, together with a record of compilation and document checks.
\end{enumerate}

The new LaTeX file compiles independently: it requires no external image,
\texttt{input} file, or bibliography database. To rebuild it with standard
TeX Live packages, run \texttt{pdflatex} repeatedly until the references and
table of contents stabilize. To reproduce an older numerical certificate,
first unpack its named source ZIP and follow that package's README; paths
in the preserved mathematical parts are relative to those unpacked
packages. An old verification record is identified as inherited evidence,
not described as a new execution performed for this edition.

The integration checks cover source inventory, section and citation
resolution, formula layout, and targeted mathematical comparisons. They do
not constitute independent external peer review of every theorem in every
unpublished source package. In particular, preserving a source file in the
archive does not endorse its superseded wording. The corrected statements
and explicit status labels in this manuscript govern the cumulative
edition.


\clearpage
\begin{thebibliography}{999}
\small
\addcontentsline{toc}{section}{References}
\bibitem{fd:DLMF}
NIST Digital Library of Mathematical Functions,
\S5.9 (Binet and digamma integral representations) and \S6.2
(exponential-integral definitions and connection formulas).
\url{https://dlmf.nist.gov/5.9}; \url{https://dlmf.nist.gov/6.2}.

\bibitem{fd:Kolchin}
E.~R.~Kolchin, \emph{Differential Algebra and Algebraic Groups},
Academic Press, 1973.  Primitive Kolchin--Ostrowski dependence theorem.

\bibitem{fd:FJ}
J.~Fres\'an and P.~Jossen, \emph{Exponential Motives},
research monograph manuscript, \S12.8, in the version cited by the
companion papers.  Euler--Mascheroni motive and non-splitting.
\url{http://javier.fresan.perso.math.cnrs.fr/expmot.pdf}.

\bibitem{fd:Prevost}
M.~Pr\'evost, A family of criteria for irrationality of Euler's constant,
\emph{J. Comput. Appl. Math.} \textbf{219} (2008), 484--492.
\url{https://arxiv.org/abs/math/0507231}.

\bibitem{fd:early}
KEIO, \emph{Gamma Transcendence Attempt: Interior-Zero Kernels and
Arithmetic Deficits}, research note, 4 September 2026.
Archived as \texttt{Gamma\_Transcendence\_Attempt\_2026-09-04.md},
including the exact saddle certificate and logarithm-cancellation data.

\bibitem{fd:multiscale}
KEIO, \emph{Gamma Stieltjes Multiscale Attack}, research note,
version 1, archived as
\texttt{KEIO\_Gamma\_Stieltjes\_Multiscale\_Attack\_v1.md}.
Read together with the corrections in \cite{fd:positivev2}.

\bibitem{fd:positivev2}
KEIO, \emph{Gamma Rational Positive Certificate}, research note,
version 2, archived as
\texttt{KEIO\_Gamma\_Rational\_Positive\_Certificate\_v2.md}.

\bibitem{fd:uniformv3}
KEIO, \emph{Gamma Uniform Family Attack}, research note,
version 3, archived as
\texttt{KEIO\_Gamma\_Uniform\_Family\_Attack\_v3.md}, together with
\texttt{v3\_geometric\_family\_proof.md} and exact primal--dual certificates.





\bibitem{pm:Baker1966}
A. Baker,
\emph{Linear forms in the logarithms of algebraic numbers. I},
Mathematika \textbf{13} (1966), 204--216.
\url{https://doi.org/10.1112/S0025579300003971}.

\bibitem{pm:Beukers2006}
F. Beukers,
\emph{A refined version of the Siegel--Shidlovskii theorem},
Ann. of Math. (2) \textbf{163} (2006), 369--379.
\url{https://arxiv.org/abs/math/0405549}.

\bibitem{pm:Blagouchine2017}
I. V. Blagouchine,
\emph{A note on some recent results for the Bernoulli numbers of the
second kind}, J. Integer Seq. \textbf{20} (2017), Article 17.3.8.
\url{https://arxiv.org/abs/1612.03292}.

\bibitem{pm:BuchstaberVeselov2024}
V. M. Buchstaber and A. P. Veselov,
\emph{Todd polynomials and Hirzebruch numbers},
Proc. Steklov Inst. Math. \textbf{325} (2024), 74--85.
\url{https://doi.org/10.1134/S0081543824020044}.

\bibitem{pm:FischlerRivoal2016}
S. Fischler and T. Rivoal,
\emph{Arithmetic theory of E-operators},
J. \'Ec. polytech. Math. \textbf{3} (2016), 31--65.
\url{https://doi.org/10.5802/jep.28}.

\bibitem{pm:FischlerRivoal2024}
S. Fischler and T. Rivoal,
\emph{Relations between values of arithmetic Gevrey series, and
applications to values of the Gamma function},
J. Number Theory \textbf{261} (2024), 36--54.
\url{https://doi.org/10.1016/j.jnt.2024.02.016}.

\bibitem{pm:KEIO001}
KEIO,
\emph{Exponential-Integral Periods across Arithmetic, Pad\'e, Hankel,
Reciprocal Geometry, and Frobenius Framing},
unpublished companion manuscript [001], Version 6.3, 31 August 2026.
Author-supplied source is included with the reproducibility archive.

\bibitem{pm:KEIO002}
KEIO,
\emph{The Polyexponential Trace-Tower Program: Moving Levels,
Faber and Inverse Monodromy, Positive Spectra, and Arithmetic Barriers},
unpublished companion manuscript [002], Version 3.1, 1 September 2026.
Author-supplied source is included with the reproducibility archive.

\bibitem{pm:RosserSchoenfeld1962}
J. B. Rosser and L. Schoenfeld,
\emph{Approximate formulas for some functions of prime numbers},
Illinois J. Math. \textbf{6} (1962), 64--94.
\url{https://doi.org/10.1215/ijm/1255631807}.

\bibitem{pm:Roth1955}
K. F. Roth,
\emph{Rational approximations to algebraic numbers},
Mathematika \textbf{2} (1955), 1--20; corrigendum, 168.
\url{https://doi.org/10.1112/S0025579300000644}.

\bibitem{pm:Sondow2003}
J. Sondow,
\emph{Criteria for irrationality of Euler's constant},
Proc. Amer. Math. Soc. \textbf{131} (2003), 3335--3344.
\url{https://arxiv.org/abs/math/0209070}.

\bibitem{pm:SondowZlobin2004}
J. Sondow and S. Zlobin,
\emph{A hypergeometric approach, via linear forms involving logarithms,
to criteria for irrationality of Euler's constant},
Math. Slovaca \textbf{59} (2009), no. 3, 307--314.
\url{https://doi.org/10.2478/s12175-009-0127-2}.

\bibitem{pm:WolfsVanAssche2025}
W. Van Assche and T. Wolfs,
\emph{Rational approximation of Euler's constant using multiple
orthogonal polynomials},
Combinatorics and Number Theory \textbf{14} (2025), no. 2, 141--162.
\url{https://doi.org/10.2140/cnt.2025.14.141}.

\bibitem{pm:WZ1992}
H. S. Wilf and D. Zeilberger,
\emph{An algorithmic proof theory for hypergeometric (ordinary and
``$q$'') multisum/integral identities},
Invent. Math. \textbf{108} (1992), 575--633.
\url{https://sites.math.rutgers.edu/~zeilberg/mamarimY/Zeilberger_y1992_p575.pdf}.

\bibitem{ec:pilehrood}
Kh. Hessami Pilehrood and T. Hessami Pilehrood,
\emph{On a continued fraction expansion for Euler's constant},
\url{https://arxiv.org/abs/1010.1420}.

\bibitem{ec:dlmf}
NIST Digital Library of Mathematical Functions, Chapters 5, 10 and 18,
\url{https://dlmf.nist.gov/}.

\bibitem{ec:source-records}
KEIO, \emph{Gamma research records of 6 September 2026},
the eight dated reports and their exact proof packages reproduced in the
source archive of this cumulative edition.  Internal research sources;
not a claim of external peer review or historical priority.
\bibitem{ar:SondowVacca}
J. Sondow, \emph{New Vacca-type rational series for Euler's constant and
its ``alternating'' analog $\log(4/\pi)$}, arXiv:math/0508042.
\url{https://arxiv.org/abs/math/0508042}.

\bibitem{ar:HardouinSinger}
C. Hardouin and M. F. Singer, \emph{Differential Galois theory of linear
difference equations}, Math. Ann. \textbf{342} (2008), 333--377.
\url{https://arxiv.org/abs/0801.1493}.

\bibitem{ar:DLMFGamma}
NIST Digital Library of Mathematical Functions, Chapter 5, especially
Sections 5.5 and 5.15 (Gamma multiplication and polygamma identities).
\url{https://dlmf.nist.gov/5.5}; \url{https://dlmf.nist.gov/5.15}.

\bibitem{ar:Ridout}
D. Ridout, \emph{The $p$-adic generalization of the Thue--Siegel--Roth
theorem}, Mathematika \textbf{5} (1958), 40--48.
\url{https://doi.org/10.1112/S0025579300001339}.

\bibitem{ar:AB}
B. Adamczewski and Y. Bugeaud, \emph{On the complexity of algebraic
numbers I. Expansions in integer bases}, Ann. of Math. \textbf{165}
(2007), 547--565.
\url{https://annals.math.princeton.edu/2007/165-2/p04}.

\bibitem{ar:Christol}
G. Christol, T. Kamae, M. Mend\`es France and G. Rauzy,
\emph{Suites alg\'ebriques, automates et substitutions},
Bull. Soc. Math. France \textbf{108} (1980), 401--419.
\url{https://www.numdam.org/articles/10.24033/bsmf.1926/}.

\bibitem{ar:PolyaCarlson}
J. Bell, R. Miles and T. Ward, \emph{Towards a P\'olya--Carlson dichotomy
for algebraic dynamics}, Indag. Math. \textbf{25} (2014), 652--668.
\url{https://arxiv.org/abs/1307.2369}. Used for its stated classical
P\'olya--Carlson input.

\bibitem{ar:Ivic}
A. Ivi\'c, \emph{On the integral of the error term in the Dirichlet
divisor problem}, arXiv:math/0510114.
\url{https://arxiv.org/abs/math/0510114}.

\bibitem{ar:LauTsang}
Y.-K. Lau and K.-M. Tsang, \emph{Mean square of the remainder term in
the Dirichlet divisor problem}, J. Th\'eor. Nombres Bordeaux
\textbf{7} (1995), 75--92.
\url{https://www.numdam.org/item/JTNB_1995__7_1_75_0.pdf}.

\bibitem{ar:DixitKumar}
A. Dixit and R. Kumar, \emph{Superimposing theta structure on a
generalized modular relation}, arXiv:2005.08316, especially Section 2.1.
\url{https://arxiv.org/abs/2005.08316}.

\bibitem{ar:DLMFBessel}
NIST Digital Library of Mathematical Functions, Sections 10.31, 10.32
and 10.38 (modified Bessel expansions, integrals, and order derivatives).
\url{https://dlmf.nist.gov/10.31}; \url{https://dlmf.nist.gov/10.32};
\url{https://dlmf.nist.gov/10.38}.

\bibitem{ar:Stanley}
R. P. Stanley, \emph{Differentiably finite power series}, European
J. Combin. \textbf{1} (1980), 175--188.
\url{https://math.mit.edu/~rstan/pubs/pubfiles/45.pdf}.

\bibitem{ar:Blagouchine}
I. V. Blagouchine and M.-A. Coppo, \emph{A note on some constants related
to the zeta-function and their relationship with the Gregory
coefficients}, arXiv:1703.08601.
\url{https://arxiv.org/abs/1703.08601}.
See also I. V. Blagouchine, \emph{Three notes on Ser's and Hasse's
representations for the zeta-functions}, Integers \textbf{18A} (2018), A3.
\url{https://arxiv.org/abs/1606.02044}.

\bibitem{ar:MZ}
T. Matala-aho and W. Zudilin, \emph{Euler's factorial series and global
relations}, J. Number Theory \textbf{186} (2018), 202--210.
\url{https://arxiv.org/abs/1703.02633}.

\bibitem{ar:Gould}
H. W. Gould and J. Quaintance, \emph{A linear binomial recurrence and
the Bell numbers and polynomials} (2007).
\url{https://doiserbia.nb.rs/ft.aspx?id=1452-86300702371G}.

\bibitem{ar:DLMFLaguerre}
NIST Digital Library of Mathematical Functions, equation 18.14.8,
Laguerre polynomial inequality.
\url{https://dlmf.nist.gov/18.14.E8}.
\bibitem{inf:Erdos}
P. Erd\H{o}s, \emph{Beweis eines Satzes von Tschebyschef}, Acta Litt. Sci. Szeged \textbf{5} (1932), 194--198.
\url{https://users.renyi.hu/~p_erdos/1932-01.pdf}.

\bibitem{inf:FischlerOsc}
S. Fischler, \emph{Nesterenko's criterion when the small linear forms oscillate}, arXiv:1201.2651.
\url{https://arxiv.org/abs/1201.2651}.

\bibitem{inf:FischlerVectors}
S. Fischler, \emph{Nesterenko's linear independence criterion for vectors}, arXiv:1202.2279.
\url{https://arxiv.org/abs/1202.2279}.

\bibitem{inf:DLMFGamma}
NIST Digital Library of Mathematical Functions, \emph{Gamma Function: Integral Representations}, Section 5.9, especially equation 5.9.10.
\url{https://dlmf.nist.gov/5.9}.

\bibitem{inf:DLMFHahn}
NIST Digital Library of Mathematical Functions, \emph{Continuous Hahn polynomials: definitions and explicit representations}, Sections 18.19--18.20, equations 18.19.2--18.19.4 and 18.20.9.
\url{https://dlmf.nist.gov/18.19}; \url{https://dlmf.nist.gov/18.20}.

\bibitem{inf:Baker}
A. Baker, \emph{Linear forms in the logarithms of algebraic numbers (III)}, Mathematika \textbf{14} (1967), 220--228.
\url{https://doi.org/10.1112/S0025579300003843}.

\bibitem{inf:Varju}
P. P. Varj\'u, \emph{On the dimension of Bernoulli convolutions for all transcendental parameters}, arXiv:1810.08905.
\url{https://arxiv.org/abs/1810.08905}.

\bibitem{inf:Hochman}
M. Hochman, \emph{On self-similar sets with overlaps and inverse theorems for entropy}, Annals of Mathematics \textbf{180} (2014), 773--822; arXiv:1212.1873.
\url{https://arxiv.org/abs/1212.1873}.

\bibitem{inf:BV}
E. Breuillard and P. P. Varj\'u, \emph{Entropy of Bernoulli convolutions and uniform exponential growth for linear groups}, arXiv:1510.04043v4, 2021.
\url{https://arxiv.org/abs/1510.04043}.
\bibitem{wr:stacks-differentials}
The Stacks Project Authors, \emph{Differentials}, Tags 00RM and following.
\url{https://stacks.math.columbia.edu/tag/00RM}.

\bibitem{wr:stacks-etale}
The Stacks Project Authors, \emph{Formally \`etale maps}, Tag 00UP.
\url{https://stacks.math.columbia.edu/tag/00UP}.

\bibitem{wr:milne}
J. S. Milne, \emph{Algebraic Number Theory}, course notes,
Section 8 (Frobenius and the Chebotarev density theorem).
\url{https://www.jmilne.org/math/CourseNotes/ANT.pdf}.

\bibitem{wr:ekt}
C. Elsholtz, B. Klahn, and M. Technau, On polynomials with roots modulo
almost all primes, \emph{Acta Arith.} \textbf{205} (2022), 251--263.
\url{https://arxiv.org/abs/2206.13466}.

\bibitem{wr:hensel}
K. Conrad, \emph{Hensel's Lemma}, expository notes, Theorem 4.1.
\url{https://kconrad.math.uconn.edu/blurbs/gradnumthy/hensel.pdf}.

\bibitem{wr:bhargava}
M. Bhargava, Generalized factorials and fixed divisors over subsets of
a Dedekind domain, \emph{J. Number Theory} \textbf{72} (1998), 67--75.
\url{https://doi.org/10.1006/jnth.1998.2220}.

\bibitem{wr:dlmf}
NIST Digital Library of Mathematical Functions, Chapter 5,
Sections 5.2, 5.5, 5.7, and 5.8 (Gamma integral, reflection,
multiplication, expansion, and product).
\url{https://dlmf.nist.gov/5}.

\bibitem{wr:turing}
A. M. Turing, On computable numbers, with an application to the
Entscheidungsproblem, \emph{Proc. London Math. Soc.} (2) \textbf{42}
(1936--1937), 230--265.
\url{https://www.cs.virginia.edu/~robins/Turing_Paper_1936.pdf}.


\bibitem{at:DLMF}
NIST Digital Library of Mathematical Functions,
Sections 5.2, 5.7 and 6.2.
\url{https://dlmf.nist.gov/5.2},
\url{https://dlmf.nist.gov/5.7},
\url{https://dlmf.nist.gov/6.2}.

\bibitem{at:Ernvall2018}
A.-M. Ernvall-Hyt\"onen, T. Matala-aho and L. Sepp\"al\"a,
\emph{Euler's divergent series in arithmetic progressions},
arXiv:1809.03859 (2018).
\url{https://arxiv.org/abs/1809.03859}.

\bibitem{at:Hoeffding1948}
W. Hoeffding,
\emph{A class of statistics with asymptotically normal distribution},
Ann. Math. Statist. \textbf{19} (1948), no.~3, 293--325.
\url{https://doi.org/10.1214/aoms/1177730196}.

\bibitem{at:Katz1982}
N. M. Katz,
\emph{A conjecture in the arithmetic theory of differential equations},
Bull. Soc. Math. France \textbf{110} (1982), 203--239.
\url{https://doi.org/10.24033/bsmf.1960}.

\bibitem{at:Rivoal2012}
T. Rivoal,
\emph{On the arithmetic nature of the values of the Gamma function,
Euler's constant and Gompertz's constant},
Michigan Math. J. \textbf{61} (2012), 239--254.
\url{https://rivoal.perso.math.cnrs.fr/articles/gammater.pdf}.

\bibitem{at:Somekh2016}
A. Somekh-Baruch, A. Leshem and V. Saligrama,
\emph{On the non-existence of unbiased estimators in constrained
estimation problems}, arXiv:1609.07415 (2016).
\url{https://arxiv.org/abs/1609.07415}.

\bibitem{at:Srinivasan2008}
V. Ravi Srinivasan,
\emph{Extensions by antiderivatives, exponentials of integrals and
by iterated logarithms}, arXiv:0811.3004 (2008).
\url{https://arxiv.org/abs/0811.3004}.


\end{thebibliography}
\end{document}
